% GENERATED FILE -- DO NOT EDIT.
%
% Produced by harness/assemble_contract.py from module.yaml and claims/*.tex.
% Edit the claim files and regenerate; `--check` fails if this file is stale.
\documentclass[11pt,a4paper]{article}
\pdftrailerid{e39ac026d212eb3be0ae6bc32f5103b8}
%% --- 1. Mathematics (order matters: physics after braket wins \ket, \tr) ---
\usepackage{amsmath}
\usepackage{amssymb}
\usepackage{amsfonts}
\usepackage{amsthm}
\usepackage{mathtools}
\usepackage{dsfont}
\usepackage{braket}
\usepackage{physics}
%% --- 2. Page furniture ------------------------------------------------------
\usepackage{xcolor}
\usepackage{array}    % >{...} column prefixes in the wide tables of some claims
\usepackage{enumitem}
\usepackage{graphicx}
\usepackage{booktabs}
\usepackage{longtable}
\usepackage{xltabular}              % front- and back-matter tables across pages
\usepackage{needspace}
\usepackage{xurl}
\usepackage{tcolorbox}
\tcbuselibrary{skins,breakable}
%% --- 3. Links and navigation (hyperref late; bookmark, refcount after it) ---
\usepackage{xr-hyper}   % before hyperref: the contract's labels
\usepackage{hyperref}
\usepackage{bookmark}
\usepackage{refcount}
%% --- 4. Colours: every distinction also has a shape (rule, caps, size, weight)
\definecolor{nnGrey}{gray}{0.40}      % non-normative text
\definecolor{metaGrey}{gray}{0.30}    % the generated meta line
\definecolor{techGrey}{gray}{0.45}    % heading of a technical claim
\definecolor{mainRule}{RGB}{0,110,140}
\definecolor{reqAccent}{RGB}{165,70,0}
\definecolor{linkBlue}{RGB}{0,62,150}
\definecolor{extLink}{RGB}{120,40,140}  % links into another PDF
\hypersetup{colorlinks=true, linkcolor=linkBlue, citecolor=linkBlue,
  urlcolor=linkBlue, filecolor=extLink, bookmarksnumbered=true,
  bookmarksopen=true, bookmarksopenlevel=1, pdfdisplaydoctitle=true}
% A link into another PDF (\href to a .pdf, \hyperref{file}..., xr-hyper) is a
% relative URI, <file>#nameddest=<destination>, in the file-link colour:
% browser viewers (Edge, Chrome) ignore hyperref's default GoToR action.
\makeatletter
\renewcommand*{\hyper@linkfile}[3]{% anchor text, file name, destination
  \begingroup
    \ifx\\#3\\\def\CQ@target{#2}\else\def\CQ@target{#2\hyper@hash nameddest=#3}\fi
    \let\@urlcolor\@filecolor
    \hyper@linkurl{#1}{\CQ@target}%
  \endgroup}
\makeatother
%% --- 5. The mapped Unicode set (harness/contract_navigation.py, UNICODE_TEX)
\DeclareUnicodeCharacter{03B1}{\ensuremath{\alpha}}
\DeclareUnicodeCharacter{03B2}{\ensuremath{\beta}}
\DeclareUnicodeCharacter{03B3}{\ensuremath{\gamma}}
\DeclareUnicodeCharacter{03B4}{\ensuremath{\delta}}
\DeclareUnicodeCharacter{03B5}{\ensuremath{\varepsilon}}
\DeclareUnicodeCharacter{03F5}{\ensuremath{\epsilon}}
\DeclareUnicodeCharacter{03B6}{\ensuremath{\zeta}}
\DeclareUnicodeCharacter{03B7}{\ensuremath{\eta}}
\DeclareUnicodeCharacter{03B8}{\ensuremath{\theta}}
\DeclareUnicodeCharacter{03D1}{\ensuremath{\vartheta}}
\DeclareUnicodeCharacter{03B9}{\ensuremath{\iota}}
\DeclareUnicodeCharacter{03BA}{\ensuremath{\kappa}}
\DeclareUnicodeCharacter{03BB}{\ensuremath{\lambda}}
\DeclareUnicodeCharacter{03BC}{\ensuremath{\mu}}
\DeclareUnicodeCharacter{03BD}{\ensuremath{\nu}}
\DeclareUnicodeCharacter{03BE}{\ensuremath{\xi}}
\DeclareUnicodeCharacter{03C0}{\ensuremath{\pi}}
\DeclareUnicodeCharacter{03C1}{\ensuremath{\rho}}
\DeclareUnicodeCharacter{03F1}{\ensuremath{\varrho}}
\DeclareUnicodeCharacter{03C3}{\ensuremath{\sigma}}
\DeclareUnicodeCharacter{03C2}{\ensuremath{\varsigma}}
\DeclareUnicodeCharacter{03C4}{\ensuremath{\tau}}
\DeclareUnicodeCharacter{03C5}{\ensuremath{\upsilon}}
\DeclareUnicodeCharacter{03C6}{\ensuremath{\varphi}}
\DeclareUnicodeCharacter{03D5}{\ensuremath{\phi}}
\DeclareUnicodeCharacter{03C7}{\ensuremath{\chi}}
\DeclareUnicodeCharacter{03C8}{\ensuremath{\psi}}
\DeclareUnicodeCharacter{03C9}{\ensuremath{\omega}}
\DeclareUnicodeCharacter{0393}{\ensuremath{\Gamma}}
\DeclareUnicodeCharacter{0394}{\ensuremath{\Delta}}
\DeclareUnicodeCharacter{0398}{\ensuremath{\Theta}}
\DeclareUnicodeCharacter{039B}{\ensuremath{\Lambda}}
\DeclareUnicodeCharacter{039E}{\ensuremath{\Xi}}
\DeclareUnicodeCharacter{03A0}{\ensuremath{\Pi}}
\DeclareUnicodeCharacter{03A3}{\ensuremath{\Sigma}}
\DeclareUnicodeCharacter{03A5}{\ensuremath{\Upsilon}}
\DeclareUnicodeCharacter{03A6}{\ensuremath{\Phi}}
\DeclareUnicodeCharacter{03A8}{\ensuremath{\Psi}}
\DeclareUnicodeCharacter{03A9}{\ensuremath{\Omega}}
\DeclareUnicodeCharacter{00E1}{\'a}
\DeclareUnicodeCharacter{00E0}{\`a}
\DeclareUnicodeCharacter{00E2}{\^a}
\DeclareUnicodeCharacter{00E4}{\"a}
\DeclareUnicodeCharacter{00E5}{\aa{}}
\DeclareUnicodeCharacter{00E9}{\'e}
\DeclareUnicodeCharacter{00E8}{\`e}
\DeclareUnicodeCharacter{00EA}{\^e}
\DeclareUnicodeCharacter{00EB}{\"e}
\DeclareUnicodeCharacter{00ED}{\'{\i}}
\DeclareUnicodeCharacter{00EF}{\"{\i}}
\DeclareUnicodeCharacter{00F3}{\'o}
\DeclareUnicodeCharacter{00F2}{\`o}
\DeclareUnicodeCharacter{00F4}{\^o}
\DeclareUnicodeCharacter{00F6}{\"o}
\DeclareUnicodeCharacter{00F8}{\o{}}
\DeclareUnicodeCharacter{00FA}{\'u}
\DeclareUnicodeCharacter{00F9}{\`u}
\DeclareUnicodeCharacter{00FB}{\^u}
\DeclareUnicodeCharacter{00FC}{\"u}
\DeclareUnicodeCharacter{00F1}{\~n}
\DeclareUnicodeCharacter{00E7}{\c{c}}
\DeclareUnicodeCharacter{00DF}{\ss{}}
\DeclareUnicodeCharacter{00C4}{\"A}
\DeclareUnicodeCharacter{00C5}{\AA{}}
\DeclareUnicodeCharacter{00C9}{\'E}
\DeclareUnicodeCharacter{00D6}{\"O}
\DeclareUnicodeCharacter{00DC}{\"U}
\DeclareUnicodeCharacter{2013}{\textendash{}}
\DeclareUnicodeCharacter{2014}{\textemdash{}}
\DeclareUnicodeCharacter{2018}{`}
\DeclareUnicodeCharacter{2019}{'}
\DeclareUnicodeCharacter{201C}{``}
\DeclareUnicodeCharacter{201D}{''}
\DeclareUnicodeCharacter{2026}{\ldots{}}
\DeclareUnicodeCharacter{00A0}{~}
\DeclareUnicodeCharacter{00A7}{\S{}}
\DeclareUnicodeCharacter{2264}{\ensuremath{\le}}
\DeclareUnicodeCharacter{2265}{\ensuremath{\ge}}
\DeclareUnicodeCharacter{2260}{\ensuremath{\neq}}
\DeclareUnicodeCharacter{2248}{\ensuremath{\approx}}
\DeclareUnicodeCharacter{2261}{\ensuremath{\equiv}}
\DeclareUnicodeCharacter{223C}{\ensuremath{\sim}}
\DeclareUnicodeCharacter{2243}{\ensuremath{\simeq}}
\DeclareUnicodeCharacter{2245}{\ensuremath{\cong}}
\DeclareUnicodeCharacter{221D}{\ensuremath{\propto}}
\DeclareUnicodeCharacter{226A}{\ensuremath{\ll}}
\DeclareUnicodeCharacter{226B}{\ensuremath{\gg}}
\DeclareUnicodeCharacter{2192}{\ensuremath{\to}}
\DeclareUnicodeCharacter{2190}{\ensuremath{\leftarrow}}
\DeclareUnicodeCharacter{2194}{\ensuremath{\leftrightarrow}}
\DeclareUnicodeCharacter{21D2}{\ensuremath{\Rightarrow}}
\DeclareUnicodeCharacter{21D0}{\ensuremath{\Leftarrow}}
\DeclareUnicodeCharacter{21D4}{\ensuremath{\Leftrightarrow}}
\DeclareUnicodeCharacter{21A6}{\ensuremath{\mapsto}}
\DeclareUnicodeCharacter{2208}{\ensuremath{\in}}
\DeclareUnicodeCharacter{2209}{\ensuremath{\notin}}
\DeclareUnicodeCharacter{2286}{\ensuremath{\subseteq}}
\DeclareUnicodeCharacter{2282}{\ensuremath{\subset}}
\DeclareUnicodeCharacter{2287}{\ensuremath{\supseteq}}
\DeclareUnicodeCharacter{2283}{\ensuremath{\supset}}
\DeclareUnicodeCharacter{2229}{\ensuremath{\cap}}
\DeclareUnicodeCharacter{222A}{\ensuremath{\cup}}
\DeclareUnicodeCharacter{2216}{\ensuremath{\setminus}}
\DeclareUnicodeCharacter{2205}{\ensuremath{\emptyset}}
\DeclareUnicodeCharacter{2200}{\ensuremath{\forall}}
\DeclareUnicodeCharacter{2203}{\ensuremath{\exists}}
\DeclareUnicodeCharacter{2115}{\ensuremath{\mathbb{N}}}
\DeclareUnicodeCharacter{2124}{\ensuremath{\mathbb{Z}}}
\DeclareUnicodeCharacter{211A}{\ensuremath{\mathbb{Q}}}
\DeclareUnicodeCharacter{211D}{\ensuremath{\mathbb{R}}}
\DeclareUnicodeCharacter{2102}{\ensuremath{\mathbb{C}}}
\DeclareUnicodeCharacter{00B7}{\ensuremath{\cdot}}
\DeclareUnicodeCharacter{00D7}{\ensuremath{\times}}
\DeclareUnicodeCharacter{2212}{\ensuremath{-}}
\DeclareUnicodeCharacter{00B1}{\ensuremath{\pm}}
\DeclareUnicodeCharacter{2213}{\ensuremath{\mp}}
\DeclareUnicodeCharacter{221A}{\ensuremath{\surd}}
\DeclareUnicodeCharacter{2211}{\ensuremath{\sum}}
\DeclareUnicodeCharacter{220F}{\ensuremath{\prod}}
\DeclareUnicodeCharacter{222B}{\ensuremath{\int}}
\DeclareUnicodeCharacter{221E}{\ensuremath{\infty}}
\DeclareUnicodeCharacter{2202}{\ensuremath{\partial}}
\DeclareUnicodeCharacter{2207}{\ensuremath{\nabla}}
\DeclareUnicodeCharacter{2297}{\ensuremath{\otimes}}
\DeclareUnicodeCharacter{2295}{\ensuremath{\oplus}}
\DeclareUnicodeCharacter{2218}{\ensuremath{\circ}}
\DeclareUnicodeCharacter{2308}{\ensuremath{\lceil}}
\DeclareUnicodeCharacter{2309}{\ensuremath{\rceil}}
\DeclareUnicodeCharacter{230A}{\ensuremath{\lfloor}}
\DeclareUnicodeCharacter{230B}{\ensuremath{\rfloor}}
\DeclareUnicodeCharacter{27E8}{\ensuremath{\langle}}
\DeclareUnicodeCharacter{27E9}{\ensuremath{\rangle}}
\DeclareUnicodeCharacter{2016}{\ensuremath{\|}}
\DeclareUnicodeCharacter{2113}{\ensuremath{\ell}}
\DeclareUnicodeCharacter{00B0}{\ensuremath{^{\circ}}}
\DeclareUnicodeCharacter{00B2}{\ensuremath{^{2}}}
\DeclareUnicodeCharacter{00B3}{\ensuremath{^{3}}}
\DeclareUnicodeCharacter{00B9}{\ensuremath{^{1}}}
\DeclareUnicodeCharacter{2080}{\ensuremath{_{0}}}
\DeclareUnicodeCharacter{2081}{\ensuremath{_{1}}}
\DeclareUnicodeCharacter{2082}{\ensuremath{_{2}}}
\DeclareUnicodeCharacter{207F}{\ensuremath{^{n}}}
\DeclareUnicodeCharacter{2096}{\ensuremath{_{k}}}
%% --- 6. Numbering: one counter per section -----------------------
\numberwithin{equation}{section}
\setcounter{tocdepth}{1}
%% --- 7. The project macro set, then the claim environments -----------------
% spec/macros.tex, read relative to this module's directory (latexmk -cd).
% =============================================================================
%  spec/macros.tex -- the one macro set every generated contract inputs.
%
%  Two parts: part A, the macros spec/notation.md catalogues, and part B,
%  the macros local to one module, grouped by module.
%  harness/assemble_contract.py emits % =============================================================================
%  spec/macros.tex -- the one macro set every generated contract inputs.
%
%  Two parts: part A, the macros spec/notation.md catalogues, and part B,
%  the macros local to one module, grouped by module.
%  harness/assemble_contract.py emits % =============================================================================
%  spec/macros.tex -- the one macro set every generated contract inputs.
%
%  Two parts: part A, the macros spec/notation.md catalogues, and part B,
%  the macros local to one module, grouped by module.
%  harness/assemble_contract.py emits \input{../../macros} for it, so the
%  file is read from spec/modules/<id>/.
%
%  RULES
%    * A macro that spec/notation.md catalogues is a verbatim copy of it
%      (part A). Changing one is a notation change: bump spec/notation.md.
%    * Part B holds the module-local macros. A macro defined by two modules
%      with the same expansion appears once, under the first module (in the
%      order below) that defines it. Two macros with the same name and
%      DIFFERENT expansions are both kept under a module prefix, and the
%      claims write the prefixed names:
%        \Chain  -> \redChain (reductions, \mathfrak{C}: the reduction chain)
%                   \pcChain  (protocol-compiler, \mathfrak{R}: the two-segment
%                             chain; \Chsrc and \Chrec are built on it)
%        \Nmsm   -> \redNmsm  (reductions, \mathcal{N}: a POVM measurement
%                             channel)
%                   \pcNmsm   (\mathcal{N}_{\mathsf A}: Alice's measurement
%                             channel, written by
%                             transmitter/thm:certificate-sufficiency,
%                             transmitter/rem:xreg and
%                             reductions/lem:source-transfer)
%    * Deliberately omitted: a \provenance{tag}{details} macro. Provenance
%      is a header field of every claim file (spec/claim-format.md section
%      3); a claim body that calls such a macro fails to compile, which is
%      the intended signal.
%
%  REQUIRED PACKAGES, loaded by the assembler's preamble before this file:
%    amsmath, amssymb, amsfonts, mathtools, dsfont, braket, physics, xcolor
%    (braket before physics, so that physics wins for \ket, \bra, \tr, ...).
%
%  This file contains only definitions: no \usepackage, no document body.
% =============================================================================


% =============================================================================
%  PART A -- notation macros, the ones spec/notation.md catalogues, each a
%  verbatim copy of its catalogue row; grouped under numbered headings.
% =============================================================================

% -----------------------------------------------------------------------------
% 1. State spaces and operator sets (a "system label" -> spaces of operators)
% -----------------------------------------------------------------------------
\newcommand{\Hs}[1]{\mathcal{H}(#1)}        % Hilbert space of a system label
\newcommand{\Lin}[1]{\mathcal{L}(#1)}       % bounded linear operators on H(#1)
\newcommand{\Linn}[2]{\mathcal{L}(#1,#2)}   % linear maps H(#1) -> H(#2)
\newcommand{\Pos}[1]{\mathrm{Pos}(#1)}      % positive semidefinite operators on H(#1)
\newcommand{\Density}{\mathrm{D}}           % density operators: \Density(A) = D(A)
\newcommand{\Seq}[1]{S_{=}(#1)}             % normalized states (trace = 1)
\newcommand{\Sleq}[1]{S_{\leq}(#1)}         % subnormalized states (0 <= trace <= 1)
\newcommand{\id}{\mathds{1}}                % identity operator
\newcommand{\idon}[1]{\mathds{1}_{#1}}      % identity on a named system, e.g. \idon{C}

% DEPRECATED: applies \mathcal to both arguments itself, so it takes bare
% letters, not system labels. Prefer \Linn; do not use in new contracts.
\providecommand{\Lop}[2]{\mathcal{L}(\mathcal{#1}, \mathcal{#2})}


% -----------------------------------------------------------------------------
% 2. Composite systems
% -----------------------------------------------------------------------------
\newcommand{\dsum}{\oplus}                  % direct sum of system labels
% Tensor product: use \otimes directly.
% n-fold direct sum of an operator space: \Lin{A}^{\oplus n}


% -----------------------------------------------------------------------------
% 3. Events, conditioning and classical registers (for STATES)
%    Convention:  \wedge  -> subnormalized "and-event" component, e.g. \rho_{A\wedge x}
%                 |        -> normalized conditional state,        e.g. \rho_{A|x}
% (Write these directly in subscripts; no extra macros needed.)
% -----------------------------------------------------------------------------


% -----------------------------------------------------------------------------
% 4. Photon-number subspaces / direct-sum blocks
%    Convention:  operators carry the restriction "on top", in brackets;
%                 states carry it as a \wedge subscript (see Sec. 3).
% -----------------------------------------------------------------------------
\newcommand{\Pnum}[1]{P_{#1}}               % photon-number projector (system from context)
\newcommand{\Pnumon}[2]{P_{#1}^{#2}}        % with explicit system superscript, e.g. P_n^C
\newcommand{\blk}[2]{#1^{(#2)}}             % operator block, e.g. \blk{M}{\leqN} = M^{(<=N)}
\newcommand{\leqN}{\leq N}                  % use as \blk{M}{\leqN}
\newcommand{\gtN}{> N}                      % use as \blk{M}{\gtN}
\newcommand{\eqn}{= n}                      % use as \blk{M}{\eqn}


% -----------------------------------------------------------------------------
% 5. System labels used in the protocol
% -----------------------------------------------------------------------------
\newcommand{\Chan}{C}                       % quantum-channel (Eve) input to Bob
\newcommand{\Prb}{P}                        % probe input
\newcommand{\Flag}{{\mathrm{Flg}}}          % flag system (EXISTS ONLY AFTER squashing)
\newcommand{\Loss}{{\mathrm{Lss}}}          % loss system
\newcommand{\Vc}{{\mathrm{Vcm}}}            % vacuum system
\newcommand{\Nout}{N_{\mathrm{out}}}
\newcommand{\Nin}{N_{\mathrm{in}}}


% -----------------------------------------------------------------------------
% 6. Alphabets
% -----------------------------------------------------------------------------
\newcommand{\Yset}{\mathcal{Y}}             % Bob's (coarse-grained) detection alphabet
\newcommand{\Yraw}{\mathcal{Y}_{\mathrm{raw}}}% raw click patterns {0,1}^4
\newcommand{\Yscset}{\mathcal{Y}_{\singlec}}% single-click outcomes
\newcommand{\Ymcset}{\mathcal{Y}_{\mc}}     % multi-click outcomes
\newcommand{\Tset}{\mathcal{T}_{\mathrm{set}}}% private setting set (the probe settings in the optical instance)
\newcommand{\Xset}{\mathcal{X}}             % Alice's alphabet
\newcommand{\Dset}{\{H,V,+,-\}}             % polarization / detector labels
\newcommand{\pass}{\top}                    % pass sifting


% -----------------------------------------------------------------------------
% 7. Mode / outcome abbreviations
% -----------------------------------------------------------------------------
\newcommand{\on}{{\mathrm{on}}}
\newcommand{\off}{{\mathrm{off}}}
\newcommand{\dc}{{\mathrm{dc}}}             % dark count
\newcommand{\sq}{\mathrm{sq}}               % squashed
\newcommand{\mc}{\mathrm{mc}}               % multi-click
\newcommand{\nc}{\mathrm{nc}}               % no-click
\newcommand{\singlec}{\mathrm{sc}}          % single-click
\newcommand{\vac}{\mathrm{vac}}             % vacuum
\newcommand{\kout}{{k_\mathrm{out}}}
\newcommand{\kin}{{k_\mathrm{in}}}


% -----------------------------------------------------------------------------
% 8. Maps, channels and miscellaneous operators
% -----------------------------------------------------------------------------
\newcommand{\Prob}{\mathrm{Prob}}           % probability (use \Prob[\cdot] or \Prob(\cdot))
\newcommand{\lmin}{\lambda_{\min}}          % minimum eigenvalue
\newcommand{\Wflag}{W_{\Flag}}              % flag-space weight
\newcommand{\transpose}{\mathsf{T}}
\newcommand{\innerp}[2]{\left\langle #1 \vert #2 \right\rangle}
\renewcommand{\vec}[1]{\boldsymbol{#1}}     % NB: overrides the LaTeX arrow accent


% -----------------------------------------------------------------------------
% 9. Key-rate numerics
% -----------------------------------------------------------------------------
\newcommand{\cG}{\mathcal{G}}               % protocol map
\newcommand{\cZ}{\mathcal{Z}}               % pinching map
\newcommand{\Amap}{\mathcal{A}_{\mathrm{ann}}} % announcement + measurement map
\newcommand{\Amapnb}{\mathcal{A}_{\mathrm{ann}|\mathrm{nb}}}
\newcommand{\simpl}{}                       % empty suffix for simplified block maps
\newcommand{\acSet}{\mathcal{S}_{\mathrm{acc}}} % acceptance set
\newcommand{\testE}{\mathcal{E}_{\mathrm{test}}} % test-statistics map
\newcommand{\ecLeak}{\mathrm{leak}^{\mathrm{EC}}_{\mathrm{obs}}}
\newcommand{\pgen}{p_{\mathrm{gen}}}
\newcommand{\askeyr}{R_{\infty}}             % asymptotic key rate
\newcommand{\ppass}{p_{\mathrm{pass}}}       % sifting pass probability


% -----------------------------------------------------------------------------
% 10. Finite-size analysis / Renyi entropies
% -----------------------------------------------------------------------------

% --- Round types: test vs. key generation (register T, alphabet \Fset) --------
\newcommand{\tst}{\mathrm{test}}            % test round
\newcommand{\gen}{\mathrm{gen}}             % key-generation round
\newcommand{\Fset}{\mathcal{T}}             % round-type alphabet {test, gen} (not \Tset)
\newcommand{\Chat}{\widehat{C}}             % acceptance-test register
\newcommand{\Chatset}{\widehat{\mathcal{C}}}% its alphabet

% --- Basis / key extraction functions ----------------------------------------
\newcommand{\bA}{b_{A}}                     % Alice's basis-extraction function
\newcommand{\bB}{b_{B}}                     % Bob's basis-extraction function
\newcommand{\kB}{k_{B}}                     % Bob's key-extraction function

% --- Maps --------------------------------------------------------------------
\newcommand{\Mprot}{\mathcal{M}}            % single-round protocol map
\newcommand{\Phimap}{\boldsymbol{\Phi}}     % test-statistics map
\newcommand{\Pdeph}{\mathcal{P}}            % dephasing (pinching) map on \Blreg

% --- Isometries --------------------------------------------------------------
\newcommand{\Vmsm}{V_{\mathrm{msm}}}        % measurement isometry
\newcommand{\Vann}{V_{\mathrm{ann}}}        % announcement isometry
\newcommand{\VI}{V_{I}}                     % Eve's copy of the announcements
\newcommand{\VChat}{V_{\Chat}}              % acceptance-data isometry
\newcommand{\VR}{V_{R}}                     % sifting + key map -> raw key R
\newcommand{\VZ}{V_{Z}}                     % raw key -> key register Z
\newcommand{\Wprot}{W}                      % full protocol isometry

% --- Renyi entropies ---------------------------------------------------------
% SECURITY-CRITICAL: see spec/notation.md §4. Never change an arrow to
% make an inequality work out; never convert silently between the sandwiched
% (tilde) and Petz (overline) families.
\newcommand{\Hsandup}{\widetilde{H}^{\uparrow}_{\alpha}}   % sandwiched, up
\newcommand{\Hsanddown}{\widetilde{H}^{\downarrow}_{\alpha}}% sandwiched, down
\newcommand{\Hpetzup}{\overline{H}^{\uparrow}_{\alpha}}    % Petz, up
\newcommand{\Hpetzdown}{\overline{H}^{\downarrow}_{\alpha}}% Petz, down
\newcommand{\hup}{h^{\uparrow}_{\alpha}}                   % single-round rate
\newcommand{\hqkd}{h^{\uparrow,\mathrm{QKD}}_{\alpha}}     % single-round QKD rate
\newcommand{\hlow}{\underline{h}_{\alpha}}                 % lower bound on \hqkd
% --- spec/notation.md 3.5 and 4.5 -------------------------------------------
% \hlb  is the first-order linearization lower bound on \hlow: the headline
%       key rate (the `first-order` tier of foundations/def:tiers). NOT \hlow
%       itself.
% \tnum is the instance's one declared numerical tolerance. Its
%       near-collision with the declared marginal \tau_A is disclosed in
%       spec/notation.md 3.5; bare \tau is reserved and denotes neither.
\newcommand{\hlb}{\underline{h}^{\mathrm{FO}}_{\alpha}}    % first-order lower bound on \hlow
\newcommand{\tnum}{\tau_{\mathrm{num}}}                    % declared numerical tolerance


% =============================================================================
%  PART B -- module-local macros, grouped by module. Deprecated aliases,
%  defined so that a text written with them compiles, used in no new contract: \Lop,
%  \testE, \Pdeph, \pass, \ecLeak (spec/notation.md section 3.9).
% =============================================================================

% --- meat-core -----------------------------------------------------------------
\newcommand{\SecSet}{\mathbb{S}}            % sector index set of Q
\newcommand{\Qsys}{Q}                       % declared sector-carrying measured system
\newcommand{\Vset}{\mathcal{V}}             % declared joint outcome alphabet of the Q-POVM
\newcommand{\CAset}{\mathcal{C}_{A}}        % Alice announcement alphabet
\newcommand{\CBset}{\mathcal{C}_{B}}        % Bob announcement alphabet
\newcommand{\Party}{\mathfrak{P}}           % declared party partition
\newcommand{\Oset}[1]{\mathcal{O}_{#1}}     % readout alphabet of party #1
\newcommand{\rd}[1]{\pi_{#1}}               % readout map of party #1
\newcommand{\Sadm}{\mathcal{S}_{A}}         % admissible Alice subspace
\newcommand{\nomarg}{\mathsf{none}}         % "no trusted marginal declared"
\newcommand{\Zset}{\mathcal{Z}_{\perp}}     % raw-key alphabet incl. \perp
\newcommand{\Zval}{\mathcal{Z}}             % key alphabet without \perp
\newcommand{\Etil}{\widetilde{E}}           % single-round purifying register
\newcommand{\Ehat}{\widehat{E}}             % Eve's n-round register
\newcommand{\OmAT}{\Omega_{\mathrm{AT}}}    % acceptance event
\newcommand{\epa}{\varepsilon_{\mathrm{PA}}}
\newcommand{\eev}{\varepsilon_{\mathrm{EV}}}
\newcommand{\lEC}{\lambda_{\mathrm{EC}}}
\newcommand{\Decl}{\mathfrak{D}}            % protocol declaration
\newcommand{\CList}{\mathfrak{L}}           % declared constraint list
\newcommand{\PinS}{\mathcal{P}_{\SecSet}}   % sector pinching on Q
\newcommand{\Feas}{\mathcal{W}}             % feasible set of block states
\newcommand{\Chkey}{\Chatset_{\mathrm{key}}}% key-live acceptance values
\newcommand{\Kann}{\mathcal{K}}             % announcement-class alphabet
\newcommand{\hold}{\mathsf{h}}              % declared key holder (A or B)
\newcommand{\sfD}{D}                        % classical KL divergence (base 2)
\newcommand{\Dtwo}{D}                       % quantum relative entropy (base 2)
\newcommand{\Dsand}{\widetilde{D}_{\alpha}} % sandwiched Renyi divergence
\newcommand{\idmap}{\mathrm{id}}            % identity channel
\newcommand{\Aev}{\mathcal{A}}              % Eve's per-round attack channel
\newcommand{\AttSet}{\mathfrak{A}}          % licensed attack-channel set
\newcommand{\Chacc}{\Chatset_{\mathrm{acc}}}% acceptance support of \acSet
\newcommand{\epsstar}{\varepsilon^{\star}}      % distance to an exact POVM
\newcommand{\dfH}{\Delta_{\mathrm{herm}}}       % Hermiticity defect
\newcommand{\dfP}{\Delta_{\mathrm{pos}}}        % negative part
\newcommand{\dfC}{\Delta_{\mathrm{cpl}}}        % completeness defect
\newcommand{\lmoq}{g}                       % acceptance-side gradient block
\newcommand{\lmoJ}{G}                       % feasible-side gradient block
\newcommand{\bacc}{\beta_{\mathrm{acc}}}    % certified acceptance-side minimum
\newcommand{\Itset}{\mathfrak{I}}           % finite index set of iterates (C22)

% --- reductions ----------------------------------------------------------------
\newcommand{\Bsys}{B}                        % physical (pre-squashing) system
\newcommand{\Tcl}[1]{\mathcal{T}(#1)}        % trace-class operators on H(#1)
\newcommand{\Ins}{\mathcal{I}}               % generalized-instrument branch map
\newcommand{\redNmsm}{\mathcal{N}}              % POVM measurement channel
\newcommand{\Eprot}{\mathcal{E}}             % n-round protocol channel
\newcommand{\Grnd}{\mathsf{G}}               % one round map of a protocol
\newcommand{\Kfin}{\mathsf{K}}               % end-of-protocol processing
\newcommand{\Crit}{\mathsf{sec}}             % value of a security criterion
\newcommand{\Pflag}{\Pi_{\Flag}}             % flag-sector projector on Q
\newcommand{\Pleq}{\Pi_{\leqN}}              % preserved-block projector
\newcommand{\Pgt}{\Pi_{\gtN}}                % above-cutoff projector
\newcommand{\Sph}{S^{\mathrm{ph}}}           % family of physical inputs
\newcommand{\Srch}{S^{\mathrm{rch}}}         % reachable squashed states
\newcommand{\Ssq}{S^{\sq}}                   % squashed admissible set
\newcommand{\Ffss}{F^{\mathrm{FSS}}}         % flag-state-squasher target
\newcommand{\Fwp}{F^{\mathrm{WP}}}           % multi-bound WPFSS target
\newcommand{\Mwp}{M^{\mathrm{WP}}}           % renormalized above-cutoff POVM
\newcommand{\Mdec}{{M'}^{\mathrm{ph}}}       % declared physical family on B
\newcommand{\Ppvm}{P^{\Flag}}                % flag PVM
\newcommand{\Lmp}{{\Lambda_{\mathrm{mp}}}}   % measure-and-prepare channel
\newcommand{\Lwp}{{\Lambda^{\mathrm{WP}}}}   % multi-bound WPFSS channel
\newcommand{\Lfss}{{\Lambda^{\mathrm{FSS}}}} % flag-state-squasher channel
\newcommand{\lmax}{\lambda_{\max}}           % maximum eigenvalue
\newcommand{\Lnoise}{\widetilde{\Lambda}}    % flag-noise channel
\newcommand{\cg}{\varpi}                     % declared coarse-graining map
\newcommand{\MCK}{\mathsf{MC}}               % multiclick raw-pattern set
\newcommand{\redChain}{\mathfrak{C}}            % reduction chain
\newcommand{\Lnk}{\ell}                      % a single typed link
\newcommand{\dsq}{\preceq}                      % the squashing relation
\newcommand{\opn}[1]{\left\lVert #1\right\rVert}           % operator (spectral) norm
\newcommand{\mxn}[1]{\left\lVert #1\right\rVert_{\max}}    % entrywise max modulus
\newcommand{\Fld}{\mathbb{F}}                   % exact ordered substrate field
\newcommand{\Fldi}{\mathbb{F}_{i}}              % its Gaussian extension F + iF
\newcommand{\lat}[1]{\big\lfloor #1\big\rfloor_{\kappa}} % round DOWN onto the lattice

% --- protocol-compiler ---------------------------------------------------------
\newcommand{\Spec}{\mathfrak{S}}             % protocol specification (compiler input)
\newcommand{\pcChain}{\mathfrak{R}}            % reduction chain (two typed segments)
\newcommand{\Chsrc}{\pcChain_{\mathrm{src}}}   % its Alice-side source-map segment
\newcommand{\Chrec}{\pcChain_{\mathrm{rec}}}   % its receiver-side squashing segment
\newcommand{\rvirt}{\rho^{\mathrm{v}}}       % virtual (PRE-source-map) signal state
\newcommand{\Fph}{F^{\mathrm{ph}}}           % physical (pre-reduction) POVM element
\newcommand{\Vph}{\mathcal{V}^{\mathrm{ph}}} % physical outcome alphabet
\newcommand{\Tpub}{\mathcal{T}_{\mathrm{pub}}}% public-choice alphabet
\newcommand{\Uset}{\mathcal{U}}              % residual outcome alphabet
\newcommand{\idxm}{\mathsf{idx}}             % label-to-index map
\newcommand{\Asig}{A'}                       % transmitted (signal) system -- the REAL one
\newcommand{\Avirt}{A''}                     % virtual transmitted system (source descriptor of Xi_1)
\newcommand{\Abar}{\bar{A}}                  % Alice's measured register
\newcommand{\Ahat}{\hat{A}}                  % Alice's shield register
\newcommand{\pcNmsm}{\mathcal{N}_{\mathsf A}}  % Alice's measurement channel (see the RULES block)
\newcommand{\Phip}{\Phi^{+}}                 % maximally entangled vector label
\newcommand{\Smap}{\mathfrak{s}}             % the emitted sector map (fraktur: a compile-time object)

% --- receiver ------------------------------------------------------------------
\newcommand{\Modes}{\mathbb{M}}              % declared ordered mode set
\newcommand{\Inp}{\mathbb{I}}                % declared input modes
\newcommand{\InpC}{\mathbb{I}_{\Chan}}       % channel input modes
\newcommand{\InpS}{\mathbb{I}_{\Prb}}        % probe (source) input modes
\newcommand{\Env}{\mathbb{E}}                % dilation (environment) modes
\newcommand{\Port}[1]{\mathbb{P}_{#1}}       % mode set of detector port #1
\newcommand{\Emd}[1]{E_{#1}}                 % coordinate projector on modes #1
\newcommand{\Gmat}{G}                        % instrument (transfer) matrix
\newcommand{\Fock}{\mathfrak{F}}             % Fock / second-quantization functor
\newcommand{\Symn}[1]{\mathrm{Sym}^{#1}}     % symmetric tensor power
\newcommand{\normord}[1]{{:}\mspace{2mu}#1\mspace{2mu}{:}}   % normal ordering
\newcommand{\rk}{\operatorname{rank}}        % rank (physics.sty owns \rank)
\newcommand{\SCK}{\mathsf{SC}}               % single-click pattern set
\newcommand{\Recv}{\mathfrak{R}}             % receiver instrument tuple
\newcommand{\Cat}{\mathfrak{C}}              % component catalogue
\newcommand{\Grph}{\mathfrak{E}}             % declared element graph
\newcommand{\defic}{\mathbb{L}}              % loss deficiency operator
\newcommand{\Pvac}[1]{\Pi_{\vac}(#1)}        % joint no-click element on ports #1
\newcommand{\pmc}{p_{\mc}}                   % uniform high-photon mc weight
\newcommand{\Bnd}{\beta}                     % the scalar bound function beta(n)
\newcommand{\Rot}{R}                         % planar rotation matrix
\newcommand{\Wplate}{W}                      % retarder (wave-plate) matrix
\newcommand{\Bsp}{B}                         % beamsplitter core matrix

% --- statistics ----------------------------------------------------------------
\newcommand{\honest}{p^{\mathrm{hon}}}          % honest per-round distribution
\newcommand{\ecomp}{\varepsilon_{\mathrm{comp}}}% completeness budget
\newcommand{\simplex}[1]{\Delta(#1)}            % probability simplex
\newcommand{\Chatcon}{\Chatset_{\mathrm{con}}}  % constrained outcomes
\newcommand{\Ireg}[3]{\mathcal{I}_{#1}(#2,#3)}  % regularized incomplete beta
\newcommand{\Ncount}[1]{\mathsf{N}_{#1}}        % per-outcome count
\newcommand{\qlo}{\mathbf{q}^{-}}               % box lower bounds
\newcommand{\qhi}{\mathbf{q}^{+}}               % box upper bounds
 for it, so the
%  file is read from spec/modules/<id>/.
%
%  RULES
%    * A macro that spec/notation.md catalogues is a verbatim copy of it
%      (part A). Changing one is a notation change: bump spec/notation.md.
%    * Part B holds the module-local macros. A macro defined by two modules
%      with the same expansion appears once, under the first module (in the
%      order below) that defines it. Two macros with the same name and
%      DIFFERENT expansions are both kept under a module prefix, and the
%      claims write the prefixed names:
%        \Chain  -> \redChain (reductions, \mathfrak{C}: the reduction chain)
%                   \pcChain  (protocol-compiler, \mathfrak{R}: the two-segment
%                             chain; \Chsrc and \Chrec are built on it)
%        \Nmsm   -> \redNmsm  (reductions, \mathcal{N}: a POVM measurement
%                             channel)
%                   \pcNmsm   (\mathcal{N}_{\mathsf A}: Alice's measurement
%                             channel, written by
%                             transmitter/thm:certificate-sufficiency,
%                             transmitter/rem:xreg and
%                             reductions/lem:source-transfer)
%    * Deliberately omitted: a \provenance{tag}{details} macro. Provenance
%      is a header field of every claim file (spec/claim-format.md section
%      3); a claim body that calls such a macro fails to compile, which is
%      the intended signal.
%
%  REQUIRED PACKAGES, loaded by the assembler's preamble before this file:
%    amsmath, amssymb, amsfonts, mathtools, dsfont, braket, physics, xcolor
%    (braket before physics, so that physics wins for \ket, \bra, \tr, ...).
%
%  This file contains only definitions: no \usepackage, no document body.
% =============================================================================


% =============================================================================
%  PART A -- notation macros, the ones spec/notation.md catalogues, each a
%  verbatim copy of its catalogue row; grouped under numbered headings.
% =============================================================================

% -----------------------------------------------------------------------------
% 1. State spaces and operator sets (a "system label" -> spaces of operators)
% -----------------------------------------------------------------------------
\newcommand{\Hs}[1]{\mathcal{H}(#1)}        % Hilbert space of a system label
\newcommand{\Lin}[1]{\mathcal{L}(#1)}       % bounded linear operators on H(#1)
\newcommand{\Linn}[2]{\mathcal{L}(#1,#2)}   % linear maps H(#1) -> H(#2)
\newcommand{\Pos}[1]{\mathrm{Pos}(#1)}      % positive semidefinite operators on H(#1)
\newcommand{\Density}{\mathrm{D}}           % density operators: \Density(A) = D(A)
\newcommand{\Seq}[1]{S_{=}(#1)}             % normalized states (trace = 1)
\newcommand{\Sleq}[1]{S_{\leq}(#1)}         % subnormalized states (0 <= trace <= 1)
\newcommand{\id}{\mathds{1}}                % identity operator
\newcommand{\idon}[1]{\mathds{1}_{#1}}      % identity on a named system, e.g. \idon{C}

% DEPRECATED: applies \mathcal to both arguments itself, so it takes bare
% letters, not system labels. Prefer \Linn; do not use in new contracts.
\providecommand{\Lop}[2]{\mathcal{L}(\mathcal{#1}, \mathcal{#2})}


% -----------------------------------------------------------------------------
% 2. Composite systems
% -----------------------------------------------------------------------------
\newcommand{\dsum}{\oplus}                  % direct sum of system labels
% Tensor product: use \otimes directly.
% n-fold direct sum of an operator space: \Lin{A}^{\oplus n}


% -----------------------------------------------------------------------------
% 3. Events, conditioning and classical registers (for STATES)
%    Convention:  \wedge  -> subnormalized "and-event" component, e.g. \rho_{A\wedge x}
%                 |        -> normalized conditional state,        e.g. \rho_{A|x}
% (Write these directly in subscripts; no extra macros needed.)
% -----------------------------------------------------------------------------


% -----------------------------------------------------------------------------
% 4. Photon-number subspaces / direct-sum blocks
%    Convention:  operators carry the restriction "on top", in brackets;
%                 states carry it as a \wedge subscript (see Sec. 3).
% -----------------------------------------------------------------------------
\newcommand{\Pnum}[1]{P_{#1}}               % photon-number projector (system from context)
\newcommand{\Pnumon}[2]{P_{#1}^{#2}}        % with explicit system superscript, e.g. P_n^C
\newcommand{\blk}[2]{#1^{(#2)}}             % operator block, e.g. \blk{M}{\leqN} = M^{(<=N)}
\newcommand{\leqN}{\leq N}                  % use as \blk{M}{\leqN}
\newcommand{\gtN}{> N}                      % use as \blk{M}{\gtN}
\newcommand{\eqn}{= n}                      % use as \blk{M}{\eqn}


% -----------------------------------------------------------------------------
% 5. System labels used in the protocol
% -----------------------------------------------------------------------------
\newcommand{\Chan}{C}                       % quantum-channel (Eve) input to Bob
\newcommand{\Prb}{P}                        % probe input
\newcommand{\Flag}{{\mathrm{Flg}}}          % flag system (EXISTS ONLY AFTER squashing)
\newcommand{\Loss}{{\mathrm{Lss}}}          % loss system
\newcommand{\Vc}{{\mathrm{Vcm}}}            % vacuum system
\newcommand{\Nout}{N_{\mathrm{out}}}
\newcommand{\Nin}{N_{\mathrm{in}}}


% -----------------------------------------------------------------------------
% 6. Alphabets
% -----------------------------------------------------------------------------
\newcommand{\Yset}{\mathcal{Y}}             % Bob's (coarse-grained) detection alphabet
\newcommand{\Yraw}{\mathcal{Y}_{\mathrm{raw}}}% raw click patterns {0,1}^4
\newcommand{\Yscset}{\mathcal{Y}_{\singlec}}% single-click outcomes
\newcommand{\Ymcset}{\mathcal{Y}_{\mc}}     % multi-click outcomes
\newcommand{\Tset}{\mathcal{T}_{\mathrm{set}}}% private setting set (the probe settings in the optical instance)
\newcommand{\Xset}{\mathcal{X}}             % Alice's alphabet
\newcommand{\Dset}{\{H,V,+,-\}}             % polarization / detector labels
\newcommand{\pass}{\top}                    % pass sifting


% -----------------------------------------------------------------------------
% 7. Mode / outcome abbreviations
% -----------------------------------------------------------------------------
\newcommand{\on}{{\mathrm{on}}}
\newcommand{\off}{{\mathrm{off}}}
\newcommand{\dc}{{\mathrm{dc}}}             % dark count
\newcommand{\sq}{\mathrm{sq}}               % squashed
\newcommand{\mc}{\mathrm{mc}}               % multi-click
\newcommand{\nc}{\mathrm{nc}}               % no-click
\newcommand{\singlec}{\mathrm{sc}}          % single-click
\newcommand{\vac}{\mathrm{vac}}             % vacuum
\newcommand{\kout}{{k_\mathrm{out}}}
\newcommand{\kin}{{k_\mathrm{in}}}


% -----------------------------------------------------------------------------
% 8. Maps, channels and miscellaneous operators
% -----------------------------------------------------------------------------
\newcommand{\Prob}{\mathrm{Prob}}           % probability (use \Prob[\cdot] or \Prob(\cdot))
\newcommand{\lmin}{\lambda_{\min}}          % minimum eigenvalue
\newcommand{\Wflag}{W_{\Flag}}              % flag-space weight
\newcommand{\transpose}{\mathsf{T}}
\newcommand{\innerp}[2]{\left\langle #1 \vert #2 \right\rangle}
\renewcommand{\vec}[1]{\boldsymbol{#1}}     % NB: overrides the LaTeX arrow accent


% -----------------------------------------------------------------------------
% 9. Key-rate numerics
% -----------------------------------------------------------------------------
\newcommand{\cG}{\mathcal{G}}               % protocol map
\newcommand{\cZ}{\mathcal{Z}}               % pinching map
\newcommand{\Amap}{\mathcal{A}_{\mathrm{ann}}} % announcement + measurement map
\newcommand{\Amapnb}{\mathcal{A}_{\mathrm{ann}|\mathrm{nb}}}
\newcommand{\simpl}{}                       % empty suffix for simplified block maps
\newcommand{\acSet}{\mathcal{S}_{\mathrm{acc}}} % acceptance set
\newcommand{\testE}{\mathcal{E}_{\mathrm{test}}} % test-statistics map
\newcommand{\ecLeak}{\mathrm{leak}^{\mathrm{EC}}_{\mathrm{obs}}}
\newcommand{\pgen}{p_{\mathrm{gen}}}
\newcommand{\askeyr}{R_{\infty}}             % asymptotic key rate
\newcommand{\ppass}{p_{\mathrm{pass}}}       % sifting pass probability


% -----------------------------------------------------------------------------
% 10. Finite-size analysis / Renyi entropies
% -----------------------------------------------------------------------------

% --- Round types: test vs. key generation (register T, alphabet \Fset) --------
\newcommand{\tst}{\mathrm{test}}            % test round
\newcommand{\gen}{\mathrm{gen}}             % key-generation round
\newcommand{\Fset}{\mathcal{T}}             % round-type alphabet {test, gen} (not \Tset)
\newcommand{\Chat}{\widehat{C}}             % acceptance-test register
\newcommand{\Chatset}{\widehat{\mathcal{C}}}% its alphabet

% --- Basis / key extraction functions ----------------------------------------
\newcommand{\bA}{b_{A}}                     % Alice's basis-extraction function
\newcommand{\bB}{b_{B}}                     % Bob's basis-extraction function
\newcommand{\kB}{k_{B}}                     % Bob's key-extraction function

% --- Maps --------------------------------------------------------------------
\newcommand{\Mprot}{\mathcal{M}}            % single-round protocol map
\newcommand{\Phimap}{\boldsymbol{\Phi}}     % test-statistics map
\newcommand{\Pdeph}{\mathcal{P}}            % dephasing (pinching) map on \Blreg

% --- Isometries --------------------------------------------------------------
\newcommand{\Vmsm}{V_{\mathrm{msm}}}        % measurement isometry
\newcommand{\Vann}{V_{\mathrm{ann}}}        % announcement isometry
\newcommand{\VI}{V_{I}}                     % Eve's copy of the announcements
\newcommand{\VChat}{V_{\Chat}}              % acceptance-data isometry
\newcommand{\VR}{V_{R}}                     % sifting + key map -> raw key R
\newcommand{\VZ}{V_{Z}}                     % raw key -> key register Z
\newcommand{\Wprot}{W}                      % full protocol isometry

% --- Renyi entropies ---------------------------------------------------------
% SECURITY-CRITICAL: see spec/notation.md §4. Never change an arrow to
% make an inequality work out; never convert silently between the sandwiched
% (tilde) and Petz (overline) families.
\newcommand{\Hsandup}{\widetilde{H}^{\uparrow}_{\alpha}}   % sandwiched, up
\newcommand{\Hsanddown}{\widetilde{H}^{\downarrow}_{\alpha}}% sandwiched, down
\newcommand{\Hpetzup}{\overline{H}^{\uparrow}_{\alpha}}    % Petz, up
\newcommand{\Hpetzdown}{\overline{H}^{\downarrow}_{\alpha}}% Petz, down
\newcommand{\hup}{h^{\uparrow}_{\alpha}}                   % single-round rate
\newcommand{\hqkd}{h^{\uparrow,\mathrm{QKD}}_{\alpha}}     % single-round QKD rate
\newcommand{\hlow}{\underline{h}_{\alpha}}                 % lower bound on \hqkd
% --- spec/notation.md 3.5 and 4.5 -------------------------------------------
% \hlb  is the first-order linearization lower bound on \hlow: the headline
%       key rate (the `first-order` tier of foundations/def:tiers). NOT \hlow
%       itself.
% \tnum is the instance's one declared numerical tolerance. Its
%       near-collision with the declared marginal \tau_A is disclosed in
%       spec/notation.md 3.5; bare \tau is reserved and denotes neither.
\newcommand{\hlb}{\underline{h}^{\mathrm{FO}}_{\alpha}}    % first-order lower bound on \hlow
\newcommand{\tnum}{\tau_{\mathrm{num}}}                    % declared numerical tolerance


% =============================================================================
%  PART B -- module-local macros, grouped by module. Deprecated aliases,
%  defined so that a text written with them compiles, used in no new contract: \Lop,
%  \testE, \Pdeph, \pass, \ecLeak (spec/notation.md section 3.9).
% =============================================================================

% --- meat-core -----------------------------------------------------------------
\newcommand{\SecSet}{\mathbb{S}}            % sector index set of Q
\newcommand{\Qsys}{Q}                       % declared sector-carrying measured system
\newcommand{\Vset}{\mathcal{V}}             % declared joint outcome alphabet of the Q-POVM
\newcommand{\CAset}{\mathcal{C}_{A}}        % Alice announcement alphabet
\newcommand{\CBset}{\mathcal{C}_{B}}        % Bob announcement alphabet
\newcommand{\Party}{\mathfrak{P}}           % declared party partition
\newcommand{\Oset}[1]{\mathcal{O}_{#1}}     % readout alphabet of party #1
\newcommand{\rd}[1]{\pi_{#1}}               % readout map of party #1
\newcommand{\Sadm}{\mathcal{S}_{A}}         % admissible Alice subspace
\newcommand{\nomarg}{\mathsf{none}}         % "no trusted marginal declared"
\newcommand{\Zset}{\mathcal{Z}_{\perp}}     % raw-key alphabet incl. \perp
\newcommand{\Zval}{\mathcal{Z}}             % key alphabet without \perp
\newcommand{\Etil}{\widetilde{E}}           % single-round purifying register
\newcommand{\Ehat}{\widehat{E}}             % Eve's n-round register
\newcommand{\OmAT}{\Omega_{\mathrm{AT}}}    % acceptance event
\newcommand{\epa}{\varepsilon_{\mathrm{PA}}}
\newcommand{\eev}{\varepsilon_{\mathrm{EV}}}
\newcommand{\lEC}{\lambda_{\mathrm{EC}}}
\newcommand{\Decl}{\mathfrak{D}}            % protocol declaration
\newcommand{\CList}{\mathfrak{L}}           % declared constraint list
\newcommand{\PinS}{\mathcal{P}_{\SecSet}}   % sector pinching on Q
\newcommand{\Feas}{\mathcal{W}}             % feasible set of block states
\newcommand{\Chkey}{\Chatset_{\mathrm{key}}}% key-live acceptance values
\newcommand{\Kann}{\mathcal{K}}             % announcement-class alphabet
\newcommand{\hold}{\mathsf{h}}              % declared key holder (A or B)
\newcommand{\sfD}{D}                        % classical KL divergence (base 2)
\newcommand{\Dtwo}{D}                       % quantum relative entropy (base 2)
\newcommand{\Dsand}{\widetilde{D}_{\alpha}} % sandwiched Renyi divergence
\newcommand{\idmap}{\mathrm{id}}            % identity channel
\newcommand{\Aev}{\mathcal{A}}              % Eve's per-round attack channel
\newcommand{\AttSet}{\mathfrak{A}}          % licensed attack-channel set
\newcommand{\Chacc}{\Chatset_{\mathrm{acc}}}% acceptance support of \acSet
\newcommand{\epsstar}{\varepsilon^{\star}}      % distance to an exact POVM
\newcommand{\dfH}{\Delta_{\mathrm{herm}}}       % Hermiticity defect
\newcommand{\dfP}{\Delta_{\mathrm{pos}}}        % negative part
\newcommand{\dfC}{\Delta_{\mathrm{cpl}}}        % completeness defect
\newcommand{\lmoq}{g}                       % acceptance-side gradient block
\newcommand{\lmoJ}{G}                       % feasible-side gradient block
\newcommand{\bacc}{\beta_{\mathrm{acc}}}    % certified acceptance-side minimum
\newcommand{\Itset}{\mathfrak{I}}           % finite index set of iterates (C22)

% --- reductions ----------------------------------------------------------------
\newcommand{\Bsys}{B}                        % physical (pre-squashing) system
\newcommand{\Tcl}[1]{\mathcal{T}(#1)}        % trace-class operators on H(#1)
\newcommand{\Ins}{\mathcal{I}}               % generalized-instrument branch map
\newcommand{\redNmsm}{\mathcal{N}}              % POVM measurement channel
\newcommand{\Eprot}{\mathcal{E}}             % n-round protocol channel
\newcommand{\Grnd}{\mathsf{G}}               % one round map of a protocol
\newcommand{\Kfin}{\mathsf{K}}               % end-of-protocol processing
\newcommand{\Crit}{\mathsf{sec}}             % value of a security criterion
\newcommand{\Pflag}{\Pi_{\Flag}}             % flag-sector projector on Q
\newcommand{\Pleq}{\Pi_{\leqN}}              % preserved-block projector
\newcommand{\Pgt}{\Pi_{\gtN}}                % above-cutoff projector
\newcommand{\Sph}{S^{\mathrm{ph}}}           % family of physical inputs
\newcommand{\Srch}{S^{\mathrm{rch}}}         % reachable squashed states
\newcommand{\Ssq}{S^{\sq}}                   % squashed admissible set
\newcommand{\Ffss}{F^{\mathrm{FSS}}}         % flag-state-squasher target
\newcommand{\Fwp}{F^{\mathrm{WP}}}           % multi-bound WPFSS target
\newcommand{\Mwp}{M^{\mathrm{WP}}}           % renormalized above-cutoff POVM
\newcommand{\Mdec}{{M'}^{\mathrm{ph}}}       % declared physical family on B
\newcommand{\Ppvm}{P^{\Flag}}                % flag PVM
\newcommand{\Lmp}{{\Lambda_{\mathrm{mp}}}}   % measure-and-prepare channel
\newcommand{\Lwp}{{\Lambda^{\mathrm{WP}}}}   % multi-bound WPFSS channel
\newcommand{\Lfss}{{\Lambda^{\mathrm{FSS}}}} % flag-state-squasher channel
\newcommand{\lmax}{\lambda_{\max}}           % maximum eigenvalue
\newcommand{\Lnoise}{\widetilde{\Lambda}}    % flag-noise channel
\newcommand{\cg}{\varpi}                     % declared coarse-graining map
\newcommand{\MCK}{\mathsf{MC}}               % multiclick raw-pattern set
\newcommand{\redChain}{\mathfrak{C}}            % reduction chain
\newcommand{\Lnk}{\ell}                      % a single typed link
\newcommand{\dsq}{\preceq}                      % the squashing relation
\newcommand{\opn}[1]{\left\lVert #1\right\rVert}           % operator (spectral) norm
\newcommand{\mxn}[1]{\left\lVert #1\right\rVert_{\max}}    % entrywise max modulus
\newcommand{\Fld}{\mathbb{F}}                   % exact ordered substrate field
\newcommand{\Fldi}{\mathbb{F}_{i}}              % its Gaussian extension F + iF
\newcommand{\lat}[1]{\big\lfloor #1\big\rfloor_{\kappa}} % round DOWN onto the lattice

% --- protocol-compiler ---------------------------------------------------------
\newcommand{\Spec}{\mathfrak{S}}             % protocol specification (compiler input)
\newcommand{\pcChain}{\mathfrak{R}}            % reduction chain (two typed segments)
\newcommand{\Chsrc}{\pcChain_{\mathrm{src}}}   % its Alice-side source-map segment
\newcommand{\Chrec}{\pcChain_{\mathrm{rec}}}   % its receiver-side squashing segment
\newcommand{\rvirt}{\rho^{\mathrm{v}}}       % virtual (PRE-source-map) signal state
\newcommand{\Fph}{F^{\mathrm{ph}}}           % physical (pre-reduction) POVM element
\newcommand{\Vph}{\mathcal{V}^{\mathrm{ph}}} % physical outcome alphabet
\newcommand{\Tpub}{\mathcal{T}_{\mathrm{pub}}}% public-choice alphabet
\newcommand{\Uset}{\mathcal{U}}              % residual outcome alphabet
\newcommand{\idxm}{\mathsf{idx}}             % label-to-index map
\newcommand{\Asig}{A'}                       % transmitted (signal) system -- the REAL one
\newcommand{\Avirt}{A''}                     % virtual transmitted system (source descriptor of Xi_1)
\newcommand{\Abar}{\bar{A}}                  % Alice's measured register
\newcommand{\Ahat}{\hat{A}}                  % Alice's shield register
\newcommand{\pcNmsm}{\mathcal{N}_{\mathsf A}}  % Alice's measurement channel (see the RULES block)
\newcommand{\Phip}{\Phi^{+}}                 % maximally entangled vector label
\newcommand{\Smap}{\mathfrak{s}}             % the emitted sector map (fraktur: a compile-time object)

% --- receiver ------------------------------------------------------------------
\newcommand{\Modes}{\mathbb{M}}              % declared ordered mode set
\newcommand{\Inp}{\mathbb{I}}                % declared input modes
\newcommand{\InpC}{\mathbb{I}_{\Chan}}       % channel input modes
\newcommand{\InpS}{\mathbb{I}_{\Prb}}        % probe (source) input modes
\newcommand{\Env}{\mathbb{E}}                % dilation (environment) modes
\newcommand{\Port}[1]{\mathbb{P}_{#1}}       % mode set of detector port #1
\newcommand{\Emd}[1]{E_{#1}}                 % coordinate projector on modes #1
\newcommand{\Gmat}{G}                        % instrument (transfer) matrix
\newcommand{\Fock}{\mathfrak{F}}             % Fock / second-quantization functor
\newcommand{\Symn}[1]{\mathrm{Sym}^{#1}}     % symmetric tensor power
\newcommand{\normord}[1]{{:}\mspace{2mu}#1\mspace{2mu}{:}}   % normal ordering
\newcommand{\rk}{\operatorname{rank}}        % rank (physics.sty owns \rank)
\newcommand{\SCK}{\mathsf{SC}}               % single-click pattern set
\newcommand{\Recv}{\mathfrak{R}}             % receiver instrument tuple
\newcommand{\Cat}{\mathfrak{C}}              % component catalogue
\newcommand{\Grph}{\mathfrak{E}}             % declared element graph
\newcommand{\defic}{\mathbb{L}}              % loss deficiency operator
\newcommand{\Pvac}[1]{\Pi_{\vac}(#1)}        % joint no-click element on ports #1
\newcommand{\pmc}{p_{\mc}}                   % uniform high-photon mc weight
\newcommand{\Bnd}{\beta}                     % the scalar bound function beta(n)
\newcommand{\Rot}{R}                         % planar rotation matrix
\newcommand{\Wplate}{W}                      % retarder (wave-plate) matrix
\newcommand{\Bsp}{B}                         % beamsplitter core matrix

% --- statistics ----------------------------------------------------------------
\newcommand{\honest}{p^{\mathrm{hon}}}          % honest per-round distribution
\newcommand{\ecomp}{\varepsilon_{\mathrm{comp}}}% completeness budget
\newcommand{\simplex}[1]{\Delta(#1)}            % probability simplex
\newcommand{\Chatcon}{\Chatset_{\mathrm{con}}}  % constrained outcomes
\newcommand{\Ireg}[3]{\mathcal{I}_{#1}(#2,#3)}  % regularized incomplete beta
\newcommand{\Ncount}[1]{\mathsf{N}_{#1}}        % per-outcome count
\newcommand{\qlo}{\mathbf{q}^{-}}               % box lower bounds
\newcommand{\qhi}{\mathbf{q}^{+}}               % box upper bounds
 for it, so the
%  file is read from spec/modules/<id>/.
%
%  RULES
%    * A macro that spec/notation.md catalogues is a verbatim copy of it
%      (part A). Changing one is a notation change: bump spec/notation.md.
%    * Part B holds the module-local macros. A macro defined by two modules
%      with the same expansion appears once, under the first module (in the
%      order below) that defines it. Two macros with the same name and
%      DIFFERENT expansions are both kept under a module prefix, and the
%      claims write the prefixed names:
%        \Chain  -> \redChain (reductions, \mathfrak{C}: the reduction chain)
%                   \pcChain  (protocol-compiler, \mathfrak{R}: the two-segment
%                             chain; \Chsrc and \Chrec are built on it)
%        \Nmsm   -> \redNmsm  (reductions, \mathcal{N}: a POVM measurement
%                             channel)
%                   \pcNmsm   (\mathcal{N}_{\mathsf A}: Alice's measurement
%                             channel, written by
%                             transmitter/thm:certificate-sufficiency,
%                             transmitter/rem:xreg and
%                             reductions/lem:source-transfer)
%    * Deliberately omitted: a \provenance{tag}{details} macro. Provenance
%      is a header field of every claim file (spec/claim-format.md section
%      3); a claim body that calls such a macro fails to compile, which is
%      the intended signal.
%
%  REQUIRED PACKAGES, loaded by the assembler's preamble before this file:
%    amsmath, amssymb, amsfonts, mathtools, dsfont, braket, physics, xcolor
%    (braket before physics, so that physics wins for \ket, \bra, \tr, ...).
%
%  This file contains only definitions: no \usepackage, no document body.
% =============================================================================


% =============================================================================
%  PART A -- notation macros, the ones spec/notation.md catalogues, each a
%  verbatim copy of its catalogue row; grouped under numbered headings.
% =============================================================================

% -----------------------------------------------------------------------------
% 1. State spaces and operator sets (a "system label" -> spaces of operators)
% -----------------------------------------------------------------------------
\newcommand{\Hs}[1]{\mathcal{H}(#1)}        % Hilbert space of a system label
\newcommand{\Lin}[1]{\mathcal{L}(#1)}       % bounded linear operators on H(#1)
\newcommand{\Linn}[2]{\mathcal{L}(#1,#2)}   % linear maps H(#1) -> H(#2)
\newcommand{\Pos}[1]{\mathrm{Pos}(#1)}      % positive semidefinite operators on H(#1)
\newcommand{\Density}{\mathrm{D}}           % density operators: \Density(A) = D(A)
\newcommand{\Seq}[1]{S_{=}(#1)}             % normalized states (trace = 1)
\newcommand{\Sleq}[1]{S_{\leq}(#1)}         % subnormalized states (0 <= trace <= 1)
\newcommand{\id}{\mathds{1}}                % identity operator
\newcommand{\idon}[1]{\mathds{1}_{#1}}      % identity on a named system, e.g. \idon{C}

% DEPRECATED: applies \mathcal to both arguments itself, so it takes bare
% letters, not system labels. Prefer \Linn; do not use in new contracts.
\providecommand{\Lop}[2]{\mathcal{L}(\mathcal{#1}, \mathcal{#2})}


% -----------------------------------------------------------------------------
% 2. Composite systems
% -----------------------------------------------------------------------------
\newcommand{\dsum}{\oplus}                  % direct sum of system labels
% Tensor product: use \otimes directly.
% n-fold direct sum of an operator space: \Lin{A}^{\oplus n}


% -----------------------------------------------------------------------------
% 3. Events, conditioning and classical registers (for STATES)
%    Convention:  \wedge  -> subnormalized "and-event" component, e.g. \rho_{A\wedge x}
%                 |        -> normalized conditional state,        e.g. \rho_{A|x}
% (Write these directly in subscripts; no extra macros needed.)
% -----------------------------------------------------------------------------


% -----------------------------------------------------------------------------
% 4. Photon-number subspaces / direct-sum blocks
%    Convention:  operators carry the restriction "on top", in brackets;
%                 states carry it as a \wedge subscript (see Sec. 3).
% -----------------------------------------------------------------------------
\newcommand{\Pnum}[1]{P_{#1}}               % photon-number projector (system from context)
\newcommand{\Pnumon}[2]{P_{#1}^{#2}}        % with explicit system superscript, e.g. P_n^C
\newcommand{\blk}[2]{#1^{(#2)}}             % operator block, e.g. \blk{M}{\leqN} = M^{(<=N)}
\newcommand{\leqN}{\leq N}                  % use as \blk{M}{\leqN}
\newcommand{\gtN}{> N}                      % use as \blk{M}{\gtN}
\newcommand{\eqn}{= n}                      % use as \blk{M}{\eqn}


% -----------------------------------------------------------------------------
% 5. System labels used in the protocol
% -----------------------------------------------------------------------------
\newcommand{\Chan}{C}                       % quantum-channel (Eve) input to Bob
\newcommand{\Prb}{P}                        % probe input
\newcommand{\Flag}{{\mathrm{Flg}}}          % flag system (EXISTS ONLY AFTER squashing)
\newcommand{\Loss}{{\mathrm{Lss}}}          % loss system
\newcommand{\Vc}{{\mathrm{Vcm}}}            % vacuum system
\newcommand{\Nout}{N_{\mathrm{out}}}
\newcommand{\Nin}{N_{\mathrm{in}}}


% -----------------------------------------------------------------------------
% 6. Alphabets
% -----------------------------------------------------------------------------
\newcommand{\Yset}{\mathcal{Y}}             % Bob's (coarse-grained) detection alphabet
\newcommand{\Yraw}{\mathcal{Y}_{\mathrm{raw}}}% raw click patterns {0,1}^4
\newcommand{\Yscset}{\mathcal{Y}_{\singlec}}% single-click outcomes
\newcommand{\Ymcset}{\mathcal{Y}_{\mc}}     % multi-click outcomes
\newcommand{\Tset}{\mathcal{T}_{\mathrm{set}}}% private setting set (the probe settings in the optical instance)
\newcommand{\Xset}{\mathcal{X}}             % Alice's alphabet
\newcommand{\Dset}{\{H,V,+,-\}}             % polarization / detector labels
\newcommand{\pass}{\top}                    % pass sifting


% -----------------------------------------------------------------------------
% 7. Mode / outcome abbreviations
% -----------------------------------------------------------------------------
\newcommand{\on}{{\mathrm{on}}}
\newcommand{\off}{{\mathrm{off}}}
\newcommand{\dc}{{\mathrm{dc}}}             % dark count
\newcommand{\sq}{\mathrm{sq}}               % squashed
\newcommand{\mc}{\mathrm{mc}}               % multi-click
\newcommand{\nc}{\mathrm{nc}}               % no-click
\newcommand{\singlec}{\mathrm{sc}}          % single-click
\newcommand{\vac}{\mathrm{vac}}             % vacuum
\newcommand{\kout}{{k_\mathrm{out}}}
\newcommand{\kin}{{k_\mathrm{in}}}


% -----------------------------------------------------------------------------
% 8. Maps, channels and miscellaneous operators
% -----------------------------------------------------------------------------
\newcommand{\Prob}{\mathrm{Prob}}           % probability (use \Prob[\cdot] or \Prob(\cdot))
\newcommand{\lmin}{\lambda_{\min}}          % minimum eigenvalue
\newcommand{\Wflag}{W_{\Flag}}              % flag-space weight
\newcommand{\transpose}{\mathsf{T}}
\newcommand{\innerp}[2]{\left\langle #1 \vert #2 \right\rangle}
\renewcommand{\vec}[1]{\boldsymbol{#1}}     % NB: overrides the LaTeX arrow accent


% -----------------------------------------------------------------------------
% 9. Key-rate numerics
% -----------------------------------------------------------------------------
\newcommand{\cG}{\mathcal{G}}               % protocol map
\newcommand{\cZ}{\mathcal{Z}}               % pinching map
\newcommand{\Amap}{\mathcal{A}_{\mathrm{ann}}} % announcement + measurement map
\newcommand{\Amapnb}{\mathcal{A}_{\mathrm{ann}|\mathrm{nb}}}
\newcommand{\simpl}{}                       % empty suffix for simplified block maps
\newcommand{\acSet}{\mathcal{S}_{\mathrm{acc}}} % acceptance set
\newcommand{\testE}{\mathcal{E}_{\mathrm{test}}} % test-statistics map
\newcommand{\ecLeak}{\mathrm{leak}^{\mathrm{EC}}_{\mathrm{obs}}}
\newcommand{\pgen}{p_{\mathrm{gen}}}
\newcommand{\askeyr}{R_{\infty}}             % asymptotic key rate
\newcommand{\ppass}{p_{\mathrm{pass}}}       % sifting pass probability


% -----------------------------------------------------------------------------
% 10. Finite-size analysis / Renyi entropies
% -----------------------------------------------------------------------------

% --- Round types: test vs. key generation (register T, alphabet \Fset) --------
\newcommand{\tst}{\mathrm{test}}            % test round
\newcommand{\gen}{\mathrm{gen}}             % key-generation round
\newcommand{\Fset}{\mathcal{T}}             % round-type alphabet {test, gen} (not \Tset)
\newcommand{\Chat}{\widehat{C}}             % acceptance-test register
\newcommand{\Chatset}{\widehat{\mathcal{C}}}% its alphabet

% --- Basis / key extraction functions ----------------------------------------
\newcommand{\bA}{b_{A}}                     % Alice's basis-extraction function
\newcommand{\bB}{b_{B}}                     % Bob's basis-extraction function
\newcommand{\kB}{k_{B}}                     % Bob's key-extraction function

% --- Maps --------------------------------------------------------------------
\newcommand{\Mprot}{\mathcal{M}}            % single-round protocol map
\newcommand{\Phimap}{\boldsymbol{\Phi}}     % test-statistics map
\newcommand{\Pdeph}{\mathcal{P}}            % dephasing (pinching) map on \Blreg

% --- Isometries --------------------------------------------------------------
\newcommand{\Vmsm}{V_{\mathrm{msm}}}        % measurement isometry
\newcommand{\Vann}{V_{\mathrm{ann}}}        % announcement isometry
\newcommand{\VI}{V_{I}}                     % Eve's copy of the announcements
\newcommand{\VChat}{V_{\Chat}}              % acceptance-data isometry
\newcommand{\VR}{V_{R}}                     % sifting + key map -> raw key R
\newcommand{\VZ}{V_{Z}}                     % raw key -> key register Z
\newcommand{\Wprot}{W}                      % full protocol isometry

% --- Renyi entropies ---------------------------------------------------------
% SECURITY-CRITICAL: see spec/notation.md §4. Never change an arrow to
% make an inequality work out; never convert silently between the sandwiched
% (tilde) and Petz (overline) families.
\newcommand{\Hsandup}{\widetilde{H}^{\uparrow}_{\alpha}}   % sandwiched, up
\newcommand{\Hsanddown}{\widetilde{H}^{\downarrow}_{\alpha}}% sandwiched, down
\newcommand{\Hpetzup}{\overline{H}^{\uparrow}_{\alpha}}    % Petz, up
\newcommand{\Hpetzdown}{\overline{H}^{\downarrow}_{\alpha}}% Petz, down
\newcommand{\hup}{h^{\uparrow}_{\alpha}}                   % single-round rate
\newcommand{\hqkd}{h^{\uparrow,\mathrm{QKD}}_{\alpha}}     % single-round QKD rate
\newcommand{\hlow}{\underline{h}_{\alpha}}                 % lower bound on \hqkd
% --- spec/notation.md 3.5 and 4.5 -------------------------------------------
% \hlb  is the first-order linearization lower bound on \hlow: the headline
%       key rate (the `first-order` tier of foundations/def:tiers). NOT \hlow
%       itself.
% \tnum is the instance's one declared numerical tolerance. Its
%       near-collision with the declared marginal \tau_A is disclosed in
%       spec/notation.md 3.5; bare \tau is reserved and denotes neither.
\newcommand{\hlb}{\underline{h}^{\mathrm{FO}}_{\alpha}}    % first-order lower bound on \hlow
\newcommand{\tnum}{\tau_{\mathrm{num}}}                    % declared numerical tolerance


% =============================================================================
%  PART B -- module-local macros, grouped by module. Deprecated aliases,
%  defined so that a text written with them compiles, used in no new contract: \Lop,
%  \testE, \Pdeph, \pass, \ecLeak (spec/notation.md section 3.9).
% =============================================================================

% --- meat-core -----------------------------------------------------------------
\newcommand{\SecSet}{\mathbb{S}}            % sector index set of Q
\newcommand{\Qsys}{Q}                       % declared sector-carrying measured system
\newcommand{\Vset}{\mathcal{V}}             % declared joint outcome alphabet of the Q-POVM
\newcommand{\CAset}{\mathcal{C}_{A}}        % Alice announcement alphabet
\newcommand{\CBset}{\mathcal{C}_{B}}        % Bob announcement alphabet
\newcommand{\Party}{\mathfrak{P}}           % declared party partition
\newcommand{\Oset}[1]{\mathcal{O}_{#1}}     % readout alphabet of party #1
\newcommand{\rd}[1]{\pi_{#1}}               % readout map of party #1
\newcommand{\Sadm}{\mathcal{S}_{A}}         % admissible Alice subspace
\newcommand{\nomarg}{\mathsf{none}}         % "no trusted marginal declared"
\newcommand{\Zset}{\mathcal{Z}_{\perp}}     % raw-key alphabet incl. \perp
\newcommand{\Zval}{\mathcal{Z}}             % key alphabet without \perp
\newcommand{\Etil}{\widetilde{E}}           % single-round purifying register
\newcommand{\Ehat}{\widehat{E}}             % Eve's n-round register
\newcommand{\OmAT}{\Omega_{\mathrm{AT}}}    % acceptance event
\newcommand{\epa}{\varepsilon_{\mathrm{PA}}}
\newcommand{\eev}{\varepsilon_{\mathrm{EV}}}
\newcommand{\lEC}{\lambda_{\mathrm{EC}}}
\newcommand{\Decl}{\mathfrak{D}}            % protocol declaration
\newcommand{\CList}{\mathfrak{L}}           % declared constraint list
\newcommand{\PinS}{\mathcal{P}_{\SecSet}}   % sector pinching on Q
\newcommand{\Feas}{\mathcal{W}}             % feasible set of block states
\newcommand{\Chkey}{\Chatset_{\mathrm{key}}}% key-live acceptance values
\newcommand{\Kann}{\mathcal{K}}             % announcement-class alphabet
\newcommand{\hold}{\mathsf{h}}              % declared key holder (A or B)
\newcommand{\sfD}{D}                        % classical KL divergence (base 2)
\newcommand{\Dtwo}{D}                       % quantum relative entropy (base 2)
\newcommand{\Dsand}{\widetilde{D}_{\alpha}} % sandwiched Renyi divergence
\newcommand{\idmap}{\mathrm{id}}            % identity channel
\newcommand{\Aev}{\mathcal{A}}              % Eve's per-round attack channel
\newcommand{\AttSet}{\mathfrak{A}}          % licensed attack-channel set
\newcommand{\Chacc}{\Chatset_{\mathrm{acc}}}% acceptance support of \acSet
\newcommand{\epsstar}{\varepsilon^{\star}}      % distance to an exact POVM
\newcommand{\dfH}{\Delta_{\mathrm{herm}}}       % Hermiticity defect
\newcommand{\dfP}{\Delta_{\mathrm{pos}}}        % negative part
\newcommand{\dfC}{\Delta_{\mathrm{cpl}}}        % completeness defect
\newcommand{\lmoq}{g}                       % acceptance-side gradient block
\newcommand{\lmoJ}{G}                       % feasible-side gradient block
\newcommand{\bacc}{\beta_{\mathrm{acc}}}    % certified acceptance-side minimum
\newcommand{\Itset}{\mathfrak{I}}           % finite index set of iterates (C22)

% --- reductions ----------------------------------------------------------------
\newcommand{\Bsys}{B}                        % physical (pre-squashing) system
\newcommand{\Tcl}[1]{\mathcal{T}(#1)}        % trace-class operators on H(#1)
\newcommand{\Ins}{\mathcal{I}}               % generalized-instrument branch map
\newcommand{\redNmsm}{\mathcal{N}}              % POVM measurement channel
\newcommand{\Eprot}{\mathcal{E}}             % n-round protocol channel
\newcommand{\Grnd}{\mathsf{G}}               % one round map of a protocol
\newcommand{\Kfin}{\mathsf{K}}               % end-of-protocol processing
\newcommand{\Crit}{\mathsf{sec}}             % value of a security criterion
\newcommand{\Pflag}{\Pi_{\Flag}}             % flag-sector projector on Q
\newcommand{\Pleq}{\Pi_{\leqN}}              % preserved-block projector
\newcommand{\Pgt}{\Pi_{\gtN}}                % above-cutoff projector
\newcommand{\Sph}{S^{\mathrm{ph}}}           % family of physical inputs
\newcommand{\Srch}{S^{\mathrm{rch}}}         % reachable squashed states
\newcommand{\Ssq}{S^{\sq}}                   % squashed admissible set
\newcommand{\Ffss}{F^{\mathrm{FSS}}}         % flag-state-squasher target
\newcommand{\Fwp}{F^{\mathrm{WP}}}           % multi-bound WPFSS target
\newcommand{\Mwp}{M^{\mathrm{WP}}}           % renormalized above-cutoff POVM
\newcommand{\Mdec}{{M'}^{\mathrm{ph}}}       % declared physical family on B
\newcommand{\Ppvm}{P^{\Flag}}                % flag PVM
\newcommand{\Lmp}{{\Lambda_{\mathrm{mp}}}}   % measure-and-prepare channel
\newcommand{\Lwp}{{\Lambda^{\mathrm{WP}}}}   % multi-bound WPFSS channel
\newcommand{\Lfss}{{\Lambda^{\mathrm{FSS}}}} % flag-state-squasher channel
\newcommand{\lmax}{\lambda_{\max}}           % maximum eigenvalue
\newcommand{\Lnoise}{\widetilde{\Lambda}}    % flag-noise channel
\newcommand{\cg}{\varpi}                     % declared coarse-graining map
\newcommand{\MCK}{\mathsf{MC}}               % multiclick raw-pattern set
\newcommand{\redChain}{\mathfrak{C}}            % reduction chain
\newcommand{\Lnk}{\ell}                      % a single typed link
\newcommand{\dsq}{\preceq}                      % the squashing relation
\newcommand{\opn}[1]{\left\lVert #1\right\rVert}           % operator (spectral) norm
\newcommand{\mxn}[1]{\left\lVert #1\right\rVert_{\max}}    % entrywise max modulus
\newcommand{\Fld}{\mathbb{F}}                   % exact ordered substrate field
\newcommand{\Fldi}{\mathbb{F}_{i}}              % its Gaussian extension F + iF
\newcommand{\lat}[1]{\big\lfloor #1\big\rfloor_{\kappa}} % round DOWN onto the lattice

% --- protocol-compiler ---------------------------------------------------------
\newcommand{\Spec}{\mathfrak{S}}             % protocol specification (compiler input)
\newcommand{\pcChain}{\mathfrak{R}}            % reduction chain (two typed segments)
\newcommand{\Chsrc}{\pcChain_{\mathrm{src}}}   % its Alice-side source-map segment
\newcommand{\Chrec}{\pcChain_{\mathrm{rec}}}   % its receiver-side squashing segment
\newcommand{\rvirt}{\rho^{\mathrm{v}}}       % virtual (PRE-source-map) signal state
\newcommand{\Fph}{F^{\mathrm{ph}}}           % physical (pre-reduction) POVM element
\newcommand{\Vph}{\mathcal{V}^{\mathrm{ph}}} % physical outcome alphabet
\newcommand{\Tpub}{\mathcal{T}_{\mathrm{pub}}}% public-choice alphabet
\newcommand{\Uset}{\mathcal{U}}              % residual outcome alphabet
\newcommand{\idxm}{\mathsf{idx}}             % label-to-index map
\newcommand{\Asig}{A'}                       % transmitted (signal) system -- the REAL one
\newcommand{\Avirt}{A''}                     % virtual transmitted system (source descriptor of Xi_1)
\newcommand{\Abar}{\bar{A}}                  % Alice's measured register
\newcommand{\Ahat}{\hat{A}}                  % Alice's shield register
\newcommand{\pcNmsm}{\mathcal{N}_{\mathsf A}}  % Alice's measurement channel (see the RULES block)
\newcommand{\Phip}{\Phi^{+}}                 % maximally entangled vector label
\newcommand{\Smap}{\mathfrak{s}}             % the emitted sector map (fraktur: a compile-time object)

% --- receiver ------------------------------------------------------------------
\newcommand{\Modes}{\mathbb{M}}              % declared ordered mode set
\newcommand{\Inp}{\mathbb{I}}                % declared input modes
\newcommand{\InpC}{\mathbb{I}_{\Chan}}       % channel input modes
\newcommand{\InpS}{\mathbb{I}_{\Prb}}        % probe (source) input modes
\newcommand{\Env}{\mathbb{E}}                % dilation (environment) modes
\newcommand{\Port}[1]{\mathbb{P}_{#1}}       % mode set of detector port #1
\newcommand{\Emd}[1]{E_{#1}}                 % coordinate projector on modes #1
\newcommand{\Gmat}{G}                        % instrument (transfer) matrix
\newcommand{\Fock}{\mathfrak{F}}             % Fock / second-quantization functor
\newcommand{\Symn}[1]{\mathrm{Sym}^{#1}}     % symmetric tensor power
\newcommand{\normord}[1]{{:}\mspace{2mu}#1\mspace{2mu}{:}}   % normal ordering
\newcommand{\rk}{\operatorname{rank}}        % rank (physics.sty owns \rank)
\newcommand{\SCK}{\mathsf{SC}}               % single-click pattern set
\newcommand{\Recv}{\mathfrak{R}}             % receiver instrument tuple
\newcommand{\Cat}{\mathfrak{C}}              % component catalogue
\newcommand{\Grph}{\mathfrak{E}}             % declared element graph
\newcommand{\defic}{\mathbb{L}}              % loss deficiency operator
\newcommand{\Pvac}[1]{\Pi_{\vac}(#1)}        % joint no-click element on ports #1
\newcommand{\pmc}{p_{\mc}}                   % uniform high-photon mc weight
\newcommand{\Bnd}{\beta}                     % the scalar bound function beta(n)
\newcommand{\Rot}{R}                         % planar rotation matrix
\newcommand{\Wplate}{W}                      % retarder (wave-plate) matrix
\newcommand{\Bsp}{B}                         % beamsplitter core matrix

% --- statistics ----------------------------------------------------------------
\newcommand{\honest}{p^{\mathrm{hon}}}          % honest per-round distribution
\newcommand{\ecomp}{\varepsilon_{\mathrm{comp}}}% completeness budget
\newcommand{\simplex}[1]{\Delta(#1)}            % probability simplex
\newcommand{\Chatcon}{\Chatset_{\mathrm{con}}}  % constrained outcomes
\newcommand{\Ireg}[3]{\mathcal{I}_{#1}(#2,#3)}  % regularized incomplete beta
\newcommand{\Ncount}[1]{\mathsf{N}_{#1}}        % per-outcome count
\newcommand{\qlo}{\mathbf{q}^{-}}               % box lower bounds
\newcommand{\qhi}{\mathbf{q}^{+}}               % box upper bounds

\makeatletter
\newcommand*{\ClaimHeadFont}{\bfseries}   % switched by prominence
\newcommand*{\ClaimNoteFont}{\mdseries}
\newcommand*{\ClaimShortId}{}
\newcommand*{\SetShortId}[1]{\renewcommand*{\ClaimShortId}{#1}}
\newcommand*{\RestatedName}{}
% Upright bodies throughout: italic is reserved for non-normative text.
% \@endparpenalty\@M keeps the meta line on the statement's page.
\newtheoremstyle{claim}{\topsep}{\topsep}{\normalfont\@endparpenalty\@M}{}%
  {\ClaimHeadFont}{.}{.5em}{\thmname{#1}\thmnumber{ #2}\thmnote{ {\ClaimNoteFont(#3)}}}
\newtheoremstyle{claimass}{\topsep}{\topsep}{\normalfont\@endparpenalty\@M}{}%
  {\ClaimHeadFont}{.}{.5em}{\thmname{#1}\thmnumber{ #2}%
  \ifx\ClaimShortId\@empty\else\ [\textsf{\ClaimShortId}]\fi%
  \thmnote{ {\ClaimNoteFont(#3)}}}
\newtheoremstyle{claimremark}{.5\topsep}{.5\topsep}{\small\@endparpenalty\@M}{}%
  {\itshape}{.}{.5em}{\thmname{#1}\thmnumber{ #2}\thmnote{ {\upshape(#3)}}}
\newtheoremstyle{claimreq}{\topsep}{\topsep}{\normalfont\@endparpenalty\@M}{}%
  {\color{reqAccent}\scshape}{.}{.5em}%
  {\thmname{#1}\thmnumber{ #2}\thmnote{ {\upshape(#3)}}}
\makeatother
\theoremstyle{claim}
\newtheorem{theorem}{Theorem}[section]
\newtheorem{lemma}[theorem]{Lemma}
\newtheorem{proposition}[theorem]{Proposition}
\newtheorem{corollary}[theorem]{Corollary}
\newtheorem{fact}[theorem]{Fact}
\newtheorem{definition}[theorem]{Definition}
\newtheorem*{restatedclaim}{\RestatedName}
\theoremstyle{claimass}
\newtheorem{assumption}[theorem]{Assumption}
\newtheorem*{restatedassumption}{\RestatedName}
\theoremstyle{claimremark}
\newtheorem{remark}[theorem]{Remark}
\newtheorem{example}[theorem]{Example}
\theoremstyle{claimreq}
\newtheorem{requirement}[theorem]{Requirement}
\newtheorem*{restatedrequirement}{\RestatedName}
%% --- 8. Prominence: main results a margin rule and a bold heading ----------
\newtcolorbox{mainclaim}{enhanced jigsaw, breakable, blanker,
  borderline west={1.5pt}{-8pt}{mainRule}, left=0pt, right=0pt, top=0pt, bottom=0pt,
  boxsep=0pt, before skip=\topsep, after skip=\topsep, parbox=false,
  before upper={\renewcommand*{\ClaimNoteFont}{\bfseries}}}
\newenvironment{technicalclaim}{\renewcommand*{\ClaimHeadFont}{\mdseries\color{techGrey}}%
  \renewcommand*{\ClaimNoteFont}{}}{}
%% --- 9. Non-normative text, the meta line, navigation ------------------------
\makeatletter
% Tie what follows to this paragraph: no page break, no list penalty before
% the next theorem head (the \@nobreak route a section heading uses).
\newcommand{\keepwithnext}{\par\nobreak\vspace{3pt}\@nobreaktrue
  \everypar{\if@nobreak\@nobreakfalse\clubpenalty\@M\else\clubpenalty\@clubpenalty\everypar{}\fi}}
\newcommand{\claimsummary}[2]{% #1 claim id (named destination), #2 summary
  \par\addvspace{\bigskipamount}\Needspace{6\baselineskip}%
  {\small\itshape\color{nnGrey}\noindent\hypertarget{claim-#1}{}#2\par}\keepwithnext}
\newcommand{\claimanchor}[1]{% a claim without a summary: the destination alone
  \par\addvspace{\bigskipamount}\Needspace{6\baselineskip}%
  \phantomsection\hypertarget{claim-#1}{}\keepwithnext}
\newcommand{\claimleadin}[1]{{\color{nnGrey}\noindent#1\par}\keepwithnext}
\newcommand{\claimmeta}[1]{\par\nobreak\vspace{2pt}%
  {\footnotesize\color{metaGrey}\raggedright\interlinepenalty\@M\noindent#1\par}%
  \addvspace{\medskipamount}}
% After an assumption: the sentence generated from its verification field; after
% any claim with one: its plain-text reading (user_statement). Both
% non-normative, both grey.
\newcommand{\claimverification}[1]{\par\nobreak\vspace{2pt}%
  {\footnotesize\color{metaGrey}\noindent#1\par}\nobreak}
\newcommand{\claimreading}[1]{\par\nobreak\vspace{2pt}%
  {\footnotesize\color{nnGrey}\noindent\emph{Reading:}\enspace#1\par}\nobreak}
\makeatother
\newcommand{\claimbookmark}[4]{% #1 id, #2 kind word, #3 label, #4 plain title
  \bookmark[dest=claim-#1,level=2]{#2 \getrefnumber{#3} #4}}
\newcommand{\metasep}{\unskip\enspace\textperiodcentered\enspace\ignorespaces}
\newcommand{\metalabel}[1]{\emph{#1}\enspace}
\newcommand{\sectionintro}[1]{{\color{nnGrey}\noindent#1\par}\smallskip}
\newcommand{\sectionimports}[1]{{\footnotesize\color{metaGrey}\raggedright\noindent\metalabel{Imports}#1\par}\smallskip}
\newenvironment{overview}{\par\color{nnGrey}\itshape\noindent\ignorespaces}{\par}
% References: the whole "Def. 2.1" is the link; ids are the fallback text of
% every link into another PDF (a relative URI to a named destination).
\newcommand{\mref}[2]{\hyperref[#2]{#1~\ref*{#2}}}
\newcommand{\aref}[2]{\hyperref[#2]{\textsf{#1}}}
\DeclareUrlCommand{\claimid}{\urlstyle{sf}\def\UrlBreaks{\do\/\do\:\do\-\do\.\do\_}%
  \def\UrlBigBreaks{}\def\UrlNoBreaks{}}
\newcommand{\extref}[2]{\hyperref{#1}{}{claim-#2}{\claimid{#2}}}
\newcommand{\extaref}[3]{\hyperref{#1}{}{claim-#2}{\textsf{#3}}}
\newcommand{\coderef}[2]{\href{#1}{\claimid{#2}}}
% Claim-body commands (spec/claim-format.md section 4): \xref{<full id>} prints
% the kind, title and id of a claim of another contract, linked to its stable
% destination in that contract's PDF (one \DeclareXref per cited claim below);
% \term[gloss]{phrase} marks a defined phrase for the index (the gloss feeds
% the Terms table); \tref[phrase]{text} is a use of a defined term. The
% assembler rewrites both in the body into \termdef (the definition's named
% destination term-<slug>) and \termref / \exttermref (a link to it, in this
% or another contract's PDF); the fallbacks below serve text it copies
% elsewhere, such as the digest. \iface is an interface key.
\makeatletter
\newcommand{\DeclareXref}[4]{% id, kind word, title, contract PDF
  \@namedef{xref@#1}{\hyperref{#4}{}{claim-#1}{#2 ``#3'' ({\small\claimid{#1}})}}}
\newcommand{\xref}[1]{\@ifundefined{xref@#1}%
  {\textbf{[unresolved \nolinkurl{#1}]}\@latex@warning{Unresolved xref #1}}%
  {\@nameuse{xref@#1}}}
\makeatother
\DeclareUrlCommand{\iface}{\urlstyle{tt}}
\newcommand{\term}[2][]{\emph{#2}\index{#2}}
\newcommand{\tref}[2][]{#2}
\newcommand{\termdef}[2]{\hypertarget{#1}{}\emph{#2}\index{#2}}
\newcommand{\termref}[2]{\hyperlink{#1}{#2}}
\newcommand{\exttermref}[3]{\hyperref{#1}{}{#2}{#3}}
\hypersetup{pdftitle={MEAT key-rate engine (meat-core) -- statements digest 1.32.0}}

% The digest: the contract's labels under the prefix C-, its numbers checked.
\externaldocument[C-]{contract}
\theoremstyle{claim}\newtheorem*{digestclaim}{\RestatedName}
\theoremstyle{claimass}\newtheorem*{digestassumption}{\RestatedName}
\theoremstyle{claimreq}\newtheorem*{digestrequirement}{\RestatedName}
\makeatletter
% #1 label, #2 the number the assembler computed: a local label carrying it, so
% that a reference between two digest items stays in the digest; the build
% stops when the contract's own number differs.
\newcommand{\digestlabel}[2]{\phantomsection
  \edef\digest@theirs{\getrefnumber{C-#1}}\def\digest@ours{#2}%
  \ifx\digest@theirs\digest@ours\else
    \PackageError{digest}{#1 is #2 here but \digest@theirs\space in the contract}{}\fi
  \def\@currentlabel{#2}\label{#1}}
\makeatother
% A display's local label (amsmath records the \tag, i.e. the contract's number).
\newcommand{\digesteqlabel}[1]{\label{#1}}

% Claims of other contracts cited here (\xref, [[...]]).
\DeclareXref{foundations/ass:snapped-measurement}{Assumption}{The device measures the snapped family where the run reads the snap}{../foundations/contract.pdf}
\DeclareXref{foundations/cor:least-eigenvalue-over-enclosure}{Corollary}{Least eigenvalue and rank over an enclosure}{../foundations/contract.pdf}
\DeclareXref{foundations/cor:rank-over-factor-enclosure}{Corollary}{Rank over a factor enclosure}{../foundations/contract.pdf}
\DeclareXref{foundations/def:binary64}{Definition}{Binary64 numbers, declared values and dyadic lattices}{../foundations/contract.pdf}
\DeclareXref{foundations/def:declared-order}{Definition}{Sums in a declared order and sums that round nowhere}{../foundations/contract.pdf}
\DeclareXref{foundations/def:evaluated-off-snap}{Definition}{A reference handed beside a snap; a block evaluated off its snap}{../foundations/contract.pdf}
\DeclareXref{foundations/def:exact-island}{Definition}{The exact island}{../foundations/contract.pdf}
\DeclareXref{foundations/def:exact-reference}{Definition}{Reference of a measurement family, exact or enclosed}{../foundations/contract.pdf}
\DeclareXref{foundations/def:factor-enclosure}{Definition}{Factor enclosure of a block}{../foundations/contract.pdf}
\DeclareXref{foundations/def:kl}{Definition}{Classical divergences and the quantum relative entropy}{../foundations/contract.pdf}
\DeclareXref{foundations/def:pinching}{Definition}{Operator conventions and pinching}{../foundations/contract.pdf}
\DeclareXref{foundations/def:renyi}{Definition}{Sandwiched Rényi divergence and conditional entropies}{../foundations/contract.pdf}
\DeclareXref{foundations/def:round-model}{Definition}{Round model, attack sequences and attack sets}{../foundations/contract.pdf}
\DeclareXref{foundations/def:sector-map}{Definition}{Sector maps and block-list operators}{../foundations/contract.pdf}
\DeclareXref{foundations/def:snap}{Definition}{Snapping declared data onto a lattice}{../foundations/contract.pdf}
\DeclareXref{foundations/def:tiers}{Definition}{The tiers of a reported key rate}{../foundations/contract.pdf}
\DeclareXref{foundations/def:tolerance}{Definition}{The declared tolerance and the defects it bounds}{../foundations/contract.pdf}
\DeclareXref{foundations/fact:concavity}{Fact}{Concavity of the pinched trace-power functional}{../foundations/contract.pdf}
\DeclareXref{foundations/fact:pinch-id}{Fact}{Pinched trace-power form of the conditional Rényi entropy}{../foundations/contract.pdf}
\DeclareXref{foundations/lem:floor-stricter}{Lemma}{The positivity floor against entrywise closeness to the cone}{../foundations/contract.pdf}
\DeclareXref{foundations/lem:kernel}{Lemma}{Extended divergence kernel}{../foundations/contract.pdf}
\DeclareXref{foundations/lem:least-eigenvalue-bound}{Lemma}{Least-eigenvalue lower bound from an approximate eigenbasis}{../foundations/contract.pdf}
\DeclareXref{foundations/lem:matrix-function-perturbation}{Lemma}{Perturbation bounds for matrix functions}{../foundations/contract.pdf}
\DeclareXref{foundations/lem:pinchmono}{Lemma}{Pinching lowers the trace power}{../foundations/contract.pdf}
\DeclareXref{foundations/lem:recentred-enclosure}{Lemma}{Zeroing the centres an enclosure cannot tell from zero}{../foundations/contract.pdf}
\DeclareXref{foundations/lem:register-compression}{Lemma}{Compressing a register onto a finite-dimensional subspace}{../foundations/contract.pdf}
\DeclareXref{foundations/rem:noise-floor}{Remark}{What domination controls, and why it needs a noise floor}{../foundations/contract.pdf}
\DeclareXref{foundations/req:consumer-decides-exactly}{Requirement}{Exact decision by every consumer}{../foundations/contract.pdf}
\DeclareXref{foundations/req:declared-tolerance}{Requirement}{One declared tolerance per instance}{../foundations/contract.pdf}
\DeclareXref{foundations/req:exact-emission}{Requirement}{Exact verification before emission}{../foundations/contract.pdf}
\DeclareXref{foundations/req:exact-island}{Requirement}{Exact arithmetic on the island, unbudgeted rounding elsewhere}{../foundations/contract.pdf}
\DeclareXref{foundations/req:norm-conversion}{Requirement}{Explicit conversion between norms}{../foundations/contract.pdf}
\DeclareXref{foundations/req:rounding-direction}{Requirement}{Rounding direction of emitted bounds}{../foundations/contract.pdf}
\DeclareXref{foundations/req:snap-reference}{Requirement}{Snaps of a measurement family: kept, or moved and disclosed}{../foundations/contract.pdf}
\DeclareXref{foundations/req:tiers}{Requirement}{Reporting of the tiers}{../foundations/contract.pdf}
\DeclareXref{foundations/req:tolerance-scope}{Requirement}{Scope of the declared tolerance}{../foundations/contract.pdf}
\DeclareXref{protocol-compiler/ass:alice-meas}{Assumption}{Alice's declared measurement where no source certificate is demanded}{../protocol-compiler/contract.pdf}
\DeclareXref{protocol-compiler/ass:atk}{Assumption}{Attack-set licence}{../protocol-compiler/contract.pdf}
\DeclareXref{protocol-compiler/ass:link}{Assumption}{Soundness of every reduction link}{../protocol-compiler/contract.pdf}
\DeclareXref{protocol-compiler/ass:rhom}{Assumption}{Round homogeneity of the declaration}{../protocol-compiler/contract.pdf}
\DeclareXref{protocol-compiler/ass:spec}{Assumption}{Specification faithful to the apparatus}{../protocol-compiler/contract.pdf}
\DeclareXref{protocol-compiler/def:compile}{Definition}{The compilation map}{../protocol-compiler/contract.pdf}
\DeclareXref{protocol-compiler/def:declaration}{Definition}{Protocol declaration}{../protocol-compiler/contract.pdf}
\DeclareXref{protocol-compiler/def:exact}{Definition}{Exact admissibility of declared operator data}{../protocol-compiler/contract.pdf}
\DeclareXref{protocol-compiler/def:feasible}{Definition}{Feasible set}{../protocol-compiler/contract.pdf}
\DeclareXref{protocol-compiler/def:homog}{Definition}{Licensed constraint list}{../protocol-compiler/contract.pdf}
\DeclareXref{protocol-compiler/def:target}{Definition}{Target conditions of a well-formed declaration}{../protocol-compiler/contract.pdf}
\DeclareXref{protocol-compiler/def:validator}{Definition}{The emitted validator}{../protocol-compiler/contract.pdf}
\DeclareXref{protocol-compiler/def:weighted}{Definition}{Choice-weighted families}{../protocol-compiler/contract.pdf}
\DeclareXref{protocol-compiler/lem:weighted}{Lemma}{Properties of the choice-weighted families}{../protocol-compiler/contract.pdf}
\DeclareXref{protocol-compiler/lem:weighted-factor}{Lemma}{A factor enclosure through a setting-dependent weighting}{../protocol-compiler/contract.pdf}
\DeclareXref{protocol-compiler/rem:block-restriction}{Remark}{Restriction of the feasible set to block-diagonal states}{../protocol-compiler/contract.pdf}
\DeclareXref{protocol-compiler/rem:engine-admission}{Remark}{Relation to the consuming engine's admission conditions}{../protocol-compiler/contract.pdf}
\DeclareXref{protocol-compiler/rem:obligations-discharged}{Remark}{Discharge of the chain obligations by the reductions module}{../protocol-compiler/contract.pdf}
\DeclareXref{protocol-compiler/rem:party-local-soundness}{Remark}{Soundness of a party-local entry}{../protocol-compiler/contract.pdf}
\DeclareXref{protocol-compiler/req:no-fixed-fraction}{Requirement}{No fixed fraction outside the measurement operators}{../protocol-compiler/contract.pdf}
\DeclareXref{protocol-compiler/req:reference-companion-compiled}{Requirement}{The reference companion compiled beside the declaration}{../protocol-compiler/contract.pdf}
\DeclareXref{protocol-compiler/thm:band-a}{Theorem}{The Band-A route}{../protocol-compiler/contract.pdf}
\DeclareXref{protocol-compiler/thm:feasibility}{Theorem}{Feasibility by an explicit or a verified candidate}{../protocol-compiler/contract.pdf}
\DeclareXref{receiver/ass:trust}{Assumption}{Trusted and exactly characterised receiver}{../receiver/contract.pdf}
\DeclareXref{receiver/req:eig-ledger-row}{Requirement}{Invoking settings on the eigenvalue-slack row of the ledger}{../receiver/contract.pdf}
\DeclareXref{receiver/req:evaluation-label}{Requirement}{The evaluation label on every artifact}{../receiver/contract.pdf}
\DeclareXref{receiver/thm:enclosed-weight}{Theorem}{The weight on enclosed bound data bounds the multiclick element}{../receiver/contract.pdf}
\DeclareXref{reductions/ass:recv}{Assumption}{Supplier artefact whose claims hold of the device}{../reductions/contract.pdf}
\DeclareXref{reductions/ass:tsm}{Assumption}{Output-state test statistics}{../reductions/contract.pdf}
\DeclareXref{reductions/cor:companion-factor-enclosure}{Corollary}{Factor enclosures carried to the reference companion}{../reductions/contract.pdf}
\DeclareXref{reductions/cor:reference-companion}{Corollary}{The reference companion of a reduction instance}{../reductions/contract.pdf}
\DeclareXref{reductions/cor:sequential}{Corollary}{Output-state equality against adaptive attack sequences}{../reductions/contract.pdf}
\DeclareXref{reductions/def:attset}{Definition}{Admissible attack sets in the consumer's format}{../reductions/contract.pdf}
\DeclareXref{reductions/def:protocol}{Definition}{Measurement-factorizing $n$-round protocol}{../reductions/contract.pdf}
\DeclareXref{reductions/def:reduction-instance}{Definition}{Reduction instance}{../reductions/contract.pdf}
\DeclareXref{reductions/def:roundmodel}{Definition}{Round model and attack sequences, typed by source mode}{../reductions/contract.pdf}
\DeclareXref{reductions/lem:row-relaxation}{Lemma}{A rounded constraint entry with its constant relaxed}{../reductions/contract.pdf}
\DeclareXref{reductions/lem:source-transfer}{Lemma}{Exactness of the source map in trusted-source mode}{../reductions/contract.pdf}
\DeclareXref{reductions/rem:why-exact}{Remark}{Exactness of the reductions and generality of the transfer}{../reductions/contract.pdf}
\DeclareXref{reductions/req:recv-row-provenance}{Requirement}{Omission of the supplier row for a receiver artefact the run built}{../reductions/contract.pdf}
\DeclareXref{reductions/req:reference-companion}{Requirement}{Emitting the reference companion beside the artefact}{../reductions/contract.pdf}
\DeclareXref{reductions/thm:licence}{Theorem}{The attack-set licence}{../reductions/contract.pdf}
\DeclareXref{reductions/thm:two-sided}{Theorem}{The two-sided reduction in trusted-source mode}{../reductions/contract.pdf}
\DeclareXref{reductions/wpfss/def:weight-route}{Definition}{The weight route}{../reductions/contract.pdf}
\DeclareXref{statistics/ass:reach}{Assumption}{An honest law produced by a feasible state}{../statistics/contract.pdf}
\DeclareXref{statistics/binomial-box/thm:completeness-bound}{Theorem}{The completeness bound}{../statistics/contract.pdf}
\DeclareXref{statistics/singleton/thm:singleton-completeness}{Theorem}{Completeness of a singleton acceptance set}{../statistics/contract.pdf}
\DeclareXref{transmitter/ass:ens}{Assumption}{Emission of the specified ensemble}{../transmitter/contract.pdf}
\DeclareXref{transmitter/ass:seq}{Assumption}{Sequential round structure}{../transmitter/contract.pdf}
\DeclareXref{transmitter/ass:src}{Assumption}{Retention of the system A in trusted-source mode}{../transmitter/contract.pdf}
\DeclareXref{transmitter/decoy/ass:wcp}{Assumption}{Emission of the declared weak-coherent-pulse states}{../transmitter/contract.pdf}
\DeclareXref{transmitter/def:source-segment}{Definition}{Source segment and its interface conditions}{../transmitter/contract.pdf}
\DeclareXref{transmitter/def:source-spec}{Definition}{Source specification}{../transmitter/contract.pdf}
\DeclareXref{transmitter/def:srcert}{Definition}{Source-replacement certificate}{../transmitter/contract.pdf}
\DeclareXref{transmitter/rem:source-mode}{Remark}{The two source modes}{../transmitter/contract.pdf}
\DeclareXref{transmitter/thm:certificate-sufficiency}{Theorem}{Checkability and sufficiency of the source-replacement certificate}{../transmitter/contract.pdf}

\begin{document}
\begin{center}
{\LARGE\bfseries MEAT key-rate engine (meat-core)}\\[4pt]
{\large Statements digest, version 1.32.0}
\end{center}
\noindent{\small Every definition, fact, assumption, result and requirement of the \href{contract.pdf}{contract}, with the contract's numbers; remarks, examples, technical steps and all proofs are left out. Each heading opens its statement in the contract.}

\part*{Mathematics}\pdfbookmark[0]{Mathematics}{digest-part-1}

\section*{1\quad The declaration and the protocol channel}\pdfbookmark[1]{1 The declaration and the protocol channel}{digest-group-meat-core/def:declaration}

% claim meat-core/def:declaration
\digestlabel{def:declaration}{1.1}\pdfbookmark[2]{Definition 1.1 The protocol declaration as the MEAT core reads it}{digest-meat-core/def:declaration}
\renewcommand*{\RestatedName}{\hyperref{contract.pdf}{}{claim-meat-core/def:declaration}{Definition~1.1}}%
\begin{digestclaim}[The protocol declaration as the MEAT core reads it]
In this module a protocol declaration is a tuple
\begin{equation*}
  \Decl=\Big(A,\ \tau_A;\ \Qsys,\ \SecSet;\ \{M_x\}_{x\in\Xset};\
  \{F_v\}_{v\in\Vset};\ \Party;\ a_A,\ a_B;\ \phi;\ r,\ \hold;\
  \CList\Big)
\end{equation*}
of \xref{protocol-compiler/def:declaration}, in its basic format or in
its general format, (DX1)--(DX6) there, whose components satisfy the
well-formedness conditions of Definition~\ref{def:admissible-declaration} and the two
conditions on its raw-key map below. The two values of its
marginal slot, $\tau_A\in\Seq{A}$ and $\tau_A=\nomarg$, are the two
source modes, trusted-source mode and
untrusted-source mode. Write $\Kann$ for the announcement-class
alphabet of that definition, $\CAset\times\CBset$ in the basic format and
the product of the message alphabets of (DX5) in the general one,
$a(x,v)\in\Kann$ for the joint announcement it assigns to an outcome pair,
and $\kappa$ for a class. \emph{Sifting is announcement-determined}: there is
$\varsigma:\Kann\to\{\mathrm{pass},\mathrm{fail}\}$ with
$r(x,v)=\perp\iff\varsigma(a(x,v))=\mathrm{fail}$. \emph{The raw key is
party-computable}: for the declared key holder
$\hold\in\{\mathsf A,\mathsf B\}$ there is
$r_{\hold}:\Oset{\hold}\times\Kann\to\Zset$ with
\begin{equation*}\tag{\getrefnumber{C-eq:partycomp}}\digesteqlabel{eq:partycomp}
  r(x,v)=r_{\hold}\big(\rd{\hold}(x,v),\,a(x,v)\big)
  \qquad\text{for all }(x,v)\in\Xset\times\Vset
\end{equation*}
(Definition~\ref{def:admissible-declaration}\,(viii)). We write
\begin{equation*}\tag{\getrefnumber{C-eq:chkey}}\digesteqlabel{eq:chkey}
  \Chkey\coloneqq\big\{\hat c\in\Chatset:\ \exists\,\kappa\in\phi^{-1}(\hat c)
  \text{ with }\varsigma(\kappa)=\mathrm{pass}\big\}
\end{equation*}
for the set of
\emph{key-live acceptance value}s,
$\varsigma$ being the predicate of the sifting condition above. That
condition fixes $\varsigma$ on every class some outcome pair realises, and
on a class $\kappa\notin a(\Xset\times\Vset)$, which no outcome pair
realises, we take $\varsigma(\kappa)=\mathrm{fail}$; so $\Chkey$ is the set
of acceptance values to which some realised class passing sifting maps,
and it is fixed by the declaration.

\emph{What this module reads of $\Party$.} The party partition
$\Party$ of that definition consists of (a) its two
readout maps $\rd{\mathsf A}$ and $\rd{\mathsf B}$,
naming what each party physically observes in a round, and (b)
\emph{system metadata}:
the party label for $A$ and, for each $s\in\SecSet$, the declared tensor
factorization $\Hs{\Qsys_s}=\bigotimes_{i}\Hs{\Qsys_{s,i}}$ with a party
label per factor (a comma-indexed \emph{tensor} factor $\Qsys_{s,i}$, never
the parenthesized superscript $X^{(s)}$, which this module reserves for
\emph{sector blocks} of an operator). The \emph{default partition} is that of
\xref{protocol-compiler/def:target}: $\rd{\mathsf A}(x,v)=x$,
$\rd{\mathsf B}(x,v)=v$ with $A$ labelled $\mathsf A$ and every $\Qsys_s$
labelled $\mathsf B$. Part~(a) is consumed by the derivations, through
Definition~\ref{def:admissible-declaration}\,(viii). Part~(b) is consumed by no
derivation; it carries an obligation of the compiler: the declaration's POVMs
must be realizable by each party measuring the factors it is labelled with,
so that $\rd{\mathsf A}$ and $\rd{\mathsf B}$ are what the parties can in fact
read. Part~(b) is read by the premises that place the parties'
laboratories: Assumption~\ref{ass:lab} takes all its labels as naming the
parties' laboratories, and A-SRC (\xref{transmitter/ass:src}) reads the label
of $A$; so a run attests that obligation there, and the claims that assume
those premises rest on it through them. This module never verifies the
obligation.

\emph{The declared summation order.} The declaration fixes, as part of the
declared alphabet $\Vset$ of outcomes, a declared order
$v_1\prec v_2\prec\dots\prec v_{|\Vset|}$ on it. For each sector $s$ the
completeness sum is the sum in the declared order of the blocks
$F^{(s)}_v$, the sequential sum
\begin{equation*}\tag{\getrefnumber{C-eq:declorder}}\digesteqlabel{eq:declorder}
  \Sigma^{(s)}_0\coloneqq0,\qquad
  \Sigma^{(s)}_i\coloneqq\Sigma^{(s)}_{i-1}+F^{(s)}_{v_i}
  \quad(1\le i\le|\Vset|),\qquad
  \Sigma^{(s)}\coloneqq\Sigma^{(s)}_{|\Vset|},
\end{equation*}
formed left to right, entrywise, with no re-ordering and no
re-association. In exact arithmetic \eqref{eq:declorder} is the unordered
sum and the order is immaterial; in binary64 it is not
(\xref{foundations/def:declared-order}), so the order is part of the datum and
this module reads it off the declared alphabet rather than choosing one.

\emph{Exact admissibility.} A declaration reaches this module with its
operator data exactly admissible: every $F_v$ exactly Hermitian and
exactly positive semidefinite per sector block, every $\hat O_k$ exactly
Hermitian per sector block (positivity is asked of no constraint operator
anywhere in this module), every $M_x$ and the declared $\tau_A$ exactly
Hermitian and exactly positive semidefinite, and the completeness sum
\eqref{eq:declorder} of each sector exactly the identity with every addition
rounding nowhere. These are the conditions the compiler decides on the
numbers it emits (\xref{protocol-compiler/def:exact}).

\emph{The general format.} The general format is that of
\xref{protocol-compiler/def:declaration}, (DX1)--(DX6) there, and the basic
format is its case of one Alice sector, one measured factor,
POVM form, two messages neither of which reads the other, and no
party-local entry. This contract refers to five of that definition's
displays under its own labels, each read as it is defined there: the
product sector map of (DX1), over the ordered set $\mathbb{F}$ of
measured factors, factor $f$ having the sectors $s_f\in\SecSet^{f}$ with
projectors $\Pi^{f}_{s_f}$ and block dimensions $d_{f,s_f}$,
\begin{equation*}\tag{\getrefnumber{C-eq:productmap}}\digesteqlabel{eq:productmap}
\begin{gathered}
  \Qsys=\bigotimes_{f\in\mathbb{F}}\Qsys_f,\qquad
  \SecSet=\prod_{f\in\mathbb{F}}\SecSet^{f},\qquad
  \Hs{\Qsys_s}=\bigotimes_{f}\Hs{\Qsys_{f,s_f}},\\
  \Pi_s=\bigotimes_f\Pi^{f}_{s_f},\qquad d_s=\prod_fd_{f,s_f};
\end{gathered}
\end{equation*}
the product family of (DX3),
\begin{equation*}\tag{\getrefnumber{C-eq:productfamily}}\digesteqlabel{eq:productfamily}
  F_v\coloneqq\bigotimes_{f\in\mathbb{F}}F^{f}_{v_f},\qquad
  F^{(s)}_v=\bigotimes_{f\in\mathbb{F}}F^{f,(s_f)}_{v_f}
  \qquad\big(v=(v_f)_f,\ s=(s_f)_f\big);
\end{equation*}
the two conditions on the declared $\{N_x\}$ of the ensemble form of
(DX4), whose blocks lie along the index set $\mathbb{T}$ of Alice's sectors,
$\Hs{A}=\bigoplus_{t\in\mathbb{T}}\Hs{A_t}$, and the family of the POVM-form
declaration an ensemble-form declaration denotes, with $P_t$ the projector
onto the support of $\tau^{(t)}_A$ and the inverse square root taken on that
support and zero off it,
\begin{equation*}\tag{\getrefnumber{C-eq:ensembleform}}\digesteqlabel{eq:ensembleform}
  \sum_{x\in\Xset}N^{(t)}_x=\tau^{(t)}_A
  \quad\text{and}\quad\tau^{(t)}_A\ne0
  \qquad\text{for every }t\in\mathbb{T},
\end{equation*}
\begin{equation*}\tag{\getrefnumber{C-eq:ensemblepovm}}\digesteqlabel{eq:ensemblepovm}
  M^{(t)}_x\coloneqq\big(\tau^{(t)}_A\big)^{-1/2}N^{(t)}_x
  \big(\tau^{(t)}_A\big)^{-1/2}
  +\frac{\tr N^{(t)}_x}{\tr\tau^{(t)}_A}\,\big(\idon{A_t}-P_t\big);
\end{equation*}
and the joint announcement function a transcript protocol of (DX5)
derives by composition from its $L$ messages, message $m$ having the sender
$\pi_m$, whose readout map is $\rd{\pi_m}$, the alphabet $\mathcal{C}_m$ and
the table $a_m$,
\begin{equation*}\tag{\getrefnumber{C-eq:transcript}}\digesteqlabel{eq:transcript}
\begin{gathered}
  \kappa_m(x,v)\coloneqq a_m\big(\rd{\pi_m}(x,v),\,\kappa_1(x,v),\dots,
  \kappa_{m-1}(x,v)\big),\\
  a(x,v)\coloneqq\big(\kappa_m(x,v)\big)_{m=1}^{L}\in
  \Kann\coloneqq\prod_{m=1}^{L}\mathcal{C}_m .
\end{gathered}
\end{equation*}

\emph{Exact admissibility in the general format.} The exact-admissibility
paragraph above holds with these readings: the conditions on $\{F_v\}$ are
decided per measured factor, on the family $\{F^{f}_{v_f}\}$, its blocks and
its declared order, and on no product; the conditions on $\tau_A$, on every
$M_x$ and on the right-hand side $\Theta_k$ of a general entry are
decided block by block along $\mathbb{T}$, the trace of $\tau_A$
accumulated over the blocks in their declared order and rounding nowhere;
and in ensemble form the conditions on $\{M_x\}$ are replaced by those of
(DX4) on $\{N_x\}$, the denoted $\{M_x\}$ being a POVM by
Lemma~\ref{lem:ensemble-form} and decided by nothing.
\end{digestclaim}

% claim meat-core/def:admissible-declaration
\digestlabel{def:admissible-declaration}{1.2}\pdfbookmark[2]{Definition 1.2 Admissible declaration}{digest-meat-core/def:admissible-declaration}
\renewcommand*{\RestatedName}{\hyperref{contract.pdf}{}{claim-meat-core/def:admissible-declaration}{Definition~1.2}}%
\begin{digestclaim}[Admissible declaration]
The tuple $\Decl$ of Definition~\ref{def:declaration} is an
\emph{admissible declaration}
when it satisfies: (i) all systems finite-dimensional, all
alphabets finite, all maps total; (ii) the marginal slot holds
$\tau_A\in\Seq{A}$ (trusted-source mode) or the value $\nomarg$
(untrusted-source mode); in particular $\tr\tau_A=1$ whenever a
marginal is declared, so $\tau_A=0$ is inadmissible;
(iii) $\{M_x\}$ and $\{F_v\}$ are POVMs and every $F_v$ and every
$\hat O_k$ is sector-block-diagonal; (iv) $r$ maps into
$\Zval\sqcup\{\perp\}$; (v) sifting is announcement-determined via
$\varsigma$; (vi) $\bowtie_k\in\{\preceq,=,\succeq\}$ and each entry
carries one of the two declared forms of Definition~\ref{def:homog}, so
that $\Theta_k$ is Hermitian on $A$, in trusted-source mode, respectively
$\bowtie_k\in\{\le,=,\ge\}$ and $\Theta_k=c_k\in\mathbb R$ in
untrusted-source mode; (vii) $\Feas\neq\varnothing$; (viii)
\emph{party computability:} the declaration names a party partition
$\Party$ (Definition~\ref{def:declaration}) and
\begin{itemize}[leftmargin=1.6em]
\item each party announces a function of its own readout: $a_A$ is
defined on $\Oset{\mathsf A}$ and $a_B$ on $\Oset{\mathsf B}$, with
$a(x,v)=(a_A(\rd{\mathsf A}(x,v)),a_B(\rd{\mathsf B}(x,v)))$;
\item the declared key holder $\hold\in\{\mathsf A,\mathsf B\}$
computes the raw key from \emph{its own readout} and the announcements,
$r(x,v)=r_{\hold}(\rd{\hold}(x,v),a(x,v))$ with $r_{\hold}$ a declared
total function, Eq.~\eqref{eq:partycomp}.
\end{itemize}
Under the default partition ($\rd{\mathsf A}(x,v)=x$,
$\rd{\mathsf B}(x,v)=v$) these conditions read:
$a_A:\Xset\to\CAset$, $a_B:\Vset\to\CBset$, and
$r(x,v)=r_{\mathsf A}(x,a(x,v))$ resp.\ $r_{\mathsf B}(v,a(x,v))$.

\emph{In the general format} of Definition~\ref{def:declaration}\,(DX1)--(DX6)
the same eight conditions read as follows, and at one Alice sector, one
measured factor, POVM form and two messages they are the
conditions above. (iii): every $M_x$, the declared $\tau_A$ and the
right-hand side of every general entry are block diagonal along
Alice's sectors $\mathbb{T}$, \emph{by their type as
block lists against a declared factor tagged $A$}, and
$\{M_x\}$ is a POVM; in ensemble form the declared $\{N_x\}$ satisfy
\eqref{eq:ensembleform} and $\{M_x\}$ is the family \eqref{eq:ensemblepovm},
a POVM block diagonal along $\mathbb{T}$ by
Lemma~\ref{lem:ensemble-form}\,(i); $\{F_v\}$ is the product family
\eqref{eq:productfamily} of per-factor POVMs, each block diagonal for its
factor's sector map, and every $\hat O_k$ is block diagonal for the
sector map of its scope. (vi): an entry may in addition carry the
party-local form of Definition~\ref{def:homog}, in either mode, with
$\bowtie_k\in\{\preceq,=,\succeq\}$ and $c_k\in\mathbb{R}$. (viii): each
message of the transcript is a function of its sender's own readout and of
the earlier messages, $a$ is the composition \eqref{eq:transcript}, and
\eqref{eq:partycomp} holds with that $a$.
\end{digestclaim}

% claim meat-core/lem:ensemble-form
\digestlabel{lem:ensemble-form}{1.3}\pdfbookmark[2]{Lemma 1.3 The measurement an ensemble-form declaration denotes}{digest-meat-core/lem:ensemble-form}
\renewcommand*{\RestatedName}{\hyperref{contract.pdf}{}{claim-meat-core/lem:ensemble-form}{Lemma~1.3}}%
\begin{digestclaim}[The measurement an ensemble-form declaration denotes]
Let $\Hs{A}=\bigoplus_{t\in\mathbb{T}}\Hs{A_t}$, and for every $t$ let
$\tau^{(t)}_A\in\Pos{A_t}$ be non-zero and $N^{(t)}_x\in\Pos{A_t}$,
$x\in\Xset$, with $\sum_{x}N^{(t)}_x=\tau^{(t)}_A$. Let $P_t$ be the
projector onto the support of $\tau^{(t)}_A$, let
$(\tau^{(t)}_A)^{-1/2}$ be the inverse square root on that support, zero off
it, and let $M^{(t)}_x$ be the operator of \eqref{eq:ensemblepovm},
$M_x\coloneqq\bigoplus_tM^{(t)}_x$. Then:
\begin{enumerate}[label=(\roman*),leftmargin=2.2em]
\item $M^{(t)}_x\succeq0$ and $\sum_xM^{(t)}_x=\idon{A_t}$ for every $t$, so
$\{M_x\}_{x\in\Xset}$ is a POVM on $A$ and every element is block diagonal
along $\mathbb{T}$;
\item $(\tau^{(t)}_A)^{1/2}M^{(t)}_x(\tau^{(t)}_A)^{1/2}=N^{(t)}_x$, and in
particular $\tr[M_x\tau_A]=\sum_t\tr N^{(t)}_x$;
\item $M^{(t)}_x=0$ if and only if $N^{(t)}_x=0$;
\item conversely, if $\tau_A\in\Pos{A}$ and $\{M_x\}$ is a POVM on $A$, both
block diagonal along $\mathbb{T}$, then
$N_x\coloneqq\tau_A^{1/2}M_x\tau_A^{1/2}$ is positive semidefinite and block
diagonal along $\mathbb{T}$ and $\sum_xN_x=\tau_A$.
\end{enumerate}
\end{digestclaim}

% claim meat-core/def:maps
\digestlabel{def:maps}{1.4}\pdfbookmark[2]{Definition 1.4 Generated registers, isometries and maps}{digest-meat-core/def:maps}
\renewcommand*{\RestatedName}{\hyperref{contract.pdf}{}{claim-meat-core/def:maps}{Definition~1.4}}%
\begin{digestclaim}[Generated registers, isometries and maps]
The classical registers are $X$ (alphabet $\Xset$), $Y$ (alphabet $\Vset$),
the announcement register $C_AC_B$ and $I$ (alphabet $\Kann$ each), $\Chat$
(alphabet $\Chatset$), $R$ and $Z$ (alphabet $\Zset$). The announcement
register is one register, named after the two messages of the basic format;
in the general format it holds the whole tuple $a(x,v)$ of messages. Define
\begin{align*}
  \Vmsm &\coloneqq \sum_{x\in\Xset}\sum_{v\in\Vset}
    \ket{x,v}_{XY}\otimes\big(\sqrt{M_x}\otimes\sqrt{F_v}\big),
    &&\Hs{A\Qsys}\to\Hs{XY\,A\Qsys},\tag{\getrefnumber{C-eq:Vmsm}}\digesteqlabel{eq:Vmsm}\\
  \Vann &\coloneqq \sum_{x,v}\ket{a(x,v)}_{C_AC_B}\otimes\dyad{x,v}_{XY},
    &&\text{controlled on }XY,\tag{\getrefnumber{C-eq:Vann}}\digesteqlabel{eq:Vann}\\
  \VI &\coloneqq \sum_{\kappa\in\Kann}\ket{\kappa}_{I}\otimes
    \dyad{\kappa}_{C_AC_B},
    &&\text{controlled on }C_AC_B,\tag{\getrefnumber{C-eq:VI}}\digesteqlabel{eq:VI}\\
  \VChat &\coloneqq \sum_{\kappa\in\Kann}\ket{\phi(\kappa)}_{\Chat}\otimes
    \dyad{\kappa}_{C_AC_B},
    &&\text{controlled on }C_AC_B,\tag{\getrefnumber{C-eq:VChat}}\digesteqlabel{eq:VChat}\\
  \VR &\coloneqq \sum_{x,v}\ket{r(x,v)}_{R}\otimes\dyad{x,v}_{XY},
    &&\text{controlled on }XY,\tag{\getrefnumber{C-eq:VR}}\digesteqlabel{eq:VR}\\
  \VZ &\coloneqq \sum_{j\in\Zset}\ket{j}_{Z}\otimes Z_j,
    \qquad Z_j\coloneqq\dyad{j}_{R},
    &&\text{controlled on }R,\tag{\getrefnumber{C-eq:VZ}}\digesteqlabel{eq:VZ}
\end{align*}
each acting as the identity on all registers not named. The full protocol
isometry and the
\emph{single-round protocol map}
are
\begin{equation*}\tag{\getrefnumber{C-eq:Mprot}}\digesteqlabel{eq:Mprot}
\begin{gathered}
  \Wprot\coloneqq\VR\,\VChat\,\VI\,\Vann\,\Vmsm,\\
  \Mprot(\omega)\coloneqq
  \tr_{A\Qsys\,XY\,C_AC_B\,R}\!\left[\VZ\,\Wprot\,\omega\,\Wprot^{\dagger}
  \VZ^{\dagger}\right],
  \qquad \Mprot: A\Qsys\to Z\Chat I .
\end{gathered}
\end{equation*}
The
\emph{key pinching}
is $\cZ(\tau)\coloneqq\sum_{j\in\Zset}Z_j\,\tau\,Z_j$ on $R$, the pinching
of \xref{foundations/def:pinching} by the projectors $Z_j$ (extended by the
identity to any ambient system), and the
\emph{statistics map}
is
\begin{equation*}\tag{\getrefnumber{C-eq:Phimap}}\digesteqlabel{eq:Phimap}
  \Phimap[\omega]\coloneqq\sum_{\hat c\in\Chatset}
  \Bigg(\sum_{(x,v):\,\phi(a(x,v))=\hat c}
  \tr\!\big[(M_x\otimes F_v)\,\omega\big]\Bigg)\,\hat e_{\hat c},
  \qquad \Phimap:\Lin{A\Qsys}\to\mathbb{R}^{\Chatset},
\end{equation*}
with $\hat e_{\hat c}$ the unit vector at position $\hat c$.
\end{digestclaim}

% claim meat-core/thm:protocol-channel
\digestlabel{thm:protocol-channel}{1.5}\pdfbookmark[2]{Theorem 1.5 The declaration generates a protocol channel}{digest-meat-core/thm:protocol-channel}
\renewcommand*{\RestatedName}{\hyperref{contract.pdf}{}{claim-meat-core/thm:protocol-channel}{Theorem~1.5}}%
\begin{digestclaim}[The declaration generates a protocol channel]
Under Definition~\ref{def:admissible-declaration}:
\begin{enumerate}[label=(\roman*),leftmargin=2.2em]
\item Each of \eqref{eq:Vmsm}--\eqref{eq:VZ}, and hence $\Wprot$, is an
isometry, and $\Mprot$ is a channel $A\Qsys\to Z\Chat I$ with $\Chat$, $I$,
$Z$ classical.
\item For every system $E$ and every $\omega\in\Sleq{A\Qsys E}$,
\begin{equation*}\tag{\getrefnumber{C-eq:povm-form}}\digesteqlabel{eq:povm-form}
\begin{split}
  (\Mprot\otimes\idmap_E)(\omega)
  ={}&\sum_{x,v}
  \dyad{r(x,v)}_{Z}\otimes\dyad{\phi(a(x,v))}_{\Chat}\otimes{}\\
  &\qquad\dyad{a(x,v)}_{I}\otimes
  \tr_{A\Qsys}\!\big[(M_x\otimes F_v\otimes\idon{E})\,\omega\big].
\end{split}
\end{equation*}
\item The $\Chat$-marginal of $(\Mprot\otimes\idmap_E)(\omega)$ is the
probability vector $\Phimap[\omega_{A\Qsys}]$; $\Phimap$ is affine, and
$\Phimap[\omega]$ is a probability distribution on $\Chatset$ for every
$\omega\in\Seq{A\Qsys}$.
\end{enumerate}
\end{digestclaim}

% claim meat-core/def:branch
\digestlabel{def:branch}{1.6}\pdfbookmark[2]{Definition 1.6 Branch maps}{digest-meat-core/def:branch}
\renewcommand*{\RestatedName}{\hyperref{contract.pdf}{}{claim-meat-core/def:branch}{Definition~1.6}}%
\begin{digestclaim}[Branch maps]
For $\hat c\in\Chatset$, a class $\kappa\in\phi^{-1}(\hat c)$ with
$\varsigma(\kappa)=\mathrm{pass}$ is a
\emph{pass class}
of $\hat c$, and one with $\varsigma(\kappa)=\mathrm{fail}$ a
\emph{fail class};
let
$\Kann_{\mathrm{pass}}(\hat c)\coloneqq
\{\kappa\in\phi^{-1}(\hat c):\varsigma(\kappa)=\mathrm{pass}\}$ and
$\Kann_{\mathrm{fail}}(\hat c)\coloneqq
\{\kappa\in\phi^{-1}(\hat c):\varsigma(\kappa)=\mathrm{fail}\}$. For each
class $\kappa$ define the operator
\begin{equation*}\tag{\getrefnumber{C-eq:kraus}}\digesteqlabel{eq:kraus}
\begin{split}
  K_{\kappa}\coloneqq{}&\sum_{(x,v)\in a^{-1}(\kappa)}
  \ket{r(x,v)}_{R}\otimes\ket{x,v}_{XY}\otimes\ket{\kappa}_{C_AC_B}
  \otimes\big(\sqrt{M_x}\otimes\sqrt{F_v}\big)\\
  &\;:\;\Hs{A\Qsys}\to\Hs{R\,XY\,C_AC_B\,A\Qsys},
\end{split}
\end{equation*}
and the CP
\emph{branch map}s,
the pass map $\cG_{\hat c}$ and the fail map $\Theta_{\hat c}$,
\begin{equation*}\tag{\getrefnumber{C-eq:GTheta}}\digesteqlabel{eq:GTheta}
  \cG_{\hat c}(\omega)\coloneqq
  \sum_{\kappa\in\Kann_{\mathrm{pass}}(\hat c)}K_{\kappa}\,\omega\,
  K_{\kappa}^{\dagger},
  \qquad
  \Theta_{\hat c}(\omega)\coloneqq
  \sum_{\kappa\in\Kann_{\mathrm{fail}}(\hat c)}K_{\kappa}\,\omega\,
  K_{\kappa}^{\dagger}.
\end{equation*}
For $s\in\SecSet$, the sector versions $K_{\kappa}^{(s)}$, $\cG_{\hat c,s}$,
$\Theta_{\hat c,s}$ are \eqref{eq:kraus}--\eqref{eq:GTheta} with
$\sqrt{F_v}$ replaced by $\sqrt{F_v^{(s)}}$, acting
$\Hs{A\Qsys_s}\to\Hs{R\,XY\,C_AC_B\,A\Qsys_s}$.
\end{digestclaim}

% claim meat-core/lem:branch-maps
\digestlabel{lem:branch-maps}{1.7}\pdfbookmark[2]{Lemma 1.7 Properties of the pass and fail branch maps}{digest-meat-core/lem:branch-maps}
\renewcommand*{\RestatedName}{\hyperref{contract.pdf}{}{claim-meat-core/lem:branch-maps}{Lemma~1.7}}%
\begin{digestclaim}[Properties of the pass and fail branch maps]
Under Definition~\ref{def:admissible-declaration}, for every $\hat c\in\Chatset$:
\begin{enumerate}[label=(\roman*),leftmargin=2.2em]
\item $\cG_{\hat c}$ and $\Theta_{\hat c}$ are CP and trace-non-increasing,
with
\begin{equation*}\tag{\getrefnumber{C-eq:traces}}\digesteqlabel{eq:traces}
  \tr\cG_{\hat c}(\omega)+\tr\Theta_{\hat c}(\omega)
  =\Phimap[\omega](\hat c)
  \qquad\text{for all }\omega\in\Pos{A\Qsys}.
\end{equation*}
\item $\cG_{\hat c}(\omega)$ is supported on
$\operatorname{span}\{\ket{j}_R:j\in\Zval\}$ and
$\Theta_{\hat c}(\omega)=\dyad{\perp}_{R}\otimes\theta_{\hat c}(\omega)$ for
a CP map $\theta_{\hat c}$; in particular their supports are orthogonal and
$\cZ(\Theta_{\hat c}(\omega)^{p})=\Theta_{\hat c}(\omega)^{p}$ for every
$p>0$.
\item (Sector decomposition.) If
$\omega=\bigoplus_{s\in\SecSet}\omega_s$ is block-diagonal with respect to
the decomposition $\Hs{A\Qsys}=\bigoplus_s\Hs{A\Qsys_s}$, then
\begin{equation*}\tag{\getrefnumber{C-eq:G-sector}}\digesteqlabel{eq:G-sector}
  \cG_{\hat c}(\omega)=\bigoplus_{s\in\SecSet}\cG_{\hat c,s}(\omega_s),
  \qquad
  \Theta_{\hat c}(\omega)=\bigoplus_{s\in\SecSet}\Theta_{\hat c,s}(\omega_s),
  \qquad
  \Phimap[\omega]=\sum_{s\in\SecSet}\Phimap_s[\omega_s],
\end{equation*}
where the direct sums refer to the (mutually orthogonal) subspaces
$\Hs{R\,XY\,C_AC_B}\otimes\Hs{A\Qsys_s}$ and
$\Phimap_s$ is \eqref{eq:Phimap} with $F_v^{(s)}$ in place of $F_v$.
\item If $\hat c\notin\Chkey$ then $\cG_{\hat c}=0$.
\end{enumerate}
The complete positivity of $\cG_{\hat c}$ and $\Theta_{\hat c}$ in (i),
and all of (iii), hold as well when $\{F_v\}$ is any family of
sector-block-diagonal positive semidefinite operators, complete or not,
the other components satisfying Definition~\ref{def:admissible-declaration}.
\end{digestclaim}

\section*{2\quad Accumulation}\pdfbookmark[1]{2 Accumulation}{digest-group-meat-core/fact:meat}

% claim meat-core/fact:meat
\digestlabel{fact:meat}{2.1}\pdfbookmark[2]{Fact 2.1 Marginal-constrained entropy accumulation, sequential form}{digest-meat-core/fact:meat}
\renewcommand*{\RestatedName}{\hyperref{contract.pdf}{}{claim-meat-core/fact:meat}{Fact~2.1}}%
\begin{digestclaim}[Marginal-constrained entropy accumulation, sequential form]
\emph{(\cite[Thm.~4.2a, Eqs.~(100)--(101)]{arqand2025meat}.)}
For each $j\in\{1,\dots,n\}$ take a state
$\sigma^{(j-1)}\in\Seq{A_{j-1}}$ and a channel
$\Mprot_j\in\mathrm{CPTP}(A_{j-1}E_{j-1},\,S_j\Chat_jE_j)$ with
$\Chat_j$ classical. Let
\begin{equation*}
  \rho_{S_1^n\Chat_1^nE_n}
  =\big(\Mprot_n\circ\dots\circ\Mprot_1\big)\big[\omega^{0}\big],
  \qquad
  \omega^{0}\in\Seq{A_0^{n-1}E_0},
\end{equation*}
subject to the product marginal constraint
$\omega^{0}_{A_0^{n-1}}=\sigma^{(0)}\otimes\dots\otimes\sigma^{(n-1)}$;
$E_0$ is unconstrained, and $\omega^{0}$ ranges over \emph{all}
normalized states compatible with the marginal constraint (purity is
not assumed). Suppose
$\rho=p_{\Omega}\,\rho_{|\Omega}+(1-p_{\Omega})\rho_{|\bar\Omega}$
with $p_{\Omega}\in(0,1]$ and $\rho_{|\Omega},\rho_{|\bar\Omega}$
normalized, all $\Chat_j$ isomorphic to one register $\Chat$ with
alphabet $\Chatset$, and let $S_{\Omega}$ be a convex set of
probability distributions on $\Chatset$ such that
$\operatorname{freq}_{\hat c_1^n}\in S_{\Omega}$ for every frequency
vector with nonzero probability in $\rho_{|\Omega}$. Then, for every
$\alpha\in(1,\infty)$ (locator note (vi): the source also covers the
endpoint $\alpha=\infty$, which no claim of this module uses),
\begin{equation*}\tag{\getrefnumber{C-eq:meat}}\digesteqlabel{eq:meat}
  \Hsandup\big(S_1^n\big|\Chat_1^nE_n\big)_{\rho_{|\Omega}}
  \;\ge\;n\,\hup-\frac{\alpha}{\alpha-1}\log_2\frac{1}{p_{\Omega}},
\end{equation*}
where
\begin{equation*}\tag{\getrefnumber{C-eq:meat-rate}}\digesteqlabel{eq:meat-rate}
  \hup=\inf_{\mathbf q\in S_{\Omega}}\ \inf_{\nu\in\Sigma}
  \left(\frac{\alpha}{\alpha-1}\,
  \sfD\!\left(\mathbf q\,\middle\|\,\boldsymbol{\nu}_{\Chat}\right)
  +\sum_{\hat c\in\operatorname{supp}(\boldsymbol{\nu}_{\Chat})}
  q(\hat c)\,\Hsandup(S|E\Etil)_{\nu_{|\hat c}}\right),
\end{equation*}
with $\sfD$ the classical (base-2) KL divergence of
\xref{foundations/def:kl} and $\Sigma=\Sigma_{S\Chat E\Etil}$ the
\emph{marginal-constrained convex range}
(\cite[Def.~4.2, Eq.~(98)]{arqand2025meat}) of the channels
$\Mprot_j\otimes\idmap_{\Etil}$ and states $\sigma^{(j-1)}$:
\begin{equation*}\tag{\getrefnumber{C-eq:meat-range}}\digesteqlabel{eq:meat-range}
  \Sigma\coloneqq\operatorname{conv}\!\Big(\bigcup_{j=1}^{n}\Sigma_j\Big),
  \qquad
  \Sigma_j\coloneqq\big\{(\Mprot_j\otimes\idmap_{\Etil})[\omega]:
  \omega\in\Seq{A_{j-1}E_{j-1}\Etil},\
  \omega_{A_{j-1}}=\sigma^{(j-1)}\big\},
\end{equation*}
where every output is embedded in common registers $S\Chat E\Etil$ (the
results hold for any choice of embeddings;
\cite[footnote~12]{arqand2025meat}) and $\Etil$ is a register of large
enough dimension to serve as a purifying register for any of the
$A_{j-1}E_{j-1}$ registers (enlargeable without loss of generality;
\cite[footnote~17]{arqand2025meat}).

\emph{Locator notes.} (i) The statement above is the source's
Thm.~4.2a with its registers written in this module's fonts; its
Eq.~(100) is \eqref{eq:meat}, its Eq.~(101) is \eqref{eq:meat-rate},
and its Def.~4.2, Eq.~(98) is \eqref{eq:meat-range}. Thm.~4.2a is the
special case of the source's Thm.~4.2b in which the per-round secret
test registers (the unbarred $C_j$ of the source) are trivial, which is
the case Theorem~\ref{thm:protocol-channel} generates. The source states
its entropies in base~2
(bits). (ii) The divergence in the penalty term is the \emph{KL}
divergence $\sfD$: the source obtains the rate objective as the convex
conjugate of a log-sum-exponential
(\cite[Lemma~4.11, Eqs.~(113)--(114)]{arqand2025meat}) and evaluates
it in its Eq.~(118), stating there explicitly that the first term ``is
the KL divergence''.
(iii) An alternative import is the tensor-product
specialization \cite[Cor.~4.2, Eq.~(137)]{arqand2025meat}, in which
the round channels carry no $E_j$ output and Eve's entire action is
the choice of one global input state. The source presents that
corollary as a register substitution of the theorems above, and its
Sec.~5 applies the resulting analysis when every public announcement
is made only after \emph{all} measurements are complete; for per-round
(``on-the-fly'') announcements --- the model
Theorem~\ref{thm:protocol-channel} generates, whose round channel
announces $I_j$ in every round --- the source's own
route (its Sec.~5, Protocols~1--2 and the reduction between them) is
the sequential form quoted here, in which Eve carries a quantum memory
$E_{j-1}$ between rounds and may adapt each round's action to all
earlier announcements. The present import therefore covers a strictly
larger attack model than that specialization would. (iv) The theorem has no
announcement register of its own: announcements enter the sequential
model inside the $E_j$ outputs (the source's Sec.~5: Eve
``immediately incorporates'' each $I_j$ into her side information).
(v) Indexing is off by one between this module's rounds and the
source's registers (\cite[Rem.~5.1]{arqand2025meat}): round $j$ here
uses the source's $A_{j-1}E_{j-1}$. (vi) The source states the
theorem for $\alpha\in(1,\infty]$, and its Def.~2.6 defines the
entropies at the endpoint; \xref{foundations/def:renyi}
defines the sandwiched family on $(1,\infty)$ only, so the Fact is
quoted, and Theorems \ref{thm:accumulated-entropy} and \ref{thm:constrained-accumulation} are stated, on
$(1,\infty)$. No claim uses the endpoint: every key-length statement
is confined to $(1,2)$.
(vii) \emph{Trivial marginal registers are admissible} --- the import
check for untrusted-source mode, read off the source at the
locators named. The hypotheses that
\cite[Thm.~4.2a]{arqand2025meat} places on the marginal registers are,
in full, that one ``take a state $\sigma^{(j-1)}\in S_=(A_{j-1})$, and
a channel $M_j\in\mathrm{CPTP}(A_{j-1}E_{j-1},S_j\widehat C_jE_j)$,
such that $\widehat C_j$ are classical'', together with the product
condition
$\omega_{A_0^{n-1}}=\sigma^{(0)}_{A_0}\otimes\dots\otimes
\sigma^{(n-1)}_{A_{n-1}}$ on the input state; and
\cite[Def.~4.2, Eq.~(98)]{arqand2025meat} likewise only ``consider[s]
channels $E_j\in\mathrm{CPTP}(A_{j-1}\bar A_{j-1},Z_j)$ and states
$\sigma^{(j-1)}\in S_=(A_{j-1})$'' (in the source's symbols, where $E_j$ is a
channel and $Z_j$ its output register, not this module's register $E_j$ and
raw key $Z_j$). \emph{No dimension bound and no
rank condition is imposed on $A_{j-1}$ at either locator}, and the
source's standing convention is explicit that none is imposed globally
either: ``we will assume that all systems are finite-dimensional, but
we will not impose any bounds on the system dimensions unless otherwise
specified'' (\cite[Sec.~2, the convention stated below its Table~1] 
{arqand2025meat}). A one-dimensional $A_{j-1}$ is therefore an
admissible instantiation, and at it the product marginal condition is
satisfied by every normalized $\omega$, since $\Seq{A_0^{n-1}}$ is then
a singleton --- the hypothesis discharges vacuously rather than
failing. Independently, the source reads the same equivalence into its
own definition of channel conditional entropy: see the sentence
following \cite[Eq.~(1)]{arqand2025meat}. Nothing else in Fact~\ref{fact:meat}
mentions $A_{j-1}$ beyond the purifier requirement, which asks $\Etil$
to purify $A_{j-1}E_{j-1}$ and is only weakened by shrinking
$A_{j-1}$.
\end{digestclaim}

% claim meat-core/def:homog
\digestlabel{def:homog}{2.2}\pdfbookmark[2]{Definition 2.2 Licensed constraint list and its two classes}{digest-meat-core/def:homog}
\renewcommand*{\RestatedName}{\hyperref{contract.pdf}{}{claim-meat-core/def:homog}{Definition~2.2}}%
\begin{digestclaim}[Licensed constraint list and its two classes]
The declared forms of an entry of $\CList$ --- homogeneous, carrying
a real scalar $c_k$, general, carrying a Hermitian $\Theta_k$ on $A$,
and, in the general format, party-local --- are those of
\xref{protocol-compiler/def:homog}, and the list is \emph{licensed} when it
is licensed for the declared source mode there. Its two classes
for an untagged entry are the following; in the general format an untagged
entry is read for its lifted operator, and the cited class also
licenses a party-local entry in either mode (the paragraph on party-local
entries).
\begin{enumerate}[label=(\alph*),leftmargin=2.2em]
\item \emph{trusted-source mode:} every entry of $\CList$ is declared
in the homogeneous form, carrying $c_k\in\mathbb{R}$; the right-hand side it
denotes is then \emph{constructed} as $\Theta_k\coloneqq c_k\,\tau_A$ from the
declared $c_k$ and the declared $\tau_A$, and such a list is a
\emph{$\tau_A$-homogeneous list}
by construction;
\item \emph{untrusted-source mode:} $\CList$ is a
\emph{scalar list},
i.e.\ every right-hand side is a real number $\Theta_k=c_k$ and the entry
reads $\tr[(\idon{A}\otimes\hat O_k)\,\omega]\bowtie_k c_k$ --- a
comparison of real numbers, with $\bowtie_k\in\{\le,=,\ge\}$.
\end{enumerate}
Licensing is therefore a \emph{type} check in both modes and compares no
numbers. A general entry remains admissible input, every structural and
numerical claim of this module being proved for it, but it is never licensed,
so a list containing one carries no key-length claim, which is the safe
direction. The two classes agree at $d_A=1$
(Proposition~\ref{prop:trivial-alice-system}\,(iii)), and (b) is what the
licensing condition \eqref{eq:attset-un} transports in untrusted-source
mode, through \eqref{C-eq:transport-un}:
\emph{the operator form $\Theta_k=c_k\idon{A}$ is not licensed}.

\emph{Party-local entries.} In the general format an entry may be declared
in the party-local form (\xref{protocol-compiler/def:homog}); it is
written $(\hat O_k,\bowtie_k,c_k;f)$, and its state-level reading, in either
source mode, is
\begin{equation*}\tag{\getrefnumber{C-eq:partylocal-entry}}\digesteqlabel{eq:partylocal-entry}
  \tr_{\Qsys_f}\!\Big[\big(\idon{}\otimes
  (\hat O_k-c_k\,\idon{\Qsys_f})\big)\,\omega\Big]\;\bowtie_k\;0
  \qquad\text{as an operator on the systems of }\omega
  \text{ other than }\Qsys_f ,
\end{equation*}
those systems being $A$ and the measured factors
$\Qsys_{f'}$, $f'\ne f$.
\end{digestclaim}

% claim meat-core/def:roundmodel
\digestlabel{def:roundmodel}{2.3}\pdfbookmark[2]{Definition 2.3 Round model and licensed attack sets}{digest-meat-core/def:roundmodel}
\renewcommand*{\RestatedName}{\hyperref{contract.pdf}{}{claim-meat-core/def:roundmodel}{Definition~2.3}}%
\begin{digestclaim}[Round model and licensed attack sets]
Fix a declaration $\Decl$ with single-round protocol channel
$\Mprot\in\mathrm{CPTP}(A\Qsys\to Z\Chat I)$ (Theorem~\ref{thm:protocol-channel}).
The round model of this contract is the round model of
\xref{foundations/def:round-model} with per-round system $\Qsys$, Alice's
system $A$ and the declared source mode, the round index of $\Qsys$ and $A$
dropped. An attack sequence is thus a family $(\Aev_j)_{j=1}^{n}$ of
channels whose output carries the systems the adversary supplies in the
declared source mode, the display being that definition's with its
per-round system $S$ taken to be $\Qsys$:
\begin{equation*}\tag{\getrefnumber{C-eq:attackchannel}}\digesteqlabel{eq:attackchannel}
  \Aev_j\in
  \begin{cases}
    \mathrm{CPTP}(E_{j-1},\ \Qsys\,E'_j), &
      \text{trusted-source mode},\\[0.3ex]
    \mathrm{CPTP}(E_{j-1},\ A\,\Qsys\,E'_j), &
      \text{untrusted-source mode},
  \end{cases}
\end{equation*}
where $E_0$ is Eve's initial register and $E'_j$ her outgoing memory. In
trusted-source mode $A$ is \emph{not} an output of $\Aev_j$: it is
supplied by Alice, inaccessible to the adversary, and carried into the
round channel as a separate input --- that, and nothing else, is what
A-SRC (\xref{transmitter/ass:src}) licenses. The generated
\emph{round channel}s
are
\begin{equation*}\tag{\getrefnumber{C-eq:roundchannel}}\digesteqlabel{eq:roundchannel}
\begin{gathered}
  \Mprot_j\coloneqq
  \begin{cases}
    (\Mprot\otimes\idmap_{E'_j})\circ(\idmap_{A}\otimes\Aev_j)
      \;\in\;\mathrm{CPTP}(A\,E_{j-1},\;Z\,\Chat\,E_j), &
      \text{trusted mode},\\[0.5ex]
    (\Mprot\otimes\idmap_{E'_j})\circ\Aev_j
      \;\in\;\mathrm{CPTP}(E_{j-1},\;Z\,\Chat\,E_j), &
      \text{untrusted mode},
  \end{cases}\\
  E_j\coloneqq E'_j\,I_j:
\end{gathered}
\end{equation*}
the attack delivers the round input on $\Qsys$ --- and, in
untrusted-source mode, on $A$ as well --- together with Eve's
new memory, the protocol channel produces $(Z_j,\Chat_j,I_j)$, and the
announcement $I_j$ --- classical and public --- is placed in Eve's
register. Content of $E_{j-1}$ that Eve wishes to retain (earlier
announcements included) she carries inside $E'_j$; every physical
attack is represented by a sequence that retains \emph{all}
announcements, and Theorem~\ref{thm:accumulated-entropy} uses that representation, so
that $E_n=I_1^n\Ehat$ with $\Ehat$ Eve's retained quantum memory. The
untrusted-mode round channel has no $A$ input at all, which is what
makes it an instance of Fact~\ref{fact:meat} with the marginal
register $A_{j-1}$ trivial (Fact~\ref{fact:meat}, locator note (vii);
Theorem~\ref{thm:accumulated-entropy}).

For a licensed list $\CList$ (Definition~\ref{def:homog}) and
finite-dimensional registers $E',E''$, the
\emph{licensed attack set}
is the attack set of $\CList$ of \xref{foundations/def:round-model}
at $S=\Qsys$; its displays at that system are, in trusted-source mode,
\begin{equation*}\tag{\getrefnumber{C-eq:attset}}\digesteqlabel{eq:attset}
  \AttSet_{\CList}(E',\,\Qsys E'')\coloneqq
  \Big\{\Aev\in\mathrm{CPTP}(E',\Qsys E''):\
  \Aev^{\dagger}\big(\hat O_k\otimes\idon{E''}\big)
  \bowtie_k c_k\,\idon{E'}\ \ \forall k\Big\},
\end{equation*}
and, with the constraint operator lifted to the larger supplied system,
\begin{equation*}\tag{\getrefnumber{C-eq:attset-un}}\digesteqlabel{eq:attset-un}
  \AttSet_{\CList}(E',\,A\Qsys E'')\coloneqq
  \Big\{\Aev\in\mathrm{CPTP}(E',A\Qsys E''):\
  \Aev^{\dagger}\big(\idon{A}\otimes\hat O_k\otimes\idon{E''}\big)
  \bowtie_k c_k\,\idon{E'}\ \ \forall k\Big\}
\end{equation*}
in untrusted-source mode; the two are distinguished by the output
system named in the argument. An attack sequence is
$\CList$-admissible (an admissible attack sequence for $\CList$) if
$\Aev_j$ lies in the set of
\eqref{eq:attset}, resp.\ \eqref{eq:attset-un}, with $E'=E_{j-1}$ and
$E''=E'_j$, for every $j$. Where $\CList$ carries entries declared
party-local (Definition~\ref{def:homog}), each such entry $(\hat O_k,\bowtie_k,c_k;f)$
contributes to the licensed attack set, in place of the scalar condition of
\eqref{eq:attset} or \eqref{eq:attset-un}, the condition of an entry tagged
with the factor $f$ in \xref{foundations/def:round-model}\,(iii) at
$S=\Qsys$. With $\Qsys_{f^c}$ the product of the measured factors other than
$f$, and $R$ equal to $A$ in untrusted-source mode and trivial in
trusted-source mode, its display at that system is
\begin{equation*}\tag{\getrefnumber{C-eq:attset-pl}}\digesteqlabel{eq:attset-pl}
\begin{gathered}
  \big(\idmap_{R'}\otimes\Aev\big)^{\dagger}
  \Big(X\otimes\big(\hat O_k-c_k\,\idon{\Qsys_f}\big)\otimes\idon{E''}\Big)
  \;\bowtie_k\;0\\
  \text{for every finite-dimensional }R'\text{ and every }
  X\in\Pos{R'R\,\Qsys_{f^c}} ,
\end{gathered}
\end{equation*}
as an operator on $R'E'$, the tensor factors of the argument written in the
order $R'R\,\Qsys_{f^c}$, $\Qsys_f$, $E''$; untagged entries, those not
declared party-local, keep their conditions, read for their
lifted operators
(\xref{protocol-compiler/def:declaration}\,(DX6)). The reference system $R'$ is what
carries Alice's retained system in trusted-source mode, where $A$ is no
output of $\Aev$. With $R'$ trivial and $X=\idon{}$ the condition is the
scalar one for the lifted operator, so the set with tagged entries is
contained in the set with the same entries untagged. For
$\CList=\varnothing$ the condition is empty:
$\AttSet_{\varnothing}=\mathrm{CPTP}$, every attack sequence is
admissible, and the model is unrestricted. The adjoint condition in
\eqref{eq:attset} is exactly the published squashed-attack-set format
of \cite[Eq.~(151)]{tupkary2026decoy}, and
restricting attention to a declared subset of channels is the standard
modelling step of \cite[Def.~8.3, Rem.~8.4]{tupkary2026decoy}. Whether
a given physical protocol reduces to this model with a given
$\AttSet_{\CList}$ is \emph{not} decided here: it is what
A-ATK (\xref{protocol-compiler/ass:atk}) certifies.
\end{digestclaim}

% claim meat-core/def:hqkd-list
\digestlabel{def:hqkd-list}{2.4}\pdfbookmark[2]{Definition 2.4 Constrained target set and single-round rate}{digest-meat-core/def:hqkd-list}
\renewcommand*{\RestatedName}{\hyperref{contract.pdf}{}{claim-meat-core/def:hqkd-list}{Definition~2.4}}%
\begin{digestclaim}[Constrained target set and single-round rate]
For a register $\Etil'$ and a (possibly empty) declared list $\CList$,
let
\begin{equation*}\tag{\getrefnumber{C-eq:sigL}}\digesteqlabel{eq:sigL}
  \Sigma^{\CList}(\Etil')\coloneqq
  \big\{(\Mprot\otimes\idmap_{\Etil'})(\sigma):\
  \sigma\in\Seq{A\Qsys\Etil'},\ \sigma\models\Decl,\
  \sigma_{A\Qsys}\models\CList\big\},
\end{equation*}
where $\sigma\models\Decl$ abbreviates the declaration's marginal
condition,
\begin{equation*}\tag{\getrefnumber{C-eq:margcond}}\digesteqlabel{eq:margcond}
  \sigma\models\Decl
  \quad:\Longleftrightarrow\quad
  \begin{cases}
    \sigma_A=\tau_A, & \text{trusted-source mode},\\
    \text{no condition beyond }\tr\sigma=1,
      & \text{untrusted-source mode},
  \end{cases}
\end{equation*}
and $\sigma_{A\Qsys}\models\CList$ abbreviates ``$\sigma_{A\Qsys}$
satisfies every entry of $\CList$'' in the form
Definition~\ref{def:homog} gives it for the declared mode, i.e.\
$\tr_{\Qsys}[(\idon{A}\otimes\hat O_k)\,\sigma_{A\Qsys}]
\bowtie_k\Theta_k$ in trusted-source mode and
$\tr[(\idon{A}\otimes\hat O_k)\,\sigma_{A\Qsys}]\bowtie_k c_k$ in
untrusted-source mode, for all $k$; an entry declared party-local
is read, in either mode, in the form \eqref{eq:partylocal-entry} of
Definition~\ref{def:homog}, an affine semidefinite condition on
$\sigma_{A\Qsys}$. $\Sigma^{\CList}(\Etil')$ is
convex in either mode: the state set is the intersection of the convex
$\Seq{A\Qsys\Etil'}$ with affine equalities and operator- or
scalar-inequality conditions, each convex, and
$\Mprot\otimes\idmap_{\Etil'}$ is affine.
Define the
\emph{constrained single-round rate}
\begin{equation*}\tag{\getrefnumber{C-eq:hqkd-list}}\digesteqlabel{eq:hqkd-list}
  \hqkd(\CList)\coloneqq\inf_{\mathbf q\in\acSet}\
  \inf_{\substack{\sigma\in\Seq{A\Qsys\Etil}\\
  \sigma\models\Decl,\ \sigma_{A\Qsys}\models\CList}}
  \left(\frac{\alpha}{\alpha-1}\,
  \sfD\!\left(\mathbf q\,\middle\|\,\Phimap[\sigma_{A\Qsys}]\right)
  +\sum_{\hat c\in\Chkey}
  q(\hat c)\,\Hsandup(Z|I\Etil)_{\nu^{\sigma}_{|\hat c}}\right),
\end{equation*}
with $\nu^{\sigma}\coloneqq(\Mprot\otimes\idmap_{\Etil})(\sigma)$ and
$\dim\Hs{\Etil}\ge\dim\Hs{A\Qsys}$; by
Lemma~\ref{lem:purifier-indep} the value does not depend on that
choice. For $\CList=\varnothing$ we write
$\hqkd\coloneqq\hqkd(\varnothing)$, the unconstrained single-round rate,
which Theorem~\ref{thm:accumulated-entropy} writes out as
\eqref{eq:hqkd-def}. Since a nonempty
$\CList$ shrinks the inner feasible set,
$\hqkd(\CList)\ge\hqkd(\varnothing)$ always.
\end{digestclaim}

% claim meat-core/def:admissible-acceptance-set
\digestlabel{def:admissible-acceptance-set}{2.5}\pdfbookmark[2]{Definition 2.5 Admissible acceptance set}{digest-meat-core/def:admissible-acceptance-set}
\renewcommand*{\RestatedName}{\hyperref{contract.pdf}{}{claim-meat-core/def:admissible-acceptance-set}{Definition~2.5}}%
\begin{digestclaim}[Admissible acceptance set]
An \emph{admissible acceptance set}
is a closed convex subset $\acSet$ of the probability simplex on $\Chatset$.
\end{digestclaim}

% claim meat-core/lem:purifier-indep
\digestlabel{lem:purifier-indep}{2.7}\pdfbookmark[2]{Lemma 2.7 Independence of the purifier size}{digest-meat-core/lem:purifier-indep}
\renewcommand*{\RestatedName}{\hyperref{contract.pdf}{}{claim-meat-core/lem:purifier-indep}{Lemma~2.7}}%
\begin{digestclaim}[Independence of the purifier size]
Fix $\CList$ (possibly empty) and write $h(\Etil')$ for the value of
\eqref{eq:hqkd-list} computed with a register $\Etil'$ of
$\dim\Hs{\Etil'}\ge\dim\Hs{A\Qsys}$ in place of $\Etil$. Then
$h(\Etil')$ is the same for every such $\Etil'$, and the inner infimum
is unchanged if restricted to pure $\sigma$.
\end{digestclaim}

% claim meat-core/lem:range
\digestlabel{lem:range}{2.8}\pdfbookmark[2]{Lemma 2.8 Constraint transport and range inclusion}{digest-meat-core/lem:range}
\renewcommand*{\RestatedName}{\hyperref{contract.pdf}{}{claim-meat-core/lem:range}{Lemma~2.8}}%
\begin{digestclaim}[Constraint transport and range inclusion]
Let $\CList$ be licensed (Definition~\ref{def:homog}) or empty, and let
$\Aev_j$ be $\CList$-admissible for the declared source mode
(Definition~\ref{def:roundmodel}).
\begin{enumerate}[label=(\roman*),leftmargin=2.2em]
\item \emph{(Transport.)} \emph{Trusted-source mode:} for every
$\omega\in\Seq{A\,E_{j-1}\,\Etil}$ with $\omega_A=\tau_A$, the state
$\varsigma\coloneqq(\idmap_A\otimes\Aev_j\otimes\idmap_{\Etil})(\omega)
\in\Seq{A\Qsys E'_j\Etil}$ satisfies $\varsigma_A=\tau_A$ and
$\varsigma_{A\Qsys}\models\CList$.
\emph{Untrusted-source mode:} for every
$\omega\in\Seq{E_{j-1}\,\Etil}$, the state
$\varsigma\coloneqq(\Aev_j\otimes\idmap_{\Etil})(\omega)
\in\Seq{A\Qsys E'_j\Etil}$ satisfies $\varsigma\models\Decl$ --- which
in that mode is only $\tr\varsigma=1$ --- and
$\varsigma_{A\Qsys}\models\CList$ in the scalar form of
Definition~\ref{def:homog}\,(b). \emph{In either mode}, every entry of
$\CList$ declared party-local holds of $\varsigma_{A\Qsys}$ in
the form \eqref{eq:partylocal-entry}, the attack being admissible for it in
the sense of \eqref{eq:attset-pl}.
\item \emph{(Range inclusion.)} With $\Mprot_j$ the round channel
\eqref{eq:roundchannel} and $\Sigma_j$ its output set
\eqref{eq:meat-range},
$\Sigma_j\subseteq\Sigma^{\CList}(E'_j\Etil)$ up to the regrouping of the
registers of the output, $(Z\Chat)(E'_jI_j)\Etil\mapsto(Z\Chat I)(E'_j\Etil)$;
hence, embedding the finitely many $E'_j\Etil$ into a common register
$\Etil^{*}$ with $\dim\Hs{\Etil^{*}}=\max_j\dim\Hs{E'_j\Etil}$,
\begin{equation*}
  \Sigma=\operatorname{conv}\Big(\bigcup_{j}\Sigma_j\Big)
  \;\subseteq\;\Sigma^{\CList}(\Etil^{*}).
\end{equation*}
\item \emph{(Antitonicity.)} If $\Sigma\subseteq\Sigma'$ then the
value of \eqref{eq:meat-rate} computed over $\Sigma$ is at least its
value computed over $\Sigma'$.
\end{enumerate}
\end{digestclaim}

% claim meat-core/lem:product-marginal
\digestlabel{lem:product-marginal}{2.9}\pdfbookmark[2]{Lemma 2.9 Product marginal of Alice's retained systems}{digest-meat-core/lem:product-marginal}
\renewcommand*{\RestatedName}{\hyperref{contract.pdf}{}{claim-meat-core/lem:product-marginal}{Lemma~2.9}}%
\begin{digestclaim}[Product marginal of Alice's retained systems]
Let the declaration be in trusted-source mode with $d_A\ge2$, and let
its $A$, $\tau_A$ and $\{M_x\}$ be those of a source specification
(\xref{transmitter/def:source-spec}) that carries a source-replacement
certificate $\Psi$ for $(A,\mathbb{T},\tau_A,\{M_x\})$ against its
virtual ensemble $\{p_x,\rvirt_x\}$ on $\Avirt$, in any of the forms of
\xref{transmitter/def:srcert}, an ensemble presentation being read as
the block-form certificate it defines; let its source segment satisfy
the source identity $\Xi(\rvirt_x)=\rho_x$ for every $x\in\Xset$, with $\Xi$
the composite channel of the segment and $\Xi=\idmap$ when the segment is
empty (\xref{transmitter/def:source-segment}\,(S3)). Let A-SRC
(\xref{transmitter/ass:src}), A-ENS (\xref{transmitter/ass:ens}), A-SEQ
(\xref{transmitter/ass:seq}), A-SEED (Assumption~\ref{ass:seed})
and A-LAB (Assumption~\ref{ass:lab}) at the party to which the
party partition assigns $A$ hold of the physical protocol, let the
adversary's initial register $E$ be finite-dimensional, and let $\sigma_E$
be the state in which it starts. Put
\begin{equation*}\tag{\getrefnumber{C-eq:product-input}}\digesteqlabel{eq:product-input}
  \omega^{0}\coloneqq\Big(\bigotimes_{j=1}^{n}\Psi_{A_j\Avirt_j}\Big)
  \otimes\sigma_E\ \in\ \Seq{A_1^nE_0},
  \qquad E_0\coloneqq(\Avirt)_1^n\,E .
\end{equation*}
Then:
\begin{enumerate}[label=(\roman*),leftmargin=2.2em]
\item \emph{(source replacement over the $n$ rounds)} against every attack,
the physical protocol and the protocol that starts from $\omega^{0}$, applies
$\Xi$ to $\Avirt_j$ in round $j$ in the place of the source's emission, and
applies $\pcNmsm$, Alice's measurement as the channel of
\xref{transmitter/thm:certificate-sufficiency}\,(iii), to $A_j$ at the point of
round $j$ at which the declared protocol's measurements are made, produce
the same state after every round, on every register that both then hold; in
the second protocol $A_j$ is an input of that measurement and of no other
operation, which is how the round model of Definition~\ref{def:roundmodel}
takes Alice's system in trusted-source mode;
\item $\omega^{0}_{A_1^n}=\tau_A^{\otimes n}$;
\item for every attack sequence of Definition~\ref{def:roundmodel}
and every $j$, the state on which the round channel $\Mprot_j$ acts, in a run
of the round model from $\omega^{0}$, has the marginal $\tau_A$ on $A_j$.
\end{enumerate}
At $d_A=1$, where $\tau_A=[1]$ and no certificate is declared, (ii) and (iii)
hold of every state on $A_1^nE_0$.
\end{digestclaim}

% claim meat-core/thm:accumulated-entropy
\digestlabel{thm:accumulated-entropy}{2.10}\pdfbookmark[2]{Theorem 2.10 Accumulated entropy bound under adaptive attacks}{digest-meat-core/thm:accumulated-entropy}
\begin{mainclaim}
\renewcommand*{\RestatedName}{\hyperref{contract.pdf}{}{claim-meat-core/thm:accumulated-entropy}{Theorem~2.10}}%
\begin{digestclaim}[Accumulated entropy bound under adaptive attacks]
Let $\Decl$ satisfy Definition~\ref{def:admissible-declaration}, let
$\acSet$ satisfy Definition~\ref{def:admissible-acceptance-set}, and let
Assumption~A-RHOM (\xref{protocol-compiler/ass:rhom}) hold; in
trusted-source mode let A-SRC
(\xref{transmitter/ass:src}), A-ENS (\xref{transmitter/ass:ens}), A-SEQ
(\xref{transmitter/ass:seq}), \ref{ass:seed} (A-SEED) and
\ref{ass:lab} (A-LAB) at the party to which the party partition
assigns $A$ hold as well, and, at $d_A\ge2$, let the declaration's $A$,
$\tau_A$ and $\{M_x\}$ be those of a source specification carrying a
source-replacement certificate against its virtual ensemble,
with a source segment satisfying the source identity, as
Lemma~\ref{lem:product-marginal} states them; together with these, the
adversary's initial register being finite-dimensional as every register of
the round model is, the five assumptions are the hypotheses of
Lemma~\ref{lem:product-marginal}, so that by that lemma the
protocol's $n$-round input in that mode has the marginal below (in
untrusted-source mode the lemma is not invoked and none of the hypotheses
named for that mode is used here). Let $(\Aev_j)_{j=1}^{n}$ be an
\emph{arbitrary} attack sequence (Definition~\ref{def:roundmodel}; no
admissibility restriction; its registers $E_0$ and $E'_j$ finite-dimensional,
of any dimensions, as every register of that model is), let
\begin{equation*}
  \rho_{Z_1^n\Chat_1^nI_1^n\Ehat}
  =\big(\Mprot_n\circ\dots\circ\Mprot_1\big)\big[\omega^{0}\big],
  \qquad
  \omega^{0}\in\Seq{A_1^nE_0}\ \text{with}\
  \omega^{0}_{A_1^n}=\tau_A^{\otimes n}
  \ \ \text{(trusted mode)},
\end{equation*}
respectively $\omega^{0}\in\Seq{E_0}$ arbitrary (untrusted mode, where
the round channels take no $A$ input),
with $\Mprot_j$ the round channels \eqref{eq:roundchannel} and Eve's
final register written $E_n=I_1^n\Ehat$
(Definition~\ref{def:roundmodel}), and let $\OmAT$, the acceptance event, be
the event that the frequency vector $\operatorname{freq}_{\hat c_1^n}$ of the
acceptance values of the $n$ rounds lies in $\acSet$, with $\Pr[\OmAT]>0$.
Then for
every $\alpha\in(1,\infty)$,
\begin{equation*}\tag{\getrefnumber{C-eq:c3}}\digesteqlabel{eq:c3}
  \Hsandup\big(Z_1^n\big|\Chat_1^nI_1^n\Ehat\big)_{\rho_{|\OmAT}}
  \;\ge\;n\,\hqkd-\frac{\alpha}{\alpha-1}\log_2\frac{1}{\Pr[\OmAT]},
\end{equation*}
with $\hqkd=\hqkd(\varnothing)$ the single-round rate of
Definition~\ref{def:hqkd-list} at the empty list, which reads
\begin{equation*}\tag{\getrefnumber{C-eq:hqkd-def}}\digesteqlabel{eq:hqkd-def}
\begin{gathered}
  \hqkd=\inf_{\mathbf q\in\acSet}\
  \inf_{\substack{\omega\in\Seq{A\Qsys\Etil}\\ \omega\models\Decl}}
  \left(\frac{\alpha}{\alpha-1}\,
  \sfD\!\left(\mathbf q\,\middle\|\,\Phimap[\omega_{A\Qsys}]\right)
  +\sum_{\hat c\in\Chkey}
  q(\hat c)\,\Hsandup(Z|I\Etil)_{\nu^{\omega}_{|\hat c}}\right),\\
  \nu^{\omega}\coloneqq(\Mprot\otimes\idmap_{\Etil})(\omega),
\end{gathered}
\end{equation*}
with $\omega\models\Decl$ as in \eqref{eq:margcond} --- so the inner
infimum is over states with $\omega_A=\tau_A$ in trusted-source mode
and over \emph{all} normalized states on $A\Qsys\Etil$ in
untrusted-source mode. It ranges over all (not necessarily pure) such
states, with
$\dim\Hs{\Etil}\ge\dim\Hs{A}\cdot\dim\Hs{\Qsys}$
(Lemma~\ref{lem:purifier-indep}: the value is independent of that
choice).
\end{digestclaim}
\end{mainclaim}

\section*{3\quad Constrained accumulation}\pdfbookmark[1]{3 Constrained accumulation}{digest-group-meat-core/thm:constrained-accumulation}

% claim meat-core/thm:constrained-accumulation
\digestlabel{thm:constrained-accumulation}{3.1}\pdfbookmark[2]{Theorem 3.1 Constrained accumulation under the attack-set licence}{digest-meat-core/thm:constrained-accumulation}
\renewcommand*{\RestatedName}{\hyperref{contract.pdf}{}{claim-meat-core/thm:constrained-accumulation}{Theorem~3.1}}%
\begin{digestclaim}[Constrained accumulation under the attack-set licence]
Let the hypotheses of Theorem~\ref{thm:accumulated-entropy} hold in either
source mode, let $\CList$ be
licensed for that mode
(Definition~\ref{def:homog}), let A-ATK
(\xref{protocol-compiler/ass:atk}) hold, and let the attack sequence be
$\CList$-admissible (Definition~\ref{def:roundmodel}). Then for every
$\alpha\in(1,\infty)$,
\begin{equation*}\tag{\getrefnumber{C-eq:c17}}\digesteqlabel{eq:c17}
  \Hsandup\big(Z_1^n\big|\Chat_1^nI_1^n\Ehat\big)_{\rho_{|\OmAT}}
  \;\ge\;n\,\hqkd(\CList)
  -\frac{\alpha}{\alpha-1}\log_2\frac{1}{\Pr[\OmAT]},
\end{equation*}
with $\hqkd(\CList)$ the constrained single-round rate of
Definition~\ref{def:hqkd-list}, and
$\hqkd(\CList)\ge\hqkd(\varnothing)=\hqkd$.
\end{digestclaim}

\section*{4\quad Branch entropy and sectors}\pdfbookmark[1]{4 Branch entropy and sectors}{digest-group-meat-core/thm:branch-formula}

% claim meat-core/thm:branch-formula
\digestlabel{thm:branch-formula}{4.1}\pdfbookmark[2]{Theorem 4.1 Branch entropy from the protocol marginal}{digest-meat-core/thm:branch-formula}
\renewcommand*{\RestatedName}{\hyperref{contract.pdf}{}{claim-meat-core/thm:branch-formula}{Theorem~4.1}}%
\begin{digestclaim}[Branch entropy from the protocol marginal]
For $\alpha\in(1,\infty)$, $\omega\in\Pos{A\Qsys}$ and $\hat c\in\Chatset$
define
\begin{equation*}\tag{\getrefnumber{C-eq:fK}}\digesteqlabel{eq:fK}
  P_{\hat c}(\omega)\coloneqq\Phimap[\omega](\hat c),\qquad
  T_{\hat c}(\omega)\coloneqq
  \tr\!\left[\left(\cZ\big(\cG_{\hat c}(\omega)^{1/\alpha}\big)
  \right)^{\alpha}\right],\qquad
  K_{\hat c}(\omega)\coloneqq
  T_{\hat c}(\omega)+P_{\hat c}(\omega)-\tr\cG_{\hat c}(\omega),
\end{equation*}
and, for block-diagonal $\omega=\bigoplus_s\omega_s$, the sector deficits
\begin{equation*}\tag{\getrefnumber{C-eq:f-def}}\digesteqlabel{eq:f-def}
  f_{\hat c,s}(\omega_s)\coloneqq
  \tr\!\left[\left(\cZ\big(\cG_{\hat c,s}(\omega_s)^{1/\alpha}\big)
  \right)^{\alpha}\right]-\tr\cG_{\hat c,s}(\omega_s).
\end{equation*}
Let $\alpha\in(1,\infty)$, let $\omega\in\Seq{A\Qsys\Etil}$ be \emph{pure},
let $\nu^{\omega}=(\Mprot\otimes\idmap_{\Etil})(\omega)$, and let
$\hat c\in\Chatset$ with $P_{\hat c}\coloneqq
P_{\hat c}(\omega_{A\Qsys})>0$. Then:
\begin{enumerate}[label=(\roman*),leftmargin=2.2em]
\item \emph{(the identity)}
\begin{equation*}\tag{\getrefnumber{C-eq:c5}}\digesteqlabel{eq:c5}
  \Hsanddown(Z|I\Etil)_{\nu^{\omega}_{|\hat c}}
  =\frac{1}{1-\alpha}\log_2\frac{K_{\hat c}(\omega_{A\Qsys})}{P_{\hat c}};
\end{equation*}
\item \emph{(the bounds)} $0\le K_{\hat c}(\omega_{A\Qsys})\le P_{\hat c}$, so
the right-hand side of \eqref{eq:c5} is non-negative;
\item \emph{(the sector form, on the whole cone)} for every
$\hat c\in\Chatset$ and \emph{every} block-diagonal
$\varrho=\bigoplus_s\varrho_s\in\Pos{A\Qsys}$, with none of the hypotheses
above on $\omega$ and $P_{\hat c}$ (no purity, no normalization and no
hypothesis $P_{\hat c}(\varrho)>0$),
\begin{equation*}\tag{\getrefnumber{C-eq:c5-sector}}\digesteqlabel{eq:c5-sector}
  K_{\hat c}(\varrho)
  =P_{\hat c}(\varrho)+\sum_{s\in\SecSet}f_{\hat c,s}(\varrho_s).
\end{equation*}
\end{enumerate}
The two bounds of (ii) hold on the whole cone, as a statement separate from
the identity (i): for every $\alpha\in(1,\infty)$, every
$\hat c\in\Chatset$ and \emph{every} $\varrho\in\Pos{A\Qsys}$, with no
purity, no normalization, no block-diagonality and no hypothesis
$P_{\hat c}(\varrho)>0$, one has $0\le K_{\hat c}(\varrho)\le
P_{\hat c}(\varrho)$. We call this the theorem's \emph{cone clause}.
\end{digestclaim}

% claim meat-core/cor:concavity
\digestlabel{cor:concavity}{4.2}\pdfbookmark[2]{Corollary 4.2 Concavity of the branch quantities on subnormalised states}{digest-meat-core/cor:concavity}
\renewcommand*{\RestatedName}{\hyperref{contract.pdf}{}{claim-meat-core/cor:concavity}{Corollary~4.2}}%
\begin{digestclaim}[Concavity of the branch quantities on subnormalised states]
For $\alpha\in(1,\infty)$, the map
$X\mapsto\tr[(\cZ(X^{1/\alpha}))^{\alpha}]=Q_{1/\alpha}(X)$ is concave on
$\Sleq{}$ of its space; consequently, for every $\hat c$, $s$, the deficit
$f_{\hat c,s}$ of \eqref{eq:f-def} is concave on
$\{\sigma\in\Pos{A\Qsys_s}:\tr\sigma\le1\}$, and
$K_{\hat c}$ is concave on $\Feas$.
\end{digestclaim}

% claim meat-core/lem:cone-concavity
\digestlabel{lem:cone-concavity}{4.3}\pdfbookmark[2]{Lemma 4.3 Concavity on the positive cone}{digest-meat-core/lem:cone-concavity}
\renewcommand*{\RestatedName}{\hyperref{contract.pdf}{}{claim-meat-core/lem:cone-concavity}{Lemma~4.3}}%
\begin{digestclaim}[Concavity on the positive cone]
Let $\alpha\in(1,\infty)$, put $p\coloneqq1/\alpha\in(0,1)$, and let $\cZ$
be a pinching built from a complete family of mutually orthogonal
projectors. On the cone $\Pos{}$ of the space $\cZ$ acts on, consider the
functional
\begin{equation*}\tag{\getrefnumber{C-eq:Tcone}}\digesteqlabel{eq:Tcone}
  T(X)\coloneqq\tr\!\big[\big(\cZ(X^{p})\big)^{\alpha}\big].
\end{equation*}
\begin{enumerate}[label=(\alph*),leftmargin=2.4em]
\item $T$ is positively $1$-homogeneous: $T(tX)=t\,T(X)$ for every
$t\ge0$.
\item $T$ is concave on all of $\Pos{}$.
\item Consider a declaration whose components other than $\{F_v\}$
satisfy Definition~\ref{def:admissible-declaration} and whose $\{F_v\}$ is any family of
sector-block-diagonal positive semidefinite operators, complete or not.
Then each deficit $f_{\hat c,s}$ of \eqref{eq:f-def} is concave on all of
$\Pos{A\Qsys_s}$, and $K_{\hat c}$ of \eqref{eq:fK} is concave on the
sector-block-diagonal part of $\Pos{A\Qsys}$, in particular on $\Feas$.
\end{enumerate}
\end{digestclaim}

% claim meat-core/cor:f-nonpos
\digestlabel{cor:f-nonpos}{4.4}\pdfbookmark[2]{Corollary 4.4 Non-positivity of the sector deficits}{digest-meat-core/cor:f-nonpos}
\renewcommand*{\RestatedName}{\hyperref{contract.pdf}{}{claim-meat-core/cor:f-nonpos}{Corollary~4.4}}%
\begin{digestclaim}[Non-positivity of the sector deficits]
For every $\hat c$, $s$ and \emph{every} $\sigma\in\Pos{A\Qsys_s}$, with no
trace hypothesis (in particular not $\tr\sigma\le1$),
$f_{\hat c,s}(\sigma)\le0$, and
$f_{\hat c,s}(\sigma)=0$ if and only if $\cG_{\hat c,s}(\sigma)$ is
block-diagonal for the key pinching $\cZ$.
\end{digestclaim}

% claim meat-core/def:keyinert
\digestlabel{def:keyinert}{4.5}\pdfbookmark[2]{Definition 4.5 Key-inert pair}{digest-meat-core/def:keyinert}
\renewcommand*{\RestatedName}{\hyperref{contract.pdf}{}{claim-meat-core/def:keyinert}{Definition~4.5}}%
\begin{digestclaim}[Key-inert pair]
For $\hat c\in\Chkey$ and $s\in\SecSet$ define the
\emph{effective outcome set}
of a pass class $\kappa\in\Kann_{\mathrm{pass}}(\hat c)$ on sector $s$ as
\begin{equation*}
  E^{(s)}_{\kappa}\coloneqq
  \big\{(x,v)\in a^{-1}(\kappa):\ M_x\ne0,\ F^{(s)}_v\ne0\big\}.
\end{equation*}
The pair $(\hat c,s)$ is
\emph{key-inert}
if for every $\kappa\in\Kann_{\mathrm{pass}}(\hat c)$ the raw-key map $r$ is
constant on $E^{(s)}_{\kappa}$. This is a finite, mechanically checkable
predicate on the declaration.
\end{digestclaim}

% claim meat-core/prop:exact-cancellation
\digestlabel{prop:exact-cancellation}{4.6}\pdfbookmark[2]{Proposition 4.6 Exact cancellation of key-inert sectors}{digest-meat-core/prop:exact-cancellation}
\renewcommand*{\RestatedName}{\hyperref{contract.pdf}{}{claim-meat-core/prop:exact-cancellation}{Proposition~4.6}}%
\begin{digestclaim}[Exact cancellation of key-inert sectors]
If $(\hat c,s)$ is key-inert, then for \emph{every}
$\sigma\in\Pos{A\Qsys_s}$ the operator $\cG_{\hat c,s}(\sigma)$ is
block-diagonal for the key pinching $\cZ$, and consequently
$f_{\hat c,s}(\sigma)=0$. In particular such sectors contribute nothing to
$K_{\hat c}$ in \eqref{eq:c5-sector}, for every input; they remain
optimization variables, entering through $\Phimap$, $P_{\hat c}$, and the
marginal constraint.
\end{digestclaim}

% claim meat-core/def:keyactive
\digestlabel{def:keyactive}{4.7}\pdfbookmark[2]{Definition 4.7 Key-active index set}{digest-meat-core/def:keyactive}
\renewcommand*{\RestatedName}{\hyperref{contract.pdf}{}{claim-meat-core/def:keyactive}{Definition~4.7}}%
\begin{digestclaim}[Key-active index set]
The
\emph{key-active index set}
of $\hat c$ is
$\SecSet_{K}(\hat c)\coloneqq\{s\in\SecSet:(\hat c,s)\text{ is not
key-inert}\}$. By Proposition~\ref{prop:exact-cancellation},
$K_{\hat c}(\omega)=P_{\hat c}(\omega)
+\sum_{s\in\SecSet_{K}(\hat c)}f_{\hat c,s}(\omega_s)$ for block-diagonal
$\omega$; the objective computed with the sum restricted to
$\SecSet_K(\hat c)$ is \emph{exactly equal} to the full sum, not a bound.
\end{digestclaim}

% claim meat-core/def:joint-sectors
\digestlabel{def:joint-sectors}{4.8}\pdfbookmark[2]{Definition 4.8 Joint-sector deficits and the joint key-active index set}{digest-meat-core/def:joint-sectors}
\renewcommand*{\RestatedName}{\hyperref{contract.pdf}{}{claim-meat-core/def:joint-sectors}{Definition~4.8}}%
\begin{digestclaim}[Joint-sector deficits and the joint key-active index set]
Let $\Decl$ be a declaration in the general format of
Definition~\ref{def:declaration}, with Alice sectors $t\in\mathbb{T}$ and
joint sectors $(t,s)\in\mathbb{T}\times\SecSet$. Regard
$\Pos{A_t\Qsys_s}$ as the subcone of $\Pos{A\Qsys_s}$ of operators supported
on $\Hs{A_t}\otimes\Hs{\Qsys_s}$. For $\hat c\in\Chatset$ the
\emph{joint-sector} maps, statistics and deficit are the restrictions of the
sector objects of Definition~\ref{def:branch} and of \eqref{eq:f-def} to that
subcone,
\begin{equation*}\tag{\getrefnumber{C-eq:joint-deficit}}\digesteqlabel{eq:joint-deficit}
  \cG_{\hat c,(t,s)}\coloneqq\cG_{\hat c,s}\big|_{\Pos{A_t\Qsys_s}},\qquad
  P_{\hat c,(t,s)}(\sigma)\coloneqq P_{\hat c}(\sigma),\qquad
  f_{\hat c,(t,s)}(\sigma)\coloneqq f_{\hat c,s}(\sigma)
  \qquad(\sigma\in\Pos{A_t\Qsys_s}),
\end{equation*}
so that on such a $\sigma$ the Kraus operators \eqref{eq:kraus} act through
$\sqrt{M^{(t)}_x}\otimes\sqrt{F^{(s)}_v}$. For
$\hat c\in\Chkey$ and a pass class $\kappa\in\Kann_{\mathrm{pass}}(\hat c)$
the \emph{effective outcome set on the joint sector} is
\begin{equation*}\tag{\getrefnumber{C-eq:joint-effective}}\digesteqlabel{eq:joint-effective}
  E^{(t,s)}_{\kappa}\coloneqq\big\{(x,v)\in a^{-1}(\kappa):\
  M^{(t)}_x\ne0,\ F^{(s)}_v\ne0\big\}\subseteq E^{(s)}_{\kappa},
\end{equation*}
and the pair $(\hat c,(t,s))$ is \emph{key-inert} if $r$ is constant on
$E^{(t,s)}_{\kappa}$ for every $\kappa\in\Kann_{\mathrm{pass}}(\hat c)$. The
\emph{joint key-active index set}
is $\mathbb{J}_K(\hat c)\coloneqq\{(t,s):(\hat c,(t,s))\text{ is not
key-inert}\}$.

Key-inertness of a joint pair is decided on declared data. In
POVM form $M^{(t)}_x\ne0$ is a comparison of the declared block with
the zero block. In ensemble form it is $N^{(t)}_x\ne0$, by
Lemma~\ref{lem:ensemble-form}\,(iii), again a comparison of a declared block
with zero. For a factorised family
$F^{(s)}_v=\bigotimes_fF^{f,(s_f)}_{v_f}$ is non-zero exactly when every
factor block $F^{f,(s_f)}_{v_f}$ is, a tensor product of matrices vanishing
exactly when one of them does; so the predicate is decided per factor and no
product is formed. With $|\mathbb{T}|=1$ one has
$E^{(t_0,s)}_{\kappa}=E^{(s)}_{\kappa}$, and the pair $(\hat c,(t_0,s))$ is
key-inert exactly when $(\hat c,s)$ is in the sense of
Definition~\ref{def:keyinert}, with
$\mathbb{J}_K(\hat c)=\{t_0\}\times\SecSet_K(\hat c)$.
\end{digestclaim}

% claim meat-core/lem:joint-sectors
\digestlabel{lem:joint-sectors}{4.9}\pdfbookmark[2]{Lemma 4.9 The objective split over joint sectors}{digest-meat-core/lem:joint-sectors}
\renewcommand*{\RestatedName}{\hyperref{contract.pdf}{}{claim-meat-core/lem:joint-sectors}{Lemma~4.9}}%
\begin{digestclaim}[The objective split over joint sectors]
Let $\Decl$ satisfy Definition~\ref{def:admissible-declaration} in the general format of
Definition~\ref{def:declaration}, let $\alpha\in(1,\infty)$ and
$\hat c\in\Chatset$.
\begin{enumerate}[label=(\roman*),leftmargin=2.2em]
\item \emph{(decomposition)} For $s\in\SecSet$ and
$\sigma=\bigoplus_{t}\sigma_t\in\Pos{A\Qsys_s}$ with
$\sigma_t\in\Pos{A_t\Qsys_s}$,
\begin{equation*}\tag{\getrefnumber{C-eq:G-joint}}\digesteqlabel{eq:G-joint}
  \cG_{\hat c,s}(\sigma)=\bigoplus_{t\in\mathbb{T}}
  \cG_{\hat c,(t,s)}(\sigma_t),\qquad
  \Theta_{\hat c,s}(\sigma)=\bigoplus_{t\in\mathbb{T}}
  \Theta_{\hat c,s}(\sigma_t),\qquad
  \Phimap_s[\sigma]=\sum_{t}\Phimap_s[\sigma_t],
\end{equation*}
the direct sums along the mutually orthogonal subspaces
$\Hs{R\,XY\,C_AC_B}\otimes\Hs{A_t\Qsys_s}$.
\item \emph{(additivity)} For such $\sigma$,
$f_{\hat c,s}(\sigma)=\sum_tf_{\hat c,(t,s)}(\sigma_t)$. Consequently, for
every $\omega=\bigoplus_{(t,s)}\omega^{(t,s)}$ block diagonal along the joint
sectors,
\begin{equation*}\tag{\getrefnumber{C-eq:K-joint}}\digesteqlabel{eq:K-joint}
  P_{\hat c}(\omega)=\sum_{(t,s)}P_{\hat c,(t,s)}\big(\omega^{(t,s)}\big),
  \qquad
  K_{\hat c}(\omega)=P_{\hat c}(\omega)
  +\sum_{(t,s)}f_{\hat c,(t,s)}\big(\omega^{(t,s)}\big),
\end{equation*}
and for every $\sigma\in\Pos{A_t\Qsys_s}$: $f_{\hat c,(t,s)}(\sigma)\le0$,
$0\le P_{\hat c,(t,s)}(\sigma)+f_{\hat c,(t,s)}(\sigma)$, and
$f_{\hat c,(t,s)}$ is concave and positively $1$-homogeneous on
$\Pos{A_t\Qsys_s}$.
\item \emph{(key-inert joint pairs)} If $\hat c\in\Chkey$ and
$(\hat c,(t,s))$ is key-inert (Definition~\ref{def:joint-sectors}), then
$f_{\hat c,(t,s)}(\sigma)=0$ for \emph{every} $\sigma\in\Pos{A_t\Qsys_s}$.
Hence the second sum of \eqref{eq:K-joint} may be restricted to
$(t,s)\in\mathbb{J}_K(\hat c)$, as an identity and not as a bound.
\item \emph{(relation to the sector notions)} If $(\hat c,s)$ is key-inert in
the sense of Definition~\ref{def:keyinert} then $(\hat c,(t,s))$ is key-inert
for every $t$; the converse fails in general, and holds when
$|\mathbb{T}|=1$.
\end{enumerate}
\end{digestclaim}

\section*{5\quad The program and its reduction}\pdfbookmark[1]{5 The program and its reduction}{digest-group-meat-core/def:feasible}

% claim meat-core/def:feasible
\digestlabel{def:feasible}{5.1}\pdfbookmark[2]{Definition 5.1 Feasible set and computable objective}{digest-meat-core/def:feasible}
\renewcommand*{\RestatedName}{\hyperref{contract.pdf}{}{claim-meat-core/def:feasible}{Definition~5.1}}%
\begin{digestclaim}[Feasible set and computable objective]
For a declaration $\Decl$, $\Feas$ is the feasible set of
\xref{protocol-compiler/def:feasible}. Read as a set
of operators on $A\Qsys$, a list of sector blocks being the block-diagonal
operator it denotes, it is
\begin{equation*}\tag{\getrefnumber{C-eq:feas}}\digesteqlabel{eq:feas}
  \Feas\coloneqq\Big\{\!\omega\in\Pos{A\Qsys}:\
  (\idon{A}\otimes\Pi_s)\,\omega\,(\idon{A}\otimes\Pi_{s'})=0
  \ \forall s\ne s',\quad
  \tr\omega=1,\quad
  \omega\models\Decl,\quad
  \omega\models\CList\Big\},
\end{equation*}
writing $\omega_s$ for the sector blocks, with $\omega\models\Decl$ the
marginal condition of \eqref{eq:margcond} and $\omega\models\CList$ the
constraint conditions in the form
Definition~\ref{def:homog} gives them for the declared mode. Written
out, $\Feas$ is
\begin{equation*}\tag{\getrefnumber{C-eq:feas-modes}}\digesteqlabel{eq:feas-modes}
\begin{aligned}
  \text{trusted:}\quad&
  \omega_A=\tau_A,\quad
  \tr_{\Qsys}[(\idon{A}\otimes\hat O_k)\,\omega]\bowtie_k\Theta_k\
  \ \forall k;\\
  \text{untrusted:}\quad&
  \text{(no marginal condition)},\quad
  \tr[(\idon{A}\otimes\hat O_k)\,\omega]\bowtie_k c_k\ \ \forall k .
\end{aligned}
\end{equation*}
The normalization $\tr\omega=1$ is written \emph{explicitly} in
\eqref{eq:feas}, and is not derived from anything: in
trusted-source mode it also follows from $\omega_A=\tau_A$ (whence
$\tr\omega=\tr\tau_A=1$), but that derivation is unavailable in
untrusted-source mode, where normalization is the \emph{only}
unconditional constraint the feasible set carries. Leaving it implied
would empty the untrusted-mode program of its one defining condition.
$\Feas$ is convex and compact in either mode and in either format
(\xref{protocol-compiler/def:feasible}). For
$\alpha\in(1,\infty)$,
$\mathbf q$ a distribution on $\Chatset$ and $\omega\in\Feas$, define the
\emph{objective}
\begin{equation*}\tag{\getrefnumber{C-eq:objective}}\digesteqlabel{eq:objective}
  O_{\alpha}(\mathbf q,\omega)\coloneqq
  \frac{\alpha}{\alpha-1}\sum_{\hat c\notin\Chkey}
  q(\hat c)\log_2\frac{q(\hat c)}{P_{\hat c}(\omega)}
  +\sum_{\hat c\in\Chkey}\left[
  q(\hat c)\log_2\frac{q(\hat c)}{P_{\hat c}(\omega)}
  +\frac{q(\hat c)}{\alpha-1}\log_2\frac{q(\hat c)}{K_{\hat c}(\omega)}
  \right],
\end{equation*}
with the conventions of \xref{foundations/def:kl} at vanishing arguments
($0\log_2\frac0v=0$; $u\log_2\frac u0=+\infty$ for $u>0$), and
\begin{equation*}\tag{\getrefnumber{C-eq:hlow-def}}\digesteqlabel{eq:hlow-def}
  \hlow\coloneqq\inf_{\mathbf q\in\acSet}\ \inf_{\omega\in\Feas}\
  O_{\alpha}(\mathbf q,\omega).
\end{equation*}

\emph{The general format.} For a declaration in the general format of
Definition~\ref{def:declaration}, $\Feas$ is the set of lists over the
joint sectors of \xref{protocol-compiler/def:feasible},
$\omega=\bigoplus_{(t,s)}\omega^{(t,s)}$, with the $s$-block
$\hat O^{\uparrow(s)}_k$ of the lifted operator of an entry scoped to $f_k$
and, for a party-local entry scoped to $f$, the split $s=(s_{f^c},s_f)$ and
the product $\Qsys_{f^c,s_{f^c}}$ of the other factors' blocks written
there, with every block positive semidefinite; its conditions read, block by
block,
\begin{equation*}\tag{\getrefnumber{C-eq:feas-general}}\digesteqlabel{eq:feas-general}
\begin{aligned}
  &\textstyle\sum_{(t,s)}\tr\omega^{(t,s)}=1
  &&\text{(normalization)};\\
  &\textstyle\sum_{s}\tr_{\Qsys_s}\omega^{(t,s)}=\tau^{(t)}_A
  \quad\forall t
  &&\text{(marginal, trusted-source mode)};\\
  &\textstyle\sum_{s}\tr_{\Qsys_s}\!\big[(\idon{A_t}\otimes
  \hat O^{\uparrow(s)}_k)\,\omega^{(t,s)}\big]\bowtie_k\Theta^{(t)}_k
  \quad\forall t
  &&\text{(trusted; }\Theta^{(t)}_k=c_k\tau^{(t)}_A
  \text{ if homogeneous)};\\
  &\textstyle\sum_{(t,s)}\tr\!\big[(\idon{A_t}\otimes
  \hat O^{\uparrow(s)}_k)\,\omega^{(t,s)}\big]\bowtie_kc_k
  &&\text{(scalar entry, untrusted-source mode)};\\
  &\textstyle\sum_{s_f\in\SecSet^{f}}\tr_{\Qsys_{f,s_f}}\!\Big[\big(\idon{}
  \otimes(\hat O^{(s_f)}_k-c_k\idon{})\big)\,
  \omega^{(t,(s_{f^c},s_f))}\Big]\bowtie_k0
  \quad\forall(t,s_{f^c})
  &&\text{(party-local, either mode)},
\end{aligned}
\end{equation*}
the last as an operator on $\Hs{A_t}\otimes\Hs{\Qsys_{f^c,s_{f^c}}}$. The
objective \eqref{eq:objective} and $\hlow$ of \eqref{eq:hlow-def} are
defined on this $\Feas$ by the same formulas.
\end{digestclaim}

% claim meat-core/prop:trivial-alice-system
\digestlabel{prop:trivial-alice-system}{5.2}\pdfbookmark[2]{Proposition 5.2 The two source modes coincide at a trivial Alice system}{digest-meat-core/prop:trivial-alice-system}
\renewcommand*{\RestatedName}{\hyperref{contract.pdf}{}{claim-meat-core/prop:trivial-alice-system}{Proposition~5.2}}%
\begin{digestclaim}[The two source modes coincide at a trivial Alice system]
Let $d_A=1$. Then trusted-source mode admits exactly one marginal,
$\tau_A=[1]$, and:
\begin{enumerate}[label=(\roman*),leftmargin=2.2em]
\item the two feasible sets of \eqref{eq:feas} coincide;
\item the $n$-round condition $\omega_{A_1^n}=\tau_A^{\otimes n}$ of
Lemma~\ref{lem:product-marginal}\,(ii) holds for every normalized state; it is
satisfied vacuously;
\item the licensed constraint classes of Definition~\ref{def:homog}
coincide: trusted-mode $\tau_A$-homogeneity is exactly the scalar typing
of untrusted-source mode;
\item A-SRC (\xref{transmitter/ass:src}) is vacuous: what it
licenses is the condition $\tr\omega=1$, satisfied by every normalized
state, so it deletes no attack.
\end{enumerate}
\end{digestclaim}

% claim meat-core/thm:reduction
\digestlabel{thm:reduction}{5.4}\pdfbookmark[2]{Theorem 5.4 Reduction of the rate to the convex program}{digest-meat-core/thm:reduction}
\begin{mainclaim}
\renewcommand*{\RestatedName}{\hyperref{contract.pdf}{}{claim-meat-core/thm:reduction}{Theorem~5.4}}%
\begin{digestclaim}[Reduction of the rate to the convex program]
Let $\Decl$ satisfy Definition~\ref{def:admissible-declaration} with $\CList=\varnothing$,
in either source mode, and let $\alpha\in(1,\infty)$. Then, with
$\hqkd=\hqkd(\varnothing)$ the single-round rate of
Definition~\ref{def:hqkd-list} at the empty list,
\begin{equation*}\tag{\getrefnumber{C-eq:c8}}\digesteqlabel{eq:c8}
  \hqkd\;\ge\;\hlow .
\end{equation*}
\end{digestclaim}
\end{mainclaim}

% claim meat-core/prop:a-neutral
\digestlabel{prop:a-neutral}{5.5}\pdfbookmark[2]{Proposition 5.5 Value neutrality along Alice's sectors}{digest-meat-core/prop:a-neutral}
\renewcommand*{\RestatedName}{\hyperref{contract.pdf}{}{claim-meat-core/prop:a-neutral}{Proposition~5.5}}%
\begin{digestclaim}[Value neutrality along Alice's sectors]
Let $\Decl$ satisfy Definition~\ref{def:admissible-declaration} in the general format of
Definition~\ref{def:declaration}, in either source mode and with an arbitrary
list $\CList$, general and party-local entries included.
Let $\Feas^{\circ}$ be the set defined as $\Feas$ is in
Definition~\ref{def:feasible} but with block diagonality required along
$\SecSet$ only --- the feasible set of the same declaration with the factor
tagged $A$ forgotten --- and let
$\mathcal{P}_{\mathbb{T}}(X)\coloneqq\sum_{t}\Pi^{A}_tX\Pi^{A}_t$. Then for
every $\alpha\in(1,\infty)$, every distribution $\mathbf q$ on $\Chatset$ and
every $\omega\in\Feas^{\circ}$, the operator
$\omega'\coloneqq(\mathcal{P}_{\mathbb{T}}\otimes\idmap_{\Qsys})(\omega)$
satisfies
\begin{equation*}\tag{\getrefnumber{C-eq:a-neutral}}\digesteqlabel{eq:a-neutral}
  \omega'\in\Feas,\qquad
  P_{\hat c}(\omega')=P_{\hat c}(\omega),\qquad
  K_{\hat c}(\omega')\ge K_{\hat c}(\omega)\quad\forall\hat c,\qquad
  O_{\alpha}(\mathbf q,\omega')\le O_{\alpha}(\mathbf q,\omega),
\end{equation*}
the last in the extended-real sense. Consequently
$\inf_{\acSet\times\Feas}O_{\alpha}=\inf_{\acSet\times\Feas^{\circ}}O_{\alpha}$,
and $\omega\mapsto\omega'$ maps minimisers over $\Feas^{\circ}$ to minimisers
over $\Feas$.
\end{digestclaim}

% claim meat-core/cor:reduction-constrained
\digestlabel{cor:reduction-constrained}{5.6}\pdfbookmark[2]{Corollary 5.6 The reduction for arbitrary constraint lists}{digest-meat-core/cor:reduction-constrained}
\renewcommand*{\RestatedName}{\hyperref{contract.pdf}{}{claim-meat-core/cor:reduction-constrained}{Corollary~5.6}}%
\begin{digestclaim}[The reduction for arbitrary constraint lists]
Let $\CList$ be arbitrary: $\Theta_k$ an arbitrary Hermitian operator on $A$
in trusted-source mode, $\Theta_k=c_k$ an arbitrary real number in
untrusted-source mode, and entries declared party-local in
either mode (Definition~\ref{def:admissible-declaration}\,(vi)); no licensing restriction
(Definition~\ref{def:homog}) is needed in either mode. Then
\begin{equation*}\tag{\getrefnumber{C-eq:c8-general}}\digesteqlabel{eq:c8-general}
  \hqkd(\CList)\;\ge\;\hlow(\CList),
\end{equation*}
where $\hqkd(\CList)$ is the constrained single-round rate of
Definition~\ref{def:hqkd-list} and $\hlow(\CList)$ denotes
\eqref{eq:hlow-def} over the feasible set $\Feas$ of \eqref{eq:feas},
which carries the same list ($\hlow(\CList)$ and the unadorned
$\hlow$ are the same quantity; the argument is written when the
dependence on the declared list matters, as in
$\hlow(\CList)\ge\hlow(\varnothing)$).
\end{digestclaim}

% claim meat-core/prop:value-range
\digestlabel{prop:value-range}{5.8}\pdfbookmark[2]{Proposition 5.8 Range of the value and its finite entropy part}{digest-meat-core/prop:value-range}
\renewcommand*{\RestatedName}{\hyperref{contract.pdf}{}{claim-meat-core/prop:value-range}{Proposition~5.8}}%
\begin{digestclaim}[Range of the value and its finite entropy part]
Let $\alpha\in(1,\infty)$.
\begin{enumerate}[label=(\roman*),leftmargin=2.2em]
\item \emph{(lower bound)} $O_{\alpha}(\mathbf q,\omega)\ge0$ for every
distribution $\mathbf q$ on $\Chatset$ and every $\omega\in\Feas$; hence
$\hlow\in[0,+\infty]$.
\item \emph{(entropy-part ceiling)} For every $\omega\in\Feas$ and every
$\hat c\in\Chkey$ with $P_{\hat c}(\omega)>0$, the entropy term
$e_{\hat c}(\omega)\coloneqq\frac{1}{1-\alpha}
\log_2\frac{K_{\hat c}(\omega)}{P_{\hat c}(\omega)}$ satisfies
$e_{\hat c}(\omega)\le\log_2|\Zset|=\log_2(d_Z+1)$.
\item \emph{(no overall ceiling)} No finite a-priori upper bound on $\hlow$
exists: the divergence penalty is unbounded already at fixed $\alpha$.
\item \emph{(infinite value)} $\hlow=+\infty$ whenever no
$(\mathbf q,\omega)$ with $\mathbf q\ll\Phimap[\omega]$ is feasible. If
$\hlow=+\infty$, then no $\mathbf q\in\acSet$ is absolutely continuous with
respect to $\Phimap[\omega]$ for any $\omega\in\Feas$.
\end{enumerate}
\end{digestclaim}

% claim meat-core/thm:convexity
\digestlabel{thm:convexity}{5.10}\pdfbookmark[2]{Theorem 5.10 Joint convexity of the program}{digest-meat-core/thm:convexity}
\renewcommand*{\RestatedName}{\hyperref{contract.pdf}{}{claim-meat-core/thm:convexity}{Theorem~5.10}}%
\begin{digestclaim}[Joint convexity of the program]
For every $\alpha\in(1,\infty)$ and every declaration satisfying
Definition~\ref{def:admissible-declaration} (any $\CList$), the extended-real objective
$O_{\alpha}$ of \eqref{eq:objective} is jointly convex and lower
semicontinuous on $\Delta(\Chatset)\times\Feas$, where $\Delta(\Chatset)$
is the probability simplex. Consequently \eqref{eq:hlow-def} is a convex
minimization over the compact convex set $\acSet\times\Feas$, and the
infimum is attained whenever it is finite.
\end{digestclaim}

\section*{6\quad Sector classes and sector removal}\pdfbookmark[1]{6 Sector classes and sector removal}{digest-group-meat-core/def:classes}

% claim meat-core/def:classes
\digestlabel{def:classes}{6.1}\pdfbookmark[2]{Definition 6.1 Sector classes and the removal checklist}{digest-meat-core/def:classes}
\renewcommand*{\RestatedName}{\hyperref{contract.pdf}{}{claim-meat-core/def:classes}{Definition~6.1}}%
\begin{digestclaim}[Sector classes and the removal checklist]
Let $\Decl$ satisfy Definition~\ref{def:admissible-declaration}. For $s\in\SecSet$:
\begin{enumerate}[label=(\alph*),leftmargin=2.2em]
\item $s$ is \emph{key-active} if $(\hat c,s)$ fails to be key-inert
(Definition~\ref{def:keyinert}) for at least one $\hat c\in\Chkey$; these
sectors carry the deficits $f_{\hat c,s}$ in the objective.
\item $s$ is \emph{constraint-active} if it is key-inert for every
$\hat c\in\Chkey$; by Proposition~\ref{prop:exact-cancellation} it contributes nothing to
any $K_{\hat c}$ but remains an optimization variable (it enters
$\Phimap$, every $P_{\hat c}$, the marginal constraint, and $\CList$).
\item $s$ is \emph{state-blind} if $F^{(s)}_v=c_{v,s}\,\idon{\Qsys_s}$ for
scalars $c_{v,s}\ge0$, for every $v\in\Vset$ (then
$\sum_vc_{v,s}=1$ by POVM completeness restricted to the sector); if
$\CList\neq\varnothing$ we additionally require each
$\hat O_k^{(s)}=o_{k,s}\,\idon{\Qsys_s}$ for scalars $o_{k,s}$.
\item $s$ is \emph{removable} if the hypotheses of
Theorem~\ref{thm:removal} hold for it.
\end{enumerate}
All four predicates are decided by finitely many linear-algebra checks on
the declaration. Items (a) and (b) reduce to key-inertness
(Definition~\ref{def:keyinert}), a finite combinatorial predicate on the
declared function tables and on which of the $M_x$, $F^{(s)}_v$ are the zero
operator, decided exactly, at no window. Items (c) and (d) are decided
exactly as well: (c) is a condition on the
declared values themselves
(\xref{foundations/def:binary64}), so that a block that is only
\emph{nearly} scalar fails it and $s$ is then not state-blind, and (d) is
decided by the checklist below.

\emph{The removal checklist.} A \emph{candidate witness} for a sector
$s^{*}$ is a target list $s_1,\dots,s_L$ as index data, states $\chi_i$ as
binary64 sector-block data, weights $\mu_i$ as binary64 data, and the
evaluation order of (W6). It is a \emph{witness} for $s^{*}$ in the sense
of~(d) when $\CList$ carries no entry declared party-local and all six
of the following hold:
\begin{enumerate}[label=(W\arabic*),leftmargin=3.0em]
\item \emph{$s^{*}$ is state-blind} in the sense of (c), the
$\hat O^{(s^{*})}_k$ half included whenever $\CList\neq\varnothing$;
\item \emph{the witness does not name $s^{*}$, and names no sector
twice}: $L\ge1$, $s_i\in\SecSet\setminus\{s^{*}\}$ for every $i$, and the
$s_i$ pairwise distinct, all compared as index data at
no window. A candidate naming $s^{*}$ fails (W2), and so does one that
repeats a target;
\item \emph{weights}: $\mu_i\ge0$ for every $i$ --- a comparison of an
exhibited binary64 datum against zero --- and $\sum_i\mu_i=1$, the sum
formed in the evaluation order the witness exhibits and verified to
round nowhere, exactly as (W6) requires of its
residuals, so that the computed sum \emph{is} the real sum;
\item \emph{states}: $\chi_i\succeq0$ and $\tr\chi_i=1$ for every $i$,
decided on the \emph{exhibited form} and never on a computed spectrum:
the witness presents each $\chi_i$ diagonal in the declared basis of
$\Qsys_{s_i}$ --- every off-diagonal entry the datum $0$ in the sense of
(c), every diagonal entry a datum $\ge0$, and their sum formed in the
exhibited order, rounding nowhere, and equal to the datum $1$ --- so that
both conditions hold of the declared numbers by construction. Any other
presentation fails (W4), which can only \emph{refuse} a removal;
\item \emph{key-inert targets}, which is (RM2): every pair
$(\hat c,s_i)$ with $\hat c\in\Chkey$ and $i\in\{1,\dots,L\}$ is key-inert
(Definition~\ref{def:keyinert});
\item \emph{both equality families close at residual exactly zero}: (RM1)
for every $v\in\Vset$, and (RM3) for every $k$ (vacuous when
$\CList=\varnothing$).
\end{enumerate}
The condition on $\CList$ and (W1)--(W6) are the hypotheses of
Theorem~\ref{thm:removal} transcribed, all of them and in the theorem's own
words --- the condition on $\CList$ from its first hypothesis, (W1) and
(W2) from its opening sentence, (W3) and (W4) from its quantifiers over
$\mu_i$ and $\chi_i$, (W5) from (RM2), (W6) from (RM1) and (RM3) --- which
is what (d) means by ``the hypotheses of Theorem~\ref{thm:removal}
hold''.
Condition (W6) is read on the residual itself: each residual is formed in
an evaluation order that the witness exhibits alongside its data, every
elementary operation of which rounds nothing, so that the computed residual
\emph{is} the real residual and its comparison against zero decides the
real identity.

\emph{Deleting a set of removable sectors needs one condition more, and it is
not part of this definition.} Removability is a property of one sector, and
Theorem~\ref{thm:removal} licenses the deletion of one sector; that two
sectors are each removable does not license deleting both.
Proposition~\ref{prop:removal-compose} states the further hypothesis on the
set $\SecSet_{\mathrm{del}}$ a run deletes.
Nothing in (W1)--(W6) implies that hypothesis.

\emph{In the general format} of Definition~\ref{def:declaration} the four
classes remain predicates of the sectors $s\in\SecSet$ of $\Qsys$, decided
as above, with $\hat O^{(s)}_k$ the block of the lifted operator of an
entry. The finer notion along Alice's sectors --- which joint
pairs $(\hat c,(t,s))$ carry a deficit --- is
Definition~\ref{def:joint-sectors}, and a sector that is key-active here may
have key-inert joint pairs. Theorem~\ref{thm:removal} is stated for a list
carrying no entry declared party-local, so for a list carrying one,
clause~(d) is unsatisfiable and no sector carries a witness in its sense.
\end{digestclaim}

% claim meat-core/thm:removal
\digestlabel{thm:removal}{6.2}\pdfbookmark[2]{Theorem 6.2 Sound removal of a state-blind sector}{digest-meat-core/thm:removal}
\renewcommand*{\RestatedName}{\hyperref{contract.pdf}{}{claim-meat-core/thm:removal}{Theorem~6.2}}%
\begin{digestclaim}[Sound removal of a state-blind sector]
Let $\CList$ carry no entry declared party-local, let
$s^{*}\in\SecSet$ be state-blind, and suppose there exist
\emph{pairwise distinct} sectors
$s_1,\dots,s_L\in\SecSet\setminus\{s^{*}\}$, states
$\chi_i\in\Seq{\Qsys_{s_i}}$, and weights $\mu_i\ge0$ with
$\sum_{i=1}^{L}\mu_i=1$, such that
\begin{enumerate}[label=(RM\arabic*),leftmargin=3.0em]
\item \emph{response match:}\quad
$c_{v,s^{*}}=\sum_{i}\mu_i\tr\big[F^{(s_i)}_v\chi_i\big]$ for every
$v\in\Vset$;
\item \emph{key-inert targets:}\quad every pair $(\hat c,s_i)$ with
$\hat c\in\Chkey$ is key-inert;
\item \emph{constraint match} (vacuous for $\CList=\varnothing$):\quad
$o_{k,s^{*}}=\sum_i\mu_i\tr\big[\hat O^{(s_i)}_k\chi_i\big]$ for every
$k$.
\end{enumerate}
Then restricting the infimum \eqref{eq:hlow-def} to
$\{\omega\in\Feas:\omega_{s^{*}}=0\}$ changes neither its value nor its
attainment: the sector $s^{*}$ may be deleted from the variable list, and
its rows deleted from the generated maps, without loss.
\end{digestclaim}

% claim meat-core/prop:retention
\digestlabel{prop:retention}{6.3}\pdfbookmark[2]{Proposition 6.3 Unsoundness of dropping a constraint-active sector}{digest-meat-core/prop:retention}
\renewcommand*{\RestatedName}{\hyperref{contract.pdf}{}{claim-meat-core/prop:retention}{Proposition~6.3}}%
\begin{digestclaim}[Unsoundness of dropping a constraint-active sector]
There exists a declaration satisfying Definition~\ref{def:admissible-declaration} and a
constraint-active sector $s_0$ such that forcing $\omega_{s_0}=0$ strictly
\emph{increases} $\hlow$. Dropping a constraint-active sector from the
variable list is therefore unsound: it can overstate the extractable key.
\end{digestclaim}

% claim meat-core/prop:no-averaging
\digestlabel{prop:no-averaging}{6.4}\pdfbookmark[2]{Proposition 6.4 Unsoundness of averaging sectors}{digest-meat-core/prop:no-averaging}
\renewcommand*{\RestatedName}{\hyperref{contract.pdf}{}{claim-meat-core/prop:no-averaging}{Proposition~6.4}}%
\begin{digestclaim}[Unsoundness of averaging sectors]
Coarse-graining the sector decomposition of a declaration --- replacing
$\Qsys=\bigoplus_s\Qsys_s$ and $\{F^{(s)}_v\}$ by a single sector carrying
averaged POVM elements --- is unsound: there exists a declaration
$\Decl$ and an averaged declaration $\bar\Decl$ reproducing the same
outcome statistics on the averaged state space for which
$\hlow(\bar\Decl)>\hlow(\Decl)$.
\end{digestclaim}

% claim meat-core/prop:removal-compose
\digestlabel{prop:removal-compose}{6.5}\pdfbookmark[2]{Proposition 6.5 Deleting a set of removable sectors}{digest-meat-core/prop:removal-compose}
\renewcommand*{\RestatedName}{\hyperref{contract.pdf}{}{claim-meat-core/prop:removal-compose}{Proposition~6.5}}%
\begin{digestclaim}[Deleting a set of removable sectors]
Let $\Decl$ satisfy Definition~\ref{def:admissible-declaration} and let
$\SecSet_{\mathrm{del}}\subseteq\SecSet$ be the set of sectors a run
deletes. Suppose \emph{every $s\in\SecSet_{\mathrm{del}}$
carries a witness in the sense of Definition~\ref{def:classes}\,(d) all of
whose targets lie outside $\SecSet_{\mathrm{del}}$}. Then restricting the
infimum \eqref{eq:hlow-def} to
$\{\omega\in\Feas:\omega_s=0\ \text{for all }s\in\SecSet_{\mathrm{del}}\}$
changes neither its value nor its attainment: the whole set
$\SecSet_{\mathrm{del}}$ may be deleted from the variable list, and its
rows from the generated maps, without loss.
\end{digestclaim}

\section*{7\quad The Choi program}\pdfbookmark[1]{7 The Choi program}{digest-group-meat-core/lem:admissible-alice-subspace}

% claim meat-core/lem:admissible-alice-subspace
\digestlabel{lem:admissible-alice-subspace}{7.1}\pdfbookmark[2]{Lemma 7.1 The admissible Alice subspace}{digest-meat-core/lem:admissible-alice-subspace}
\renewcommand*{\RestatedName}{\hyperref{contract.pdf}{}{claim-meat-core/lem:admissible-alice-subspace}{Lemma~7.1}}%
\begin{digestclaim}[The admissible Alice subspace]
Let $\Decl$ be a declaration, and write
\begin{equation*}\tag{\getrefnumber{C-eq:SA}}\digesteqlabel{eq:SA}
  \Sadm\coloneqq
  \begin{cases}
    \operatorname{supp}\tau_A, & \text{trusted-source mode},\\[0.3ex]
    \Hs{A}, & \text{untrusted-source mode},
  \end{cases}
\end{equation*}
the \emph{admissible Alice subspace}; it is determined by the
marginal slot alone and never by $\CList$ or by a particular feasible
point. Then:
\begin{enumerate}[label=(\roman*),leftmargin=2.2em]
\item every $\omega\in\Feas$ satisfies
$\operatorname{supp}\omega\subseteq\Sadm\otimes\Hs{\Qsys}$;
\item the bound~(i) is tight for $\CList=\varnothing$, and for a nonempty
$\CList$ it remains a valid upper bound.
\end{enumerate}
\end{digestclaim}

% claim meat-core/def:support-frame
\digestlabel{def:support-frame}{7.2}\pdfbookmark[2]{Definition 7.2 Support frame and declared isometry}{digest-meat-core/def:support-frame}
\renewcommand*{\RestatedName}{\hyperref{contract.pdf}{}{claim-meat-core/def:support-frame}{Definition~7.2}}%
\begin{digestclaim}[Support frame and declared isometry]
Let the declaration be in trusted-source mode, with declared marginal
$\tau_A$. Let $A_0$ be a register of dimension $d_{A_0}$ and let
$\iota\in\Linn{A_0}{A}$ satisfy
\begin{equation*}\tag{\getrefnumber{C-eq:iota}}\digesteqlabel{eq:iota}
  \iota^{\dagger}\iota=\idon{A_0},
  \qquad
  \iota\,\iota^{\dagger}\,\tau_A=\tau_A .
\end{equation*}
Write $\Pi_{\iota}\coloneqq\iota\iota^{\dagger}$, which is the orthogonal
projector onto $\operatorname{ran}\iota$
($\Pi_{\iota}^{\dagger}=\Pi_{\iota}$ and
$\Pi_{\iota}^{2}=\iota(\iota^{\dagger}\iota)\iota^{\dagger}=\Pi_{\iota}$
by the first condition). The first condition says $\iota$ is an isometry;
the second says its range \emph{contains} the support of $\tau_A$, since
$\Pi_{\iota}\tau_A=\tau_A$ holds exactly when every column of $\tau_A$
lies in $\operatorname{ran}\iota$, i.e.\ when
$\operatorname{ran}\tau_A\subseteq\operatorname{ran}\iota$. We call
$\operatorname{ran}\iota$ the
\emph{support frame}
of the run and $\iota$ the
\emph{declared isometry};
a support frame is given by any $\iota$ satisfying \eqref{eq:iota}. The
frame is
\emph{tight}
when $\operatorname{ran}\iota=\operatorname{supp}\tau_A$, equivalently when
$d_{A_0}=\operatorname{rank}\tau_A$; in general
$\operatorname{rank}\tau_A\le d_{A_0}\le d_A$. The images under $\iota$
of a fixed orthonormal basis of $A_0$ are the
\emph{declared basis of $A_0$}.

Two constructions satisfy \eqref{eq:iota} exactly:
\begin{enumerate}[label=(\alph*),leftmargin=2.2em]
\item \emph{The
  \emph{full frame}.}
  $A_0=A$, $d_{A_0}=d_A$, $\iota=\idon{A}$.
  Then $\iota^{\dagger}\iota=\idon{A}$ and
  $\Pi_{\iota}\tau_A=\idon{A}\,\tau_A=\tau_A$: both conditions of
  \eqref{eq:iota} hold for \emph{any} admissible $\tau_A$, and both hold
  \emph{exactly} as binary64 data, a product with the identity rounding
  nothing. This construction is always available and uses no property of
  $\tau_A$ beyond admission. It is tight exactly when $\tau_A$ is
  nonsingular.
\item \emph{The
  \emph{diagonal frame}.}
  If the declared $\tau_A$ is diagonal in
  the declared basis of $A$, $\tau_A=\sum_it_i\dyad{i}$, then
  $\operatorname{supp}\tau_A=\operatorname{span}\{\ket i:t_i\neq0\}$, and
  one may take $d_{A_0}=\#\{i:t_i\neq0\}$ and
  $\iota=\sum_j\ket{i_j}\bra{j}$, the coordinate injection onto those
  indices in increasing order. Both conditions of \eqref{eq:iota} are then
  arithmetic in $\{0,1\}$ and hold exactly, and the frame is tight.
  Whether $t_i\neq0$ is a bit-level comparison of a declared binary64
  datum against zero and not a threshold: $\tau_A$ is admitted only with
  $\tau_A\succeq0$ and $\tr\tau_A=1$ decided exactly
  (Definition~\ref{def:declaration}, the exact-admissibility paragraph), so
  its diagonal entries \emph{are}
  its eigenvalues and \emph{are} the declared numbers.
  The construction is available only for a $\tau_A$ that is diagonal in the
  declared basis, every off-diagonal entry being the datum $0$.
  Construction~(a) is always available
  and exhibits no tightness at all.
\end{enumerate}
Set
\begin{equation*}\tag{\getrefnumber{C-eq:Wtau}}\digesteqlabel{eq:Wtau}
  W_{\tau}\coloneqq\tau_A^{1/2}\,\iota\ \in\ \Linn{A_0}{A}.
\end{equation*}
\end{digestclaim}

% claim meat-core/lem:choi-marginal
\digestlabel{lem:choi-marginal}{7.3}\pdfbookmark[2]{Lemma 7.3 Fixed-marginal parametrisation on a support frame}{digest-meat-core/lem:choi-marginal}
\renewcommand*{\RestatedName}{\hyperref{contract.pdf}{}{claim-meat-core/lem:choi-marginal}{Lemma~7.3}}%
\begin{digestclaim}[Fixed-marginal parametrisation on a support frame]
Let the declaration be in trusted-source mode, let $\iota$ be the
declared isometry of a support frame, satisfying
\eqref{eq:iota}, and let $W_{\tau}$ be \eqref{eq:Wtau}
(Definition~\ref{def:support-frame}). Then:
\begin{enumerate}[label=(\roman*),leftmargin=2.2em]
\item $\Pi_{\iota}\tau_A^{1/2}=\tau_A^{1/2}$.
\item For every support frame, tight or not,
\begin{equation*}\tag{\getrefnumber{C-eq:WWtau}}\digesteqlabel{eq:WWtau}
  W_{\tau}W_{\tau}^{\dagger}
  =\tau_A^{1/2}\,\Pi_{\iota}\,\tau_A^{1/2}
  =\tau_A^{1/2}\,\tau_A^{1/2}=\tau_A .
\end{equation*}
\item $\operatorname{ran}W_{\tau}=\operatorname{supp}\tau_A$, and the
Moore--Penrose inverse $W_{\tau}^{+}$ satisfies
\begin{equation*}\tag{\getrefnumber{C-eq:Wplus}}\digesteqlabel{eq:Wplus}
  W_{\tau}W_{\tau}^{+}=\Pi_{\operatorname{supp}\tau_A},
  \qquad
  W_{\tau}^{+}W_{\tau}=P_{\iota}
  \coloneqq\Pi_{(\ker W_{\tau})^{\perp}} ,
\end{equation*}
where $\ker W_{\tau}=\{\xi\in\Hs{A_0}:
\iota\xi\in(\operatorname{supp}\tau_A)^{\perp}\}$ has dimension
$d_{A_0}-\operatorname{rank}\tau_A$. The frame is tight if and only if
$\ker W_{\tau}=0$, i.e.\ $P_{\iota}=\idon{A_0}$, and then $W_{\tau}$ has
full column rank and
$W_{\tau}^{+}=(W_{\tau}^{\dagger}W_{\tau})^{-1}W_{\tau}^{\dagger}$.
\item For any finite-dimensional system $B'$, the map
$\chi_{B'}(J)\coloneqq(W_{\tau}\otimes\idon{B'})\,J\,
(W_{\tau}^{\dagger}\otimes\idon{B'})$ maps
$\{J\in\Pos{A_0B'}:\tr_{B'}J=\idon{A_0}\}$ \emph{onto}
$\{\omega\in\Pos{AB'}:\tr_{B'}\omega=\tau_A\}$. It is \emph{injective},
hence a bijection with inverse
$\omega\mapsto(W_{\tau}^{+}\otimes\idon{B'})\,\omega\,
(W_{\tau}^{+\dagger}\otimes\idon{B'})$, when the support frame is tight,
and, if $\dim\Hs{B'}\ge2$, only then; if $\dim\Hs{B'}=1$ its domain is the
single point $\idon{A_0}$. In general $\chi_{B'}(J)=\chi_{B'}(J')$ if and only if
$(P_{\iota}\otimes\idon{B'})J(P_{\iota}\otimes\idon{B'})
=(P_{\iota}\otimes\idon{B'})J'(P_{\iota}\otimes\idon{B'})$.
\end{enumerate}
\end{digestclaim}

% claim meat-core/lem:choi-factor
\digestlabel{lem:choi-factor}{7.4}\pdfbookmark[2]{Lemma 7.4 Fixed-marginal parametrisation with any factor}{digest-meat-core/lem:choi-factor}
\renewcommand*{\RestatedName}{\hyperref{contract.pdf}{}{claim-meat-core/lem:choi-factor}{Lemma~7.4}}%
\begin{digestclaim}[Fixed-marginal parametrisation with any factor]
Let $\Decl$ satisfy Definition~\ref{def:admissible-declaration} in trusted-source mode, in the
general format of Definition~\ref{def:declaration}.
\begin{enumerate}[label=(\roman*),leftmargin=2.2em]
\item \emph{(any factor serves)} Let $\tau\in\Pos{A'}$ on a
finite-dimensional system $A'$, of any trace, and let $W\in\Linn{A_0}{A'}$
satisfy $WW^{\dagger}=\tau$. Let $P_W$ be the projector onto
$(\ker W)^{\perp}$. For every finite-dimensional $B'$ the map
$\chi^{W}_{B'}(J)\coloneqq(W\otimes\idon{B'})J(W^{\dagger}\otimes\idon{B'})$
maps $\{J\in\Pos{A_0B'}:\tr_{B'}J=\idon{A_0}\}$ onto
$\{\omega\in\Pos{A'B'}:\tr_{B'}\omega=\tau\}$; it is injective when
$\ker W=0$, and, if $\dim\Hs{B'}\ge2$, only then; and
$\chi^{W}_{B'}(J)=\chi^{W}_{B'}(J')$ exactly when
$(P_W\otimes\idon{})J(P_W\otimes\idon{})=(P_W\otimes\idon{})J'(P_W\otimes\idon{})$.
The operator $W_{\tau}$ of \eqref{eq:Wtau} is such a factor of $\tau_A$ for
every admissible frame.
\item \emph{(per Alice sector)} For every $t\in\mathbb{T}$ let $W_t\in
\Linn{A_{0,t}}{A_t}$ satisfy $W_tW_t^{\dagger}=\tau^{(t)}_A$, and put
$A_0\coloneqq\bigoplus_tA_{0,t}$, $W\coloneqq\bigoplus_tW_t$. Then
$WW^{\dagger}=\tau_A$, and the assignment
$(J_{t,s})\mapsto\big(\omega^{(t,s)}\big)$,
$\omega^{(t,s)}\coloneqq(W_t\otimes\idon{\Qsys_s})J_{t,s}
(W_t^{\dagger}\otimes\idon{\Qsys_s})$, maps
\begin{equation*}\tag{\getrefnumber{C-eq:choi-coupling-t}}\digesteqlabel{eq:choi-coupling-t}
  \Big\{(J_{t,s})_{(t,s)}:\ J_{t,s}\in\Pos{A_{0,t}\Qsys_s},\quad
  \sum_{s\in\SecSet}\tr_{\Qsys_s}J_{t,s}=\idon{A_{0,t}}
  \ \ \forall t\in\mathbb{T}\Big\}
\end{equation*}
onto the set of lists over the joint sectors with positive semidefinite
blocks and $\sum_s\tr_{\Qsys_s}\omega^{(t,s)}=\tau^{(t)}_A$ for every $t$,
bijectively when $\ker W_t=0$ for every $t$, and, if
$\dim\Hs{\Qsys}\ge2$, only then. The set
\eqref{eq:choi-coupling-t} is a Cartesian product over $t$. If every
$\iota_t$ satisfies \eqref{eq:iota} for $\tau^{(t)}_A$ on $A_t$, then
$\iota\coloneqq\bigoplus_t\iota_t$ satisfies \eqref{eq:iota} for $\tau_A$ and
$W_{\tau}=\bigoplus_t(\tau^{(t)}_A)^{1/2}\iota_t$.
\item \emph{(the objective reads Alice's side through $N_x$)} For a factor
$W=\bigoplus_tW_t$ as in (ii) put
\begin{equation*}\tag{\getrefnumber{C-eq:NfromW}}\digesteqlabel{eq:NfromW}
  N^{W}_x\coloneqq W^{\dagger}M_xW=\bigoplus_tW_t^{\dagger}M^{(t)}_xW_t
  \in\Pos{A_0},
\end{equation*}
let $K^{N}_{\kappa}$ be \eqref{eq:kraus} with $\sqrt{N^{W}_x}$ in place of
$\sqrt{M_x}$, an operator
$\Hs{A_0\Qsys}\to\Hs{R\,XY\,C_AC_B\,A_0\Qsys}$, and let $\cG^{N}_{\hat c}$ be
\eqref{eq:GTheta} with $K^{N}_{\kappa}$. Then for every $J\in\Pos{A_0\Qsys}$,
every $(x,v)$ and every $\hat c$, with
$\omega\coloneqq\chi^{W}_{\Qsys}(J)$,
\begin{equation*}\tag{\getrefnumber{C-eq:choi-N}}\digesteqlabel{eq:choi-N}
  \tr[(M_x\otimes F_v)\,\omega]=\tr[(N^{W}_x\otimes F_v)\,J],\qquad
  \tr\cG_{\hat c}(\omega)=\tr\cG^{N}_{\hat c}(J),\qquad
  T\big(\cG_{\hat c}(\omega)\big)=T\big(\cG^{N}_{\hat c}(J)\big),
\end{equation*}
$T$ the functional \eqref{eq:Tcone}; hence $P_{\hat c}$, $K_{\hat c}$ and
$O_{\alpha}$ at $\omega$ are functions of $J$, $\{N^{W}_x\}$, $\{F_v\}$ and
the classical maps alone, and the same holds sector by sector for the
joint-sector quantities of Definition~\ref{def:joint-sectors}.
\item \emph{(ensemble form)} On the full frame, $A_{0,t}=A_t$ and
$W_t=(\tau^{(t)}_A)^{1/2}$, one has
$N^{W}_x=\tau_A^{1/2}M_x\tau_A^{1/2}=N_x$, the operators of
Definition~\ref{def:declaration}\,(DX4): for an ensemble-form declaration
they are the \emph{declared} numbers. So the program of an ensemble-form
declaration is the program \eqref{eq:choi-coupling-t}, (iii) in the declared
$\{N_x\}$ and $\{F_v\}$, with every constraint row in the conjugated form of
\eqref{eq:choi-feas}; it is defined, and equals $\hlow$ of the denoted
declaration, at every rank of every $\tau^{(t)}_A$, and no rank is computed.
\end{enumerate}
\end{digestclaim}

% claim meat-core/thm:choi
\digestlabel{thm:choi}{7.5}\pdfbookmark[2]{Theorem 7.5 The Choi program}{digest-meat-core/thm:choi}
\renewcommand*{\RestatedName}{\hyperref{contract.pdf}{}{claim-meat-core/thm:choi}{Theorem~7.5}}%
\begin{digestclaim}[The Choi program]
Let the declaration be in trusted-source mode, fix a
support frame (Definition~\ref{def:support-frame}), apply
Lemma~\ref{lem:choi-marginal} with $B'=\Qsys$, and restrict to
sector-block-diagonal operators. Then the assignment
$J=\bigoplus_{s}J_s\mapsto\omega=\chi_{\Qsys}(J)$, with
$\omega_s=\chi_{\Qsys_s}(J_s)\coloneqq
(W_{\tau}\otimes\idon{\Qsys_s})J_s(W_{\tau}^{\dagger}\otimes\idon{\Qsys_s})$,
maps
\begin{equation*}\tag{\getrefnumber{C-eq:choi-feas}}\digesteqlabel{eq:choi-feas}
  \Feas^{J}\coloneqq\Big\{\!(J_s)_{s\in\SecSet}:\ J_s\in\Pos{A_0\Qsys_s},\
  \sum_{s\in\SecSet}\tr_{\Qsys_s}J_s=\idon{A_0},\
  W_{\tau}\Big(\textstyle\sum_s
  \tr_{\Qsys_s}[(\idon{A_0}\otimes\hat O^{(s)}_k)J_s]\Big)
  W_{\tau}^{\dagger}\bowtie_k\Theta_k\ \forall k\!\Big\}
\end{equation*}
\emph{onto} $\Feas$ of \eqref{eq:feas}, bijectively when the support frame
is tight (Lemma~\ref{lem:choi-marginal}), and, if $\Feas\ne\varnothing$ and
$\dim\Hs{\Qsys}\ge2$, only then. Writing
$\Phimap_{J}[J]\coloneqq\Phimap[\chi_{\Qsys}(J)]$,
$\cG^{J}_{\hat c,s}\coloneqq\cG_{\hat c,s}\circ\chi_{\Qsys_s}$,
$f^{J}_{\hat c,s}\coloneqq f_{\hat c,s}\circ\chi_{\Qsys_s}$,
$P^{J}_{\hat c}(J)\coloneqq\Phimap_J[J](\hat c)$ and
$K^{J}_{\hat c}(J)\coloneqq P^{J}_{\hat c}(J)
+\sum_{s\in\SecSet_K(\hat c)}f^{J}_{\hat c,s}(J_s)$, the value
\eqref{eq:hlow-def} equals
\begin{equation*}\tag{\getrefnumber{C-eq:choi-program}}\digesteqlabel{eq:choi-program}
\boxed{\ \
\begin{aligned}
  \hlow=\inf_{\mathbf q,\ (J_s)_s}\quad
  &\frac{\alpha}{\alpha-1}\sum_{\hat c\notin\Chkey}
  q(\hat c)\log_2\frac{q(\hat c)}{P^{J}_{\hat c}(J)}
  +\sum_{\hat c\in\Chkey}\left[
  q(\hat c)\log_2\frac{q(\hat c)}{P^{J}_{\hat c}(J)}
  +\frac{q(\hat c)}{\alpha-1}
  \log_2\frac{q(\hat c)}{K^{J}_{\hat c}(J)}\right]\\
  \text{s.t.}\quad
  &\mathbf q\in\acSet,\qquad (J_s)_s\in\Feas^{J},
\end{aligned}\ \ }
\end{equation*}
an \emph{exact} reformulation (not a relaxation), jointly convex in
$(\mathbf q,(J_s)_s)$, with all constraints affine or conic.

\emph{In the general format} of Definition~\ref{def:declaration}, let the
frame be a direct sum $\iota=\bigoplus_{t\in\mathbb{T}}\iota_t$ with every
$\iota_t\in\Linn{A_{0,t}}{A_t}$ satisfying \eqref{eq:iota} for
$\tau^{(t)}_A$ --- either construction of Definition~\ref{def:support-frame}
applied sector by sector --- and put $W_t\coloneqq(\tau^{(t)}_A)^{1/2}\iota_t$. Let $\Feas^{J}$
be the set of lists $(J_{t,s})$ over the joint sectors with
$J_{t,s}\in\Pos{A_{0,t}\Qsys_s}$, the coupling constraints
$\sum_s\tr_{\Qsys_s}J_{t,s}=\idon{A_{0,t}}$ for every $t$, and for every
entry the row of \eqref{eq:feas-general} evaluated at
$\omega^{(t,s)}=(W_t\otimes\idon{})J_{t,s}(W_t^{\dagger}\otimes\idon{})$: for
an entry read through its lifted operator,
$W_t\big(\sum_s\tr_{\Qsys_s}[(\idon{A_{0,t}}\otimes\hat O^{\uparrow(s)}_k)
J_{t,s}]\big)W_t^{\dagger}\bowtie_k\Theta^{(t)}_k$ for every $t$, and for a
party-local entry the last line of \eqref{eq:feas-general} at $J$,
conjugated by $W_t\otimes\idon{}$. Then $(J_{t,s})\mapsto(\omega^{(t,s)})$
maps $\Feas^{J}$ onto $\Feas$ of \eqref{eq:feas-general}, bijectively when
every sector's frame is tight, and, if $\Feas\ne\varnothing$ and
$\dim\Hs{\Qsys}\ge2$, only then, and \eqref{eq:choi-program} holds with
$K^{J}_{\hat c}(J)=P^{J}_{\hat c}(J)+\sum_{(t,s)\in\mathbb{J}_K(\hat c)}
f_{\hat c,(t,s)}(\omega^{(t,s)})$. Every constraint of $\Feas^{J}$ involves
the blocks of one Alice sector only, so $\Feas^{J}=\prod_t\Feas^{J}_t$, and a
linear function of $(J_{t,s})$ is minimised over $\Feas^{J}$ by minimising,
independently for each $t$, its part on the blocks of $t$ over
$\Feas^{J}_t$.
\end{digestclaim}

% claim meat-core/prop:frame-free-rows
\digestlabel{prop:frame-free-rows}{7.6}\pdfbookmark[2]{Proposition 7.6 Frame-free constraint rows on a tight frame}{digest-meat-core/prop:frame-free-rows}
\renewcommand*{\RestatedName}{\hyperref{contract.pdf}{}{claim-meat-core/prop:frame-free-rows}{Proposition~7.6}}%
\begin{digestclaim}[Frame-free constraint rows on a tight frame]
In (i)--(iii) let the declaration be in trusted-source mode, let $W_{\tau}$ be
\eqref{eq:Wtau} for a support frame (Definition~\ref{def:support-frame}),
and for an entry $k$ of $\CList$ and $(J_s)_{s\in\SecSet}$ with
$J_s\in\Pos{A_0\Qsys_s}$ write
$Y_k\coloneqq\sum_s\tr_{\Qsys_s}[(\idon{A_0}\otimes\hat O^{(s)}_k)J_s]$.
For a homogeneous entry, whose right-hand side is
$\Theta_k=c_k\tau_A$ (Definition~\ref{def:homog}\,(a)), call the $k$-th row
of \eqref{eq:choi-feas}, $W_{\tau}Y_kW_{\tau}^{\dagger}\bowtie_kc_k\tau_A$,
its \emph{conjugated form}, and
\begin{equation*}\tag{\getrefnumber{C-eq:frame-free-row}}\digesteqlabel{eq:frame-free-row}
  Y_k\ \bowtie_k\ c_k\,\idon{A_0}
\end{equation*}
its \emph{frame-free form}.
\begin{enumerate}[label=(\roman*),leftmargin=2.2em]
\item If $\ker W_{\tau}=0$, that is, if the frame is tight, the
frame-free form of a homogeneous row holds for $(J_s)_s$ if and only if its
conjugated form does; so imposing some homogeneous rows in the frame-free
form leaves $\Feas^{J}$ unchanged.
\item On every support frame the frame-free form of a homogeneous row
implies its conjugated form. So the set obtained from $\Feas^{J}$ by
imposing some homogeneous rows in the frame-free form is contained in
$\Feas^{J}$, and the infimum in \eqref{eq:choi-program} taken over that set
in place of $\Feas^{J}$ is at least $\hlow$.
\item In the general format, with $\iota=\bigoplus_t\iota_t$ and $W_t$ as in
Theorem~\ref{thm:choi}:
\begin{enumerate}[label=(\alph*),leftmargin=2.2em]
\item for a homogeneous entry read through its lifted operator and
every $t\in\mathbb{T}$, (i) and (ii) hold for the row of sector $t$, with
$W_t$ and $A_{0,t}$ in place of $W_{\tau}$ and $A_0$ and with the
frame-free form
$\sum_s\tr_{\Qsys_s}[(\idon{A_{0,t}}\otimes\hat O^{\uparrow(s)}_k)
J_{t,s}]\bowtie_kc_k\,\idon{A_{0,t}}$;
\item for a party-local entry, with $Y_k$ the expression of the last
line of \eqref{eq:feas-general} at $J$ in place of $\omega$, the
unconjugated row $Y_k\bowtie_k0$ is equivalent to the conjugated row
$(W_t\otimes\idon{})\,Y_k\,(W_t^{\dagger}\otimes\idon{})\bowtie_k0$ if
$\ker W_t=0$, and implies it on every frame.
\end{enumerate}
\item In untrusted-source mode there is no conjugation to remove: every
row is a condition on the state variable as written.
\end{enumerate}
\end{digestclaim}

% claim meat-core/def:untrusted-choi
\digestlabel{def:untrusted-choi}{7.10}\pdfbookmark[2]{Definition 7.10 The program in untrusted-source mode}{digest-meat-core/def:untrusted-choi}
\renewcommand*{\RestatedName}{\hyperref{contract.pdf}{}{claim-meat-core/def:untrusted-choi}{Definition~7.10}}%
\begin{digestclaim}[The program in untrusted-source mode]
In untrusted-source mode the Choi step is the identity. This is the
degenerate case of Theorem~\ref{thm:choi}, not a second program:
\begin{enumerate}[label=(\alph*),leftmargin=2.2em]
\item \emph{The parametrization.} $\chi_{B'}=\idmap$, $J=\omega$,
$J_s=\omega_s$; no $A_0$, no $\iota$, no $W_{\tau}$.
\item \emph{The variables and constraints.} $\Feas$ of
\eqref{eq:feas} is already the feasible set, and by
\eqref{eq:feas-modes} it reads
\begin{equation*}\tag{\getrefnumber{C-eq:choi-feas-un}}\digesteqlabel{eq:choi-feas-un}
  \Feas=\Big\{(\omega_s)_{s\in\SecSet}:\
  \omega_s\in\Pos{A\Qsys_s},\
  \sum_{s\in\SecSet}\tr\omega_s=1,\
  \sum_{s\in\SecSet}
  \tr\big[(\idon{A}\otimes\hat O^{(s)}_k)\,\omega_s\big]
  \bowtie_k c_k\ \forall k\Big\},
\end{equation*}
so the coupling constraint that binds the sector blocks is the single
\emph{scalar} equation $\sum_s\tr\omega_s=1$, in place of the operator
equation $\sum_s\tr_{\Qsys_s}J_s=\idon{A_0}$ of
\eqref{eq:choi-feas}.
\item \emph{The objective.} \eqref{eq:choi-program} holds verbatim with
$J_s$ replaced by $\omega_s$ and $P^{J}_{\hat c}$, $f^J_{\hat c,s}$,
$K^{J}_{\hat c}$ by their state-form originals
$P_{\hat c}$, $f_{\hat c,s}$, $K_{\hat c}$ of \eqref{eq:fK}: the
substitution $\omega_s=\chi_{\Qsys_s}(J_s)$ that
Theorem~\ref{thm:choi} performs is the identity here. Hence the value
is $\hlow$ of \eqref{eq:hlow-def} by definition rather than by a
theorem, the program is jointly convex by
Theorem~\ref{thm:convexity} directly, and the feasible set is convex
and compact by Definition~\ref{def:feasible}.
\item \emph{Consistency at $d_A=1$.} There trusted-source mode admits only
$\tau_A=[1]$, the one state on a one-dimensional $A$, and its Choi
parametrization degenerates: $d_{A_0}=\operatorname{rank}\tau_A=1$,
$\iota=[1]$, $W_{\tau}=[1]$, $\chi_{B'}=\idmap$ and $J_s=\omega_s$. So
$\idon{A_0}=[1]$ and $\sum_s\tr_{\Qsys_s}J_s=\idon{A_0}$ \emph{is}
$\sum_s\tr\omega_s=1$, the constraint rows agree because every $1\times1$
Hermitian $\Theta_k$ is $c_k\tau_A$ with $c_k=\Theta_k$
(Proposition~\ref{prop:trivial-alice-system}\,(iii)), and the two feasible
sets of \eqref{eq:feas} coincide
(Proposition~\ref{prop:trivial-alice-system}\,(i)): \eqref{eq:choi-feas}
and \eqref{eq:choi-feas-un} are the same set, and the two modes give
literally the same program.
\end{enumerate}
\end{digestclaim}

% claim meat-core/lem:estimate
\digestlabel{lem:estimate}{7.11}\pdfbookmark[2]{Lemma 7.11 Upper estimate of the value at any feasible point}{digest-meat-core/lem:estimate}
\renewcommand*{\RestatedName}{\hyperref{contract.pdf}{}{claim-meat-core/lem:estimate}{Lemma~7.11}}%
\begin{digestclaim}[Upper estimate of the value at any feasible point]
Let $\Decl$ satisfy Definition~\ref{def:admissible-declaration}, let $\acSet$ be any set
of probability vectors on the acceptance alphabet, and let
$\alpha\in(1,\infty)$.
\begin{enumerate}[label=(\roman*),leftmargin=2.2em]
\item \emph{(State form.)} For every $(\mathbf q,\omega)\in\acSet\times\Feas$,
\begin{equation*}\tag{\getrefnumber{C-eq:estimate}}\digesteqlabel{eq:estimate}
  O_{\alpha}(\mathbf q,\omega)\;\ge\;\hlow .
\end{equation*}
\item \emph{(Choi form.)} Write
$O^{J}_{\alpha}(\mathbf q,J)\coloneqq O_{\alpha}(\mathbf q,\chi_{\Qsys}(J))$
for the objective of the Choi program \eqref{eq:choi-program} of
Theorem~\ref{thm:choi}, with $\chi_{\Qsys}$ the map of that theorem. For
every $x=(\mathbf q,J)\in\acSet\times\Feas^{J}$,
\begin{equation*}\tag{\getrefnumber{C-eq:estimate-choi}}\digesteqlabel{eq:estimate-choi}
  O^{J}_{\alpha}(x)\;\ge\;\hlow .
\end{equation*}
\end{enumerate}
\end{digestclaim}

\section*{8\quad Gradients and the limit of order one}\pdfbookmark[1]{8 Gradients and the limit of order one}{digest-group-meat-core/def:branch-support}

% claim meat-core/def:branch-support
\digestlabel{def:branch-support}{8.1}\pdfbookmark[2]{Definition 8.1 Maximal branch support}{digest-meat-core/def:branch-support}
\renewcommand*{\RestatedName}{\hyperref{contract.pdf}{}{claim-meat-core/def:branch-support}{Definition~8.1}}%
\begin{digestclaim}[Maximal branch support]
For $\hat c\in\Chkey$ and $s\in\SecSet$ the
\emph{maximal branch support}
is the subspace
\begin{equation*}\tag{\getrefnumber{C-eq:Wcs}}\digesteqlabel{eq:Wcs}
  W_{\hat c,s}\coloneqq
  \sum_{\kappa\in\Kann_{\mathrm{pass}}(\hat c)}
  K^{(s)}_{\kappa}\big(\Sadm\otimes\Hs{\Qsys_s}\big)
  \ \subseteq\ \Hs{R\,XY\,C_AC_B\,A\Qsys_s},
\end{equation*}
the span of the images, under the Kraus operators of $\cG_{\hat c,s}$
(Definition~\ref{def:branch}), of the admissible Alice subspace
$\Sadm$ of \eqref{eq:SA} tensored with the sector, that is, of the
subspace the declaration admits: $\operatorname{supp}\tau_A$ in
trusted-source mode, all of $\Hs{A}$ in
untrusted-source mode.

Listing the operators of \eqref{eq:Wcs} side by side, with $W_{\tau}$ of
\eqref{eq:Wtau} on the support frame of
Definition~\ref{def:support-frame}, put
\begin{equation*}\tag{\getrefnumber{C-eq:Kstack}}\digesteqlabel{eq:Kstack}
  \mathbf K_{\hat c,s}\coloneqq
  \Big[\ K^{(s)}_{\kappa}\,(W_{\tau}\otimes\idon{\Qsys_s})\ \Big]
  _{\kappa\in\Kann_{\mathrm{pass}}(\hat c)}
  \;:\;\bigoplus_{\kappa\in\Kann_{\mathrm{pass}}(\hat c)}\Hs{A_0\Qsys_s}
  \;\longrightarrow\;\Hs{R\,XY\,C_AC_B\,A\Qsys_s},
\end{equation*}
with $W_{\tau}\otimes\idon{\Qsys_s}$ replaced by $\idon{A\Qsys_s}$, and
the domain by $\bigoplus_{\kappa}\Hs{A\Qsys_s}$, in untrusted-source mode.
The matrix $\mathbf K_{\hat c,s}$ depends on the declaration and on
$W_{\tau}$ alone and on no iterate.
\end{digestclaim}

% claim meat-core/lem:branch-support
\digestlabel{lem:branch-support}{8.2}\pdfbookmark[2]{Lemma 8.2 Properties of the maximal branch support}{digest-meat-core/lem:branch-support}
\renewcommand*{\RestatedName}{\hyperref{contract.pdf}{}{claim-meat-core/lem:branch-support}{Lemma~8.2}}%
\begin{digestclaim}[Properties of the maximal branch support]
Let $\hat c\in\Chkey$ and $s\in\SecSet$, and let $W_{\hat c,s}$ and
$\mathbf K_{\hat c,s}$ be as in Definition~\ref{def:branch-support}.
\begin{enumerate}[label=(\roman*),leftmargin=2.2em]
\item In trusted-source mode every $\omega_s=\chi_{\Qsys_s}(\Delta)$ with
  $\Delta\in\Pos{A_0\Qsys_s}$, in particular every feasible $\omega_s$
  and the image of the positive and the negative part of every Hermitian
  direction, satisfies
  $\operatorname{supp}\omega_s
  =(W_{\tau}\otimes\idon{})(\operatorname{supp}\Delta)
  \subseteq\operatorname{supp}\tau_A\otimes\Hs{\Qsys_s}
  =\Sadm\otimes\Hs{\Qsys_s}$; in untrusted-source mode
  $\Sadm=\Hs{A}$ and the inclusion holds for every
  $\omega_s\in\Pos{A\Qsys_s}$. In either mode, therefore,
  \begin{equation*}\tag{\getrefnumber{C-eq:Wcs-bound}}\digesteqlabel{eq:Wcs-bound}
    \operatorname{supp}\cG_{\hat c,s}(\omega_s)\subseteq W_{\hat c,s},
  \end{equation*}
  with equality whenever $\omega_s$ attains that maximal support,
  $\operatorname{supp}\omega_s=\Sadm\otimes\Hs{\Qsys_s}$. In
  trusted-source mode $J_s\succ0$ for the Choi block suffices for it, and
  on a tight support frame is equivalent to it; in untrusted-source mode it
  is $\omega_s\succ0$ on $\Hs{A\Qsys_s}$, which is $J_s\succ0$, as
  $J_s=\omega_s$ there (Definition~\ref{def:untrusted-choi}\,(a)).
\item $W_{\hat c,s}$ is the range of one matrix, and its dimension is
  bounded:
  \begin{equation*}\tag{\getrefnumber{C-eq:Kstack-rank}}\digesteqlabel{eq:Kstack-rank}
    W_{\hat c,s}=\operatorname{ran}\mathbf K_{\hat c,s},
    \qquad
    \dim W_{\hat c,s}=\operatorname{rank}\mathbf K_{\hat c,s}
    \;\le\;\big|\Kann_{\mathrm{pass}}(\hat c)\big|\cdot
    \dim\Sadm\cdot d_s ,
  \end{equation*}
  where $\dim\Sadm=\operatorname{rank}\tau_A$ in trusted-source mode and
  $d_A$ in untrusted-source mode.
\item $\dim W_{\hat c,s}$ is \emph{strictly} smaller than
  $\dim\Hs{R\,XY\,C_AC_B\,A\Qsys_s}$, so $0$ always lies in the ambient
  spectrum of a branch operator $X_{\hat c,s}=\cG_{\hat c,s}(\omega_s)$ of~(i)
  and ``maximal support'' is never the statement that $X_{\hat c,s}\succ0$
  on the ambient space.
\item $W_{\hat c,s}$ is \emph{not} in general $\cZ$-invariant.
\end{enumerate}
\end{digestclaim}

% claim meat-core/thm:gradient
\digestlabel{thm:gradient}{8.4}\pdfbookmark[2]{Theorem 8.4 Gradient of the pinched trace power on a face}{digest-meat-core/thm:gradient}
\renewcommand*{\RestatedName}{\hyperref{contract.pdf}{}{claim-meat-core/thm:gradient}{Theorem~8.4}}%
\begin{digestclaim}[Gradient of the pinched trace power on a face]
Let $\alpha\in(1,\infty)$, $p=1/\alpha$, let $T$ be the pinched trace
power \eqref{eq:Tcone} on the positive semidefinite cone of a branch space
(Lemma~\ref{C-lem:dk}), and let $X\succ0$. Then $T$ is Fr\'echet
differentiable at $X$ with gradient
\begin{equation*}\tag{\getrefnumber{C-eq:gradT}}\digesteqlabel{eq:gradT}
  \nabla T(X)
  =\alpha\,D X^{p}\Big[\cZ\big(Y^{\alpha-1}\big)\Big],
  \qquad Y\coloneqq\cZ(X^{p}),
\end{equation*}
i.e.\ $\frac{\mathrm d}{\mathrm dt}T(X+tH)\big|_0=\tr[\nabla T(X)H]$,
where $DX^{p}[\cdot]$ is the self-adjoint map \eqref{C-eq:dk} with
$f(x)=x^{p}$ and $Y^{\alpha-1}$ is defined by functional calculus
($0^{\alpha-1}\coloneqq0$; $Y$ may be singular). Explicitly, in the
eigenbasis $X=\sum_i\lambda_iP_i$:
\begin{equation*}\tag{\getrefnumber{C-eq:gradT-kernel}}\digesteqlabel{eq:gradT-kernel}
  \nabla T(X)=\alpha\sum_{i,j}
  \frac{\lambda_i^{p}-\lambda_j^{p}}{\lambda_i-\lambda_j}\,
  P_i\,\cZ\big(Y^{\alpha-1}\big)\,P_j
  \qquad(\text{diagonal kernel }p\lambda_i^{p-1}).
\end{equation*}
Moreover:
\begin{enumerate}[label=(\roman*),leftmargin=2.2em]
\item \emph{(faces)} Let $X\succeq0$, write
$S\coloneqq\operatorname{supp}X$, and let $H$ be Hermitian with
$\operatorname{supp}H\subseteq S$. Then $t\mapsto T(X+tH)$ is
differentiable at $t=0$, and \eqref{eq:gradT}--\eqref{eq:gradT-kernel}
hold in the following \emph{split} form:
\begin{enumerate}[label=(\alph*),leftmargin=2.2em]
\item the divided-difference map $DX^{p}[\cdot]$ is compressed to $S$,
equivalently, the sum in \eqref{eq:gradT-kernel} runs over the strictly
positive eigenvalues of $X$ only, where $X$ is positive definite and the
kernel $f^{[1]}$ is finite;
\item $Y=\cZ(X^{p})$, $Y^{\alpha-1}$ and $\cZ(Y^{\alpha-1})$ are formed on
the \emph{ambient} branch space and are \emph{not} compressed to $S$.
\end{enumerate}
The resulting $\nabla T(X)$ determines the derivative in every direction
supported in $S$; its components outside $S$ are neither determined by
this nor needed, support-enlarging directions being the subject of (ii).
Compressing the \emph{whole} computation to a subspace $S'$ is licensed
when $\sum_{j\in\Zset}Z_jS\subseteq S'$, in particular when $S\subseteq S'$
and $S'$ is $\cZ$-invariant ($Z_jS'\subseteq S'$ for every $j\in\Zset$):
then the ambient complement contributes $0$ to $T$, and the compressed and
ambient computations agree. Compressing to $S$ itself is licensed
\emph{only} when $S$ is $\cZ$-invariant; otherwise $\cZ$ transports weight
out of $S$ and the compressed value of $T$ is strictly smaller (the proof
gives the argument and a rank-one witness). For the branch operators $X=\cG_{\hat c,s}(\omega_s)$, whose
faces are in general not $\cZ$-invariant, the relevant face is
$S=W_{\hat c,s}$ of \eqref{eq:Wcs} at maximal support, and (a)--(b) is
what makes the gradient available there.
\item \emph{(one-sided derivatives)} For $X\succeq0$ and any Hermitian
$H$ with $X+tH\succeq0$ for small $t>0$, the one-sided directional
derivative $\lim_{t\downarrow0}\frac{T(X+tH)-T(X)}{t}$ exists in
$(-\infty,+\infty]$, and can equal $+\infty$ only when $H$ strictly
enlarges the support of $X$.
\end{enumerate}
\end{digestclaim}

% claim meat-core/cor:full-gradient
\digestlabel{cor:full-gradient}{8.5}\pdfbookmark[2]{Corollary 8.5 Gradient of the objective along feasible directions}{digest-meat-core/cor:full-gradient}
\renewcommand*{\RestatedName}{\hyperref{contract.pdf}{}{claim-meat-core/cor:full-gradient}{Corollary~8.5}}%
\begin{digestclaim}[Gradient of the objective along feasible directions]
Let $\alpha\in(1,\infty)$ and let $x=(\mathbf q,(J_s)_s)$ be a point with
$\mathbf q\in\acSet$ and $J_s\in\Pos{A_0\Qsys_s}$ for every
$s\in\SecSet$ --- per-block positive semidefiniteness of the $J$ factor
and nothing else: no normalization, no coupling condition
$\sum_s\tr_{\Qsys_s}J_s=\idon{A_0}$ and no declared row $\bowtie_k$ is
assumed of it, so that every point of $\acSet\times\Feas^{J}(\varepsilon)$
qualifies, at every $\varepsilon\ge0$ and the exact $\Feas^{J}$ included
(Definition~\ref{def:relaxed-feasible-set}: per-block positive
semidefiniteness is the one condition that family does not relax) ---
and such that $q(\hat c)>0$, $P^{J}_{\hat c}>0$ and $K^{J}_{\hat c}>0$
for all $\hat c$, and such that each branch operator
$X_{\hat c,s}\coloneqq\cG^{J}_{\hat c,s}(J_s)$
($\hat c\in\Chkey$, $s\in\SecSet_K(\hat c)$) has \emph{maximal support},
$\operatorname{supp}X_{\hat c,s}=W_{\hat c,s}$ of \eqref{eq:Wcs}, for
which $J_s\succ0$ suffices, so that Theorem~\ref{thm:gradient}\,(i)
applies on that face. The displayed formulas are written for
trusted-source mode; the untrusted-source-mode statement is at the end of
this corollary.

Then the objective $O^{J}_{\alpha}$ of \eqref{eq:choi-program} is
differentiable at that point \emph{along every feasible direction}, and
that is the whole of the conclusion: for every
$y\in\acSet\times\Feas^{J}$ the one-variable map
$t'\mapsto O^{J}_{\alpha}\big(x+t'(y-x)\big)$ is differentiable at
$t'=0$, with derivative the pairing of $y-x$ against the gradient
\begin{align*}
  \frac{\partial O^{J}_{\alpha}}{\partial q(\hat c)}
  &=\begin{cases}
  \dfrac{\alpha}{\alpha-1}\Big(\log_2\dfrac{q(\hat c)}{P^{J}_{\hat c}}
  +\log_2\mathrm e\Big), & \hat c\notin\Chkey,\\[1.2ex]
  \log_2\dfrac{q(\hat c)}{P^{J}_{\hat c}}+\log_2\mathrm e
  +\dfrac{1}{\alpha-1}\Big(\log_2\dfrac{q(\hat c)}{K^{J}_{\hat c}}
  +\log_2\mathrm e\Big), & \hat c\in\Chkey,
  \end{cases}\tag{\getrefnumber{C-eq:grad-q}}\digesteqlabel{eq:grad-q}\\[0.6ex]
  \nabla_{J_{s'}}O^{J}_{\alpha}
  &=-\log_2\mathrm e\!\!\sum_{\hat c\notin\Chkey}\!
  \frac{\alpha}{\alpha-1}\,\frac{q(\hat c)}{P^{J}_{\hat c}}\,
  \Gamma^{(s')}_{\hat c}
  \;-\;\log_2\mathrm e\!\!\sum_{\hat c\in\Chkey}\!\left[
  \frac{q(\hat c)}{P^{J}_{\hat c}}\,\Gamma^{(s')}_{\hat c}
  +\frac{1}{\alpha-1}\,\frac{q(\hat c)}{K^{J}_{\hat c}}
  \Big(\Gamma^{(s')}_{\hat c}+\Delta^{(s')}_{\hat c}\Big)\right],
  \tag{\getrefnumber{C-eq:grad-J}}\digesteqlabel{eq:grad-J}
\end{align*}
where
\begin{equation*}\tag{\getrefnumber{C-eq:grad-pieces}}\digesteqlabel{eq:grad-pieces}
  \Gamma^{(s')}_{\hat c}\coloneqq
  \sum_{(x,v):\,\phi(a(x,v))=\hat c}
  \big(W_{\tau}^{\dagger}M_xW_{\tau}\big)\otimes F^{(s')}_v,
  \qquad
  \Delta^{(s')}_{\hat c}\coloneqq
  \mathbb 1[s'\in\SecSet_K(\hat c)]\;
  \chi_{\Qsys_{s'}}^{\dagger}\Big(
  \cG_{\hat c,s'}^{\dagger}\big[\nabla T(X_{\hat c,s'})
  -\idon{}\big]\Big),
\end{equation*}
with $\cG^{\dagger}_{\hat c,s'}(Y)=\sum_{\kappa}K^{(s')\dagger}_{\kappa}
YK^{(s')}_{\kappa}$ the Kraus adjoint and
$\chi^{\dagger}_{\Qsys_{s'}}(A)=(W_{\tau}^{\dagger}\otimes\idon{})A
(W_{\tau}\otimes\idon{})$. Each $\nabla T(X_{\hat c,s'})$ is read in the
split form of Theorem~\ref{thm:gradient}\,(i): divided differences over
the strictly positive spectrum of $X_{\hat c,s'}$, the pinched operators
$\cZ(X^{p}_{\hat c,s'})$ and $\cZ(Y^{\alpha-1})$ on the ambient branch
space.

\emph{Untrusted-source mode.} There the \emph{state form} is primary,
the Choi form being its pullback rather than the other way round: the
variables are the sector blocks $\omega_{s}$ themselves
(Definition~\ref{def:untrusted-choi}), $\chi_{\Qsys_{s'}}=\idmap$, and
\eqref{eq:grad-q}--\eqref{eq:grad-J} hold verbatim with
$J_{s'}\to\omega_{s'}$, $P^{J}_{\hat c}\to P_{\hat c}$,
$K^{J}_{\hat c}\to K_{\hat c}$, $X_{\hat c,s}=\cG_{\hat c,s}(\omega_s)$
and the pieces \eqref{eq:grad-pieces} losing their $W_{\tau}$
conjugation:
\begin{equation*}\tag{\getrefnumber{C-eq:grad-pieces-un}}\digesteqlabel{eq:grad-pieces-un}
  \Gamma^{(s')}_{\hat c}=
  \sum_{(x,v):\,\phi(a(x,v))=\hat c} M_x\otimes F^{(s')}_v,
  \qquad
  \Delta^{(s')}_{\hat c}=
  \mathbb 1[s'\in\SecSet_K(\hat c)]\;
  \cG_{\hat c,s'}^{\dagger}\big[\nabla T(X_{\hat c,s'})-\idon{}\big],
\end{equation*}
Hermitian operators on $A\Qsys_{s'}$. The derivation is the proof
below with $\chi_{\Qsys_{s'}}$ deleted; the positivity hypothesis on the
point reads $\omega_{s}\in\Pos{A\Qsys_s}$ per block, there being no $A_0$
in that mode; and the maximal-support hypothesis
reads $\operatorname{supp}X_{\hat c,s'}=W_{\hat c,s'}$ with $\Sadm$
now all of $\Hs{A}$, for which $\omega_{s'}\succ0$ suffices. In
trusted-source mode the same formulas conjugated by $\chi$ give the
state-form gradients, and at $d_A=1$ the two agree because
$W_{\tau}=[1]$.
\end{digestclaim}

% claim meat-core/lem:kraus-freedom
\digestlabel{lem:kraus-freedom}{8.6}\pdfbookmark[2]{Lemma 8.6 Kraus freedom in the branch maps}{digest-meat-core/lem:kraus-freedom}
\renewcommand*{\RestatedName}{\hyperref{contract.pdf}{}{claim-meat-core/lem:kraus-freedom}{Lemma~8.6}}%
\begin{digestclaim}[Kraus freedom in the branch maps]
Let $\alpha\in(1,\infty)$, $p\coloneqq1/\alpha$ and $\hat c\in\Chkey$, and let
$\mathcal{H}_{\mathrm{in}}$ be a finite-dimensional Hilbert space. For every
outcome pair $(x,v)$ of a pass class of $\hat c$, that is
$a(x,v)\in\Kann_{\mathrm{pass}}(\hat c)$, let $S_{xv}$ and $S'_{xv}$ be linear
maps from $\mathcal{H}_{\mathrm{in}}$ into finite-dimensional Hilbert spaces
$\mathcal{K}$ and $\mathcal{K}'$ with
\begin{equation*}\tag{\getrefnumber{C-eq:kraus-gram}}\digesteqlabel{eq:kraus-gram}
  S'^{\dagger}_{xv}S'_{xv}\;=\;S^{\dagger}_{xv}S_{xv} .
\end{equation*}
For $\kappa\in\Kann_{\mathrm{pass}}(\hat c)$ put
\begin{equation*}\tag{\getrefnumber{C-eq:kraus-factored}}\digesteqlabel{eq:kraus-factored}
  K_{\kappa}\coloneqq\sum_{(x,v)\in a^{-1}(\kappa)}\ket{r(x,v)}_R\ket{x,v}_{XY}
  \ket{\kappa}_{C_AC_B}\otimes S_{xv},
\end{equation*}
a map into $\mathcal{B}\coloneqq\Hs{R\,XY\,C_AC_B}\otimes\mathcal{K}$, and
$K'_{\kappa}$ likewise with $S'_{xv}$, into
$\mathcal{B}'\coloneqq\Hs{R\,XY\,C_AC_B}\otimes\mathcal{K}'$. For every
positive semidefinite operator $\omega$ on $\mathcal{H}_{\mathrm{in}}$ put
$X(\omega)\coloneqq\sum_{\kappa}K_{\kappa}\omega K_{\kappa}^{\dagger}$ and
$X'(\omega)\coloneqq\sum_{\kappa}K'_{\kappa}\omega K'^{\dagger}_{\kappa}$,
the sums over $\Kann_{\mathrm{pass}}(\hat c)$. Let $\cZ$ and $\cZ'$ be the
key pinchings on $\mathcal{B}$ and $\mathcal{B}'$, by
the projectors $Z_j=\dyad{j}_R\otimes\id$, $j\in\Zset$, and $T$ and $T'$ the
functionals \eqref{eq:Tcone} with these pinchings. For $X\succeq0$ write
$\nabla T(X)$ for the operator of \eqref{eq:gradT-kernel} with the sum over
the strictly positive eigenvalues of $X$ only, $Y=\cZ(X^{p})$ and
$0^{\alpha-1}\coloneqq0$, which is the split form of
Theorem~\ref{thm:gradient}\,(i), and $\nabla T'$ likewise.
\begin{enumerate}[label=(\roman*),leftmargin=2.4em]
\item \emph{(A partial isometry.)} Let $\mathcal{W}\subseteq\mathcal{B}$ be
  the span of the subspaces
  $\ket{r(x,v)}_R\ket{x,v}_{XY}\ket{a(x,v)}_{C_AC_B}\otimes
  \operatorname{ran}S_{xv}$. There is a linear map
  $\mathbf U:\mathcal{B}\to\mathcal{B}'$ that is isometric on $\mathcal{W}$
  and vanishes on $\mathcal{W}^{\perp}$, with $\mathbf UK_{\kappa}=K'_{\kappa}$
  for every $\kappa$ and $(Z_j\otimes\idon{\mathcal{K}'})\,\mathbf U
  =\mathbf U\,(Z_j\otimes\idon{\mathcal{K}})$ for every $j$; and
  $\mathcal{W}$ is invariant under every $Z_j\otimes\idon{\mathcal{K}}$.
  Hence $X'(\omega)=\mathbf UX(\omega)\mathbf U^{\dagger}$, with
  $\operatorname{supp}X(\omega)\subseteq\mathcal{W}$, for every
  $\omega\succeq0$.
\item \emph{(Values.)} For every $\omega\succeq0$:
  $T'(X'(\omega))=T(X(\omega))$, $\tr X'(\omega)=\tr X(\omega)$,
  $\operatorname{supp}X'(\omega)=\mathbf U\operatorname{supp}X(\omega)$ and
  $\operatorname{rank}X'(\omega)=\operatorname{rank}X(\omega)$.
\item \emph{(Gradients.)} For every $\omega\succeq0$, with $X\coloneqq
  X(\omega)$ and $X'\coloneqq X'(\omega)$,
  \begin{equation*}\tag{\getrefnumber{C-eq:kraus-gradient}}\digesteqlabel{eq:kraus-gradient}
    \nabla T'(X')=\mathbf U\,\nabla T(X)\,\mathbf U^{\dagger},
    \qquad
    \sum_{\kappa}K'^{\dagger}_{\kappa}\big(\nabla T'(X')-\id\big)K'_{\kappa}
    =\sum_{\kappa}K^{\dagger}_{\kappa}\big(\nabla T(X)-\id\big)K_{\kappa}.
  \end{equation*}
\item \emph{(Compression to the support.)} Let $A\succeq0$ be an operator
  on $\mathcal{B}$.
  \begin{enumerate}[label=(\alph*),leftmargin=2.2em]
  \item \emph{(Containment.)} For every $\omega\succeq0$,
    $\operatorname{supp}X(\omega)\subseteq\sum_{\kappa}\operatorname{ran}K_{\kappa}$;
    and $\operatorname{supp}\big((Z_j\otimes\idon{\mathcal{K}})A
    (Z_j\otimes\idon{\mathcal{K}})\big)\subseteq(Z_j\otimes\idon{\mathcal{K}})
    \operatorname{supp}A$ for every $j$; and
    $\operatorname{supp}A^{q}=\operatorname{supp}A$ for every $q>0$.
  \item \emph{(A floor proves equality.)} If $Q:\mathbb{C}^{m}\to\mathcal{B}$
    is an isometry with $\operatorname{ran}Q\supseteq\operatorname{supp}A$
    and $Q^{\dagger}AQ$ is positive definite, then
    $\operatorname{ran}Q=\operatorname{supp}A$.
  \item \emph{(The compression.)} If $Q:\mathbb{C}^{m}\to\mathcal{B}$ is an
    isometry with $\operatorname{ran}Q=\operatorname{supp}A$ and $m\ge1$, then
    $N\coloneqq Q^{\dagger}AQ$ is positive definite, $A=QNQ^{\dagger}$,
    $A^{q}=QN^{q}Q^{\dagger}$ for every $q>0$, and for every operator $H$ on
    $\mathcal{B}$ the divided-difference sum of \eqref{C-eq:face-dpow} at $A$,
    over the strictly positive eigenvalues of $A$, equals
    $Q\,\mathrm{D}_N[Q^{\dagger}HQ]\,Q^{\dagger}$, with $\mathrm{D}_N$ the
    map of \xref{foundations/lem:matrix-function-perturbation}\,(iii) for
    $f(x)=x^{p}$.
  \end{enumerate}
\item \emph{(The branch maps.)} With $\mathcal{H}_{\mathrm{in}}=\Hs{A\Qsys_s}$
  and $S_{xv}=\sqrt{M_x}\otimes\sqrt{F^{(s)}_v}$, the $K_{\kappa}$ are the
  $K^{(s)}_{\kappa}$ of Definition~\ref{def:branch} and
  $X(\omega)=\cG_{\hat c,s}(\omega)$. In trusted-source mode, with
  $\mathcal{H}_{\mathrm{in}}=\Hs{A_0\Qsys_s}$, a support frame and
  $S_{xv}=\sqrt{M_x}\,W_{\tau}\otimes\sqrt{F^{(s)}_v}$, the $K_{\kappa}$ are
  the $K^{(s)}_{\kappa}(W_{\tau}\otimes\idon{\Qsys_s})$,
  $X(J_s)=\cG^{J}_{\hat c,s}(J_s)$, and
  $\sum_{\kappa}K^{\dagger}_{\kappa}(\nabla T(X)-\id)K_{\kappa}
  =\chi^{\dagger}_{\Qsys_s}\big(\cG^{\dagger}_{\hat c,s}[\nabla T(X)-\id]\big)$,
  the operator of \eqref{eq:grad-pieces}; the same holds on a
  joint sector
  $(t,s)$ with $M^{(t)}_x$ and the frame of the Alice sector $t$
  (Definition~\ref{def:joint-sectors}). In each case
  $S^{\dagger}_{xv}S_{xv}$ is $M_x\otimes F^{(s)}_v$, respectively
  $(W_{\tau}^{\dagger}M_xW_{\tau})\otimes F^{(s)}_v$. So for any $S'_{xv}$
  with that Gram matrix, the deficit $f_{\hat c,s}$, the branch operator's
  rank and the operator $\Delta^{(s)}_{\hat c}$ of \eqref{eq:grad-pieces}
  are given, by (ii) and (iii), by $X'$ in place of $X$.
\end{enumerate}
\end{digestclaim}

% claim meat-core/thm:oracle
\digestlabel{thm:oracle}{8.9}\pdfbookmark[2]{Theorem 8.9 Limit of the value as the Rényi order tends to one}{digest-meat-core/thm:oracle}
\renewcommand*{\RestatedName}{\hyperref{contract.pdf}{}{claim-meat-core/thm:oracle}{Theorem~8.9}}%
\begin{digestclaim}[Limit of the value as the Rényi order tends to one]
Let $\Decl$ satisfy Definition~\ref{def:admissible-declaration} and let
$\acSet$ satisfy Definition~\ref{def:admissible-acceptance-set}. Assume that
there is $p_0>0$ with $q(\hat c)\ge p_0$ for every $\mathbf q\in\acSet$ and
$\hat c\in\Chkey$, and that the asymptotic feasible set
$\Feas_1\coloneqq\{\omega\in\Feas:\Phimap[\omega]\in\acSet\}$ is
nonempty. Then
\begin{equation*}\tag{\getrefnumber{C-eq:oracle}}\digesteqlabel{eq:oracle}
  \lim_{\alpha\downarrow1}\hlow(\alpha)
  \;=\;h_1\coloneqq
  \inf_{\omega\in\Feas_1}\ \sum_{\hat c\in\Chkey}\
  \sum_{s\in\SecSet_K(\hat c)}
  \Dtwo\Big(\cG_{\hat c,s}(\omega_s)\,\Big\|\,
  \cZ\big(\cG_{\hat c,s}(\omega_s)\big)\Big).
\end{equation*}
\end{digestclaim}

\section*{9\quad The dropped program}\pdfbookmark[1]{9 The dropped program}{digest-group-meat-core/lem:rank-one-sector}

% claim meat-core/lem:rank-one-sector
\digestlabel{lem:rank-one-sector}{9.1}\pdfbookmark[2]{Lemma 9.1 Linear deficit on a rank-one Alice sector}{digest-meat-core/lem:rank-one-sector}
\renewcommand*{\RestatedName}{\hyperref{contract.pdf}{}{claim-meat-core/lem:rank-one-sector}{Lemma~9.1}}%
\begin{digestclaim}[Linear deficit on a rank-one Alice sector]
Let $\Decl$ satisfy Definition~\ref{def:admissible-declaration} in trusted-source mode
and the general format, and let $t\in\mathbb{T}$ with
$\tau^{(t)}_A=p_t\dyad{\hat v}$, $p_t>0$, $\ket{\hat v}\in\Hs{A_t}$ a unit
vector.
\begin{enumerate}[label=(\roman*),leftmargin=2.2em]
\item \emph{(product blocks)} Every list $(\omega^{(t,s)})_{s\in\SecSet}$ of
positive semidefinite blocks with
$\sum_s\tr_{\Qsys_s}\omega^{(t,s)}=\tau^{(t)}_A$, in particular the sector-$t$
part of every $\omega\in\Feas$, has the form
$\omega^{(t,s)}=\dyad{\hat v}\otimes\sigma_s$ with
$\sigma_s\in\Pos{\Qsys_s}$ and $\sum_s\tr\sigma_s=p_t$.
\item \emph{(pass classes are product sets)} If the
readout maps are the default ones,
$\rd{\mathsf A}(x,v)=x$ and $\rd{\mathsf B}(x,v)=v$, then for
every transcript protocol and every $\kappa\in\Kann$ there are sets
$\Xset_{\kappa}\subseteq\Xset$, $\Vset_{\kappa}\subseteq\Vset$ with
$a^{-1}(\kappa)=\Xset_{\kappa}\times\Vset_{\kappa}$; if moreover
$\hold=\mathsf A$, then $r(x,v)=r_{\mathsf A}(x,\kappa)$ on $a^{-1}(\kappa)$.
\item \emph{(linear deficit)} Under the hypotheses of (ii), let
$\hat c\in\Chkey$, $s\in\SecSet$ and $\alpha\in(1,\infty)$, put
$\mu_x\coloneqq\bra{\hat v}M^{(t)}_x\ket{\hat v}$ and, for
$\kappa\in\Kann_{\mathrm{pass}}(\hat c)$ and $\mathsf z\in\Zval$,
\begin{equation*}\tag{\getrefnumber{C-eq:rank-one-weights}}\digesteqlabel{eq:rank-one-weights}
  m_{\kappa}\coloneqq\sum_{x\in\Xset_{\kappa}}\mu_x,\qquad
  m_{\kappa,\mathsf z}\coloneqq\sum_{x\in\Xset_{\kappa}:\,
  r_{\mathsf A}(x,\kappa)=\mathsf z}\mu_x,\qquad
  \eta_{\kappa}\coloneqq
  \begin{cases}
   1-\sum_{\mathsf z}\big(m_{\kappa,\mathsf z}/m_{\kappa}\big)^{\alpha},
   & m_{\kappa}>0,\\ 0,& m_{\kappa}=0 .
  \end{cases}
\end{equation*}
Then for every $\sigma_s\in\Pos{\Qsys_s}$
\begin{equation*}\tag{\getrefnumber{C-eq:rank-one-deficit}}\digesteqlabel{eq:rank-one-deficit}
  f_{\hat c,(t,s)}\big(\dyad{\hat v}\otimes\sigma_s\big)
  =-\sum_{\kappa\in\Kann_{\mathrm{pass}}(\hat c)}\eta_{\kappa}\,m_{\kappa}
  \sum_{v\in\Vset_{\kappa}}\tr\big[F^{(s)}_v\sigma_s\big],
\end{equation*}
a \emph{linear} function of $\sigma_s$, with $0\le\eta_{\kappa}<1$. In
ensemble form $\mu_x=\tr N^{(t)}_x/p_t$ and $p_t=\tr\tau^{(t)}_A$, declared
numbers.
\item \emph{(two values)} If within every pass class of $\hat c$ the key
takes $d$ values of equal weight, $m_{\kappa,\mathsf z}=m_{\kappa}/d$, then
$\eta_{\kappa}=1-d^{1-\alpha}$ and the sector contributes to
$\frac{1}{1-\alpha}\log_2$ of the ratio $K/P$ \emph{as if the conditional
entropy were $\log_2d$ on the pass weight}: one bit per unit of pass weight
for a uniform key bit. The qualifier matters: what the
expression assigns to the pair is its fail weight plus $d^{1-\alpha}$ times
its pass weight, so with a non-zero fail weight in the same joint pair the
contribution is strictly less than that of a sector of conditional entropy
$\log_2d$, and equals it exactly when the pass weight is the whole of the
pair's weight. If the key is
constant on every $\Xset_{\kappa}\cap\{x:\mu_x>0\}$ then $\eta_{\kappa}=0$
and $f_{\hat c,(t,s)}$ vanishes on these blocks.
\end{enumerate}
\end{digestclaim}

% claim meat-core/def:dropped
\digestlabel{def:dropped}{9.2}\pdfbookmark[2]{Definition 9.2 Dropped set and dropped program}{digest-meat-core/def:dropped}
\renewcommand*{\RestatedName}{\hyperref{contract.pdf}{}{claim-meat-core/def:dropped}{Definition~9.2}}%
\begin{digestclaim}[Dropped set and dropped program]
Let $\Decl$ be a declaration in the general format of
Definition~\ref{def:declaration}. A
\emph{dropped set}
is a set
$\mathcal{D}\subseteq\Chkey\times(\mathbb{T}\times\SecSet)$ of pairs of a
key-live acceptance value and a joint sector. For
$\omega=\bigoplus_{(t,s)}\omega^{(t,s)}$ block diagonal along the joint
sectors put
\begin{equation*}\tag{\getrefnumber{C-eq:KD}}\digesteqlabel{eq:KD}
  K^{\mathcal{D}}_{\hat c}(\omega)\coloneqq P_{\hat c}(\omega)
  +\sum_{(t,s):\,(\hat c,(t,s))\notin\mathcal{D}}
  f_{\hat c,(t,s)}\big(\omega^{(t,s)}\big)
  \qquad(\hat c\in\Chkey),
\end{equation*}
with $f_{\hat c,(t,s)}$ of \eqref{eq:joint-deficit}; let
$O^{\mathcal{D}}_{\alpha}$ be \eqref{eq:objective} with
$K^{\mathcal{D}}_{\hat c}$ in place of $K_{\hat c}$, with the same
conventions at vanishing arguments; and let
\begin{equation*}\tag{\getrefnumber{C-eq:hlowD}}\digesteqlabel{eq:hlowD}
  \underline{h}^{\mathcal{D}}_{\alpha}\coloneqq
  \inf_{\mathbf q\in\acSet}\ \inf_{\omega\in\Feas}\
  O^{\mathcal{D}}_{\alpha}(\mathbf q,\omega).
\end{equation*}
The minimisation in \eqref{eq:hlowD} is the
\emph{dropped program}.
A joint sector $(t,s)$ is
\emph{fully dropped}
when $\Chkey\times\{(t,s)\}\subseteq\mathcal{D}$. For
$\mathcal{D}=\varnothing$ these are $K_{\hat c}$ (\eqref{eq:K-joint}),
$O_{\alpha}$ and $\hlow$.

\emph{What is not changed.} The feasible set is $\Feas$, unchanged: a
dropped sector's block remains a variable, enters $\Phimap$ and every
$P_{\hat c}$ with its full weight --- the term
$P_{\hat c,(t,s)}(\omega^{(t,s)})$ stays inside $P_{\hat c}$ in
\eqref{eq:KD} --- and enters the marginal condition and every constraint
row. Setting a dropped block to zero, omitting its statistics, or removing
its weight from $P_{\hat c}$ in \eqref{eq:KD} are different operations, and
each can overstate the value (Proposition~\ref{prop:retention} for the
first). \emph{The dropped set is a formulation choice of a run and not
declared data}.
\end{digestclaim}

% claim meat-core/prop:drop-bound
\digestlabel{prop:drop-bound}{9.3}\pdfbookmark[2]{Proposition 9.3 The dropped program bounds the value from below}{digest-meat-core/prop:drop-bound}
\renewcommand*{\RestatedName}{\hyperref{contract.pdf}{}{claim-meat-core/prop:drop-bound}{Proposition~9.3}}%
\begin{digestclaim}[The dropped program bounds the value from below]
Let $\Decl$ satisfy Definition~\ref{def:admissible-declaration} in the general format, in
either source mode and with an arbitrary list $\CList$; let
$\alpha\in(1,\infty)$ and let $\mathcal{D}\subseteq\mathcal{D}'$ be dropped
sets (Definition~\ref{def:dropped}). For every $\omega$ block diagonal along
the joint sectors, every $\hat c\in\Chkey$ and every distribution $\mathbf q$
on $\Chatset$:
\begin{enumerate}[label=(\roman*),leftmargin=2.2em]
\item $0\le K_{\hat c}(\omega)\le K^{\mathcal{D}}_{\hat c}(\omega)
\le K^{\mathcal{D}'}_{\hat c}(\omega)\le P_{\hat c}(\omega)$;
\item $O^{\mathcal{D}'}_{\alpha}(\mathbf q,\omega)\le
O^{\mathcal{D}}_{\alpha}(\mathbf q,\omega)\le O_{\alpha}(\mathbf q,\omega)$
in the extended-real sense, and $O^{\mathcal{D}'}_{\alpha}(\mathbf
q,\omega)\ge0$ whenever $\tr\omega=1$;
\item consequently
$0\le\underline{h}^{\mathcal{D}'}_{\alpha}\le
\underline{h}^{\mathcal{D}}_{\alpha}\le\hlow\le\hqkd(\CList)$;
\item if every pair of $\mathcal{D}$ is key-inert
(Definition~\ref{def:joint-sectors}) then
$K^{\mathcal{D}}_{\hat c}=K_{\hat c}$ and
$O^{\mathcal{D}}_{\alpha}=O_{\alpha}$ identically.
\end{enumerate}
\end{digestclaim}

% claim meat-core/cor:drop-program
\digestlabel{cor:drop-program}{9.4}\pdfbookmark[2]{Corollary 9.4 Convexity, gradient and limit of the dropped program}{digest-meat-core/cor:drop-program}
\renewcommand*{\RestatedName}{\hyperref{contract.pdf}{}{claim-meat-core/cor:drop-program}{Corollary~9.4}}%
\begin{digestclaim}[Convexity, gradient and limit of the dropped program]
Under the hypotheses of Proposition~\ref{prop:drop-bound}, with $\acSet$
convex and compact (Definition~\ref{def:admissible-acceptance-set}):
\begin{enumerate}[label=(\roman*),leftmargin=2.2em]
\item $K^{\mathcal{D}}_{\hat c}$ is concave and continuous on the cone of
operators block diagonal along the joint sectors;
\item $O^{\mathcal{D}}_{\alpha}$ is jointly convex and lower semicontinuous on
$\Delta(\Chatset)\times\Feas$ with values in $[0,+\infty]$, and the infimum
\eqref{eq:hlowD} is attained whenever it is finite;
\item if $(t,s)$ is fully dropped, then $O^{\mathcal{D}}_{\alpha}$ depends on
the block $\omega^{(t,s)}$ only through the real numbers
$\tr[(M^{(t)}_x\otimes F^{(s)}_v)\,\omega^{(t,s)}]$, $(x,v)\in\Xset\times\Vset$;
at a point where every $P_{\hat c}$ with $q(\hat c)>0$ and every
$K^{\mathcal{D}}_{\hat c}$ with $q(\hat c)>0$, $\hat c\in\Chkey$, is
positive, its derivative in the direction of a Hermitian $\Delta$ supported
on that joint sector is
\begin{equation*}\tag{\getrefnumber{C-eq:drop-gradient}}\digesteqlabel{eq:drop-gradient}
  -\log_2\!\mathrm{e}\sum_{\hat c\in\Chatset}\beta_{\hat c}\,
  \tr\big[\Gamma^{(t,s)}_{\hat c}\Delta\big],\qquad
  \Gamma^{(t,s)}_{\hat c}\coloneqq\!\!\sum_{(x,v):\,\phi(a(x,v))=\hat c}\!\!
  M^{(t)}_x\otimes F^{(s)}_v,\qquad
  \beta_{\hat c}\coloneqq
  \begin{cases}
   \frac{\alpha}{\alpha-1}\frac{q(\hat c)}{P_{\hat c}},&\hat c\notin\Chkey,\\[0.6ex]
   \frac{q(\hat c)}{P_{\hat c}}+\frac{1}{\alpha-1}
   \frac{q(\hat c)}{K^{\mathcal{D}}_{\hat c}},&\hat c\in\Chkey,
  \end{cases}
\end{equation*}
with $\beta_{\hat c}\coloneqq0$ where $q(\hat c)=0$;
\item under the additional hypotheses of Theorem~\ref{thm:oracle},
\begin{equation*}\tag{\getrefnumber{C-eq:drop-oracle}}\digesteqlabel{eq:drop-oracle}
  \lim_{\alpha\downarrow1}\underline{h}^{\mathcal{D}}_{\alpha}
  =\inf_{\omega\in\Feas_1}\ \sum_{\hat c\in\Chkey}\
  \sum_{\substack{(t,s)\in\mathbb{J}_K(\hat c)\\(\hat c,(t,s))\notin\mathcal{D}}}
  \Dtwo\Big(\cG_{\hat c,(t,s)}\big(\omega^{(t,s)}\big)\,\Big\|\,
  \cZ\big(\cG_{\hat c,(t,s)}\big(\omega^{(t,s)}\big)\big)\Big).
\end{equation*}
\end{enumerate}
\end{digestclaim}

% claim meat-core/prop:minimal-divergence
\digestlabel{prop:minimal-divergence}{9.5}\pdfbookmark[2]{Proposition 9.5 The minimal divergence as the value with every pair dropped}{digest-meat-core/prop:minimal-divergence}
\renewcommand*{\RestatedName}{\hyperref{contract.pdf}{}{claim-meat-core/prop:minimal-divergence}{Proposition~9.5}}%
\begin{digestclaim}[The minimal divergence as the value with every pair dropped]
Let $\Decl$ satisfy Definition~\ref{def:admissible-declaration} in the
general format, in either source mode and with an arbitrary list $\CList$,
let $\acSet$ satisfy Definition~\ref{def:admissible-acceptance-set}, and let
$\alpha\in(1,\infty)$. Let
$\mathcal{D}_{\mathrm{all}}\coloneqq\Chkey\times(\mathbb{T}\times\SecSet)$
be the dropped set of every pair (Definition~\ref{def:dropped}), and put
\begin{equation*}\tag{\getrefnumber{C-eq:dmin}}\digesteqlabel{eq:dmin}
  \sfD_{\min}\coloneqq\inf_{\mathbf q\in\acSet}\ \inf_{\omega\in\Feas}\
  \sfD\big(\mathbf q\,\big\|\,\Phimap[\omega]\big)\in[0,+\infty],
  \qquad
  m_{\mathrm{key}}\coloneqq\sup_{\mathbf q\in\acSet}\
  \sum_{\hat c\in\Chkey}q(\hat c)\in[0,1],
\end{equation*}
with $\inf\varnothing=+\infty$ and $\sup\varnothing=0$.
\begin{enumerate}[label=(\roman*),leftmargin=2.2em]
\item \emph{(The value with every pair dropped.)} For every distribution
  $\mathbf q$ on $\Chatset$ and every $\omega\in\Feas$,
  $O^{\mathcal{D}_{\mathrm{all}}}_{\alpha}(\mathbf q,\omega)
  =\frac{\alpha}{\alpha-1}\,\sfD(\mathbf q\|\Phimap[\omega])$; hence
  $\underline{h}^{\mathcal{D}_{\mathrm{all}}}_{\alpha}
  =\frac{\alpha}{\alpha-1}\,\sfD_{\min}$.
\item \emph{(The value between its floor and its ceiling.)}
  \begin{equation*}\tag{\getrefnumber{C-eq:dmin-sandwich}}\digesteqlabel{eq:dmin-sandwich}
    \frac{\alpha}{\alpha-1}\,\sfD_{\min}\;\le\;\hlow\;\le\;
    \frac{\alpha}{\alpha-1}\,\sfD_{\min}+\log_2(d_Z+1)\,m_{\mathrm{key}} .
  \end{equation*}
\item \emph{(When the floor is positive.)} $\sfD_{\min}=0$ if and only if
  $\mathbf q=\Phimap[\omega]$ for some $\mathbf q\in\acSet$ and some
  $\omega\in\Feas$. Equivalently, $\sfD_{\min}>0$ if and only if no
  distribution of $\acSet$ is the statistics $\Phimap[\omega]$ of a state
  $\omega\in\Feas$.
\item \emph{(Bounds at one order serve every order.)} Let $b$ be a real
  number with $b\le\underline{h}^{\mathcal{D}_{\mathrm{all}}}_{\alpha}$, and
  let $(\mathbf q,\omega)\in\acSet\times\Feas$. Then
  \begin{equation*}\tag{\getrefnumber{C-eq:dmin-bracket}}\digesteqlabel{eq:dmin-bracket}
    \frac{\alpha-1}{\alpha}\,b\;\le\;\sfD_{\min}\;\le\;
    \frac{\alpha-1}{\alpha}\,
    O^{\mathcal{D}_{\mathrm{all}}}_{\alpha}(\mathbf q,\omega),
  \end{equation*}
  and for every $\alpha'\in(1,\infty)$ the value $\underline{h}_{\alpha'}$
  of \eqref{eq:hlow-def} at the order $\alpha'$ satisfies
  $\underline{h}_{\alpha'}\ge\frac{\alpha'}{\alpha'-1}\cdot
  \frac{\alpha-1}{\alpha}\,b$. If $b>0$, then $\sfD_{\min}>0$, and no
  distribution of $\acSet$ is the statistics of a state of $\Feas$.
  For every $\alpha'\ge\alpha$,
  \begin{equation*}\tag{\getrefnumber{C-eq:dmin-ceiling}}\digesteqlabel{eq:dmin-ceiling}
    \frac{\alpha'}{\alpha'-1}\,\sfD_{\min}\;\le\;
    \frac{\alpha'}{\alpha'-1}\cdot\frac{\alpha-1}{\alpha}\,
    O^{\mathcal{D}_{\mathrm{all}}}_{\alpha}(\mathbf q,\omega)\;\le\;
    O^{\mathcal{D}_{\mathrm{all}}}_{\alpha}(\mathbf q,\omega),
  \end{equation*}
  so that $\underline{h}_{\alpha'}\le
  \frac{\alpha'}{\alpha'-1}\cdot\frac{\alpha-1}{\alpha}\,
  O^{\mathcal{D}_{\mathrm{all}}}_{\alpha}(\mathbf q,\omega)
  +\log_2(d_Z+1)\,m_{\mathrm{key}}$.
\end{enumerate}
\end{digestclaim}

\section*{10\quad Key length}\pdfbookmark[1]{10 Key length}{digest-group-meat-core/def:postprocessing-model}

% claim meat-core/def:postprocessing-model
\digestlabel{def:postprocessing-model}{10.1}\pdfbookmark[2]{Definition 10.1 Fixed-length post-processing model}{digest-meat-core/def:postprocessing-model}
\renewcommand*{\RestatedName}{\hyperref{contract.pdf}{}{claim-meat-core/def:postprocessing-model}{Definition~10.1}}%
\begin{digestclaim}[Fixed-length post-processing model]
Fix a declaration $\Decl$ with key holder $\hold$
(Definition~\ref{def:declaration}), a round count $n$, an acceptance set
$\acSet$, an output length $\ell$, a transcript length $\lEC$ and error
budgets $\epa,\eev\in(0,1]$. A hash family is a set of functions with a
probability law on it \cite[Def.~3.5]{tupkary2026decoy}. The
\emph{fixed-length post-processing model}
is the following classical processing, applied after the $n$
round channels of Definition~\ref{def:roundmodel}.
The raw string $Z_1^n$ is the content of the classical registers
$Z_1,\dots,Z_n$ of the round channels, which the single-round protocol channel
$\Mprot$ retains (Theorem~\ref{thm:protocol-channel}); the key holder computes
it from its own outcome record and the announcements $I_1^n$ through the
function $r_{\hold}$ of \eqref{eq:partycomp}, and we call the other party the
partner. The acceptance values $\Chat_1^n$ are announced. The processing
appends, in order:
\begin{enumerate}[label=(\roman*),leftmargin=2.4em]
\item the \emph{acceptance test}: the event
  $\OmAT\coloneqq\{\operatorname{freq}_{\hat c_1^n}\in\acSet\}$, a function
  of $\Chat_1^n$;
\item \emph{error correction}, one-way from the holder: a public seed
  $S_{\mathrm{pub}}$, drawn with the law the error-correction scheme
  prescribes, and the public transcript
  $L_{\mathrm{EC}}=g_{\mathrm{EC}}(Z_1^n,S_{\mathrm{pub}})$, which occupies
  a fixed-length register of $\lEC$ bits: a shorter transcript is padded to
  length $\lEC$, so that $\dim(L_{\mathrm{EC}})=2^{\lEC}$ exactly; from
  $L_{\mathrm{EC}}$ and its own data the partner forms a guess $\hat Z_1^n$
  of $Z_1^n$;
\item \emph{error verification}: a public seed $S_{\mathrm{EV}}$ selects a
  function $h_{\mathrm{EV}}$, with the law of the error-verification family,
  from a hash family of functions with
  $\lceil\log_2\frac1{\eev}\rceil$-bit values; the holder announces
  $T_{\mathrm{EV}}\coloneqq h_{\mathrm{EV}}(Z_1^n)$, the partner compares it
  with $h_{\mathrm{EV}}(\hat Z_1^n)$ and announces the one-bit outcome
  $F_{\mathrm{acc}}\in\{\mathrm{acc},\mathrm{rej}\}$, which is
  $\mathrm{acc}$ exactly when the two values agree;
\item \emph{privacy amplification}: a function $G$ is drawn, with the law of
  the privacy-amplification family, from a hash family of functions with
  $\ell$-bit values, and announced; the holder outputs
  $K\coloneqq G(Z_1^n)$ and the partner $\hat K\coloneqq G(\hat Z_1^n)$.
\end{enumerate}
We write $K_A$ and $K_B$ for Alice's and Bob's outputs, which are $K$ and
$\hat K$ in the order the key holder fixes, and
\begin{equation*}\tag{\getrefnumber{C-eq:pp-registers}}\digesteqlabel{eq:pp-registers}
  S\coloneqq(S_{\mathrm{pub}},S_{\mathrm{EV}}),\qquad
  L\coloneqq(L_{\mathrm{EC}},T_{\mathrm{EV}},F_{\mathrm{acc}}),\qquad
  E_{\mathrm{tot}}\coloneqq\Chat_1^nI_1^n\Ehat\,S\,L\,G,
\end{equation*}
where $L$ has at most
$\lambda\coloneqq\lEC+\lceil\log_2\frac1{\eev}\rceil+1$ bits and $\Ehat$ is
the adversary's retained quantum memory of Definition~\ref{def:roundmodel}.
The protocol \emph{accepts} on
$\Omega_{\mathrm{acc}}\coloneqq\OmAT\wedge\{F_{\mathrm{acc}}=\mathrm{acc}\}$;
on its complement both parties abort and output the symbol $\perp$ in place
of a key. The model is stated in the
\emph{honest-authentication setting}:
every classical message of the protocol, the round announcements included,
reaches its recipient exactly as it was sent, in the order sent and at some
time after it was sent, the adversary receives a copy of it, and no message
fails authentication.
\end{digestclaim}

% claim meat-core/def:authenticated-channel
\digestlabel{def:authenticated-channel}{10.2}\pdfbookmark[2]{Definition 10.2 Authenticated classical channel and authentication post-processing}{digest-meat-core/def:authenticated-channel}
\renewcommand*{\RestatedName}{\hyperref{contract.pdf}{}{claim-meat-core/def:authenticated-channel}{Definition~10.2}}%
\begin{digestclaim}[Authenticated classical channel and authentication post-processing]
Label each classical message of a party by its position $i$ in that party's
sequence of sent messages, and the messages a party receives by their
position in its sequence of received messages; times are global times, which
the parties need not know.
\begin{enumerate}[label=(\alph*),leftmargin=2.4em]
\item The
  \emph{authenticated classical channel}
  between the two parties gives the adversary a copy of every message sent
  over it and lets her choose when each message is delivered, subject to:
  (A1) if the $i$-th message a party receives arrives before the other
  party sent its $i$-th message, the received symbol is
  $\mathrm{auth\text{-}abort}$; (A2) otherwise the received symbol is either
  the $i$-th sent message, unaltered, or $\mathrm{auth\text{-}abort}$.
\item An \emph{authentication scheme} is a pair of procedures, one per party,
  that use a pre-shared key and an unauthenticated classical channel, on
  which the adversary may read, alter, inject, delay and suppress messages.
  It \emph{$\varepsilon$-realises} the authenticated classical channel if
  there is a simulator at the adversary's interface such that the
  distinguishing advantage between the scheme run on the key and the
  unauthenticated channel, and the authenticated classical channel with that
  simulator, is at most $\varepsilon$
  \cite[Sec.~II.E, Def.~1]{portmann2022security}. Its \emph{key consumption}
  is the number of bits of the pre-shared key it uses in one run of the
  protocol.
\item The
  \emph{authentication post-processing}
  appended to the fixed-length post-processing model
  (Definition~\ref{def:postprocessing-model}) is: a party that receives
  $\mathrm{auth\text{-}abort}$ for any message aborts, replaces its key by
  $\perp$ and sends $\mathrm{abort}$ in place of each later message; after
  privacy amplification the partner sends $\mathrm{accept}$ if it holds a key
  and $\mathrm{abort}$ otherwise; the key holder sends $\mathrm{accept}$ if
  it received $\mathrm{accept}$ and holds a key, and otherwise sends
  $\mathrm{abort}$ and replaces its key by $\perp$; the partner keeps its key
  if it receives $\mathrm{accept}$ and otherwise replaces it by $\perp$.
\end{enumerate}
\end{digestclaim}

% claim meat-core/lem:ev-correctness
\digestlabel{lem:ev-correctness}{10.4}\pdfbookmark[2]{Lemma 10.4 Correctness from error verification}{digest-meat-core/lem:ev-correctness}
\renewcommand*{\RestatedName}{\hyperref{contract.pdf}{}{claim-meat-core/lem:ev-correctness}{Lemma~10.4}}%
\begin{digestclaim}[Correctness from error verification]
Let the protocol run the fixed-length post-processing model
(Definition~\ref{def:postprocessing-model}) in the
honest-authentication setting, let Assumptions~\ref{ass:hash},
\ref{ass:seed} and~\ref{ass:conform} hold, and let $\eev\in(0,1]$. Then for every attack and
every output length $\ell$ the protocol is $\eev$-correct on the accepted
state, in the sense of \cite[Sec.~2.3.1, Def.~2.2, Eq.~(4)]{tupkary2025survey}:
\begin{equation*}\tag{\getrefnumber{C-eq:ev-correct}}\digesteqlabel{eq:ev-correct}
  \Pr\big[\Omega_{\mathrm{acc}}\wedge K_A\ne K_B\big]\le\eev .
\end{equation*}
\end{digestclaim}

% claim meat-core/lem:memory-lift
\digestlabel{lem:memory-lift}{10.5}\pdfbookmark[2]{Lemma 10.5 Approximating an attack by all-at-once attacks with finite memory}{digest-meat-core/lem:memory-lift}
\renewcommand*{\RestatedName}{\hyperref{contract.pdf}{}{claim-meat-core/lem:memory-lift}{Lemma~10.5}}%
\begin{digestclaim}[Approximating an attack by all-at-once attacks with finite memory]
Read a run of the protocol and its attacks as follows.
\begin{enumerate}[label=(\alph*),leftmargin=2.4em]
\item \emph{(the run)} Its quantum phase is an $n$-round protocol of
  \xref{reductions/def:protocol}, whose measured system $\Bsys_j$ of round $j$
  is separable and of any dimension, followed by the
  fixed-length post-processing model of
  Definition~\ref{def:postprocessing-model} in the
  honest-authentication setting. In trusted-source mode its source
  releases, for $i=1,\dots,n$, a separable transmitted system
  $\Asig_i$ of any dimension, prepared together with Alice's letter register
  $X_i$ in a state $\rho_i$ on $X_i\Asig_i$ that depends on nothing else, and
  the run reads $X_i$ from round $i$ on; in untrusted-source mode there
  are no $X_i$ and no $\Asig_i$. Every system of the run other than the
  $\Bsys_j$ and the $\Asig_i$ is finite-dimensional, and its public messages
  are classical registers it keeps.
\item \emph{(attacks)} An \emph{attack} consists of a release schedule
  $0=k_0\le k_1\le\dots\le k_n=n$ with $k_j\ge j$ for every $j$, a state
  $\zeta$ on the adversary's initial register $E_0$, and channels, in the
  sense of \xref{foundations/lem:register-compression},
  \begin{equation*}\tag{\getrefnumber{C-eq:lift-attack}}\digesteqlabel{eq:lift-attack}
    \Aev_j\in
    \begin{cases}
      \mathrm{CPTP}\big(\Asig_{k_{j-1}+1}\cdots\Asig_{k_j}E_{j-1},\
        \Bsys_jE'_j\big), & \text{trusted-source mode},\\[0.3ex]
      \mathrm{CPTP}\big(E_{j-1},\ A_j\Bsys_jE'_j\big),
        & \text{untrusted-source mode},
    \end{cases}
  \end{equation*}
  for $j=1,\dots,n$, where $E'_j$ is the adversary's memory after round $j$,
  $E_j$ is $E'_j$ together with a copy of the public messages of round $j$,
  and $\Hs{E_0}$ and every $\Hs{E'_j}$ are Hilbert spaces of any dimension.
  Under the attack the run starts from $\zeta$ and its own registers, prepares
  $\rho_i$ in round $j$ for $k_{j-1}<i\le k_j$ before $\Aev_j$ acts, takes
  $\Bsys_j$, and in untrusted-source mode Alice's system $A_j$, from the
  output of $\Aev_j$, and places a copy of the public messages of round $j$ in
  $E_j$. Its \emph{output state} is the state of the run's
  output registers, its public messages and $E'_n$. The attack is
  \emph{all-at-once} if $k_1=n$, and has \emph{finite registers} if $E_0$ and
  every $E'_j$ are finite-dimensional.
\end{enumerate}
Then:
\begin{enumerate}[label=(\roman*),leftmargin=2.4em]
\item \emph{(approximation)} for every attack and every $\eta>0$ there is an
  all-at-once attack with finite registers whose memory after round $n$ is a
  subspace of $\Hs{E'_n}$, such that the output states $\rho$ and
  $\rho^{\circ}$ of the two attacks satisfy
  $\|\iota(\rho^{\circ})-\rho\|_1\le\eta$, with $\iota$ the inclusion of that
  subspace into $\Hs{E'_n}$, acting as the identity on every other register;
\item \emph{(trace-distance criteria)} let $L_1$ and $L_2$ be completely
  positive maps, not increasing the trace, from the run's output registers
  and public messages to a finite-dimensional system, put $L\coloneqq
  L_1-L_2$, and for a state $\sigma$ of those registers and the adversary's
  memory after round $n$ put
  $f(\sigma)\coloneqq\frac12\|(L\otimes\idmap)(\sigma)\|_1$; if
  $f(\rho^{\circ})\le\varepsilon$ at the output state $\rho^{\circ}$ of every
  all-at-once attack with finite registers, then $f(\rho)\le\varepsilon$ at
  the output state $\rho$ of every attack;
\item \emph{(processing by the adversary)} let, in (ii), each $L_t$ keep the
  public messages and act on the other registers given their value $m$,
  $L_t(\dyad{m}\otimes x)=\dyad{m}\otimes L_{t,m}(x)$, and let
  $\Phi=\sum_m\dyad{m}\otimes\Phi_m$ act on a state that is classical on the
  public messages by a channel $\Phi_m$ on the adversary's register, given
  their value $m$, keeping them. Then $f((\idmap\otimes\Phi)(\sigma))\le
  f(\sigma)$ for every state $\sigma$ classical on the public messages.
\end{enumerate}
\end{digestclaim}

% claim meat-core/lem:honest-attacks
\digestlabel{lem:honest-attacks}{10.6}\pdfbookmark[2]{Lemma 10.6 Adversaries of the honest-authentication setting as round attacks}{digest-meat-core/lem:honest-attacks}
\renewcommand*{\RestatedName}{\hyperref{contract.pdf}{}{claim-meat-core/lem:honest-attacks}{Lemma~10.6}}%
\begin{digestclaim}[Adversaries of the honest-authentication setting as round attacks]
Let the run be as in Lemma~\ref{lem:memory-lift}\,(a), with A-ALICE-MEAS
(\xref{protocol-compiler/ass:alice-meas})\,(i) for Alice's measurement where
no source-replacement certificate is demanded, and let A-SEQ
(\xref{transmitter/ass:seq}), its clause (v) included,
Assumption~\ref{ass:lab} (A-LAB) at each of the two parties and
Assumption~\ref{ass:conform} (A-CONFORM) hold. Let an adversary act on the run in the honest-authentication setting
(Definition~\ref{def:postprocessing-model}) by any operations, at any times,
each chosen given the classical outcomes of her earlier ones, on her own
registers, of any dimension, on each transmitted system once the source has
released it and on her copies of the public messages once she holds them,
delivering the measured system of each round, and in untrusted-source mode
Alice's system, for the round's measurements. Then there are an all-at-once
attack in the sense of Lemma~\ref{lem:memory-lift}\,(b), with output state
$\rho$, and a channel $\Phi=\sum_m\dyad{m}\otimes\Phi_m$ that applies a
channel $\Phi_m$ to her register $E'_n$ given the value $m$ of the run's
public messages and keeps them, such that the state of the run's output
registers, its public messages and the adversary's final register is
$(\idmap\otimes\Phi)(\rho)$.
\end{digestclaim}

% claim meat-core/thm:keylength
\digestlabel{thm:keylength}{10.7}\pdfbookmark[2]{Theorem 10.7 Secrecy of the key from the accumulated entropy}{digest-meat-core/thm:keylength}
\begin{mainclaim}
\renewcommand*{\RestatedName}{\hyperref{contract.pdf}{}{claim-meat-core/thm:keylength}{Theorem~10.7}}%
\begin{digestclaim}[Secrecy of the key from the accumulated entropy]
Let the protocol run the fixed-length post-processing model
(Definition~\ref{def:postprocessing-model}) in the
honest-authentication setting, with the registers
\eqref{eq:pp-registers}. Let one of the two accumulation branches hold:
\emph{(unconstrained)} the hypotheses of
Theorem~\ref{thm:accumulated-entropy} with $\CList=\varnothing$; or
\emph{(constrained)} the hypotheses of
Theorem~\ref{thm:constrained-accumulation} ($\CList$
licensed for the declared source mode, Definition~\ref{def:homog};
A-ATK (\xref{protocol-compiler/ass:atk})), in either branch save the
condition $\Pr[\OmAT]>0$ that the accumulation theorem places on the
attack, which the attack below need not meet. Let the run be as in
Lemma~\ref{lem:memory-lift}\,(a), as A-SPEC (\xref{protocol-compiler/ass:spec})\,(i)
attests for its measurement, with A-ALICE-MEAS
(\xref{protocol-compiler/ass:alice-meas})\,(i) for Alice's measurement where
no source-replacement certificate is demanded, and A-ENS
(\xref{transmitter/ass:ens}), or the
premise a source model declares in its place, for its source states;
let it be attacked by an attack in the sense of
Lemma~\ref{lem:memory-lift}\,(b) that is all-at-once and has finite
registers; and let the reduction chain from which the declaration was
compiled carry it onto the round model, as A-LINK
(\xref{protocol-compiler/ass:link})\,(i)(a), (i)(b), (ii)(a) and (iv) attest,
the existence of each link's channel, the pull-back identity, the source
identity and complete positivity, with A-ATK in the constrained branch:
\begin{itemize}[leftmargin=2.4em]
\item[(H-CHAIN)] the marginal on $Z_1^n\Chat_1^nI_1^nE_n$ of the run's state
  after the $n$ rounds is the state the
  round channels of Definition~\ref{def:roundmodel}
  produce under an attack sequence of that definition,
  $\CList$-admissible in the constrained branch, from an input state as in
  the branch's accumulation theorem; the run's state is an extension of it,
  holding the parties' further records.
\end{itemize}
Let Assumptions~\ref{ass:keyreg}, \ref{ass:hash}, \ref{ass:seed},
\ref{ass:fixed}, \ref{ass:conform} and~\ref{ass:lab}, the last at each of
the two parties, and A-SEQ (\xref{transmitter/ass:seq}) hold, let
$\alpha\in(1,2)$ and $\epa,\eev\in(0,1]$. If the output length satisfies
\begin{equation*}\tag{\getrefnumber{C-eq:keylength}}\digesteqlabel{eq:keylength}
  \ell\;\le\;n\,\hlow-\lEC-\Big\lceil\log_2\frac1{\eev}\Big\rceil
  -\frac{\alpha}{\alpha-1}\log_2\frac1{\epa}+1,
\end{equation*}
then, under that attack, the protocol is $\epa$-secret on the accepted
state:
\begin{equation*}
  \frac12\Big\|\rho_{K\,E_{\mathrm{tot}}\wedge\Omega_{\mathrm{acc}}}
  -\frac{\idon{K}}{2^{\ell}}\otimes
  \rho_{E_{\mathrm{tot}}\wedge\Omega_{\mathrm{acc}}}\Big\|_1\le\epa .
\end{equation*}
\end{digestclaim}
\end{mainclaim}

% claim meat-core/cor:composed-security
\digestlabel{cor:composed-security}{10.8}\pdfbookmark[2]{Corollary 10.8 Security of the protocol and its net key}{digest-meat-core/cor:composed-security}
\begin{mainclaim}
\renewcommand*{\RestatedName}{\hyperref{contract.pdf}{}{claim-meat-core/cor:composed-security}{Corollary~10.8}}%
\begin{digestclaim}[Security of the protocol and its net key]
Let the hypotheses of Theorem~\ref{thm:keylength}, the run as in
Lemma~\ref{lem:memory-lift}\,(a) and (H-CHAIN) among them, hold at every
attack in the sense of Lemma~\ref{lem:memory-lift}\,(b) that is all-at-once
and has finite registers, with $\ell\ge1$ satisfying \eqref{eq:keylength},
A-SEQ with its clause (v) among them; let those of
Lemma~\ref{lem:ev-correctness} hold; and put
$\tau_{K_AK_B}\coloneqq2^{-\ell}\sum_{k\in\{0,1\}^{\ell}}
\dyad{k}_{K_A}\otimes\dyad{k}_{K_B}$.
\begin{enumerate}[label=(\roman*),leftmargin=2.4em]
\item In the honest-authentication setting, at every output state the
  setting admits, whatever the dimension of the adversary's registers and the
  order in which the source releases its signals, the protocol is
  $\epa$-secret on the accepted state as Theorem~\ref{thm:keylength} states
  it, and $(\epa+\eev)$-secure in the sense of
  \cite[Sec.~2.3, Def.~2.1, Eq.~(1)]{tupkary2025survey}:
  \begin{equation*}\tag{\getrefnumber{C-eq:composed-honest}}\digesteqlabel{eq:composed-honest}
    \frac12\Big\|\rho_{K_AK_B\,E_{\mathrm{tot}}\wedge\Omega_{\mathrm{acc}}}
    -\tau_{K_AK_B}\otimes
    \rho_{E_{\mathrm{tot}}\wedge\Omega_{\mathrm{acc}}}\Big\|_1
    \le\epa+\eev ,
  \end{equation*}
  and on the complement of $\Omega_{\mathrm{acc}}$ both parties output
  $\perp$.
\item Run with the authentication post-processing
  (Definition~\ref{def:authenticated-channel}\,(c)) over the
  authenticated classical channel
  (Definition~\ref{def:authenticated-channel}\,(a)), the protocol is
  $(\epa+\eev)$-secure in the sense of
  \cite[Sec.~3.3, Def.~3.8]{tupkary2026decoy}, which admits asymmetric
  aborts, for every output state that channel admits.
\item Let Assumptions~\ref{ass:auth}, \ref{ass:fixed}
  and~\ref{ass:conform} hold, with $\varepsilon_{\mathrm{auth}}$ and
  $k_{\mathrm{auth}}$ the parameters of Assumption~\ref{ass:auth}. Run with
  the authentication post-processing over the classical channel of
  Assumption~\ref{ass:auth}, the protocol is
  $(\epa+\eev+2\varepsilon_{\mathrm{auth}})$-secure in the sense of
  \cite[Sec.~3.3, Def.~3.8]{tupkary2026decoy}: for every attack on that run,
  with output state $\rho'$ on the key registers and the adversary's final
  information,
  \begin{equation*}\tag{\getrefnumber{C-eq:composed-auth}}\digesteqlabel{eq:composed-auth}
    \frac12\big\|\rho'-\mathcal R_{\mathrm{ideal}}(\rho')\big\|_1
    \le\epa+\eev+2\varepsilon_{\mathrm{auth}},
  \end{equation*}
  with $\mathcal R_{\mathrm{ideal}}$ the map of
  \cite[Sec.~3.3, Eqs.~(13)--(17)]{tupkary2026decoy}. It consumes at most
  $k_{\mathrm{auth}}$ bits of the pre-shared key, and when both parties
  output an $\ell$-bit key its \emph{net key gain} is
  $\ell_{\mathrm{net}}\coloneqq\ell-k_{\mathrm{auth}}$ bits; for
  $\ell_{\mathrm{net}}\le0$ the run yields no net key.
\end{enumerate}
\end{digestclaim}
\end{mainclaim}

% claim meat-core/cor:separate-executions
\digestlabel{cor:separate-executions}{10.9}\pdfbookmark[2]{Corollary 10.9 Security parameters of separate executions add}{digest-meat-core/cor:separate-executions}
\renewcommand*{\RestatedName}{\hyperref{contract.pdf}{}{claim-meat-core/cor:separate-executions}{Corollary~10.9}}%
\begin{digestclaim}[Security parameters of separate executions add]
Let executions $1,\dots,N$ of the protocol run against one adversary, in
sequence, concurrently or interleaved, the adversary carrying information
from each to the others. For each $i$ let the hypotheses of
Corollary~\ref{cor:composed-security}\,(iii) hold for execution~$i$, with
its parameters $\epa^{(i)}$, $\eev^{(i)}$ and
$\varepsilon_{\mathrm{auth}}^{(i)}$. Write
$K^{(i)}\coloneqq K_A^{(i)}K_B^{(i)}$ for its key registers,
$\mathcal R^{(i)}$ for the map $\mathcal R_{\mathrm{ideal}}$ of
\cite[Sec.~3.3, Eqs.~(13)--(17)]{tupkary2026decoy} on $K^{(i)}$, and
$\varepsilon_i\coloneqq\epa^{(i)}+\eev^{(i)}+2\varepsilon_{\mathrm{auth}}^{(i)}$.
Let Assumption~\ref{ass:separate-executions} hold. Then for every attack on
the $N$ executions under which each of them reaches its
end of an execution (Definition~\ref{C-def:execution-end}) on every
branch, with $\rho$ the joint state of $K^{(1)}\cdots K^{(N)}$ and the
adversary's information $E$ at the end of the last of them,
\begin{equation*}\tag{\getrefnumber{C-eq:separate-executions}}\digesteqlabel{eq:separate-executions}
  \frac12\Big\|\rho-\big(\mathcal R^{(1)}\otimes\cdots\otimes
  \mathcal R^{(N)}\otimes\idmap_E\big)(\rho)\Big\|_1
  \le\sum_{i=1}^{N}\varepsilon_i .
\end{equation*}
\end{digestclaim}

% claim meat-core/prop:model-pass-bound
\digestlabel{prop:model-pass-bound}{10.12}\pdfbookmark[2]{Proposition 10.12 The pass probability of devices the declared model admits}{digest-meat-core/prop:model-pass-bound}
\renewcommand*{\RestatedName}{\hyperref{contract.pdf}{}{claim-meat-core/prop:model-pass-bound}{Proposition~10.12}}%
\begin{digestclaim}[The pass probability of devices the declared model admits]
Let $\Decl$ and $\acSet$ satisfy the hypotheses of
Proposition~\ref{prop:minimal-divergence}, let $\sfD_{\min}$ be as in
\eqref{eq:dmin} and $\mathcal{D}_{\mathrm{all}}$ the dropped set of every
pair there, and let $n\ge1$ be an integer. Let $\Pr$ be a probability
distribution on $\Chatset^{\,n}$, write $\Chat_1,\dots,\Chat_n$ for the
coordinates and $\Chat_1^{t}\coloneqq(\Chat_1,\dots,\Chat_t)$, and suppose:
\begin{description}
\item[(H)] for every $t\in\{1,\dots,n\}$ and every $h\in\Chatset^{\,t-1}$
  with $\Pr[\Chat_1^{t-1}=h]>0$ there is a state $\omega_t(h)\in\Feas$ with
  \begin{equation*}\tag{\getrefnumber{C-eq:pass-hypothesis}}\digesteqlabel{eq:pass-hypothesis}
    \Pr\big[\Chat_t=\hat c\ \big|\ \Chat_1^{t-1}=h\big]
    =\Phimap[\omega_t(h)](\hat c)\qquad\text{for every }\hat c\in\Chatset ,
  \end{equation*}
  where $\Phimap[\omega](\hat c)$ is the entry of $\Phimap[\omega]$ at
  $\hat c$ and, for $t=1$, $h$ is the empty sequence and the left-hand side
  is $\Pr[\Chat_1=\hat c]$.
\end{description}
The case of independent rounds with one state $\omega\in\Feas$ is
$\omega_t(h)=\omega$ for all $t$ and $h$. Let
$\OmAT=\{\operatorname{freq}_{\Chat_1^n}\in\acSet\}$ be the event of
Definition~\ref{def:postprocessing-model}\,(i), where
$\operatorname{freq}_{\hat c_1^n}(\hat c)\coloneqq
\frac1n\,|\{t:\hat c_t=\hat c\}|$.
\begin{enumerate}[label=(\roman*),leftmargin=2.2em]
\item \emph{(The bound.)} For every real number $\delta\le\sfD_{\min}$,
  \begin{equation*}\tag{\getrefnumber{C-eq:pass-bound}}\digesteqlabel{eq:pass-bound}
    \Pr[\OmAT]\;\le\;2^{-n\delta} .
  \end{equation*}
  In particular $\Pr[\OmAT]=0$ if $\sfD_{\min}=+\infty$.
\item \emph{(From a bound at one order.)} Let $\alpha_0\in(1,\infty)$ and let
  $b$ be a real number with
  $b\le\underline{h}^{\mathcal{D}_{\mathrm{all}}}_{\alpha_0}$. Then
  $\Pr[\OmAT]\le2^{-x}$ for every real number
  $x\le n\,\frac{\alpha_0-1}{\alpha_0}\,b$. For every
  $\alpha'\in(1,\infty)$ the number
  $f(\alpha')\coloneqq\frac{\alpha'}{\alpha'-1}\cdot\frac{\alpha_0-1}{\alpha_0}\,b$
  satisfies $\frac{\alpha'-1}{\alpha'}\,n\,f(\alpha')
  =n\,\frac{\alpha_0-1}{\alpha_0}\,b$.
\end{enumerate}
\end{digestclaim}

% claim meat-core/def:max-key-length
\digestlabel{def:max-key-length}{10.13}\pdfbookmark[2]{Definition 10.13 Maximal key length}{digest-meat-core/def:max-key-length}
\renewcommand*{\RestatedName}{\hyperref{contract.pdf}{}{claim-meat-core/def:max-key-length}{Definition~10.13}}%
\begin{digestclaim}[Maximal key length]
The
\emph{maximal key length}
is
\[
  \ell_{\max}\coloneqq\max\{0,\lfloor
  n\hlow-\lEC-\lceil\log_2\tfrac1{\eev}\rceil
  -\tfrac{\alpha}{\alpha-1}\log_2\tfrac1{\epa}+1\rfloor\},
\]
the right-hand side of \eqref{eq:keylength} under a floor and a clamp at
$0$; the floor and the clamp are both in the safe direction: an integer
at most the floor of a number is at most that number, and a key of length
$0$ is trivially secret and correct and claims no key.
\end{digestclaim}

% claim meat-core/prop:monotone
\digestlabel{prop:monotone}{10.17}\pdfbookmark[2]{Proposition 10.17 Monotonicity of the maximal key length}{digest-meat-core/prop:monotone}
\renewcommand*{\RestatedName}{\hyperref{contract.pdf}{}{claim-meat-core/prop:monotone}{Proposition~10.17}}%
\begin{digestclaim}[Monotonicity of the maximal key length]
The maximal key length $\ell_{\max}$ of
Definition~\ref{def:max-key-length}
is
\begin{enumerate}[label=(\roman*),leftmargin=2.4em]
\item non-increasing in $\lEC$;
\item non-decreasing in $\eev$;
\item non-decreasing in $\epa$;
\item unchanged under any modification of the declaration that leaves
  $\hlow$ unchanged.
\end{enumerate}
Consequently increasing leakage or tightening either error budget can never
increase $\ell_{\max}$.
\end{digestclaim}

% claim meat-core/cor:keylength-bound
\digestlabel{cor:keylength-bound}{10.18}\pdfbookmark[2]{Corollary 10.18 Key length from any lower bound on the value}{digest-meat-core/cor:keylength-bound}
\begin{mainclaim}
\renewcommand*{\RestatedName}{\hyperref{contract.pdf}{}{claim-meat-core/cor:keylength-bound}{Corollary~10.18}}%
\begin{digestclaim}[Key length from any lower bound on the value]
Let $\alpha\in(1,2)$, let the hypotheses of Theorem~\ref{thm:keylength}
hold in its unconstrained or its constrained branch, the run as in
Lemma~\ref{lem:memory-lift}\,(a) and hypothesis (H-CHAIN) included, at every
all-at-once attack with finite registers, A-SEQ with its clause (v) among
them, and let those of Lemma~\ref{lem:ev-correctness} hold. Let
$h\in\mathbb{R}\cup\{-\infty\}$ satisfy
\begin{equation*}\tag{\getrefnumber{C-eq:lower-bound-h}}\digesteqlabel{eq:lower-bound-h}
  h\;\le\;\hlow ,
\end{equation*}
with $\hlow$ the exact value \eqref{eq:hlow-def} of the declaration (for
the constrained branch: with $\Feas$ carrying the declared $\CList$). If the
output length satisfies
\begin{equation*}\tag{\getrefnumber{C-eq:keylength-h}}\digesteqlabel{eq:keylength-h}
  \ell\;\le\;n\,\max\{0,h\}-\lEC-\Big\lceil\log_2\frac1{\eev}\Big\rceil
  -\frac{\alpha}{\alpha-1}\log_2\frac1{\epa}+1,
\end{equation*}
then the protocol is $\epa$-secret and $\eev$-correct on the accepted
state, hence, for $\ell\ge1$, $(\epa+\eev)$-secure in the
honest-authentication setting; under the further hypotheses of
Corollary~\ref{cor:composed-security}\,(iii) it is
$(\epa+\eev+2\varepsilon_{\mathrm{auth}})$-secure over the classical channel
of that part, with net key gain $\ell-k_{\mathrm{auth}}$ when both parties
output a key. The secrecy and the security statements hold at every output
state, whatever the dimension of the adversary's registers and the order in
which the source releases its signals, as
Corollary~\ref{cor:composed-security} states them, and the correctness
statement at every attack. Moreover:
\begin{enumerate}[label=(\alph*),leftmargin=2.2em]
\item if $h$ depends on none of $(\lEC,\eev,\epa)$, the four conclusions of
  Proposition~\ref{prop:monotone} hold verbatim for the maximal key length
  (Definition~\ref{def:max-key-length}) built from \eqref{eq:keylength-h},
  with $\max\{0,h\}$ in place of $\hlow$;
\item if $h\le0$ and $\eev<1$ then no positive $\ell$ satisfies
  \eqref{eq:keylength-h} and no key is claimed (at $\eev=1$ the correctness
  statement is empty and the corner carries no security claim).
\end{enumerate}
\end{digestclaim}
\end{mainclaim}

% claim meat-core/cor:constrained-key-length
\digestlabel{cor:constrained-key-length}{10.19}\pdfbookmark[2]{Corollary 10.19 Key length from the constrained value}{digest-meat-core/cor:constrained-key-length}
\renewcommand*{\RestatedName}{\hyperref{contract.pdf}{}{claim-meat-core/cor:constrained-key-length}{Corollary~10.19}}%
\begin{digestclaim}[Key length from the constrained value]
Under the hypotheses of Theorem~\ref{thm:constrained-accumulation}, with
Corollary~\ref{cor:reduction-constrained},
\begin{equation*}
  \Hsandup\big(Z_1^n\big|\Chat_1^nI_1^n\Ehat\big)_{\rho_{|\OmAT}}
  \;\ge\;n\,\hlow(\CList)
  -\frac{\alpha}{\alpha-1}\log_2\frac{1}{\Pr[\OmAT]},
\end{equation*}
which is the form Theorem~\ref{thm:keylength} consumes; under
A-ATK (\xref{protocol-compiler/ass:atk}) and the remaining
hypotheses of Theorem~\ref{thm:keylength}, the key-length statement
of Theorem~\ref{thm:keylength} therefore holds with the constrained $\hlow$
(its constrained branch).
\end{digestclaim}

\section*{11\quad Perturbations of the declared measurement}\pdfbookmark[1]{11 Perturbations of the declared measurement}{digest-group-meat-core/prop:optol-nobound}

% claim meat-core/prop:optol-nobound
\digestlabel{prop:optol-nobound}{11.1}\pdfbookmark[2]{Proposition 11.1 Instability of the rate under measurement perturbations}{digest-meat-core/prop:optol-nobound}
\renewcommand*{\RestatedName}{\hyperref{contract.pdf}{}{claim-meat-core/prop:optol-nobound}{Proposition~11.1}}%
\begin{digestclaim}[Instability of the rate under measurement perturbations]
Let $\alpha\in(1,\infty)$ and $B>0$, write $N\coloneqq\lceil B\rceil$, and
let $\epsilon\in(0,\tfrac12]$ satisfy $2^{N}\le\epsilon\,2^{53}$.
Write $\Decl^{(i)}$ for the declaration obtained from a given one by
replacing its measurement $\{F_v\}$ on $\Qsys$ by $\{F^{(i)}_v\}$, and
$\hlow[\,\cdot\,]$ for the value \eqref{eq:hlow-def} of a declaration.
\begin{enumerate}[label=(\roman*),leftmargin=2.4em]
\item There is a declaration $\Decl$ satisfying Definition~\ref{def:admissible-declaration}
in full, with $\{F_v\}$ itself exactly admissible, and there are two
exactly admissible families $\{F^{(1)}_v\}$, $\{F^{(2)}_v\}$ with
$\max_{v}\|F^{(i)}_v-F_v\|\le\epsilon$ for $i=1,2$, such that
\begin{equation*}\tag{\getrefnumber{C-eq:optol-gap}}\digesteqlabel{eq:optol-gap}
  \hlow\big[\Decl^{(1)}\big]-\hlow\big[\Decl^{(2)}\big]\;\ge\;B .
\end{equation*}
The dimensions and alphabets of $\Decl$ are $d_A=d_{\Qsys}=1$, one
sector, $|\Xset|=1$, $|\Vset|=|\Chatset|=2$ and $d_Z=1$, the same for
every $B$ and every $\epsilon$.
\item Let $\mathrm{err}$ be a function of the declared dimensions,
alphabet sizes and $\alpha$ whose value on the data of (i) is less than
$\frac{\alpha}{\alpha-1}\,N$. Then the bound
$\big|\hlow[\Decl^{(1)}]-\hlow[\Decl^{(2)}]\big|\le\mathrm{err}$, with
$\mathrm{err}$ evaluated on the data of the declaration, fails for some
declaration and some pair of exactly admissible families within
$\epsilon$ of each other.
\item For $\alpha\in(1,2)$ and every number $n$ of rounds, the
maximal admissible key lengths $\ell_{\max}$ of
$\Decl^{(1)}$ and $\Decl^{(2)}$
(Definition~\ref{def:max-key-length}), computed with the same $\lEC$, $\eev$ and
$\epa$, differ by at least $nB-1$ bits whenever the smaller of the two is
positive.
\end{enumerate}
\end{digestclaim}

\section*{12\quad The reference companion}\pdfbookmark[1]{12 The reference companion}{digest-group-meat-core/prop:reference-companion}

% claim meat-core/prop:reference-companion
\digestlabel{prop:reference-companion}{12.1}\pdfbookmark[2]{Proposition 12.1 The compiled reference companion is an admissible declaration}{digest-meat-core/prop:reference-companion}
\renewcommand*{\RestatedName}{\hyperref{contract.pdf}{}{claim-meat-core/prop:reference-companion}{Proposition~12.1}}%
\begin{digestclaim}[The compiled reference companion is an admissible declaration]
Let $\Decl$ be an admissible declaration
(Definition~\ref{def:admissible-declaration}) that the compilation map
(\xref{protocol-compiler/def:compile}) emits from a specification, and let
$\mathbb{F}^{\mathrm{ref}}$ be a non-empty set of its
measured factors (the one measured system in
the basic format) such that, for each $f\in\mathbb{F}^{\mathrm{ref}}$,
the last link of the receiver segment of $f$ is the emission of a
reduction instance $\mathcal{J}_f$ satisfying (H-REF) of
\xref{reductions/cor:reference-companion}. Let $\Decl^{\mathrm{ref}}$ be the
tuple the compilation map gives from the same specification with, for every
$f\in\mathbb{F}^{\mathrm{ref}}$, that link's target family replaced by
the target $\{F_v[A^{\mathrm{ref}}_f]\}$ of the reference companion of
$\mathcal{J}_f$ of that corollary, $A^{\mathrm{ref}}_f$ its datum, and that
link's entries replaced by the relaxed entries
(\xref{reductions/lem:row-relaxation}) of the companion's entries, each
rounded operator Hermitian and with the sector blocks of the entry it rounds.
The companion's entries of such a link are the instance's own entries at the
companion's datum and the companion forms
(\xref{reductions/cor:reference-companion}\,(vi)) of those entries requested
of the link whose companion forms carry their certificates at that datum; a
requested entry whose companion form is not established to carry one has no
entry in $\Decl^{\mathrm{ref}}$. Then:
\begin{enumerate}[label=(\roman*),leftmargin=2.4em]
\item \emph{(what is shared)} $\Decl^{\mathrm{ref}}$ has the systems,
  sectors and dimensions, the alphabets, Alice's marginal slot and
  measurement, the party partition, the announcement and sifting maps, the
  raw-key map and the index map of $\Decl$, and in the general format its
  factor list and Alice's block lists; it differs from
  $\Decl$ at most in the measured families of the factors of
  $\mathbb{F}^{\mathrm{ref}}$ and in the entries of their last links. In
  particular the two have the same joint sectors
  (Definition~\ref{def:joint-sectors}) with the same block dimensions and the
  same coupling condition;
\item \emph{(admissible)} $\Decl^{\mathrm{ref}}$ satisfies conditions
  (i)--(vi) and (viii) of Definition~\ref{def:admissible-declaration};
\item \emph{(the relaxed entries enlarge the feasible set)} let
  $\Decl^{\mathrm{ref}}_{\mathrm{ex}}$ be formed as $\Decl^{\mathrm{ref}}$
  but with the companions' entries themselves in place of their relaxed
  entries; then the feasible set of $\Decl^{\mathrm{ref}}_{\mathrm{ex}}$
  (Definition~\ref{def:feasible}) is contained in that of
  $\Decl^{\mathrm{ref}}$, and for every acceptance set
  $\hlow[\Decl^{\mathrm{ref}}]\le\hlow[\Decl^{\mathrm{ref}}_{\mathrm{ex}}]$.
\end{enumerate}
\end{digestclaim}

% claim meat-core/cor:reference-companion-factors
\digestlabel{cor:reference-companion-factors}{12.2}\pdfbookmark[2]{Corollary 12.2 Factor enclosures of the compiled reference companion}{digest-meat-core/cor:reference-companion-factors}
\renewcommand*{\RestatedName}{\hyperref{contract.pdf}{}{claim-meat-core/cor:reference-companion-factors}{Corollary~12.2}}%
\begin{digestclaim}[Factor enclosures of the compiled reference companion]
Let $\Decl$, $\mathbb{F}^{\mathrm{ref}}$, the instances $\mathcal{J}_f$ and
$\Decl^{\mathrm{ref}}$ be as in Proposition~\ref{prop:reference-companion},
and for $f\in\mathbb{F}^{\mathrm{ref}}$ let $G_f$ be the family of
\xref{reductions/cor:reference-companion} from which the datum
$A^{\mathrm{ref}}_f$ of the reference companion of $\mathcal{J}_f$ is
formed. Let $f\in\mathbb{F}^{\mathrm{ref}}$, and read the block on a
data sector $s$ of $\mathcal{J}_f$'s target of an element of the
measured family of factor $f$ of $\Decl^{\mathrm{ref}}$ formed from the
target element $F_v[A^{\mathrm{ref}}_f]$ with the triple that part (i) of
\xref{reductions/cor:companion-factor-enclosure} forms for the element
$A^{\mathrm{ref},(s)}_{f,v}$ of $A^{\mathrm{ref}}_f$ on $s$, where the rule
of (I1) of $\mathcal{J}_f$ is a non-negative combination (that corollary)
and that part forms one from factor enclosures
(\xref{foundations/def:factor-enclosure}) that blocks of $G_f$ carry, its
core multiplied by $P_f(\varphi\,|\,t(v))$ where the family is weighted after
the chain by factor $f$'s weighting $P_f$ and the element is
$P_f(\varphi\,|\,t(v))\,F_v[A^{\mathrm{ref}}_f]$ with
$(v,\varphi)\in\Vset_{P_f}$ (\xref{protocol-compiler/def:weighted}, with
$P\coloneqq P_f$); and read every other block of the
measured families of $\Decl^{\mathrm{ref}}$ with none. Then every block so
read with a triple is a member of it, the triple being a factor enclosure of
that block.
\end{digestclaim}

\section*{13\quad Assumptions}\pdfbookmark[1]{13 Assumptions}{digest-group-meat-core/ass:keyreg}

% claim meat-core/ass:keyreg
\digestlabel{ass:keyreg}{13.1}\pdfbookmark[2]{Assumption 13.1 Dedicated key register and one-way reconciliation}{digest-meat-core/ass:keyreg}
\SetShortId{A-KEYREG}
\renewcommand*{\RestatedName}{\hyperref{contract.pdf}{}{claim-meat-core/ass:keyreg}{Assumption~13.1}}%
\begin{digestassumption}[Dedicated key register and one-way reconciliation]
(i) Error correction, error verification and privacy amplification are
applied to $Z_1^n$, the string of raw-key values $r(x,v)$ that the rounds
write to the registers $Z_1,\dots,Z_n$
(Definition~\ref{def:postprocessing-model});
(ii) the declared key holder $\hold$ (Definition~\ref{def:declaration})
computes $Z_1^n$ from its own outcome record and the public announcements;
(iii) reconciliation is one-way from $\hold$: its transcript
$L_{\mathrm{EC}}$ is a function of $Z_1^n$ and public seeds only;
(iv) $L_{\mathrm{EC}}$ occupies a fixed-length register of $\lEC$ bits,
shorter transcripts being padded, so that $\dim(L_{\mathrm{EC}})=2^{\lEC}$.
\end{digestassumption}

% claim meat-core/ass:hash
\digestlabel{ass:hash}{13.2}\pdfbookmark[2]{Assumption 13.2 Ideal two-universal privacy amplification}{digest-meat-core/ass:hash}
\SetShortId{A-HASH}
\renewcommand*{\RestatedName}{\hyperref{contract.pdf}{}{claim-meat-core/ass:hash}{Assumption~13.2}}%
\begin{digestassumption}[Ideal two-universal privacy amplification]
(i) The privacy-amplification family of
Definition~\ref{def:postprocessing-model}, applied to $Z_1^n$, is ideal
two-universal: for any two distinct inputs, the probability under the
family's law that their $\ell$-bit values coincide is exactly $2^{-\ell}$;
(ii) the error-verification family of
Definition~\ref{def:postprocessing-model} is two-universal: for any two
distinct inputs, the probability under the family's law that their
$\lceil\log_2\frac1{\eev}\rceil$-bit values coincide is at most
$2^{-\lceil\log_2(1/\eev)\rceil}$.
\end{digestassumption}

% claim meat-core/ass:seed
\digestlabel{ass:seed}{13.3}\pdfbookmark[2]{Assumption 13.3 Private, independent and fresh local randomness}{digest-meat-core/ass:seed}
\SetShortId{A-SEED}
\renewcommand*{\RestatedName}{\hyperref{contract.pdf}{}{claim-meat-core/ass:seed}{Assumption~13.3}}%
\begin{digestassumption}[Private, independent and fresh local randomness]
(i) The error-correction seed $S_{\mathrm{pub}}$, the error-verification seed
$S_{\mathrm{EV}}$ and the privacy-amplification hash $G$ of
Definition~\ref{def:postprocessing-model} are drawn with the laws that
definition gives them, $S_{\mathrm{EV}}$ and $G$ with the laws of their hash
families;
(ii) each is drawn privately: its value is available to the adversary only
once the protocol announces it;
(iii) each is independent of every other register of the protocol and of the
adversary's information at the time it is drawn;
(iv) each is drawn afresh in every run;
(v) in trusted-source mode, the letters Alice's source draws in the
$n$ rounds, and the systems it prepares from them, are jointly independent of
the adversary's registers at the start of the protocol: the joint state of
the letters, the prepared systems and those registers is the product of its
marginal on the letters and the prepared systems together and its marginal on
those registers; in particular the adversary holds no purification of a
prepared system, and no purification of a prepared system is made available
to the adversary during or after the protocol.
\end{digestassumption}

% claim meat-core/ass:fixed
\digestlabel{ass:fixed}{13.4}\pdfbookmark[2]{Assumption 13.4 Protocol parameters fixed before the run}{digest-meat-core/ass:fixed}
\SetShortId{A-FIXED}
\renewcommand*{\RestatedName}{\hyperref{contract.pdf}{}{claim-meat-core/ass:fixed}{Assumption~13.4}}%
\begin{digestassumption}[Protocol parameters fixed before the run]
(i) The acceptance set $\acSet$, the round count $n$, the output length
$\ell$, the transcript length $\lEC$, the error budgets $\epa$ and $\eev$ of
Definition~\ref{def:postprocessing-model}, and the authentication parameters
$\varepsilon_{\mathrm{auth}}$ and $k_{\mathrm{auth}}$, the declared
authentication error bound and key consumption
(Definition~\ref{def:authenticated-channel}\,(b)), are fixed before the
first round of the run;
(ii) the run applies them unchanged: none of them is chosen, adjusted or
recomputed from data observed in the run;
(iii) the R\'enyi order $\alpha\in(1,2)$ at which $\ell$ is computed is
fixed before the first round of the run, from data available before it, and
is not chosen, adjusted or recomputed from data observed in the run.
\end{digestassumption}

% claim meat-core/ass:conform
\digestlabel{ass:conform}{13.5}\pdfbookmark[2]{Assumption 13.5 The implementation runs the analysed post-processing}{digest-meat-core/ass:conform}
\SetShortId{A-CONFORM}
\renewcommand*{\RestatedName}{\hyperref{contract.pdf}{}{claim-meat-core/ass:conform}{Assumption~13.5}}%
\begin{digestassumption}[The implementation runs the analysed post-processing]
(i) The implementation evaluates the acceptance test of
Definition~\ref{def:postprocessing-model}\,(i) and aborts unless $\OmAT$
holds;
(ii) after the rounds, the only classical messages either party sends are
those of Definition~\ref{def:postprocessing-model} and of the
authentication post-processing
(Definition~\ref{def:authenticated-channel}\,(c));
(iii) both parties abort when the verification flag is $\mathrm{rej}$;
(iv) each party outputs the $\ell$-bit key that
Definition~\ref{def:postprocessing-model}\,(iv) assigns to it, or $\perp$;
(v) neither party uses any part of its key before the
end of an execution (Definition~\ref{C-def:execution-end});
(vi) a party that receives an authentication failure acts as the
authentication post-processing prescribes.
\end{digestassumption}

% claim meat-core/ass:auth
\digestlabel{ass:auth}{13.6}\pdfbookmark[2]{Assumption 13.6 Authentication of the classical channel by a keyed scheme}{digest-meat-core/ass:auth}
\SetShortId{A-AUTH}
\renewcommand*{\RestatedName}{\hyperref{contract.pdf}{}{claim-meat-core/ass:auth}{Assumption~13.6}}%
\begin{digestassumption}[Authentication of the classical channel by a keyed scheme]
(i) All classical messages of the protocol, the round announcements
included, are sent through one authentication scheme using a pre-shared key
and an unauthenticated classical channel, or all over a classical channel
that is itself the authenticated classical channel of
Definition~\ref{def:authenticated-channel}\,(a), in which case
$\varepsilon_{\mathrm{auth}}=k_{\mathrm{auth}}=0$;
(ii) in the first case, the scheme's key consumption is at most
$k_{\mathrm{auth}}$ bits;
(iii) run on two identical copies of a uniformly distributed key independent
of every other register of the protocol and of the adversary's information,
the scheme $\varepsilon_0$-realises the authenticated classical channel
(Definition~\ref{def:authenticated-channel}\,(b)), for some
$\varepsilon_0\ge0$;
(iv) before the run starts, the joint state of the two parties' copies of
the pre-shared key the run consumes with every other register of the
protocol and with the adversary's information is within trace distance
$\delta\ge0$ of the state in which those copies are replaced by two
identical copies of a uniformly distributed key independent of them;
(v) $\varepsilon_0+\delta\le\varepsilon_{\mathrm{auth}}$.
\end{digestassumption}

% claim meat-core/ass:separate-executions
\digestlabel{ass:separate-executions}{13.7}\pdfbookmark[2]{Assumption 13.7 Executions that share no key and no device state}{digest-meat-core/ass:separate-executions}
\SetShortId{A-SEPARATE}
\renewcommand*{\RestatedName}{\hyperref{contract.pdf}{}{claim-meat-core/ass:separate-executions}{Assumption~13.7}}%
\begin{digestassumption}[Executions that share no key and no device state]
For executions $1,\dots,N$ of the protocol, run against one adversary:
(i) $N$ and, for every execution, the parameters fixed before its run (its
acceptance set, round count, output length, transcript length, error
budgets and authentication parameters) are fixed before the first of the
$N$ executions starts;
(ii) no operation of the honest parties of one execution acts on a register
that the honest parties of another execution use, hold or output: no
execution uses any part of another execution's key, local randomness,
seeds or pre-shared key, no device, memory or random source carries a
state from one execution into another, and the executions interact only
through systems and classical channels the adversary has access to;
(iii) from the end of an execution
(Definition~\ref{C-def:execution-end}) until the end of the last of the $N$
executions, no part of that execution's key is used, read or otherwise
acted on, and no information about any other register of its parties
leaves their laboratories;
(iv) for every execution, every assumption on which its key and its
security parameter rest holds when every register of the honest parties of
every other execution is counted as part of the adversary's information,
and every operation of those parties as part of the adversary's actions.
\end{digestassumption}

% claim meat-core/ass:lab
\digestlabel{ass:lab}{13.8}\pdfbookmark[2]{Assumption 13.8 Each party's laboratory is sealed}{digest-meat-core/ass:lab}
\SetShortId{A-LAB}
\renewcommand*{\RestatedName}{\hyperref{contract.pdf}{}{claim-meat-core/ass:lab}{Assumption~13.8}}%
\begin{digestassumption}[Each party's laboratory is sealed]
Let $\mathsf P\in\{\mathsf A,\mathsf B\}$ be a party of the
protocol declaration (Definition~\ref{def:declaration}). Its
\emph{records} are its letters, where the party partition assigns
Alice's system $A$ to $\mathsf P$; the private settings and the outcomes of
the measurements made in its laboratory, the public choices it draws before
they are announced, and its readout $\rd{\mathsf P}$; the inputs of its
announcements; its raw or final key; and its local randomness. The \emph{systems handed out of its laboratory} are, in
trusted-source mode and where the party partition assigns $A$ to
$\mathsf P$, the systems its source emits into the quantum channel, which the
round model (Definition~\ref{def:roundmodel}) gives to the adversary,
$A$ being no output of the attack in that mode \eqref{eq:attackchannel} and
staying with $\mathsf P$; in every other case there are none. In particular,
in untrusted-source mode, where the attack supplies $A$ and $\Qsys$
\eqref{eq:attackchannel}, no laboratory hands out any system, and a source the declaration treats as untrusted lies
outside every laboratory wherever it physically stands. The \emph{systems
delivered into its laboratory} are, in untrusted-source mode and where the
party partition assigns $A$ to $\mathsf P$, the system $A$; and, in either
mode and in each sector $s$, the tensor factors $\Qsys_{s,i}$ of the declared
factorisation of $\Hs{\Qsys_s}$ that the party partition labels $\mathsf P$
(in the general format, every measured factor labelled $\mathsf P$),
physically the input modes those factors model.
\begin{enumerate}[label=(\roman*),leftmargin=2.2em]
\item During the protocol, no information about $\mathsf P$'s records leaves
its laboratory, other than its key as the protocol's output, except through
the systems handed out of its laboratory and through the public messages of
the protocol (the announcements of the
round channels of Definition~\ref{def:roundmodel}, the
messages of Definition~\ref{def:postprocessing-model} and of the
authentication post-processing
(Definition~\ref{def:authenticated-channel}\,(c)), and, where an
authentication scheme carries them, the scheme's own transmissions over the
unauthenticated classical channel): no light emitted by its devices other
than those systems, no electromagnetic, acoustic, power or timing emanation
and no side channel of its classical processing carries such information, and
neither the timing of its messages nor any auxiliary or synchronisation
signal does;
\item the adversary acts on $\mathsf P$'s laboratory only through the systems
delivered into its laboratory, as the declared measurements receive them, and
through the public messages the protocol delivers to $\mathsf P$ and, where a
scheme carries them, the transmissions it receives over the unauthenticated
channel: no light or signal sent into the laboratory, through the systems
handed out of it or otherwise, reads out any of its records, and nothing else
influences the state, configuration or calibration of its devices (its source
included, where it has one), its measurements or its classical processing.
\end{enumerate}
\end{digestassumption}

% claim meat-core/ass:eve-dim
\digestlabel{ass:eve-dim}{13.9}\pdfbookmark[2]{Assumption 13.9 The adversary's registers are finite-dimensional}{digest-meat-core/ass:eve-dim}
\SetShortId{A-EVE-DIM}
\renewcommand*{\RestatedName}{\hyperref{contract.pdf}{}{claim-meat-core/ass:eve-dim}{Assumption~13.9}}%
\begin{digestassumption}[The adversary's registers are finite-dimensional]
(i) Every register of the adversary in the round model
(\xref{foundations/def:round-model}), her initial register $E_0$, her
outgoing memories $E'_j$ and her registers $E_j$ for $j=1,\dots,n$, is
finite-dimensional, of arbitrary finite dimension.
\end{digestassumption}

\part*{Engine requirements}\pdfbookmark[0]{Engine requirements}{digest-part-2}

\section*{14\quad Requirements: admission and the declared structure}\pdfbookmark[1]{14 Requirements: admission and the declared structure}{digest-group-meat-core/req:exact-admission}

% claim meat-core/req:exact-admission
\digestlabel{req:exact-admission}{14.1}\pdfbookmark[2]{Requirement 14.1 Exact admission of declared data}{digest-meat-core/req:exact-admission}
\renewcommand*{\RestatedName}{\hyperref{contract.pdf}{}{claim-meat-core/req:exact-admission}{Requirement~14.1}}%
\begin{digestrequirement}[Exact admission of declared data]
The engine takes no tolerance on declared data and decides no
condition on it more loosely than exactly
(\xref{foundations/req:consumer-decides-exactly}): it admits a declaration
only if the declaration is exactly admissible in the sense of
Definition~\ref{def:declaration}, and refuses it otherwise. No
perturbation of a declared family is licensed to carry the key length of
the family it approximates. Proposition~\ref{prop:optol-nobound} is the
reason: it shows, for the measurement on $\Qsys$, that no tolerance could
license a key-length claim. A compiled reference companion is not admitted
under this requirement but under Requirement~\ref{req:companion-admission},
which decides exactly what it shares with the declaration and decides none
of its family; the binary64 numbers it is evaluated at are no declared
family whose key length is claimed, since the bound evaluated there is that
of the reference, at the unbudgeted arithmetic of
Remark~\ref{C-rem:fo-arithmetic}.
\end{digestrequirement}

% claim meat-core/req:declaration-admission
\digestlabel{req:declaration-admission}{14.2}\pdfbookmark[2]{Requirement 14.2 Admission of the declaration}{digest-meat-core/req:declaration-admission}
\renewcommand*{\RestatedName}{\hyperref{contract.pdf}{}{claim-meat-core/req:declaration-admission}{Requirement~14.2}}%
\begin{digestrequirement}[Admission of the declaration]
The engine decides conditions (i)--(vi) and (viii) of
Definition~\ref{def:admissible-declaration}, in the general format as that
definition reads them there, exactly on the declared data, and
computes nothing from a declaration that fails one of them, which it
refuses. It does not decide condition~(vii) on a declaration whose
constraint list is non-empty. A compiled reference companion is admitted
instead under Requirement~\ref{req:companion-admission}, which decides
exactly the components it shares with the declaration it accompanies and
none of its measured family.
\end{digestrequirement}

% claim meat-core/req:companion-admission
\digestlabel{req:companion-admission}{14.3}\pdfbookmark[2]{Requirement 14.3 Admission of the reference companion beside its declaration}{digest-meat-core/req:companion-admission}
\renewcommand*{\RestatedName}{\hyperref{contract.pdf}{}{claim-meat-core/req:companion-admission}{Requirement~14.3}}%
\begin{digestrequirement}[Admission of the reference companion beside its declaration]
Where the protocol compiler emits a compiled reference companion
$\Decl^{\mathrm{ref}}$ beside a declaration $\Decl$
(\xref{protocol-compiler/req:reference-companion-compiled}), the engine
admits it only beside $\Decl$, after admitting $\Decl$
(Requirements~\ref{req:exact-admission} and~\ref{req:declaration-admission}),
and only as follows.
\begin{enumerate}[label=(\roman*),leftmargin=2.4em]
\item \emph{Decided equal.} It decides exactly, on the declared data, that
  every component Proposition~\ref{prop:reference-companion}\,(i) names as
  shared is equal to $\Decl$'s, entry for entry and label for label: the
  shapes, the sector maps and the joint sectors, the
  alphabets and their orders, Alice's slots, the party partition, the
  announcement, sifting and raw-key maps; that the companion's entries are
  those of $\Decl$ save the entries of the links the companion replaced, one
  per measured factor whose family it replaces, each of which carries the relaxed mark of
  \xref{reductions/req:reference-companion}\,(iii) and a Hermitian binary64
  operator with the blocks of its scope; and it reads the acceptance set of
  $\Decl$, the acceptance set being one per run. A companion that fails a
  decision of this item is not admitted, the reason naming the component,
  and the run proceeds without it. The statistics map, the
  single-round protocol map, the branch maps, the
  objective and its gradient of the companion are generated from the
  companion's own slots, its measured family among them
  (Definition~\ref{def:maps}), as for any declaration; no map generated from
  $\Decl$'s measured family is used for the companion. The announcement,
  sifting and raw-key maps decided equal above are declared data, not
  generated maps.
\item \emph{Imported, never decided.} It decides none of the conditions of
  Definition~\ref{def:admissible-declaration} on the companion's measured
  family. That the family it denotes is a POVM, block diagonal for its
  sector map, is \xref{reductions/cor:reference-companion}\,(i), which rests
  on the claims that establish the handed reference
  (\xref{foundations/def:exact-reference}) on the blocks where one is
  handed, on the exact decisions of the declarer's snap on the blocks where
  the companion carries the snapped block, and on (H-REF) for each instance
  whose link it replaces;
  and with (i) above it gives the companion conditions (i)--(vi) and (viii)
  by Proposition~\ref{prop:reference-companion}\,(ii). Condition (vii) is not
  decided on it. Nor does it decide that a block of the companion's measured
  family is a member of a factor enclosure the companion carries with
  it, which is Corollary~\ref{cor:reference-companion-factors}, on the claims
  that establish the factor enclosures of the handed reference.
\item \emph{Evaluated and measured.} It reads the companion's family at the
  binary64 numbers nearest to the centres of its enclosure, rounded once,
  entrywise to nearest, when the family is docked, and records, per block of
  the measured family, the completeness defect and the least eigenvalue of
  those numbers, measured and not decided, and the largest radius of the
  enclosure. A defect so measured refuses nothing and marks nothing
  inadmissible; what these numbers mean for the bound evaluated on them is
  Remark~\ref{C-rem:fo-arithmetic}.
\item \emph{Nothing else.} It computes from the companion only what
  Requirements~\ref{req:bracket-on-reference}
  and~\ref{req:certified-evaluation} ask. The one program solved over it is
  the polish of Requirement~\ref{req:bracket-on-reference}\,(i), a search
  continued from a point of the search on $\Decl$, whose points serve only
  as points at which those requirements form brackets; no estimate is
  reported from it, and no key length is reported from it save through
  those requirements.
\end{enumerate}
Requirements~\ref{req:exact-admission} and~\ref{req:declaration-admission}
bind every declaration but a companion admitted under this requirement.
\end{digestrequirement}

% claim meat-core/req:acceptance-set-admission
\digestlabel{req:acceptance-set-admission}{14.4}\pdfbookmark[2]{Requirement 14.4 Admission of the acceptance set}{digest-meat-core/req:acceptance-set-admission}
\renewcommand*{\RestatedName}{\hyperref{contract.pdf}{}{claim-meat-core/req:acceptance-set-admission}{Requirement~14.4}}%
\begin{digestrequirement}[Admission of the acceptance set]
The engine admits an acceptance set only if it decides, exactly and on the
declared data, that the set is an admissible acceptance set
(Definition~\ref{def:admissible-acceptance-set}): the set carries a geometry
tag every well-formed datum of which denotes such a set, and its data are
well formed for that tag. It refuses every other declared acceptance set
(\iface{acceptance_set_ill_formed} for data ill formed for the geometry they
claim).
\end{digestrequirement}

% claim meat-core/req:sector-report
\digestlabel{req:sector-report}{14.6}\pdfbookmark[2]{Requirement 14.6 The sector report}{digest-meat-core/req:sector-report}
\renewcommand*{\RestatedName}{\hyperref{contract.pdf}{}{claim-meat-core/req:sector-report}{Requirement~14.6}}%
\begin{digestrequirement}[The sector report]
The engine emits, for every sector, the four predicates of
Definition~\ref{def:classes} in \iface{sector_report}, and justifies there
every sector it deletes by the six conditions (W1)--(W6) checked on the
exhibited witness, by that witness with its evaluation order of (W6),
declared with the witness and recorded with it as
\iface{witness_evaluation_order}, and by the retention of the witness's
targets.
\end{digestrequirement}

% claim meat-core/req:state-blind-test
\digestlabel{req:state-blind-test}{14.7}\pdfbookmark[2]{Requirement 14.7 Bit-level test of state-blindness}{digest-meat-core/req:state-blind-test}
\renewcommand*{\RestatedName}{\hyperref{contract.pdf}{}{claim-meat-core/req:state-blind-test}{Requirement~14.7}}%
\begin{digestrequirement}[Bit-level test of state-blindness]
The engine decides whether a sector $s$ is state-blind
(Definition~\ref{def:classes}\,(c)) by a bit-level identity of declared
data. It reads $c_{v,s}$ off the declared block as its first diagonal
entry, $c_{v,s}\coloneqq(F^{(s)}_v)_{11}$, and requires
$F^{(s)}_v-c_{v,s}\idon{\Qsys_s}=0$ entrywise and exactly: every
off-diagonal entry of the block is the datum $0$, and every diagonal entry
is the same datum. Likewise, whenever $\CList\neq\varnothing$, it requires
$\hat O^{(s)}_k-o_{k,s}\idon{\Qsys_s}=0$ with
$o_{k,s}\coloneqq(\hat O^{(s)}_k)_{11}$. No subtraction, no norm and no
window enters: two binary64 data are equal or they are not, equality being
equality of the values they denote and not of their bit patterns
(\xref{foundations/def:binary64}). The engine decides the sign condition
$c_{v,s}\ge0$ by comparing that same datum with zero as a real number. A
block that is only \emph{nearly} scalar fails, and $s$ is then not
state-blind.
\end{digestrequirement}

% claim meat-core/req:removal-checklist
\digestlabel{req:removal-checklist}{14.8}\pdfbookmark[2]{Requirement 14.8 Setting the removable flag}{digest-meat-core/req:removal-checklist}
\renewcommand*{\RestatedName}{\hyperref{contract.pdf}{}{claim-meat-core/req:removal-checklist}{Requirement~14.8}}%
\begin{digestrequirement}[Setting the removable flag]
The engine sets the flag \iface{removable} for a sector $s^{*}$ of a
declaration with one measured factor when, and only when, a candidate
witness exhibited by the implementation is a witness for $s^{*}$ in the
sense of Definition~\ref{def:classes}\,(d): the list $\CList$ carries no
entry declared party-local and the candidate meets all six of
(W1)--(W6); it sets the flag on no other ground. The
witness may come from any search, and it is never read as declared input: a
declaration supplies no witness. Each condition is decided exactly, at no
window: (W1) by the bit-level identity of
Requirement~\ref{req:state-blind-test}, (W2) as a comparison of sector
indices, (W5) as the finite combinatorial predicate of key-inertness, (W3)
and (W4) on the exhibited data as those conditions state, and (W6) on the
residual formed in the exhibited evaluation order; a witness whose
evaluation rounds anywhere, or whose residual is not the datum $0$, sets no
flag. Every other outcome --- a nonzero residual, a rounding operation, an
absent witness, a witness that names $s^{*}$, a target that is not
key-inert, a block that is not exactly scalar --- retains the sector. The
flag is set for no sector of a declaration whose list carries an entry
declared party-local, nor of a declaration with more than one measured
factor. No tolerance licenses a removal:
\iface{sector_criterion} is decided at exactly zero and not at the
declared tolerance $\tnum$.
\end{digestrequirement}

% claim meat-core/req:structure-operations
\digestlabel{req:structure-operations}{14.9}\pdfbookmark[2]{Requirement 14.9 Operations on the declared sectors}{digest-meat-core/req:structure-operations}
\renewcommand*{\RestatedName}{\hyperref{contract.pdf}{}{claim-meat-core/req:structure-operations}{Requirement~14.9}}%
\begin{digestrequirement}[Operations on the declared sectors]
The engine treats the declared sector list as input it never edits; sectors
leave the variable list only by the deletion of~(ii). The operations
it may perform on the declared sectors are exactly: (i) evaluate the
objective over the key-active sectors only (Definition~\ref{def:keyactive};
an identity, by Proposition~\ref{prop:exact-cancellation}); (ii) delete
sectors carrying the flag \iface{removable} of
Definition~\ref{def:classes}\,(d), set as
Requirement~\ref{req:removal-checklist} states --- and only those, and among those only
a set meeting the hypothesis of Proposition~\ref{prop:removal-compose},
every deleted sector's witness naming targets outside the deleted set (an
identity, by Theorem~\ref{thm:removal} and that proposition); (iii) nothing
else. In particular no simplification step merges sectors, and the only
structure-changing simplification is the deletion of~(ii).
\end{digestrequirement}

\section*{15\quad Requirements: the Choi program and gradients}\pdfbookmark[1]{15 Requirements: the Choi program and gradients}{digest-group-meat-core/req:frame-constructed}

% claim meat-core/req:frame-constructed
\digestlabel{req:frame-constructed}{15.1}\pdfbookmark[2]{Requirement 15.1 Construction and exact check of the support frame}{digest-meat-core/req:frame-constructed}
\renewcommand*{\RestatedName}{\hyperref{contract.pdf}{}{claim-meat-core/req:frame-constructed}{Requirement~15.1}}%
\begin{digestrequirement}[Construction and exact check of the support frame]
In trusted-source mode the engine constructs the
declared isometry $\iota$ of every support frame it uses and
never takes $\iota$ as declared input:
by the full frame, which is always available, by the
diagonal frame, or by another construction for which it exhibits
\eqref{eq:iota} exactly. It decides both identities of \eqref{eq:iota} on
the isometry it constructed, entrywise as binary64 data at exactly zero
(\iface{choi_isometry}), and refuses a frame that fails either of them
with \iface{choi_isometry_inexact}.
\end{digestrequirement}

% claim meat-core/req:diagonal-frame-test
\digestlabel{req:diagonal-frame-test}{15.2}\pdfbookmark[2]{Requirement 15.2 Exact diagonality test for the diagonal frame}{digest-meat-core/req:diagonal-frame-test}
\renewcommand*{\RestatedName}{\hyperref{contract.pdf}{}{claim-meat-core/req:diagonal-frame-test}{Requirement~15.2}}%
\begin{digestrequirement}[Exact diagonality test for the diagonal frame]
The engine uses the diagonal frame of
Definition~\ref{def:support-frame}\,(b) only when every off-diagonal entry
of the declared $\tau_A$ is the datum $0$. It decides this by comparing
each off-diagonal entry with $0$ as a binary64 value, at no window, as
Requirement~\ref{req:state-blind-test} compares the entries of a sector
block, and the comparison is a clause of \iface{alice_marginal_exact}. A
$\tau_A$ that is only \emph{nearly} diagonal is not diagonal: the diagonal
frame is unavailable for it, and the full frame always is.
\end{digestrequirement}

% claim meat-core/req:frame-recorded
\digestlabel{req:frame-recorded}{15.3}\pdfbookmark[2]{Requirement 15.3 Record of the support frame}{digest-meat-core/req:frame-recorded}
\renewcommand*{\RestatedName}{\hyperref{contract.pdf}{}{claim-meat-core/req:frame-recorded}{Requirement~15.3}}%
\begin{digestrequirement}[Record of the support frame]
In trusted-source mode every run records, beside any Choi block it
reports, the support frame it used: the construction that produced
it, the dimension $d_{A_0}$ and the declared isometry $\iota$, per
Alice sector in the general format, as \iface{choi_input_basis} in
\iface{optimizer_certificate}. In untrusted-source mode no $A_0$ and no
$\iota$ exist (Definition~\ref{def:untrusted-choi}\,(a)), and the field is
absent.
\end{digestrequirement}

% claim meat-core/req:row-form-by-construction
\digestlabel{req:row-form-by-construction}{15.5}\pdfbookmark[2]{Requirement 15.5 Frame-free rows only where the construction exhibits tightness}{digest-meat-core/req:row-form-by-construction}
\renewcommand*{\RestatedName}{\hyperref{contract.pdf}{}{claim-meat-core/req:row-form-by-construction}{Requirement~15.5}}%
\begin{digestrequirement}[Frame-free rows only where the construction exhibits tightness]
In trusted-source mode the engine imposes a homogeneous row of
\eqref{eq:choi-feas} in its frame-free form \eqref{eq:frame-free-row}, and
in the general format the row of an Alice sector $t$ of a homogeneous or a
party-local entry in its unconjugated form, only where the construction
that produced the frame (of sector $t$) exhibits $\ker W_{\tau}=0$
($\ker W_t=0$) exactly, with no rank computed and no window declared.
Of the two constructions of Definition~\ref{def:support-frame} that is the
diagonal frame; on a frame produced by the full frame
construction the conjugated form is imposed, whatever $\tau_A$ is, and so
it is on every frame whose construction does not exhibit
$\ker W_{\tau}=0$. The engine licenses the frame-free form by no numerical
test on $\tau_A$.
\end{digestrequirement}

% claim meat-core/req:row-form-record
\digestlabel{req:row-form-record}{15.6}\pdfbookmark[2]{Requirement 15.6 Record of the form each constraint row was imposed in}{digest-meat-core/req:row-form-record}
\renewcommand*{\RestatedName}{\hyperref{contract.pdf}{}{claim-meat-core/req:row-form-record}{Requirement~15.6}}%
\begin{digestrequirement}[Record of the form each constraint row was imposed in]
In trusted-source mode a run records, for every declared constraint row,
which of the two forms of Proposition~\ref{prop:frame-free-rows} it
imposed, per Alice sector in the general format, as
\iface{constraint_row_form} in \iface{optimizer_certificate}, beside the
frame it records under Requirement~\ref{req:frame-recorded}.
\end{digestrequirement}

% claim meat-core/req:split-form-gradient
\digestlabel{req:split-form-gradient}{15.7}\pdfbookmark[2]{Requirement 15.7 Split-form evaluation of the trace-power gradient}{digest-meat-core/req:split-form-gradient}
\renewcommand*{\RestatedName}{\hyperref{contract.pdf}{}{claim-meat-core/req:split-form-gradient}{Requirement~15.7}}%
\begin{digestrequirement}[Split-form evaluation of the trace-power gradient]
The engine evaluates the gradient of the pinched trace power at a branch
operator $X$ in the split form of Theorem~\ref{thm:gradient}\,(i): the
divided-difference map $DX^{p}[\cdot]$ compressed to
$\operatorname{supp}X$, its kernel summed over the strictly positive
eigenvalues of $X$ only, and $\cZ(X^{p})$, $Y^{\alpha-1}$ and
$\cZ(Y^{\alpha-1})$ formed on the ambient branch space. It compresses the
whole computation to a subspace $S'$ only where $S'$ contains the maximal
branch support $W_{\hat c,s}$ of Definition~\ref{def:branch-support} and
$Z_jS'\subseteq S'$ for every $j$.
\end{digestrequirement}

% claim meat-core/req:gradient-at-maximal-support
\digestlabel{req:gradient-at-maximal-support}{15.8}\pdfbookmark[2]{Requirement 15.8 Gradient evaluation at maximal support}{digest-meat-core/req:gradient-at-maximal-support}
\renewcommand*{\RestatedName}{\hyperref{contract.pdf}{}{claim-meat-core/req:gradient-at-maximal-support}{Requirement~15.8}}%
\begin{digestrequirement}[Gradient evaluation at maximal support]
The engine evaluates a gradient of the pinched trace power only at an
iterate whose support is (kept) maximal, maximality being that of
\eqref{eq:Wcs} for the branch operators, for instance after an explicit
perturbation.
\end{digestrequirement}

% claim meat-core/req:limit-not-a-bound
\digestlabel{req:limit-not-a-bound}{15.9}\pdfbookmark[2]{Requirement 15.9 The order-one limit as a validation reference only}{digest-meat-core/req:limit-not-a-bound}
\renewcommand*{\RestatedName}{\hyperref{contract.pdf}{}{claim-meat-core/req:limit-not-a-bound}{Requirement~15.9}}%
\begin{digestrequirement}[The order-one limit as a validation reference only]
The engine never reports the limit $h_1$ of \eqref{eq:oracle}, however
computed, as a bound of any kind; in particular it never reports it as the
first-order bound $\hlb$ of Theorem~\ref{thm:firstorder}.
\end{digestrequirement}

% claim meat-core/req:dropped-set-record
\digestlabel{req:dropped-set-record}{15.10}\pdfbookmark[2]{Requirement 15.10 Recording the dropped set}{digest-meat-core/req:dropped-set-record}
\renewcommand*{\RestatedName}{\hyperref{contract.pdf}{}{claim-meat-core/req:dropped-set-record}{Requirement~15.10}}%
\begin{digestrequirement}[Recording the dropped set]
A run records its dropped set $\mathcal{D}$ (Definition~\ref{def:dropped})
with the run, and reports the number it computes from \eqref{eq:hlowD} as
$\underline{h}^{\mathcal{D}}_{\alpha}$ together with $\mathcal{D}$, never,
when $\mathcal{D}\neq\varnothing$, under the name $\hlow$.
\end{digestrequirement}

\section*{16\quad Requirements: values and key lengths}\pdfbookmark[1]{16 Requirements: values and key lengths}{digest-group-meat-core/req:infinite-value}

% claim meat-core/req:infinite-value
\digestlabel{req:infinite-value}{16.1}\pdfbookmark[2]{Requirement 16.1 Propagation of an infinite value}{digest-meat-core/req:infinite-value}
\renewcommand*{\RestatedName}{\hyperref{contract.pdf}{}{claim-meat-core/req:infinite-value}{Requirement~16.1}}%
\begin{digestrequirement}[Propagation of an infinite value]
The engine propagates $\hlow=+\infty$, the infeasibility of the declared
acceptance set, rather than clamping it to any finite range; its output
\iface{h_low} has the extended-real domain $[0,+\infty]$. If $\hlow=+\infty$
reaches Theorem~\ref{thm:keylength}, the engine emits no key-length claim and
no numeric key length, and reports the mismatch between the declaration and
the acceptance set (\iface{infeasible_acceptance_set}).
\end{digestrequirement}

% claim meat-core/req:unreachable-acceptance-set
\digestlabel{req:unreachable-acceptance-set}{16.2}\pdfbookmark[2]{Requirement 16.2 The check of the acceptance set against the feasible statistics}{digest-meat-core/req:unreachable-acceptance-set}
\renewcommand*{\RestatedName}{\hyperref{contract.pdf}{}{claim-meat-core/req:unreachable-acceptance-set}{Requirement~16.2}}%
\begin{digestrequirement}[The check of the acceptance set against the feasible statistics]
A run is \emph{checked} when its acceptance set is built from an honest law
or an honest envelope that the caller gives as numbers, such as one measured
on a device or read from a table, rather than one the engine computes from
the declared devices, or when its declaration carries a constraint entry
that the caller declares rather than one the compiler or a reduction
derives, unless the run sets the option \iface{statistics_check} to
\iface{off} (default \iface{on}). On every run, checked or not, the engine
records under \iface{statistics_consistency} whether the run is checked and,
where it is not, why (out of scope, or opted out), and the report says so
beside every key length; a run that is not checked names the premise
\xref{statistics/ass:reach} as unchecked, the caller's to vouch for. On
every checked run the engine does the following once, before it evaluates
the run at any R\'enyi order and before it reports a key length.
\begin{enumerate}[label=(\roman*),leftmargin=2.4em]
\item \emph{The order and the solve.} It takes $\alpha_0$ to be fixed before
  any order is evaluated: the run's fixed order, or the lowest order its
  order search evaluates, where the search fixes that order before it
  evaluates any and evaluates none below it, and otherwise the lower end of
  the domain the search declares, which every order the search evaluates is
  at least. The engine evaluates the program only at an order in
  $(1,\infty)$ (Requirement~\ref{C-req:objective-alpha-domain}), so the solve
  runs only where $\alpha_0>1$, and at any other $\alpha_0$ it is refused,
  as (ii) records; a key length is reported only at an order in $(1,2)$
  (Requirement~\ref{req:key-length-alpha-domain}), so a fixed order at
  which one is reported lies there, and that the orders a search evaluates
  and the lower end of its declared domain lie in $(1,2)$ is the bound the
  run tier places on the domain it accepts, which no claim of this contract
  states. On the
  declaration the run searches and with the run's
  acceptance set, it runs the first-order search of the
  dropped program for the dropped set
  $\mathcal{D}_{\mathrm{all}}$ of every pair
  (Proposition~\ref{prop:minimal-divergence}) at $\alpha_0$, with the
  method, the iteration cap and the gap target the run applies to one
  evaluation of its own, and with the share
  \iface{statistics_check_budget_share} of the wall budget of one
  evaluation of the run, where it has one, as its wall budget. A smaller
  share can leave a mismatch unproven: $b$ smaller, or not positive, so
  that (iv) may leave unrefused a run that a longer solve would have
  refused, and the
  cap of (iv) looser; no bound is made unsound by it.
\item \emph{The bound.} From that solve it forms the first-order number of
  the dropped program and, unless the run sets \iface{certified_evaluation}
  to \iface{off}, the certified evaluation of that number, on the
  declaration the run searches and at that solve's points, by parts
  (ii)--(v) of Requirement~\ref{req:certified-evaluation}, whatever
  declaration the run's key length is formed on. The number $b$ is the
  certified number where a point is certified and the first-order number
  otherwise, with its tier. The solve is not refused on the sign of
  its estimate $O_0$ of (v), which is no estimate of the run and
  which Requirement~\ref{req:negative-estimate-refused} does not read: $b$
  is formed from the solve's points whatever that sign, and a negative
  $O_0$ is recorded as a diagnostic and refuses nothing. Where no bound is
  formed (the first-order number withdrawn, the solve not completed, or the
  solve refused at an order outside $(1,\infty)$), the run proceeds with
  the reason recorded.
\item \emph{The proven part.} Where $b>0$, for every order $\alpha'$ the run
  evaluates, the proven floor
  $f(\alpha')\coloneqq\frac{\alpha'}{\alpha'-1}\cdot
  \frac{\alpha_0-1}{\alpha_0}\,b$ satisfies
  $f(\alpha')\le\frac{\alpha'}{\alpha'-1}\sfD_{\min}\le\underline{h}_{\alpha'}$
  (Proposition~\ref{prop:minimal-divergence}\,(ii), (iv)), with $\sfD_{\min}$
  and $\underline{h}_{\alpha'}$ those of the declaration the run searches: a
  part of that declaration's value at $\alpha'$ that the mismatch alone
  accounts for. Where $b\le0$, $f(\alpha')\coloneqq0$. Where no bound was
  formed, $f(\alpha')$ is not defined.
\item \emph{The cap and the refusal.} Let $\tau\in(0,1]$ be the run's option
  \iface{min_model_pass_probability}, an exact rational number, default
  $\frac12$. Where $b$ was formed and $b>0$, the engine forms the exponent
  $x\coloneqq n\,\frac{\alpha_0-1}{\alpha_0}\,b$ exactly from the run's
  round count $n$, $\alpha_0$ and $b$, with the tier of $b$. By
  Proposition~\ref{prop:model-pass-bound}\,(ii), devices the declaration the
  run searches admits, run so that each round's acceptance value, given the
  earlier ones, follows the statistics $\Phimap[\omega]$ of some $\omega$ in
  that declaration's $\Feas$, pass the acceptance test with probability at
  most $2^{-x}$, the \emph{cap}. Where $2^{-x}\le\tau$, that is
  $x\ge\log_2\frac1\tau$, the run is refused as a whole with
  \iface{acceptance_set_unreachable}: it is evaluated at no order and
  reports no key length. The engine decides $x\ge\log_2\frac1\tau$ exactly
  where $\tau=2^{-k}$ for an integer $k\ge0$, by comparing $x$ with $k$;
  for any other $\tau$ it forms a rigorous enclosure
  $[\lambda_-,\lambda_+]\ni\log_2\frac1\tau$ and refuses where
  $x\ge\lambda_+$, refuses nothing where $x<\lambda_-$, and refuses where
  $\lambda_-\le x<\lambda_+$, the comparison being undecided there. The
  decision carries the tier of $b$: where $b$ is a first-order number not
  certified, the refusal or its absence is of the first-order tier. Where
  $b\le0$ or no bound was formed there is no cap, and this item refuses
  nothing. With $\tau=1$ every run with $b>0$ is refused. Where a reference
  companion carries the key length, the cap is still that of the
  declaration the run searches, which no claim relates to the companion's
  feasible set; the refusal only withholds a key.
\item \emph{The record.} On a checked run the engine records under
  \iface{statistics_consistency}: the declaration the solve ran on;
  $\alpha_0$ and which rule of (i) fixed it; $b$ with its tier and the lower bound
  $\frac{\alpha_0-1}{\alpha_0}\,b$ on $\sfD_{\min}$, or the reason none was
  formed; $\tau$; where $b>0$, the exponent $x$ of (iv) with the tier of
  $b$, rounded downward where the record cannot carry it exactly, so that
  $2^{-x}$ at the recorded $x$ stays an upper bound on the pass
  probability, and whether the comparison of (iv) was decided exactly, by
  the enclosure, or left undecided and refused; where $b\le0$ or no bound
  was formed, that there is no cap; whether (iv) refused the run, as (iv)
  decided it on the exact $x$ and never as a comparison of the recorded or
  displayed values with $\tau$; where the
  solve ran, the share of the wall budget and the wall
  budget it was given, the wall time it took and the cause it stopped on, as the solver
  reports it (such as the gap target, the iteration cap or the budget);
  where the solve returned a point, the
  estimate $O_0\coloneqq O^{\mathcal{D}_{\mathrm{all}}}_{\alpha_0}$ at that
  point, whether it is negative, and whether its membership in the relaxed set was established;
  where $b$ was formed, the gap $O_0-b$; and, where a bound was formed,
  $n\,f(\alpha')$ in bits at each order $\alpha'$ the run evaluates, a
  diagnostic of the part of the value the mismatch accounts for, which
  satisfies $\frac{\alpha'-1}{\alpha'}\,n\,f(\alpha')
  =n\,\frac{\alpha_0-1}{\alpha_0}\,b$
  (Proposition~\ref{prop:model-pass-bound}\,(ii)). The report shows, beside
  every key length it reports and, on a run refused under (iv), in place of
  one, the cap as a probability: $2^{-x}$ at the recorded $x$ together with
  a decimal approximation of it rounded upward, with the tier of $b$, as
  the probability with which devices the declaration the run searches
  admits, run as (iv) says, pass the acceptance test, both labelled as
  rounded upward; and whether (iv) refused the run, with $\tau$, as the
  record carries that decision from the exact comparison, so that a
  displayed cap a little above $\tau$ beside a refusal reads as the
  upward rounding it is. Where there is no cap it says so and why
  ($b\le0$, or no bound formed and the reason). It states the cap as a
  probability and not as a number of bits, and it does not show
  $n\,f(\alpha')$. No key length and no other number of the run reads
  these entries.
\end{enumerate}
\end{digestrequirement}

% claim meat-core/req:reported-key-length
\digestlabel{req:reported-key-length}{16.4}\pdfbookmark[2]{Requirement 16.4 Reported key length}{digest-meat-core/req:reported-key-length}
\renewcommand*{\RestatedName}{\hyperref{contract.pdf}{}{claim-meat-core/req:reported-key-length}{Requirement~16.4}}%
\begin{digestrequirement}[Reported key length]
The engine reports as \iface{key_length} the maximal key length
built from the lower bound $h\le\hlow$ it uses, that is $\ell_{\max}$ of
Definition~\ref{def:max-key-length} with $\max\{0,h\}$ in place of $\hlow$
(Corollary~\ref{cor:keylength-bound}\,(a)), with its floor and its clamp at
$0$: that is its conservative integer-output rule.
\end{digestrequirement}

% claim meat-core/req:hash-penalty
\digestlabel{req:hash-penalty}{16.5}\pdfbookmark[2]{Requirement 16.5 Key-length penalty for a merely two-universal family}{digest-meat-core/req:hash-penalty}
\renewcommand*{\RestatedName}{\hyperref{contract.pdf}{}{claim-meat-core/req:hash-penalty}{Requirement~16.5}}%
\begin{digestrequirement}[Key-length penalty for a merely two-universal family]
For a privacy-amplification family that is only two-universal, collision
probability at most $2^{-\ell}$, \eqref{eq:keylength} must be tightened by
one further bit, $\ell\to\ell-1$.
\end{digestrequirement}

% claim meat-core/req:key-length-alpha-domain
\digestlabel{req:key-length-alpha-domain}{16.6}\pdfbookmark[2]{Requirement 16.6 Rényi order of a reported key length}{digest-meat-core/req:key-length-alpha-domain}
\renewcommand*{\RestatedName}{\hyperref{contract.pdf}{}{claim-meat-core/req:key-length-alpha-domain}{Requirement~16.6}}%
\begin{digestrequirement}[Rényi order of a reported key length]
The engine reports a key length only at a R\'enyi order $\alpha\in(1,2)$;
at any other order it reports no key length
(\iface{alpha_in_the_profile_domain}).
\end{digestrequirement}

% claim meat-core/req:compiled-declaration-only
\digestlabel{req:compiled-declaration-only}{16.7}\pdfbookmark[2]{Requirement 16.7 Key lengths only on declarations the compiler emitted}{digest-meat-core/req:compiled-declaration-only}
\renewcommand*{\RestatedName}{\hyperref{contract.pdf}{}{claim-meat-core/req:compiled-declaration-only}{Requirement~16.7}}%
\begin{digestrequirement}[Key lengths only on declarations the compiler emitted]
The engine reports a key length, whether as a key-length claim of
Theorem~\ref{thm:keylength}, Corollary~\ref{cor:keylength-bound} or
Corollary~\ref{cor:fo-keylength} or as a numeric key length, only on a
declaration that the protocol compiler emitted; on any other declaration it
reports none.
\end{digestrequirement}

% claim meat-core/req:security-statement-not-vacuous
\digestlabel{req:security-statement-not-vacuous}{16.8}\pdfbookmark[2]{Requirement 16.8 Refusal of budgets with an empty security statement}{digest-meat-core/req:security-statement-not-vacuous}
\renewcommand*{\RestatedName}{\hyperref{contract.pdf}{}{claim-meat-core/req:security-statement-not-vacuous}{Requirement~16.8}}%
\begin{digestrequirement}[Refusal of budgets with an empty security statement]
The engine refuses a run whose budgets satisfy
$\epa+\eev+2\varepsilon_{\mathrm{auth}}\ge1$, deciding the sum exactly on the
three declared numbers before any program is built, and reports
\iface{security_statement_vacuous}: the
$(\epa+\eev+2\varepsilon_{\mathrm{auth}})$-security statement of
Corollary~\ref{cor:composed-security}\,(iii), at
$\varepsilon_{\mathrm{auth}}=0$ equal to $\epa+\eev$, the parameter of (ii)
(and numerically of (i)), is then empty, and Corollary~\ref{cor:keylength-bound}\,(b) does
not claim the corner.
\end{digestrequirement}

% claim meat-core/req:secrecy-convention
\digestlabel{req:secrecy-convention}{16.9}\pdfbookmark[2]{Requirement 16.9 Declared choice of the source's secrecy convention}{digest-meat-core/req:secrecy-convention}
\renewcommand*{\RestatedName}{\hyperref{contract.pdf}{}{claim-meat-core/req:secrecy-convention}{Requirement~16.9}}%
\begin{digestrequirement}[Declared choice of the source's secrecy convention]
An instance that adds the one bit of \eqref{eq:keylength} back adopts
the secrecy definition of the $\wedge\OmAT$ state, the convention of the
derivation Remark~\ref{C-rem:plus2} compares, and declares explicitly that it
does so. A key length so obtained is not one of
Theorem~\ref{thm:keylength} and carries no claim of this module.
\end{digestrequirement}

% claim meat-core/req:constrained-key-gate
\digestlabel{req:constrained-key-gate}{16.10}\pdfbookmark[2]{Requirement 16.10 Two-part gate on a key length with declared constraints}{digest-meat-core/req:constrained-key-gate}
\renewcommand*{\RestatedName}{\hyperref{contract.pdf}{}{claim-meat-core/req:constrained-key-gate}{Requirement~16.10}}%
\begin{digestrequirement}[Two-part gate on a key length with declared constraints]
A key-length output may be combined with $\CList\neq\varnothing$
\emph{only if} (i) the list lies in the licensed class of the declared
source mode (Definition~\ref{def:homog}), a condition that is mechanically
checkable and degenerate in the cases of Remark~\ref{C-rem:cgate}, and
(ii) the declaration carries the attack-set licence
\iface{attack_set_license} (Assumption~A-ATK,
\xref{protocol-compiler/ass:atk}), an attestation no program can verify.
The engine refuses any request that couples $\CList\neq\varnothing$ with a
key-length output unless \emph{both} parts are present; without them the
numerical machinery may be used on $\CList$ for exploration, but its output
does not reach \eqref{eq:keylength}.
\end{digestrequirement}

% claim meat-core/req:degenerate-gate-report
\digestlabel{req:degenerate-gate-report}{16.11}\pdfbookmark[2]{Requirement 16.11 A degenerate licensed-class check reported as no evidence}{digest-meat-core/req:degenerate-gate-report}
\renewcommand*{\RestatedName}{\hyperref{contract.pdf}{}{claim-meat-core/req:degenerate-gate-report}{Requirement~16.11}}%
\begin{digestrequirement}[A degenerate licensed-class check reported as no evidence]
Where part~(i) of the two-part gate of
Requirement~\ref{req:constrained-key-gate} is degenerate, at $d_A=1$ and in
untrusted-source mode (Remark~\ref{C-rem:cgate}), the engine never
reports a passing part~(i) as evidence of anything, and it never reports a
passing part~(i) as evidence of anything about a party-local entry either.
\end{digestrequirement}

% claim meat-core/req:marginal-licence-gate
\digestlabel{req:marginal-licence-gate}{16.12}\pdfbookmark[2]{Requirement 16.12 Source licence for a declared marginal}{digest-meat-core/req:marginal-licence-gate}
\renewcommand*{\RestatedName}{\hyperref{contract.pdf}{}{claim-meat-core/req:marginal-licence-gate}{Requirement~16.12}}%
\begin{digestrequirement}[Source licence for a declared marginal]
The refusal discipline of the two-part gate of
Requirement~\ref{req:constrained-key-gate} governs the marginal as well: the engine refuses a declared $\tau_A$
without \iface{source_replacement_license} (Assumption~A-SRC,
\xref{transmitter/ass:src}) before any key-length output, and at $d_A=1$,
where the licence is vacuous, it does not demand it
(Proposition~\ref{prop:trivial-alice-system}\,(iv)).
\end{digestrequirement}

\part*{Frank-Wolfe first-order bound}\pdfbookmark[0]{Frank-Wolfe first-order bound}{digest-part-3}

\section*{17\quad The relaxed feasible set}\pdfbookmark[1]{17 The relaxed feasible set}{digest-group-meat-core/frank-wolfe/def:relaxed-feasible-set}

% claim meat-core/frank-wolfe/def:relaxed-feasible-set
\digestlabel{def:relaxed-feasible-set}{17.1}\pdfbookmark[2]{Definition 17.1 The relaxed feasible set}{digest-meat-core/frank-wolfe/def:relaxed-feasible-set}
\renewcommand*{\RestatedName}{\hyperref{contract.pdf}{}{claim-meat-core/frank-wolfe/def:relaxed-feasible-set}{Definition~17.1}}%
\begin{digestclaim}[The relaxed feasible set]
For $\varepsilon\ge0$ the
\emph{relaxed feasible set}
$\Feas(\varepsilon)$ is the set of sector-block-diagonal
$\omega\in\Pos{A\Qsys}$ satisfying, in trusted-source mode,
\begin{equation*}\tag{\getrefnumber{C-eq:feas-eps}}\digesteqlabel{eq:feas-eps}
  \big\|\omega_A-\tau_A\big\|_{\max}\le\varepsilon,
  \quad
  \big|\tr\omega-1\big|\le d_A\varepsilon,
  \quad
  \begin{cases}
  \lmin\big(\Theta_k-A_k(\omega)\big)\ge-\varepsilon, & \bowtie_k\ =\ \preceq,\\
  \lmin\big(A_k(\omega)-\Theta_k\big)\ge-\varepsilon, & \bowtie_k\ =\ \succeq,\\
  \big\|A_k(\omega)-\Theta_k\big\|_{\max}\le\varepsilon, & \bowtie_k\ =\ =,
  \end{cases}
\end{equation*}
where $A_k(\omega)\coloneqq\tr_{\Qsys}[(\idon{A}\otimes\hat O_k)\,\omega]$
is the left-hand side of the $k$-th declared relation of
\eqref{eq:feas-modes}, and, in untrusted-source mode, the normalisation
$\big|\sum_s\tr\omega_s-1\big|\le\varepsilon$ together with the scalar
forms
$a_k(\omega)\le c_k+\varepsilon$, $a_k(\omega)\ge c_k-\varepsilon$,
$|a_k(\omega)-c_k|\le\varepsilon$, with
$a_k(\omega)\coloneqq\tr[(\idon{A}\otimes\hat O_k)\,\omega]$ and no
marginal condition. In the Choi form $\Feas^{J}(\varepsilon)$ is the set of
$(J_s)_s$ with $J_s\in\Pos{A_0\Qsys_s}$,
$\big|\tr[\chi_{\Qsys}(J)]-1\big|\le d_{A_0}\varepsilon$,
$\big\|\sum_s\tr_{\Qsys_s}J_s-\idon{A_0}\big\|_{\max}\le\varepsilon$ and
the same relaxation of the $\bowtie_k$ rows.

The normalisation clause of each form is implied by the clause written
before it, by the marginal in trusted-source mode and by the coupling
residual in the Choi form, and is written out only so that the set can be
read without re-deriving it; Proposition~\ref{prop:relaxed-set-properties}\,(c$'$)
proves both implications and derives the two dimension factors. In
untrusted-source mode there is no clause before it: there the
normalisation \emph{is} the relaxed coupling condition
(Definition~\ref{def:untrusted-choi}\,(b)), which is why it carries no
dimension factor.

Which predicate is which is part of the definition and is never
converted: an inequality between operators is relaxed by a \emph{spectral
floor} $\lmin(\cdot)\ge-\varepsilon$, an operator equality by an
\emph{entrywise} ball $\|\cdot\|_{\max}\le\varepsilon$, a scalar equality
by an interval, and these are different quantities against the same number
(\xref{foundations/lem:floor-stricter}: the entrywise reading and the
spectral floor). The factors $d_A$ and $d_{A_0}$ are not such a
conversion: each is the dimension factor of an explicit trace bound
derived in Proposition~\ref{prop:relaxed-set-properties}\,(c$'$), written
into the clause, and neither reads an entrywise defect as a spectral one.

At $\varepsilon=0$ both sets reduce to \eqref{eq:feas} and
\eqref{eq:choi-feas} respectively: each normalisation clause is then the
exact $\tr\omega=1$, respectively $\tr[\chi_{\Qsys}(J)]=1$, which the
coupling condition already forces, so re-imposing it changes nothing. We
write $\Feas$ for $\Feas(0)$ and $\Feas^{J}$ for $\Feas^{J}(0)$
throughout.

\emph{In the general format} of Definition~\ref{def:declaration} both sets
are read on lists over the joint sectors, with every
line of \eqref{eq:feas-general}, respectively of the general-format Choi
set of Theorem~\ref{thm:choi}, relaxed by the predicate its relation calls
for, as in \eqref{eq:feas-eps}, the party-local rows included. The
marginal condition and the coupling condition are then read per Alice
sector:
$\|\sum_s\tr_{\Qsys_s}\omega^{(\mathsf t,s)}-\tau^{(\mathsf t)}_A\|_{\max}
\le\varepsilon$, respectively
$\|R_{\mathsf t}\|_{\max}\le\varepsilon$ with
$R_{\mathsf t}\coloneqq\sum_s\tr_{\Qsys_s}J_{(\mathsf t,s)}
-\idon{A_{0,\mathsf t}}$, for every $\mathsf t\in\mathbb{T}$. The residual
of the single marginal clause of \eqref{eq:feas-eps}, respectively of the
single coupling clause of the Choi form, is then the block-diagonal operator
$\omega_A-\tau_A=\bigoplus_{\mathsf t}(\sum_s\tr_{\Qsys_s}
\omega^{(\mathsf t,s)}-\tau^{(\mathsf t)}_A)$, respectively
$\bigoplus_{\mathsf t}R_{\mathsf t}$ with
$J_s\coloneqq\bigoplus_{\mathsf t}J_{(\mathsf t,s)}$, whose entrywise
maximum is the largest over $\mathsf t$ of the blocks' ones; so the clauses
per Alice sector together are that single clause. The normalisation clauses
are those of the first paragraph of this definition. At $\varepsilon=0$ the
two sets are $\Feas$ of \eqref{eq:feas-general} and the general-format
$\Feas^{J}$ of Theorem~\ref{thm:choi}.
\end{digestclaim}

% claim meat-core/frank-wolfe/prop:relaxed-set-properties
\digestlabel{prop:relaxed-set-properties}{17.2}\pdfbookmark[2]{Proposition 17.2 Properties of the relaxed feasible set}{digest-meat-core/frank-wolfe/prop:relaxed-set-properties}
\renewcommand*{\RestatedName}{\hyperref{contract.pdf}{}{claim-meat-core/frank-wolfe/prop:relaxed-set-properties}{Proposition~17.2}}%
\begin{digestclaim}[Properties of the relaxed feasible set]
Let $\Decl$ satisfy Definition~\ref{def:admissible-declaration}, let $\varepsilon\ge0$, and
let $\Feas(\varepsilon)$ and $\Feas^{J}(\varepsilon)$ be as in
Definition~\ref{def:relaxed-feasible-set}.
\begin{enumerate}[label=(\alph*),leftmargin=2.4em]
\item \emph{(Monotone.)} For $0\le\varepsilon\le\varepsilon'$,
  $\Feas(\varepsilon)\subseteq\Feas(\varepsilon')$ and
  $\Feas^{J}(\varepsilon)\subseteq\Feas^{J}(\varepsilon')$; in particular
  $\Feas\subseteq\Feas(\varepsilon)$ and $\Feas^{J}\subseteq\Feas^{J}(\varepsilon)$
  for every $\varepsilon\ge0$.
\item \emph{(Convex and compact.)} $\Feas(\varepsilon)$ and
  $\Feas^{J}(\varepsilon)$ are convex and compact. Every
  $\omega\in\Feas(\varepsilon)$ has $\tr\omega\le1+d_A\varepsilon$ in
  trusted-source mode and $\tr\omega\le1+\varepsilon$ in untrusted-source
  mode, and every $J\in\Feas^{J}(\varepsilon)$ has
  $\sum_s\tr J_s\le d_{A_0}(1+\varepsilon)$, a bound that reads the coupling
  clause and never the normalisation clause.
\item \emph{(What positivity carries.)} $\Feas(\varepsilon)\subseteq\Pos{A\Qsys}$
  and $\chi_{\Qsys}\big(\Feas^{J}(\varepsilon)\big)\subseteq\Pos{A\Qsys}$.
  \emph{First}, the bounds $0\le K_{\hat c}(\omega)\le P_{\hat c}(\omega)$
  hold on the whole cone. \emph{Second}, $K_{\hat c}$ is concave on the
  sector-block-diagonal part of $\Pos{A\Qsys}$. \emph{Third},
  Theorem~\ref{thm:convexity} holds verbatim with $\Feas(\varepsilon)$ in
  place of $\Feas$, and $O^{J}_{\alpha}$ is convex and lower semicontinuous
  on $\acSet\times\Feas^{J}(\varepsilon)$. \emph{Fourth}, none of this
  survives a relaxation of positivity: the objective factors through
  $X\mapsto X^{p}$, which is undefined off the cone. One fact of the exact
  set does not survive: for $\varepsilon>0$ and $\omega\in\Feas(\varepsilon)$,
  $\sum_{\hat c}P_{\hat c}(\omega)=\tr\omega$ for an exactly admissible
  $\{F_v\}$, which the relaxed
  normalisation does not pin to $1$, so $\Phimap[\omega]$ need not be a
  probability vector and the range fact $\hlow\ge0$ of
  Proposition~\ref{prop:value-range}\,(i) is not available on
  $\Feas(\varepsilon)$ or on $\Feas^{J}(\varepsilon)$; it holds on the exact
  sets $\Feas$ and $\Feas^{J}$.
\item[(c$'$)] \emph{(Each normalisation clause is implied by the clause
  before it.)} Let $J=(J_s)_s$ be per-block positive semidefinite with
  $\|R\|_{\max}\le\varepsilon$, where
  $R\coloneqq\sum_s\tr_{\Qsys_s}J_s-\idon{A_0}$, and put
  $G\coloneqq\iota^{\dagger}\tau_A\iota=W_{\tau}^{\dagger}W_{\tau}$. Then
  $G$ is a state,
  \begin{equation*}\tag{\getrefnumber{C-eq:choi-trace-residual}}\digesteqlabel{eq:choi-trace-residual}
    \tr[\chi_{\Qsys}(J)]
    =\tr\Big[W_{\tau}^{\dagger}W_{\tau}\,
      \textstyle\sum_s\tr_{\Qsys_s}J_s\Big]
    =\tr\big[G(\idon{A_0}+R)\big]
    =1+\tr[GR],
  \end{equation*}
  and $\big|\tr[\chi_{\Qsys}(J)]-1\big|\le d_{A_0}\varepsilon$; in
  particular $\tr[\chi_{\Qsys}(J)]=1$ holds at a point of nonzero residual
  $R$ exactly where $\tr[GR]=0$. In the state form in trusted-source mode,
  $\|\omega_A-\tau_A\|_{\max}\le\varepsilon$ implies
  $|\tr\omega-1|\le d_A\varepsilon$. In untrusted-source mode the
  normalisation clause is the relaxed coupling condition itself.
\item[(c$''$)] \emph{(What relaxing the normalisation costs.)} At a point
  of $\Feas^{J}(\varepsilon)$ put
  $\varsigma\coloneqq\tr[\chi_{\Qsys}(J)]-1$, so that
  $|\varsigma|\le d_{A_0}\varepsilon$ and
  $\sum_{\hat c}P^{J}_{\hat c}(J)=1+\varsigma$, and, where $1+\varsigma>0$,
  $\bar{\mathbf P}\coloneqq\mathbf P/(1+\varsigma)$, a probability vector.
  Then for every probability vector $\mathbf q$
  \begin{equation*}\tag{\getrefnumber{C-eq:penaltyshift-fw}}\digesteqlabel{eq:penaltyshift-fw}
    \sum_{\hat c}q(\hat c)\log_2\frac{q(\hat c)}{P^{J}_{\hat c}}
    =\sum_{\hat c}q(\hat c)\log_2\frac{q(\hat c)}{(1+\varsigma)\bar P_{\hat c}}
    =\sfD(\mathbf q\|\bar{\mathbf P})-\log_2(1+\varsigma),
  \end{equation*}
  and, provided $d_{A_0}\varepsilon<1$,
  \begin{equation*}\tag{\getrefnumber{C-eq:relaxed-objective-floor}}\digesteqlabel{eq:relaxed-objective-floor}
    O^{J}_{\alpha}(\mathbf q,J)\;\ge\;
    -\frac{\alpha}{\alpha-1}\,\log_2\!\big(1+d_{A_0}\varepsilon\big)
    \qquad\text{on }\acSet\times\Feas^{J}(\varepsilon),
  \end{equation*}
  whose right-hand side is $0$ at $\varepsilon=0$ and negative otherwise.
  The state form carries the same statement with $d_A\varepsilon$ in
  trusted-source mode and $\varepsilon$ in untrusted-source mode in place of
  $d_{A_0}\varepsilon$. At $d_{A_0}\varepsilon\ge1$ the identity
  \eqref{eq:penaltyshift-fw} still holds wherever $1+\varsigma>0$ but
  carries no bound.
\item[(d)] \emph{(Reading the set at a tolerance is conservative.)} With
  $\hlow[\Feas(\varepsilon)]\coloneqq\inf_{\mathbf q\in\acSet}
  \inf_{\omega\in\Feas(\varepsilon)}O_{\alpha}(\mathbf q,\omega)$,
  \begin{equation*}\tag{\getrefnumber{C-eq:feas-tol-mono}}\digesteqlabel{eq:feas-tol-mono}
    \varepsilon\le\varepsilon'\ \Longrightarrow\
    \hlow\big[\Feas(\varepsilon')\big]\;\le\;\hlow\big[\Feas(\varepsilon)\big]
    \;\le\;\hlow\big[\Feas(0)\big]=\hlow\;\le\;\hqkd .
  \end{equation*}
\item[(d$'$)] \emph{(The Choi-form quantity.)} With
  \begin{equation*}\tag{\getrefnumber{C-eq:hlow-feasJ}}\digesteqlabel{eq:hlow-feasJ}
    \hlow\big[\Feas^{J}(\varepsilon)\big]\coloneqq
    \inf_{\mathbf q\in\acSet}\ \inf_{J\in\Feas^{J}(\varepsilon)}
    O^{J}_{\alpha}(\mathbf q,J),
  \end{equation*}
  the two relaxed quantities agree at $\varepsilon=0$, and
  \begin{equation*}\tag{\getrefnumber{C-eq:feas-tol-mono-J}}\digesteqlabel{eq:feas-tol-mono-J}
    \varepsilon\le\varepsilon'\ \Longrightarrow\
    \hlow\big[\Feas^{J}(\varepsilon')\big]\;\le\;
    \hlow\big[\Feas^{J}(\varepsilon)\big]\;\le\;
    \hlow\big[\Feas^{J}(0)\big]\;=\;\hlow\;\le\;\hqkd .
  \end{equation*}
  At $\varepsilon>0$ no relation between $\hlow[\Feas^{J}(\varepsilon)]$
  and $\hlow[\Feas(\varepsilon)]$ is asserted.
\item[(e)] \emph{(The general format.)} For a declaration in the general
  format, with both sets read as Definition~\ref{def:relaxed-feasible-set}
  reads them there and $\chi_{\Qsys}$ the assignment
  $(J_{(\mathsf t,s)})\mapsto(\omega^{(\mathsf t,s)})$ of
  Theorem~\ref{thm:choi}, (a) to (d$'$) hold as stated. In (c$'$) the
  residual is $R=\bigoplus_{\mathsf t}R_{\mathsf t}$, $G=\bigoplus_{\mathsf
  t}G_{\mathsf t}$ with $G_{\mathsf t}\coloneqq W_{\mathsf t}^{\dagger}
  W_{\mathsf t}$, and
  $\tr[\chi_{\Qsys}(J)]=1+\sum_{\mathsf t}\tr[G_{\mathsf t}R_{\mathsf t}]$.
\end{enumerate}
\end{digestclaim}

\section*{18\quad Admissible directions}\pdfbookmark[1]{18 Admissible directions}{digest-group-meat-core/frank-wolfe/def:first-order-setting}

% claim meat-core/frank-wolfe/def:first-order-setting
\digestlabel{def:first-order-setting}{18.1}\pdfbookmark[2]{Definition 18.1 Ambient program, pairings, acceptance support and linear minimisation}{digest-meat-core/frank-wolfe/def:first-order-setting}
\renewcommand*{\RestatedName}{\hyperref{contract.pdf}{}{claim-meat-core/frank-wolfe/def:first-order-setting}{Definition~18.1}}%
\begin{digestclaim}[Ambient program, pairings, acceptance support and linear minimisation]
From this definition on, the ambient program of this submodule is the
Choi program \eqref{eq:choi-program} of Theorem~\ref{thm:choi}: the
variable is $x=(\mathbf q,J)$ with $\mathbf q\in\acSet$ and
$J=(J_s)_{s\in\SecSet}\in\Feas^{J}$ of \eqref{eq:choi-feas}, and the
objective is the function $O^{J}_{\alpha}$ displayed in
\eqref{eq:choi-program}, so that
$\hlow=\inf\{O^{J}_{\alpha}(x):x\in\acSet\times\Feas^{J}\}$ by that
theorem. In untrusted-source mode the state form is primary: there
$J_s$, $\Feas^{J}$, $\chi_{\Qsys_s}$ and $O^{J}_{\alpha}$ denote $\omega_s$,
$\Feas$, the identity and $O_{\alpha}$
(Definition~\ref{def:untrusted-choi}), and every later statement of this
submodule is read with these substitutions, as
Corollary~\ref{cor:full-gradient} reads the gradients.

The
\emph{acceptance support}
of the acceptance set $\acSet$ (Definition~\ref{def:admissible-acceptance-set}) is
\begin{equation*}\tag{\getrefnumber{C-eq:chacc-def}}\digesteqlabel{eq:chacc-def}
  \Chacc\coloneqq\big\{\hat c\in\Chatset:\ q(\hat c)>0
  \ \text{for some }\mathbf q\in\acSet\big\}.
\end{equation*}
Pairings are the real inner products
$\langle\mathbf a,\mathbf b\rangle=
\sum_{\hat c\in\Chacc}a(\hat c)b(\hat c)$ on $\mathbb{R}^{\Chacc}$, the
acceptance support being where the acceptance-side gradient lives
(Lemma~\ref{lem:fo-khyp}), and
$\langle A,B\rangle=\tr[A B]$ on Hermitian operators, the latter real and,
on a tuple $(A_s)_s$, summed over $s$.

For linear forms $\lmoq_t\in\mathbb{R}^{\Chacc}$ and
$\lmoJ_t=(\lmoJ_{t,s})_{s\in\SecSet}$, a tuple of Hermitian operators of the
block shape of $\Feas^{J}$, a pair $(v_t,S_t)$ of a vector in
$\mathbb{R}^{\Chacc}$ and such a tuple satisfies the
\emph{acceptance-side inequality}
of \eqref{eq:lmo} when the first inequality below holds, and the
\emph{feasible-side inequality}
when the second holds:
\begin{equation*}\tag{\getrefnumber{C-eq:lmo}}\digesteqlabel{eq:lmo}
  \langle\lmoq_t,v_t\rangle\le
  \min_{\mathbf q'\in\acSet}\langle\lmoq_t,\mathbf q'\rangle
  \qquad\text{and}\qquad
  \langle\lmoJ_t,S_t\rangle\le
  \min_{J'\in\Feas^{J}}\langle\lmoJ_t,J'\rangle ,
\end{equation*}
the second minimum being taken over the exact $\Feas^{J}$ of
\eqref{eq:choi-feas}.
\end{digestclaim}

% claim meat-core/frank-wolfe/lem:fo-khyp
\digestlabel{lem:fo-khyp}{18.2}\pdfbookmark[2]{Lemma 18.2 Positivity needed only on the acceptance support}{digest-meat-core/frank-wolfe/lem:fo-khyp}
\renewcommand*{\RestatedName}{\hyperref{contract.pdf}{}{claim-meat-core/frank-wolfe/lem:fo-khyp}{Lemma~18.2}}%
\begin{digestclaim}[Positivity needed only on the acceptance support]
Let $\Chacc$ be the acceptance support \eqref{eq:chacc-def} of the
acceptance set $\acSet$ of Definition~\ref{def:admissible-acceptance-set}.
\begin{enumerate}[label=(\roman*),leftmargin=2.2em]
\item If $\acSet\neq\varnothing$, $\Chacc$ is the support of a single
element of $\acSet$; and, whether or not $\acSet$ is empty, every
$\mathbf q\in\acSet$ vanishes outside $\Chacc$.
\item On $\acSet\times\Feas^{J}(\varepsilon)$, for every
$\varepsilon\ge0$ and in particular on the exact
$\acSet\times\Feas^{J}$, every term of $O^{J}_{\alpha}$ indexed by
a $\hat c\notin\Chacc$ vanishes identically, so $O^{J}_{\alpha}$ restricted
to that set is a function of $(\mathbf q|_{\Chacc},J)$ alone; and for
$\mathbf q,\mathbf q'\in\acSet$ every feasible direction has
$q'(\hat c)-q(\hat c)=0$ at every $\hat c\notin\Chacc$. The only property
of the $J$ factor the first clause reads is per-block positive
semidefiniteness, so the first clause holds at every per-block positive
semidefinite point whatever, of which $\Feas^{J}(\varepsilon)$ is one
family.
\item Corollary~\ref{cor:full-gradient} holds on
$\acSet\times\Feas^{J}(\varepsilon)$, for every $\varepsilon\ge0$,
with its hypothesis ``$q(\hat c)>0$, $P^{J}_{\hat c}>0$ and
$K^{J}_{\hat c}>0$ for all $\hat c$'' weakened to: $q(\hat c)>0$ and
$P^{J}_{\hat c}>0$ for every $\hat c\in\Chacc$, and $K^{J}_{\hat c}>0$ for
every $\hat c\in\Chkey\cap\Chacc$. Then \eqref{eq:grad-q} is read as a
vector in $\mathbb{R}^{\Chacc}$ and the sums of \eqref{eq:grad-J} are
restricted to $\Chacc$, respectively to $\Chkey\cap\Chacc$; both
restrictions are identities on that set, not truncations.
\end{enumerate}
\end{digestclaim}

% claim meat-core/frank-wolfe/lem:fo-face
\digestlabel{lem:fo-face}{18.3}\pdfbookmark[2]{Lemma 18.3 Every direction is admissible at maximal support}{digest-meat-core/frank-wolfe/lem:fo-face}
\renewcommand*{\RestatedName}{\hyperref{contract.pdf}{}{claim-meat-core/frank-wolfe/lem:fo-face}{Lemma~18.3}}%
\begin{digestclaim}[Every direction is admissible at maximal support]
In the setting of Definition~\ref{def:first-order-setting}, let
$\hat c\in\Chkey$ and $s\in\SecSet_K(\hat c)$, let
$J_s\in\Pos{A_0\Qsys_s}$, and put $X\coloneqq\cG^{J}_{\hat c,s}(J_s)$.
Then for \emph{every} Hermitian $\Delta$ on $\Hs{A_0\Qsys_s}$ the
operator $H\coloneqq\cG^{J}_{\hat c,s}(\Delta)$ satisfies
\begin{equation*}\tag{\getrefnumber{C-eq:fo-face}}\digesteqlabel{eq:fo-face}
  \operatorname{supp}H\;\subseteq\;W_{\hat c,s}
\end{equation*}
with $W_{\hat c,s}$ the maximal branch support \eqref{eq:Wcs}.
Consequently, if $X$ attains maximal support,
$\operatorname{supp}X=W_{\hat c,s}$, for which $J_s\succ0$ suffices,
then $\operatorname{supp}H\subseteq\operatorname{supp}X$ for every such
$\Delta$, so that:
\begin{enumerate}[label=(\alph*),leftmargin=2.2em]
\item Theorem~\ref{thm:gradient}\,(i) applies to the pair $(X,H)$, and
  $t\mapsto T(X+tH)$ is differentiable at $t=0$, two-sidedly, the
  hypothesis being symmetric in $H\mapsto-H$;
\item the $+\infty$ case of Theorem~\ref{thm:gradient}\,(ii) is
  \emph{unreachable} from $X$ along the image of any Hermitian direction
  $\Delta$ for which $X+t H\succeq0$ for small $t>0$. That is the class
  that theorem is stated on and the class along which $T$ is defined at
  all, and it contains every direction $y-x$ that
  Theorem~\ref{thm:firstorder} pairs against, since
  $X+t\cG^{J}_{\hat c,s}(J'_s-J_s)=\cG^{J}_{\hat c,s}\big((1-t)J_s+tJ'_s\big)
  \succeq0$ for $t\in[0,1]$ whenever $J_s$ and $J'_s$ are positive
  semidefinite. No such direction strictly enlarges the support of $X$,
  because by \eqref{eq:fo-face} none leaves
  $W_{\hat c,s}=\operatorname{supp}X$.
\end{enumerate}
\end{digestclaim}

\section*{19\quad The first-order lower bound}\pdfbookmark[1]{19 The first-order lower bound}{digest-group-meat-core/frank-wolfe/thm:firstorder}

% claim meat-core/frank-wolfe/thm:firstorder
\digestlabel{thm:firstorder}{19.1}\pdfbookmark[2]{Theorem 19.1 The estimate and the first-order lower bound}{digest-meat-core/frank-wolfe/thm:firstorder}
\begin{mainclaim}
\renewcommand*{\RestatedName}{\hyperref{contract.pdf}{}{claim-meat-core/frank-wolfe/thm:firstorder}{Theorem~19.1}}%
\begin{digestclaim}[The estimate and the first-order lower bound]
Let $\alpha\in(1,\infty)$, let $\Decl$ satisfy
Definition~\ref{def:admissible-declaration}, in the basic format or in the general
format of Definition~\ref{def:declaration}, and let $\acSet$ satisfy
Definition~\ref{def:admissible-acceptance-set}. In the general format $\Feas^{J}$ is the set of
lists $(J_{(\mathsf t,s)})$ over the joint sectors of
Theorem~\ref{thm:choi}, with $J_s\coloneqq\bigoplus_{\mathsf t}
J_{(\mathsf t,s)}$ for every sector $s$ of the measured system, and
$\Feas^{J}(\varepsilon)$ is read as
Definition~\ref{def:relaxed-feasible-set} reads it there.
\begin{enumerate}[label=(\roman*),leftmargin=2.2em]
\item \emph{(The estimate over-estimates; Lemma~\ref{lem:estimate}\,(ii)
of the core, restated.)} For every $x\in\acSet\times\Feas^{J}$,
\begin{equation*}\tag{\getrefnumber{C-eq:fo-estimate}}\digesteqlabel{eq:fo-estimate}
  O^{J}_{\alpha}(x)\;\ge\;\hlow .
\end{equation*}
\item \emph{(The first-order lower bound.)} Let $\Itset$ be a finite
non-empty index set, and let
$x_t=(\mathbf q_t,J_t)\in\acSet\times\Feas^{J}$, $t\in\Itset$, satisfy
for each $t$:
\begin{itemize}[leftmargin=2.6em]
\item[(H-POS)] $q_t(\hat c)>0$ and $P^{J}_{\hat c}(J_t)>0$ for every
  $\hat c\in\Chacc$, and $K^{J}_{\hat c}(J_t)>0$ for every
  $\hat c\in\Chkey\cap\Chacc$, with $\Chacc$ the acceptance support
  \eqref{eq:chacc-def};
\item[(H-SUPP)] every key-active branch operator
  $X_{\hat c,s}\coloneqq\cG^{J}_{\hat c,s}(J_{t,s})$,
  $\hat c\in\Chkey$, $s\in\SecSet_K(\hat c)$, attains maximal support:
  $\operatorname{supp}X_{\hat c,s}=W_{\hat c,s}$, the
  maximal branch support \eqref{eq:Wcs}, equivalently
  $\operatorname{rank}X_{\hat c,s}=\dim W_{\hat c,s}
  =\operatorname{rank}\mathbf K_{\hat c,s}$ by \eqref{eq:Kstack-rank},
  the inclusion $\subseteq$ holding always by \eqref{eq:Wcs-bound};
  $J_{t,s}\succ0$ suffices for it (Lemma~\ref{lem:branch-support}\,(i));
\item[(H-NOPERT)] the pair
  $(\lmoq_t,\lmoJ_t)\coloneqq\nabla O^{J}_{\alpha}(x_t)$ is the
  split-form gradient of Theorem~\ref{thm:gradient}\,(i)(a)--(b) assembled
  by Corollary~\ref{cor:full-gradient} \emph{at $x_t$ itself}: no
  perturbation of $x_t$, of any $J_{t,s}$, or of any branch operator
  built from them is applied before the gradient is taken.
\end{itemize}
Let $(v_t,S_t)$ satisfy both inequalities of \eqref{eq:lmo} at
$(\lmoq_t,\lmoJ_t)$ (Definition~\ref{def:first-order-setting}): the
acceptance-side inequality and the
feasible-side inequality.
Then, \emph{in exact real arithmetic},
\begin{equation*}\tag{\getrefnumber{C-eq:hlb}}\digesteqlabel{eq:hlb}
  \hlb\;\coloneqq\;\max_{t\in\Itset}\Big[
  O^{J}_{\alpha}(x_t)
  +\langle\lmoq_t,v_t-\mathbf q_t\rangle
  +\langle\lmoJ_t,S_t-J_t\rangle\Big]
  \;\le\;\hlow ,
\end{equation*}
and $\max\{0,\hlb\}\le\hlow$ as well.
\item \emph{(The iterate read at a tolerance, the minima over the exact
set.)} Let $\varepsilon\ge0$ and let $\Feas^{J}(\varepsilon)$ be the
relaxed feasible set of Definition~\ref{def:relaxed-feasible-set}.
\begin{enumerate}[label=(iii\alph*),leftmargin=2.4em]
\item \emph{The estimate.} For every
$x\in\acSet\times\Feas^{J}(\varepsilon)$,
$O^{J}_{\alpha}(x)\ge\hlow[\Feas^{J}(\varepsilon)]$, the Choi-form
relaxed infimum \eqref{eq:hlow-feasJ}; and
$\hlow[\Feas^{J}(\varepsilon)]\le\hlow\le\hqkd$ by
\eqref{eq:feas-tol-mono-J}.
\item \emph{The bound.} Let $x_t\in\acSet\times\Feas^{J}(\varepsilon)$,
$t\in\Itset$, satisfy (H-POS), (H-SUPP) and (H-NOPERT) as stated in (ii),
and let $(v_t,S_t)$ satisfy \eqref{eq:lmo} \emph{as written}, its
acceptance-side inequality and its feasible-side inequality, with the two
minima over $\acSet$ and over the \emph{exact} $\Feas^{J}$. Then, in exact
real arithmetic,
$\hlb\le\hlow$ and $\max\{0,\hlb\}\le\hlow$, with $\hlb$ of
\eqref{eq:hlb}.
\end{enumerate}
\item \emph{(Lower bounds in place of the pairings.)} The conclusions of
(ii) and of (iiib) hold as well when, at some or all $t\in\Itset$, the
pairing $\langle\lmoq_t,v_t\rangle$ in \eqref{eq:hlb} is replaced by a real
number $b^{\mathrm A}_t$, or the pairing $\langle\lmoJ_t,S_t\rangle$ by a
real number $b^{\mathrm F}_t$, or both, with
\begin{equation*}\tag{\getrefnumber{C-eq:lmo-bounds}}\digesteqlabel{eq:lmo-bounds}
  b^{\mathrm A}_t\;\le\;\min_{\mathbf q'\in\acSet}\langle\lmoq_t,\mathbf q'\rangle
  \qquad\text{and}\qquad
  b^{\mathrm F}_t\;\le\;\min_{J'\in\Feas^{J}}\langle\lmoJ_t,J'\rangle ,
\end{equation*}
each inequality asked only where its number is used, the second minimum
taken over the \emph{exact} $\Feas^{J}$, and the pairings
$\langle\lmoq_t,\mathbf q_t\rangle$ and $\langle\lmoJ_t,J_t\rangle$ being
kept. At such a $t$ the inequality of \eqref{eq:lmo} on a replaced side is
not asked of $(v_t,S_t)$, and where both sides are replaced no pair
$(v_t,S_t)$ is asked at all: the bracket of \eqref{eq:hlb} at that $t$ is
then
\begin{equation*}
  O^{J}_{\alpha}(x_t)+b^{\mathrm A}_t-\langle\lmoq_t,\mathbf q_t\rangle
  +b^{\mathrm F}_t-\langle\lmoJ_t,J_t\rangle .
\end{equation*}
\item \emph{(An empty feasible set.)} In place of the hypothesis on $\Decl$
above, let $\Decl$ be a tuple of the form of
Definition~\ref{def:declaration} that satisfies conditions (i)--(vi) and
(viii) of Definition~\ref{def:admissible-declaration}, and form $\Feas$, the
objective $O_{\alpha}$ and $\hlow$ from it by the formulas of
Definition~\ref{def:feasible}. Let $\Feas=\varnothing$. Then $\hlow=+\infty$,
so every real number is a lower bound on $\hlow$, the value of
\eqref{eq:hlb} that a run forms at its iterates included. Consequently the
conclusions $\hlb\le\hlow$ and $\max\{0,\hlb\}\le\hlow$ of (ii), (iiib) and
(iv), for a real $\hlb$, hold of every tuple of the form of
Definition~\ref{def:declaration} that satisfies conditions (i)--(vi) and
(viii) of Definition~\ref{def:admissible-declaration}, whether or not it
satisfies condition~(vii): under the hypotheses of those parts where it does,
and because $\hlow=+\infty$ where it does not.
\end{enumerate}
\end{digestclaim}
\end{mainclaim}

% claim meat-core/frank-wolfe/cor:fo-keylength
\digestlabel{cor:fo-keylength}{19.5}\pdfbookmark[2]{Corollary 19.5 Key length from the first-order bound}{digest-meat-core/frank-wolfe/cor:fo-keylength}
\renewcommand*{\RestatedName}{\hyperref{contract.pdf}{}{claim-meat-core/frank-wolfe/cor:fo-keylength}{Corollary~19.5}}%
\begin{digestclaim}[Key length from the first-order bound]
Let $\alpha\in(1,2)$, let the hypotheses of Theorem~\ref{thm:keylength}
hold in its unconstrained or constrained branch, the run as in
Lemma~\ref{lem:memory-lift}\,(a) and hypothesis (H-CHAIN) included, at every
all-at-once attack with finite registers, with $\Feas^{J}$ carrying
the declared $\CList$ in the constrained branch, let those of
Lemma~\ref{lem:ev-correctness} hold, A-SEQ with its clause (v) among the
former, and let $\hlb$ be as in
Theorem~\ref{thm:firstorder}\,(ii) or~(iiib), at some or all iterates
with the lower bounds of its part~(iv) in place of the pairings. If the
output length
satisfies
\begin{equation*}\tag{\getrefnumber{C-eq:keylength-fo}}\digesteqlabel{eq:keylength-fo}
  \ell\;\le\;n\,\max\{0,\hlb\}-\lEC-\Big\lceil\log_2\frac1{\eev}\Big\rceil
  -\frac{\alpha}{\alpha-1}\log_2\frac1{\epa}+1,
\end{equation*}
then the protocol is $\epa$-secret and $\eev$-correct on the accepted
state, hence, for $\ell\ge1$, $(\epa+\eev)$-secure in the
honest-authentication setting, and under the further hypotheses of
Corollary~\ref{cor:composed-security}\,(iii)
$(\epa+\eev+2\varepsilon_{\mathrm{auth}})$-secure over the classical channel
of that part, with net key gain $\ell-k_{\mathrm{auth}}$ when both parties
output a key, the secrecy and the security statements holding at every
output state, whatever the dimension of the adversary's registers and the
order in which the source releases its signals, as
Corollary~\ref{cor:composed-security} states them, and the correctness
statement at every attack; the four conclusions of
Proposition~\ref{prop:monotone} hold verbatim for the
maximal key length built from \eqref{eq:keylength-fo}; and if
$\hlb\le0$ and $\eev<1$ no key is claimed.
\end{digestclaim}

% claim meat-core/frank-wolfe/cor:fo-dropped
\digestlabel{cor:fo-dropped}{19.6}\pdfbookmark[2]{Corollary 19.6 First-order bound for the dropped program}{digest-meat-core/frank-wolfe/cor:fo-dropped}
\begin{mainclaim}
\renewcommand*{\RestatedName}{\hyperref{contract.pdf}{}{claim-meat-core/frank-wolfe/cor:fo-dropped}{Corollary~19.6}}%
\begin{digestclaim}[First-order bound for the dropped program]
Let $\alpha\in(1,\infty)$, let $\Decl$ satisfy Definition~\ref{def:admissible-declaration} in
the general format of Definition~\ref{def:declaration}, in either source
mode, in trusted-source mode with a frame per Alice sector as in
Theorem~\ref{thm:choi}, let $\acSet$ satisfy
Definition~\ref{def:admissible-acceptance-set}, and let $\mathcal{D}$ be a dropped set
(Definition~\ref{def:dropped}). Write an Alice sector as
$\mathsf t\in\mathbb{T}$ in this claim, the letter $t$ indexing iterates as
in Theorem~\ref{thm:firstorder}. Put
$O^{J,\mathcal{D}}_{\alpha}(\mathbf q,J)\coloneqq
O^{\mathcal{D}}_{\alpha}(\mathbf q,\chi(J))$ and
$K^{J,\mathcal{D}}_{\hat c}\coloneqq K^{\mathcal{D}}_{\hat c}\circ\chi$, with
$\chi$ the assignment of that theorem in trusted-source mode. In
untrusted-source mode the program is read as
Definition~\ref{def:untrusted-choi} reads it, with no frame: $\chi$ is the
identity, $J=\omega$ is the list of the blocks
$J_{(\mathsf t,s)}=\omega^{(\mathsf t,s)}$ over the joint sectors,
$A_{0,\mathsf t}$ and $A_0$ stand for $A_{\mathsf t}$ and $A$, and
$\Feas^{J}$ is $\Feas$ of
Definition~\ref{def:feasible} in the general format, whose condition
coupling the blocks is the normalisation $\sum_{(\mathsf t,s)}
\tr\omega^{(\mathsf t,s)}=1$; so $O^{J,\mathcal{D}}_{\alpha}$ and
$K^{J,\mathcal{D}}_{\hat c}$ are $O^{\mathcal{D}}_{\alpha}$ and
$K^{\mathcal{D}}_{\hat c}$. Let
$x_t=(\mathbf q_t,J_t)\in\acSet\times\Feas^{J}$, $t\in\Itset$, satisfy
\begin{itemize}[leftmargin=3.4em]
\item[(H-POS$^{\mathcal{D}}$)] (H-POS) of Theorem~\ref{thm:firstorder} with
  $K^{J,\mathcal{D}}_{\hat c}$ in place of $K^{J}_{\hat c}$;
\item[(H-SUPP$^{\mathcal{D}}$)] for every $\hat c\in\Chkey$ and every joint
  pair $(\mathsf t,s)\in\mathbb{J}_K(\hat c)$ with
  $(\hat c,(\mathsf t,s))\notin\mathcal{D}$, the branch operator
  $\cG^{J}_{\hat c,s}(J_{t,(\mathsf t,s)})$ of the block
  $J_{t,(\mathsf t,s)}\in\Pos{A_{0,\mathsf t}\Qsys_s}$ attains the maximal
  support of that joint pair, the support it has at
  $J_{t,(\mathsf t,s)}\succ0$; nothing is asked of a pair that is key-inert
  or dropped;
\item[(H-NOPERT$^{\mathcal{D}}$)] $(\lmoq_t,\lmoJ_t)$ is the gradient of
  $O^{J,\mathcal{D}}_{\alpha}$ at $x_t$ itself, assembled as in (H-NOPERT)
  with the kept summands only and with $K^{J,\mathcal{D}}_{\hat c}$ in the
  scalar prefactors, its blocks on fully dropped joint sectors being
  \eqref{eq:drop-gradient} in Choi coordinates;
\end{itemize}
and let $(v_t,S_t)$ satisfy \eqref{eq:lmo} for these linear forms, or let,
at some or all $t$, lower bounds $b^{\mathrm A}_t$, $b^{\mathrm F}_t$ on
one or both of its minima, as in \eqref{eq:lmo-bounds} for these linear
forms, stand in for the pairings as Theorem~\ref{thm:firstorder}\,(iv)
lets them. Then, in
exact real arithmetic, with $\underline{h}^{\mathrm{FO},\mathcal{D}}_{\alpha}$
the expression \eqref{eq:hlb} for $O^{J,\mathcal{D}}_{\alpha}$, with those
replacements where they are made,
\begin{equation*}\tag{\getrefnumber{C-eq:hlb-dropped}}\digesteqlabel{eq:hlb-dropped}
  \max\{0,\underline{h}^{\mathrm{FO},\mathcal{D}}_{\alpha}\}\;\le\;\underline{h}^{\mathcal{D}}_{\alpha}
  \;\le\;\hlow\;\le\;\hqkd(\CList).
\end{equation*}
(H-POS$^{\mathcal{D}}$) is implied by (H-POS), and
(H-SUPP$^{\mathcal{D}}$) by (H-SUPP) read on joint blocks. For
$\mathcal{D}=\varnothing$ and $|\mathbb{T}|=1$ the statement is
Theorem~\ref{thm:firstorder}\,(ii), with its part~(iv) where lower bounds
stand in.

\emph{The relaxed reading.} Let $\varepsilon\ge0$ and let
$\Feas^{J}(\varepsilon)$ be the relaxed feasible set of
Definition~\ref{def:relaxed-feasible-set}, read in the general format, with
one coupling clause per Alice sector in trusted-source mode and its
normalisation clause in untrusted-source mode. Let
$x_t\in\acSet\times\Feas^{J}(\varepsilon)$, $t\in\Itset$, satisfy
(H-POS$^{\mathcal{D}}$), (H-SUPP$^{\mathcal{D}}$) and
(H-NOPERT$^{\mathcal{D}}$) as stated, and let $(v_t,S_t)$ satisfy
\eqref{eq:lmo} \emph{as written}, or the lower bounds standing in for its
pairings satisfy \eqref{eq:lmo-bounds}, with the two minima over $\acSet$
and over the \emph{exact} $\Feas^{J}$. Then \eqref{eq:hlb-dropped} holds in
full, and the case $\varepsilon=0$ is the statement above.

\emph{An empty feasible set.} In place of the hypothesis on $\Decl$ above,
let $\Decl$ be a tuple of the form of Definition~\ref{def:declaration}, in
the general format and in either source mode, that satisfies conditions
(i)--(vi) and (viii) of Definition~\ref{def:admissible-declaration}, and form
$\Feas$, $O_{\alpha}$, $\hlow$ and $\underline{h}^{\mathcal{D}}_{\alpha}$ from
it by the formulas of Definitions~\ref{def:feasible} and~\ref{def:dropped}.
Let $\Feas=\varnothing$. Then
$\underline{h}^{\mathcal{D}}_{\alpha}=\hlow=+\infty$, so every real number is
a lower bound on both, the value of \eqref{eq:hlb} for
$O^{J,\mathcal{D}}_{\alpha}$ that a run forms at its iterates included.
Consequently the first two inequalities of \eqref{eq:hlb-dropped}, for a real
$\underline{h}^{\mathrm{FO},\mathcal{D}}_{\alpha}$, hold of every such
tuple, whether or not it satisfies condition~(vii): under the hypotheses above where it does,
and by the preceding sentence where it does not.
\end{digestclaim}
\end{mainclaim}

% claim meat-core/frank-wolfe/cor:fo-reference
\digestlabel{cor:fo-reference}{19.7}\pdfbookmark[2]{Corollary 19.7 The first-order bound on the reference companion}{digest-meat-core/frank-wolfe/cor:fo-reference}
\begin{mainclaim}
\renewcommand*{\RestatedName}{\hyperref{contract.pdf}{}{claim-meat-core/frank-wolfe/cor:fo-reference}{Corollary~19.7}}%
\begin{digestclaim}[The first-order bound on the reference companion]
Let $\Decl$ and $\Decl^{\mathrm{ref}}$ be as in
Proposition~\ref{prop:reference-companion}, let $\alpha\in(1,\infty)$, and let
$\acSet$ be an admissible acceptance set
(Definition~\ref{def:admissible-acceptance-set}), the one of the run, read
for both. Let $\varepsilon\ge0$ and let $x_t=(\mathbf q_t,J_t)$,
$t\in\Itset$, be a finite non-empty set of points, produced by any
computation, a search on $\Decl$ among them, such that every $x_t$ lies in
$\acSet\times\Feas^{J}(\varepsilon)$ of $\Decl^{\mathrm{ref}}$
(Definition~\ref{def:relaxed-feasible-set}) and satisfies (H-POS), (H-SUPP)
and (H-NOPERT) of Theorem~\ref{thm:firstorder} for the objective, the branch
maps and the gradient of $\Decl^{\mathrm{ref}}$, and let $(v_t,S_t)$
satisfy \eqref{eq:lmo} at the gradient of $\Decl^{\mathrm{ref}}$ at $x_t$,
its feasible-side minimum over the exact $\Feas^{J}$ of $\Decl^{\mathrm{ref}}$
(Definition~\ref{def:first-order-setting}), or let, at some or all $t$,
lower bounds on one or both of those minima stand in for the pairings as
Theorem~\ref{thm:firstorder}\,(iv) lets them, \eqref{eq:lmo-bounds} read
at that gradient and over those sets. Write
$\hlb[\Decl^{\mathrm{ref}}]$ for \eqref{eq:hlb} so formed. Then:
\begin{enumerate}[label=(\roman*),leftmargin=2.4em]
\item in exact real arithmetic,
  $\max\{0,\hlb[\Decl^{\mathrm{ref}}]\}\le\hlow[\Decl^{\mathrm{ref}}]$, and
  $\hlow[\Decl^{\mathrm{ref}}]\le\hlow[\Decl^{\mathrm{ref}}_{\mathrm{ex}}]$
  for the companion with its unrelaxed entries;
\item the Choi variables of $\Decl$ and $\Decl^{\mathrm{ref}}$ are lists over
  the same joint sectors with the same block dimensions,
  so every point of the one program is a point of the other; where no link
  the companion replaces carries entries, as on the encoded route,
  $\acSet\times\Feas^{J}(\varepsilon)$ is the same set for both; and a point with $J_{t,s}\succ0$ for every sector
  $s$ satisfies (H-SUPP) for $\Decl^{\mathrm{ref}}$, whatever its measured
  family;
\item where the hypotheses of Corollary~\ref{cor:fo-keylength} hold for
  $\Decl^{\mathrm{ref}}$, its conclusions hold with
  $\hlb[\Decl^{\mathrm{ref}}]$ in \eqref{eq:keylength-fo}: the key length so
  formed is secret and correct for the protocol whose declaration is
  $\Decl^{\mathrm{ref}}$.
\end{enumerate}
\end{digestclaim}
\end{mainclaim}

\section*{20\quad The verified feasible side}\pdfbookmark[1]{20 The verified feasible side}{digest-group-meat-core/frank-wolfe/lem:feasible-side-dual}

% claim meat-core/frank-wolfe/lem:feasible-side-dual
\digestlabel{lem:feasible-side-dual}{20.1}\pdfbookmark[2]{Lemma 20.1 A lower bound on the feasible-side minimum from any multipliers}{digest-meat-core/frank-wolfe/lem:feasible-side-dual}
\renewcommand*{\RestatedName}{\hyperref{contract.pdf}{}{claim-meat-core/frank-wolfe/lem:feasible-side-dual}{Lemma~20.1}}%
\begin{digestclaim}[A lower bound on the feasible-side minimum from any multipliers]
\begin{enumerate}[label=(\roman*),leftmargin=2.2em]
\item \emph{(Weak duality with a trace identity.)} Let
  $\mathcal{H}_{\mathrm{in}}$ be a Hilbert space of finite dimension
  $d_{\mathrm{in}}\ge1$, let $P$ be a finite non-empty set, and for every
  $p\in P$ let $\mathcal{H}_p$ be a finite-dimensional Hilbert space. Let
  $\mathcal{R}$ be a finite, possibly empty, set of \emph{rows}; a row $r$
  consists of a finite-dimensional Hilbert space $\mathcal{K}_r$, a
  Hermitian operator $\Theta_r$ on it, a relation
  $\bowtie_r\in\{\preceq,\succeq,=\}$ and, for every $p\in P$, a linear map
  $L_{r,p}$ from operators on $\mathcal{H}_{\mathrm{in}}\otimes\mathcal{H}_p$
  to operators on $\mathcal{K}_r$ that maps Hermitian operators to Hermitian
  ones. Write $L^{*}_{r,p}$ for its adjoint, the real-linear map from
  Hermitian operators on $\mathcal{K}_r$ to Hermitian operators on
  $\mathcal{H}_{\mathrm{in}}\otimes\mathcal{H}_p$ with
  $\tr[\Lambda\,L_{r,p}(X)]=\tr[L^{*}_{r,p}(\Lambda)\,X]$ for all Hermitian
  $\Lambda$ and $X$. Let $\mathcal{J}$ be the set of lists
  $J=(J_p)_{p\in P}$ of positive semidefinite operators $J_p$ on
  $\mathcal{H}_{\mathrm{in}}\otimes\mathcal{H}_p$ with
  \begin{equation*}\tag{\getrefnumber{C-eq:abstract-feasible}}\digesteqlabel{eq:abstract-feasible}
    \sum_{p\in P}\tr_{\mathcal{H}_p}J_p=\id_{\mathcal{H}_{\mathrm{in}}}
    \qquad\text{and}\qquad
    L_r(J)\coloneqq\sum_{p\in P}L_{r,p}(J_p)\ \bowtie_r\ \Theta_r
    \quad\text{for every }r\in\mathcal{R}.
  \end{equation*}
  Let $G=(G_p)_{p\in P}$ be a list of Hermitian operators $G_p$ on
  $\mathcal{H}_{\mathrm{in}}\otimes\mathcal{H}_p$, let $Y$ be a Hermitian
  operator on $\mathcal{H}_{\mathrm{in}}$, and for every row $r$ let
  $\Lambda_r$ be a Hermitian operator on $\mathcal{K}_r$, with
  $\Lambda_r\succeq0$ when $\bowtie_r$ is $\preceq$ or $\succeq$. Put
  $\sigma_r\coloneqq-1$ when $\bowtie_r$ is $\preceq$ and
  $\sigma_r\coloneqq+1$ otherwise,
  \begin{equation*}\tag{\getrefnumber{C-eq:dual-residual}}\digesteqlabel{eq:dual-residual}
    R_p\;\coloneqq\;G_p-Y\otimes\id_{\mathcal{H}_p}
    -\sum_{r\in\mathcal{R}}\sigma_r\,L^{*}_{r,p}(\Lambda_r)
    \qquad(p\in P),
  \end{equation*}
  and let $\lambda$ be a real number with $\lambda\le\lmin(R_p)$ for every
  $p\in P$. Then for every $J\in\mathcal{J}$
  \begin{equation*}\tag{\getrefnumber{C-eq:feasible-side-dual}}\digesteqlabel{eq:feasible-side-dual}
    \langle G,J\rangle\coloneqq\sum_{p\in P}\tr[G_pJ_p]
    \;\ge\;b(G;Y,\Lambda,\lambda)\coloneqq
    d_{\mathrm{in}}\,\lambda+\tr Y
    +\sum_{r\in\mathcal{R}}\sigma_r\tr[\Lambda_r\Theta_r],
  \end{equation*}
  so that $b(G;Y,\Lambda,\lambda)\le\inf_{J\in\mathcal{J}}\langle G,J\rangle$.
  No other condition is placed on $Y$, on the $\Lambda_r$ or on $\lambda$:
  $Y$ and the $\Lambda_r$ need not be feasible for any dual program,
  $\lambda$ may be negative, and $Y=0$ with every $\Lambda_r=0$ is admitted.
\item \emph{(The Choi feasible set is of that form.)} Let $\Decl$ be a
  declaration in the general format of Definition~\ref{def:declaration}, the
  basic format included, with its Alice sectors
  $\mathsf t\in\mathbb{T}$.
  \begin{enumerate}[label=(\alph*),leftmargin=2.2em]
  \item In trusted-source mode, with a frame and the factors
    $W_{\mathsf t}$ as in Theorem~\ref{thm:choi}, the factor
    $\Feas^{J}_{\mathsf t}$ of $\Feas^{J}=\prod_{\mathsf t}\Feas^{J}_{\mathsf t}$
    (Theorem~\ref{thm:choi}, the general format) is the set $\mathcal{J}$ of
    (i) with $\mathcal{H}_{\mathrm{in}}=\Hs{A_{0,\mathsf t}}$, $P=\SecSet$,
    $\mathcal{H}_s=\Hs{\Qsys_s}$ and $J_s=J_{(\mathsf t,s)}$, and with the
    following rows. For every entry $k$ read through its
    lifted operator, one row with $\mathcal{K}=\Hs{A_{\mathsf t}}$,
    $\Theta=\Theta^{(\mathsf t)}_k$, the relation $\bowtie_k$ and
    \begin{equation*}
      L_{k,s}(X)=W_{\mathsf t}\,\tr_{\Qsys_s}\!\big[(\idon{A_{0,\mathsf t}}
      \otimes\hat O^{\uparrow(s)}_k)X\big]\,W_{\mathsf t}^{\dagger},
      \qquad
      L^{*}_{k,s}(\Lambda)=(W_{\mathsf t}^{\dagger}\Lambda W_{\mathsf t})
      \otimes\hat O^{\uparrow(s)}_k .
    \end{equation*}
    For every party-local entry $k$, scoped to the measured factor $f$,
    and every $s_{f^c}$, one row with
    $\mathcal{K}=\Hs{A_{\mathsf t}}\otimes\Hs{\Qsys_{f^c,s_{f^c}}}$,
    $\Theta=0$, the relation $\bowtie_k$ and, at $s=(s_{f^c},s_f)$,
    \begin{gather*}
      L_{(k,s_{f^c}),s}(X)=(W_{\mathsf t}\otimes\id)\,
      \tr_{\Qsys_{f,s_f}}\!\big[(\id\otimes D^{(s_f)}_k)X\big]\,
      (W_{\mathsf t}^{\dagger}\otimes\id),\\
      L^{*}_{(k,s_{f^c}),s}(\Lambda)=\big((W_{\mathsf t}^{\dagger}\otimes\id)
      \Lambda(W_{\mathsf t}\otimes\id)\big)\otimes D^{(s_f)}_k ,
    \end{gather*}
    where $D^{(s_f)}_k\coloneqq\hat O^{(s_f)}_k-c_k\idon{}$ acts on
    $\Hs{\Qsys_{f,s_f}}$ at the place of that factor in $\Hs{\Qsys_s}$, while
    $L_{(k,s_{f^c}),s'}=0$ and $L^{*}_{(k,s_{f^c}),s'}=0$ at every sector
    $s'$ whose component on the other factors is not $s_{f^c}$. The two
    kinds of map, and the formulas for their adjoints, are defined, and the
    formulas hold, for any linear map $W$ from $\Hs{A_{0,\mathsf t}}$ to
    $\Hs{A_{\mathsf t}}$ in place of $W_{\mathsf t}$; we write
    $L^{W}_{r,s}$ for a row's map so built. A
    homogeneous or party-local row in the frame-free form of
    Proposition~\ref{prop:frame-free-rows}, on a tight frame, is a row
    of the same kind with $\idon{A_{0,\mathsf t}}$ in place of
    $W_{\mathsf t}$, hence with $\Hs{A_{0,\mathsf t}}$ in place of
    $\Hs{A_{\mathsf t}}$ in its $\mathcal{K}$, and with
    $\Theta=c_k\idon{A_{0,\mathsf t}}$, respectively $\Theta=0$; and
    $\Feas^{J}_{\mathsf t}$ is unchanged when it replaces the conjugated
    row. Hence, for $Y_{\mathsf t}$, $\Lambda_{\mathsf t}=(\Lambda_r)_r$ and
    $\lambda_{\mathsf t}$ as in (i) for each $\mathsf t$, and with
    $G_{\mathsf t}\coloneqq(G_{(\mathsf t,s)})_{s\in\SecSet}$,
    \begin{equation*}\tag{\getrefnumber{C-eq:feasible-side-sum}}\digesteqlabel{eq:feasible-side-sum}
      b_{\mathrm F}\;\coloneqq\;\sum_{\mathsf t\in\mathbb{T}}
      b(G_{\mathsf t};Y_{\mathsf t},\Lambda_{\mathsf t},\lambda_{\mathsf t})
      \;\le\;\langle G,J'\rangle
      \qquad\text{for every }J'\in\Feas^{J}.
    \end{equation*}
  \item In untrusted-source mode, $\Feas^{J}$, which is $\Feas$
    (Definition~\ref{def:first-order-setting}), is the set $\mathcal{J}$ of
    (i) with $\mathcal{H}_{\mathrm{in}}$ one-dimensional, so that $Y=y$ is a
    real number, $Y\otimes\id=y\,\id$, $\tr Y=y$ and $d_{\mathrm{in}}=1$;
    with $P$ the set of joint sectors,
    $\mathcal{H}_{(\mathsf t,s)}=\Hs{A_{\mathsf t}\Qsys_s}$ and
    $J_{(\mathsf t,s)}=\omega^{(\mathsf t,s)}$; and with these rows. For every
    scalar entry $k$, one row with $\mathcal{K}=\mathbb{C}$, $\Theta=c_k$,
    the relation $\bowtie_k$,
    $L_{k,(\mathsf t,s)}(X)=\tr[(\idon{A_{\mathsf t}}\otimes
    \hat O^{\uparrow(s)}_k)X]$ and
    $L^{*}_{k,(\mathsf t,s)}(\mu)=\mu\,\idon{A_{\mathsf t}}\otimes
    \hat O^{\uparrow(s)}_k$; for every party-local entry $k$ and every pair
    $(\mathsf t,s_{f^c})$, one row as in (a) with $\idon{A_{\mathsf t}}$ in
    place of $W_{\mathsf t}$, its maps vanishing outside the joint sectors
    $(\mathsf t,(s_{f^c},s_f))$. Hence
    $b(G;y,\Lambda,\lambda)\le\langle G,J'\rangle$ for every
    $J'\in\Feas^{J}$, the whole list being one set of the form of (i).
  \end{enumerate}
\item \emph{(A scalar coupling and no rows.)} If $d_{\mathrm{in}}=1$ and
  $\mathcal{R}=\varnothing$ in (i), then, with $Y=0$ and
  $\lambda=\min_{p\in P}\lmin(G_p)$, the infimum in (i) is attained and
  $b(G;0,\Lambda,\lambda)=\min_{J\in\mathcal{J}}\langle G,J\rangle$.
\item \emph{(Rows built from an inexact frame.)} In trusted-source mode fix
  an Alice sector $\mathsf t$, a linear map $\widetilde W_{\mathsf t}$ from
  $\Hs{A_{0,\mathsf t}}$ to $\Hs{A_{\mathsf t}}$ and a real number
  $\varepsilon_{\mathsf t}\ge\opn{\widetilde W_{\mathsf t}-W_{\mathsf t}}$.
  For every row $r$ of $\Feas^{J}_{\mathsf t}$ in (ii)(a) let
  $\widetilde L_{r,s}\coloneqq L^{\widetilde W_{\mathsf t}}_{r,s}$, a row in
  frame-free form keeping its map, every other datum of the row
  ($\Theta_r$, $\hat O^{\uparrow(s)}_k$, $D^{(s_f)}_k$, $c_k$) unchanged. Let
  $G_{\mathsf t}$, $Y$ and the $\Lambda_r$ be as in (i), form
  $\widetilde R_s$ by \eqref{eq:dual-residual} with
  $\widetilde L^{*}_{r,s}$ in place of $L^{*}_{r,s}$, and let
  $\widetilde\lambda\le\lmin(\widetilde R_s)$ for every $s\in\SecSet$. Then
  for every $J\in\Feas^{J}_{\mathsf t}$, the factor cut out by the rows
  built from the exact $W_{\mathsf t}$,
  \begin{equation*}\tag{\getrefnumber{C-eq:frame-deficit}}\digesteqlabel{eq:frame-deficit}
    \begin{gathered}
    \langle G_{\mathsf t},J\rangle\;\ge\;
    b(G_{\mathsf t};Y,\Lambda,\widetilde\lambda)-\sum_{r}\kappa_r,\\
    \kappa_r\coloneqq d_{A_{0,\mathsf t}}\,\varepsilon_{\mathsf t}
    \big(2\opn{\widetilde W_{\mathsf t}}+\varepsilon_{\mathsf t}\big)\,
    \opn{\Lambda_r}\,\max_{s\in\SecSet}\opn{O_{r,s}},
    \end{gathered}
  \end{equation*}
  where $O_{r,s}$ is $\hat O^{\uparrow(s)}_k$ for a lifted-operator row,
  $D^{(s_f)}_k$ for a party-local row at $s=(s_{f^c},s_f)$ and $0$ at the
  other sectors, and $\kappa_r\coloneqq0$ for a row in frame-free form.
  Summed over $\mathsf t$ as in (ii)(a), the right-hand sides bound
  $\langle G,J'\rangle$ from below for every $J'\in\Feas^{J}$.
\end{enumerate}
\end{digestclaim}

\section*{21\quad The certified bracket}\pdfbookmark[1]{21 The certified bracket}{digest-group-meat-core/frank-wolfe/lem:certified-bracket}

% claim meat-core/frank-wolfe/lem:certified-bracket
\digestlabel{lem:certified-bracket}{21.1}\pdfbookmark[2]{Lemma 21.1 The certified bracket from enclosures of the gradient}{digest-meat-core/frank-wolfe/lem:certified-bracket}
\renewcommand*{\RestatedName}{\hyperref{contract.pdf}{}{claim-meat-core/frank-wolfe/lem:certified-bracket}{Lemma~21.1}}%
\begin{digestclaim}[The certified bracket from enclosures of the gradient]
Let $\alpha\in(1,\infty)$, let $\Decl$ be a tuple of the form of
Definition~\ref{def:declaration}, in the basic format or the general
format, that satisfies conditions (i)--(vi) and (viii) of
Definition~\ref{def:admissible-declaration}, and let $\acSet$ satisfy
Definition~\ref{def:admissible-acceptance-set}. Let $x=(\mathbf q,J)$ with
$\mathbf q\in\acSet$ and every block of $J$ positive definite, on every
joint sector in the general format, satisfy (H-POS) of
Theorem~\ref{thm:firstorder}, and let $(\lmoq,\lmoJ)$ be the gradient of
$O^{J}_{\alpha}$ at $x$ assembled by Corollary~\ref{cor:full-gradient},
read with Lemma~\ref{lem:fo-khyp}\,(iii) and, in the general format, on the
joint sectors as Theorem~\ref{thm:firstorder} reads it.
\begin{enumerate}[label=(\roman*),leftmargin=2.4em]
\item \emph{(The Euler form.)}
  \begin{equation*}\tag{\getrefnumber{C-eq:euler-form}}\digesteqlabel{eq:euler-form}
    \langle\lmoq,\mathbf q\rangle+\langle\lmoJ,J\rangle=O^{J}_{\alpha}(x).
  \end{equation*}
  Consequently, for real numbers $b^{\mathrm A}$ and $b^{\mathrm F}$ with
  $b^{\mathrm A}\le\min_{\mathbf q'\in\acSet}\langle\lmoq,\mathbf q'\rangle$
  and $b^{\mathrm F}\le\min_{J'\in\Feas^{J}}\langle\lmoJ,J'\rangle$, the
  minimum over the exact $\Feas^{J}$, the bracket of
  Theorem~\ref{thm:firstorder}\,(iv) at $x$ with both pairings replaced is
  $b^{\mathrm A}+b^{\mathrm F}$, and
  \begin{equation*}\tag{\getrefnumber{C-eq:certified-bracket}}\digesteqlabel{eq:certified-bracket}
    \max\{0,\,b^{\mathrm A}+b^{\mathrm F}\}\;\le\;\hlow ,
  \end{equation*}
  in exact real arithmetic.
\item \emph{(The acceptance side from an enclosure.)} If
  $\hat g\in\mathbb{R}^{\Chacc}$ and $\delta\ge0$ satisfy
  $|\lmoq(\hat c)-\hat g(\hat c)|\le\delta$ for every $\hat c\in\Chacc$, then
  $\min_{\mathbf q'\in\acSet}\langle\lmoq,\mathbf q'\rangle\ge
  \min_{\mathbf q'\in\acSet}\langle\hat g,\mathbf q'\rangle-\delta$.
\item \emph{(The feasible side from enclosed residuals.)} For each group of
  Lemma~\ref{lem:feasible-side-dual}\,(ii), an Alice sector $\mathsf t$ in
  trusted-source mode, its rows built from the exact frame
  $W_{\mathsf t}$ of Theorem~\ref{thm:choi}, or the one group in
  untrusted-source mode, let $Y$ and the $\Lambda_r$ be as in
  Lemma~\ref{lem:feasible-side-dual}\,(i), let $R_p$ be the residual
  \eqref{eq:dual-residual} formed with $\lmoJ$, and for every block $p$ of the
  group let $M_p$ be a Hermitian matrix, $\rho_p\ge0$ with
  $\|R_p-M_p\|_{\mathrm F}\le\rho_p$, the Frobenius norm of
  \xref{foundations/cor:least-eigenvalue-over-enclosure}, and $\ell_p$ a
  lower bound on $\lmin(M_p)$. Then the group's
  $\lambda\coloneqq\min_p(\ell_p-\rho_p)$ satisfies $\lambda\le\lmin(R_p)$
  for every $p$, and the sum over the groups of $b(\lmoJ;Y,\Lambda,\lambda)$
  of \eqref{eq:feasible-side-dual} is a number
  $b^{\mathrm F}\le\min_{J'\in\Feas^{J}}\langle\lmoJ,J'\rangle$.
\item \emph{(The dropped program.)} In the setting of
  Corollary~\ref{cor:fo-dropped}, with a dropped set $\mathcal{D}$,
  (i)--(iii) hold with $O^{J,\mathcal{D}}_{\alpha}$ in place of
  $O^{J}_{\alpha}$, (H-POS$^{\mathcal{D}}$) in place of (H-POS), its gradient
  of (H-NOPERT$^{\mathcal{D}}$) in place of $(\lmoq,\lmoJ)$, the kept pairs in
  place of the pairs that are not key-inert, and
  $\underline{h}^{\mathcal{D}}_{\alpha}$ in place of $\hlow$; and
  $\underline{h}^{\mathcal{D}}_{\alpha}\le\hlow$.
\end{enumerate}
\end{digestclaim}

% claim meat-core/frank-wolfe/lem:certified-bracket-enclosed
\digestlabel{lem:certified-bracket-enclosed}{21.2}\pdfbookmark[2]{Lemma 21.2 The certified bracket over a measured family known through enclosures}{digest-meat-core/frank-wolfe/lem:certified-bracket-enclosed}
\renewcommand*{\RestatedName}{\hyperref{contract.pdf}{}{claim-meat-core/frank-wolfe/lem:certified-bracket-enclosed}{Lemma~21.2}}%
\begin{digestclaim}[The certified bracket over a measured family known through enclosures]
Let $\Decl^{\mathrm{ref}}$ be a compiled reference companion of
Proposition~\ref{prop:reference-companion}, with $\mathbb{F}^{\mathrm{ref}}$
the set of the measured factors it replaces (the one
measured system in the basic format), let $\alpha$, $\acSet$ and
$x=(\mathbf q,J)$ be as in Lemma~\ref{lem:certified-bracket} with
$\Decl\coloneqq\Decl^{\mathrm{ref}}$, save that (H-POS) of
Theorem~\ref{thm:firstorder} is not assumed of $x$, and, where (H-POS) holds
for $\Decl^{\mathrm{ref}}$ at $x$, let $(\lmoq,\lmoJ)$ be the gradient that
lemma reads there, the blocks of the measured families of the factors of
$\mathbb{F}^{\mathrm{ref}}$ each known through an entrywise enclosure
(\xref{foundations/def:exact-reference}), some of them read, besides, with a
triple $(L_c,t,W)$ of a matrix box and a core as in
\xref{foundations/def:factor-enclosure}, the block's \emph{factor
enclosure}, and every other datum exactly. Call a tuple
$F=(F^{f})_{f\in\mathbb{F}^{\mathrm{ref}}}$ of families, each $F^{f}$
indexed as the measured family of factor $f$ is, a \emph{family of the
replaced factors}, its elements $F^{(s)}_v$ being the products
\eqref{eq:productfamily} of its blocks and of the blocks of the other factors
of $\Decl^{\mathrm{ref}}$; call it an \emph{entrywise member} when each of
its blocks is a Hermitian matrix that the entrywise enclosure of that block
contains, and a \emph{member} when, moreover, each of its blocks that is read
with a factor enclosure is a member of it. Suppose
\begin{itemize}[leftmargin=4.6em]
\item[(H-FAC)] each block of the companion's own family of the replaced
  factors that is read with a factor enclosure is a member of it, the triple
  being a factor enclosure of that block, so that the companion's family of
  the replaced factors is a member.
\end{itemize}
Call a block \emph{decided} when every entry of its entrywise enclosure has
centre and radius $0$; or some entry's centre $C_{il}$ and radius $r_{il}$
satisfy $|C_{il}|>r_{il}$; or its factor enclosure has no column; or its
factor enclosure has $k\ge1$ columns and there are $G$, $\rho$ and $\ell$ as
in \xref{foundations/cor:rank-over-factor-enclosure}\,(ii) with
$\ell-\rho>0$. Call a block \emph{of fixed rank} when it is of one of three
kinds: of the first kind when it is zero for every member; of the second kind
when there are a matrix $E$ of full row rank and, for every member, a
positive definite $\widetilde X$ with the member's block equal to
$E^{\dagger}\widetilde XE$, $E$ the same for every member
(\xref{foundations/cor:least-eigenvalue-over-enclosure}\,(iv) gives such an
$E$); and of the third kind when there are an integer $k\ge1$, a matrix box
of complex $k\times n$ matrices, $n$ the dimension of the block, and a
positive definite $k\times k$ matrix $\widetilde X$ and, for every member, a
matrix $E_{*}$ of rank $k$ that the box contains, with the member's block
equal to $E_{*}^{\dagger}\widetilde XE_{*}$
(\xref{foundations/cor:rank-over-factor-enclosure}\,(ii) gives such a box,
that of the adjoints of the factors, with $\widetilde X$ the core). Suppose
that every block of a factor of $\mathbb{F}^{\mathrm{ref}}$ that is a factor
of an element $F^{(s)}_v$ with $(x',v)$ an outcome pair of a pass class
of some $\hat c\in\Chkey$ and $M^{(\mathsf t)}_{x'}\ne0$, no other block of
which is zero for every member, is decided, and that every such block is of
fixed rank at the pairs $(\hat c,(\mathsf t,s))$ that are not
key-inert. Then:
  \begin{enumerate}[label=(\alph*),leftmargin=2.2em]
  \item every member gives the same
    effective outcome sets on every joint
    sector, the same key-inert pairs and the same
    joint key-active index sets, those
    read off the enclosures;
  \item fix, for every block that the hypothesis asks to be of fixed rank,
    one kind it is of, with the data of that kind, and, where that kind is
    the third, at every member one matrix $E_{*}$ as that kind asks, one for
    the block and not one for each element it is a factor of; below, the
    kind of a block is the one so fixed for it. At every pair
    $(\hat c,(\mathsf t,s))$ that is not key-inert, every
    element $F^{(s)}_v$ of an outcome pair $(x',v)$ of a pass class of
    $\hat c$ with $M^{(\mathsf t)}_{x'}\ne0$ is, for every member,
    $E^{(s)\dagger}_v\widetilde F^{(s)}_vE^{(s)}_v$ with
    $\widetilde F^{(s)}_v$ positive definite and $E^{(s)}_v$ of full row
    rank, the case $F^{(s)}_v=0$ included as $E^{(s)}_v$ with no rows, where
    $E^{(s)}_v$ is the same for every member if no block of $F^{(s)}_v$ is of
    the third kind or some block of it is zero for every member, and
    otherwise its rank is the same for every member and it is the tensor
    product, in the order of $\mathbb{F}$, of one matrix for each block of
    $F^{(s)}_v$: the matrix $E_{*}$ fixed for the block at that member where
    the block is of the third kind, and a matrix the same for every member
    otherwise. For an element of the latter case, call its \emph{tuple}
    $\beta(v)$ the list, in the order of $\mathbb{F}$, of its blocks of the
    third kind, each block $F^{f,(s_f)}_{v_f}$ named by the triple
    $(f,s_f,v_f)$, so that two entries of tuples agree exactly when they
    name the same block, whatever their boxes; and for a tuple $\beta$ let
    $E^{\mathrm{box}}_{\beta}$ be the tensor product, in that order, of the
    matrices $E_{*}$ fixed for its blocks, which at each member is one
    matrix for every element whose tuple is $\beta$.
    Moreover, for every pass class of $\hat c$ and any factor forms
    $M^{(\mathsf t)}_{x}=E^{(\mathsf t)\dagger}_{x}\widetilde M_{x}
    E^{(\mathsf t)}_{x}$ of Alice's elements with $E^{(\mathsf t)}_{x}$ of
    full row rank and $\widetilde M_{x}$ positive definite, the matrix
    $\widetilde E$ that stacks the matrices
    $E^{(\mathsf t)}_{x}\otimes E^{(s)}_v$ over the class's outcome pairs
    $(x,v)$ with $M^{(\mathsf t)}_{x}\ne0$, in a fixed order, is, for every
    member, $\widetilde E=\widetilde B\,C$ with $\widetilde B$ of full column
    rank and the same for every member, and $C$ a stack of a matrix $C_0$
    and, one for each tuple $\beta$ of an element of the latter case in the
    class, of a matrix $(C_{\beta}\otimes E^{\mathrm{box}}_{\beta})U_{\beta}$,
    with $C_0$, every $C_{\beta}$ and every permutation matrix $U_{\beta}$
    the same for every member; where the $E^{(\mathsf t)}_{x}$, the
    $E^{(s)}_v$ that are the same for every member and the matrices the same
    for every member of which the other $E^{(s)}_v$ are tensor products have
    entries in $\mathbb{Q}(i)$, $\widetilde B$, $C_0$ and every $C_{\beta}$
    can be taken with entries in $\mathbb{Q}(i)$;
  \item call a pair $(F',F'')$ of families of the replaced factors, the
    first read by the statistics and the second by the deficits,
    \emph{admissible} when every block of $F''$ that is a factor of an
    element $F^{(s)}_v$ of an outcome pair $(x',v)$ of a pass class of
    $\hat c$ with $M^{(\mathsf t)}_{x'}\ne0$, at a pair
    $(\hat c,(\mathsf t,s))$ that is not key-inert, no other block of which
    is zero for every member, is a positive semidefinite matrix. For such a pair form, at $x$, $P^{J}_{\hat c}$ and
    every $\Gamma^{(s')}_{\hat c}$ with $F'$, every deficit
    $f^{J}_{\hat c,s}$ and every $\Delta^{(s')}_{\hat c}$ with $F''$, an
    outcome pair whose element of Alice's is zero, or whose element
    $F^{(s)}_v$ has a block that is zero for every member, contributing
    nothing to a branch map or to the statistics, over the key-inert pairs and the joint key-active index sets
    of (a), and
    $K^{J}_{\hat c}$ as $P^{J}_{\hat c}$ plus those deficits; where
    $q(\hat c)$, $P^{J}_{\hat c}$ and $K^{J}_{\hat c}$ so formed are positive
    as (H-POS) of Theorem~\ref{thm:firstorder} asks, let
    $(\lmoq,\lmoJ)[F',F'']$ be the pair given by \eqref{eq:grad-q} and
    \eqref{eq:grad-J} with these pieces, read as
    Lemma~\ref{lem:certified-bracket} reads them. Let $\mathcal P$ be a set
    of admissible pairs that contains $(F,F)$ for every member $F$, such as
    the set of the admissible pairs of an entrywise member $F'$ and a family
    $F''$ of the replaced factors each of whose blocks lies in its entrywise
    enclosure or is a member of its factor enclosure, the one or the other
    chosen block by block. Then
    a set that contains $P^{J}_{\hat c}$, or $K^{J}_{\hat c}$, formed with
    every pair of $\mathcal P$ contains that of $\Decl^{\mathrm{ref}}$ at
    $x$; and where (H-POS) holds for $\Decl^{\mathrm{ref}}$, enclosures
    $(\hat g,\delta)$ and $(M_p,\rho_p)$ that hold, as in parts (ii) and
    (iii) of Lemma~\ref{lem:certified-bracket}, at
    $(\lmoq,\lmoJ)[F',F'']$ for every pair of $\mathcal P$ at which it is
    formed hold at the gradient of $\Decl^{\mathrm{ref}}$, and parts
    (i)--(iii) of that lemma then give
    $\max\{0,b^{\mathrm A}+b^{\mathrm F}\}\le\hlow[\Decl^{\mathrm{ref}}]$;
  \item in the setting of
    Corollary~\ref{cor:fo-dropped}, with a dropped set
    $\mathcal{D}$, (a)--(c) hold, the hypotheses read the same way, with
    $O^{J,\mathcal{D}}_{\alpha}$ in place of $O^{J}_{\alpha}$,
    (H-POS$^{\mathcal{D}}$) in place of (H-POS), its gradient of
    (H-NOPERT$^{\mathcal{D}}$) in place of $(\lmoq,\lmoJ)$ and
    $K^{J,\mathcal{D}}_{\hat c}$ in place of $K^{J}_{\hat c}$, the kept
    pairs in place of the pairs that are not key-inert, part (iv) of
    Lemma~\ref{lem:certified-bracket} in place of its parts (i)--(iii),
    and $\underline{h}^{\mathcal{D}}_{\alpha}$ in place of $\hlow$.
  \end{enumerate}
\end{digestclaim}

% claim meat-core/frank-wolfe/cor:certified-keylength
\digestlabel{cor:certified-keylength}{21.3}\pdfbookmark[2]{Corollary 21.3 Key length from the certified bracket}{digest-meat-core/frank-wolfe/cor:certified-keylength}
\begin{mainclaim}
\renewcommand*{\RestatedName}{\hyperref{contract.pdf}{}{claim-meat-core/frank-wolfe/cor:certified-keylength}{Corollary~21.3}}%
\begin{digestclaim}[Key length from the certified bracket]
Let $\alpha$, $\Decl$ and $\acSet$ be as in
Lemma~\ref{lem:certified-bracket}, let $\Itset$ be a finite non-empty set of
points $x_t=(\mathbf q_t,J_t)$ each satisfying the hypotheses of that lemma,
and let $b^{\mathrm A}_t$ and $b^{\mathrm F}_t$ be real numbers with
\eqref{eq:lmo-bounds} at the gradient at $x_t$, as parts (ii) and (iii) of
that lemma form them. Put
\begin{equation*}\tag{\getrefnumber{C-eq:hcert}}\digesteqlabel{eq:hcert}
  \underline{h}^{\mathrm{cert}}_{\alpha}\coloneqq
  \max_{t\in\Itset}\big(b^{\mathrm A}_t+b^{\mathrm F}_t\big).
\end{equation*}
\begin{enumerate}[label=(\roman*),leftmargin=2.4em]
\item $\underline{h}^{\mathrm{cert}}_{\alpha}$ is the value of \eqref{eq:hlb}
  at the points of $\Itset$ with both pairings replaced as
  Theorem~\ref{thm:firstorder}\,(iv) replaces them, and
  $\max\{0,\underline{h}^{\mathrm{cert}}_{\alpha}\}\le\hlow$ in exact real
  arithmetic.
\item Where $\Decl$ satisfies Definition~\ref{def:admissible-declaration} and
  the hypotheses of Corollary~\ref{cor:fo-keylength} hold, its conclusions
  hold with $\underline{h}^{\mathrm{cert}}_{\alpha}$ in place of $\hlb$ in
  \eqref{eq:keylength-fo}: the key length so formed is $\epa$-secret and
  $\eev$-correct, with the composed-security and monotonicity conclusions of
  that corollary, the secrecy and security statements at every output
  state, whatever the dimension of the adversary's registers and the order
  in which the source releases its signals.
\item Under a dropped set, with the lemma's part (iv),
  $\max\{0,\underline{h}^{\mathrm{cert}}_{\alpha}\}\le
  \underline{h}^{\mathcal{D}}_{\alpha}\le\hlow$; and where $\Decl$
  satisfies Definition~\ref{def:admissible-declaration} and the hypotheses
  of Corollary~\ref{cor:keylength-bound} other than
  \eqref{eq:lower-bound-h} hold, the conclusions of that corollary hold at
  $h=\underline{h}^{\mathrm{cert}}_{\alpha}$.
\end{enumerate}
\end{digestclaim}
\end{mainclaim}

% claim meat-core/frank-wolfe/cor:certified-keylength-reference
\digestlabel{cor:certified-keylength-reference}{21.4}\pdfbookmark[2]{Corollary 21.4 Key length on a reference companion from the certified bracket}{digest-meat-core/frank-wolfe/cor:certified-keylength-reference}
\renewcommand*{\RestatedName}{\hyperref{contract.pdf}{}{claim-meat-core/frank-wolfe/cor:certified-keylength-reference}{Corollary~21.4}}%
\begin{digestclaim}[Key length on a reference companion from the certified bracket]
Let $\Decl^{\mathrm{ref}}$ be a compiled reference companion of
Proposition~\ref{prop:reference-companion}, and let $\alpha$, $\acSet$ and
$\Itset$, with its points $x_t$, be as in
Corollary~\ref{cor:certified-keylength} with
$\Decl\coloneqq\Decl^{\mathrm{ref}}$. At each $t\in\Itset$ let
$b^{\mathrm A}_t$ and $b^{\mathrm F}_t$ be real numbers formed in one of two
ways: as in that corollary, with \eqref{eq:lmo-bounds} at the gradient of
$\Decl^{\mathrm{ref}}$ at $x_t$; or, where blocks of the measured family of
$\Decl^{\mathrm{ref}}$ are known through enclosures, entrywise or by factor
enclosures, as parts (ii) and (iii) of Lemma~\ref{lem:certified-bracket} form
them from enclosures $(\hat g,\delta)$ and $(M_p,\rho_p)$ that hold at every
pair of a set $\mathcal P$ as in
Lemma~\ref{lem:certified-bracket-enclosed}\,(c), at a point $x_t$ at which
the hypotheses of that lemma hold and (H-POS) holds for
$\Decl^{\mathrm{ref}}$, the blocks it reads with factor enclosures being read
with those of Corollary~\ref{cor:reference-companion-factors}, under that
corollary's hypotheses (for each replaced factor whose blocks are so read,
the rule of (I1) of the reduction instance of its last link a non-negative
combination, and the blocks of its handed reference carrying those factor
enclosures), so that its conclusion is (H-FAC). Put $\underline{h}^{\mathrm{cert}}_{\alpha}\coloneqq
\max_{t\in\Itset}(b^{\mathrm A}_t+b^{\mathrm F}_t)$, as in \eqref{eq:hcert}.
\begin{enumerate}[label=(\roman*),leftmargin=2.4em]
\item Where $\Decl^{\mathrm{ref}}$ satisfies
  Definition~\ref{def:admissible-declaration} and the hypotheses of
  Corollary~\ref{cor:fo-keylength} hold, the conclusions of
  Corollary~\ref{cor:fo-reference}\,(iii) hold with
  $\underline{h}^{\mathrm{cert}}_{\alpha}$ in place of
  $\hlb[\Decl^{\mathrm{ref}}]$, for the protocol whose declaration is
  $\Decl^{\mathrm{ref}}$.
\item In the setting of Corollary~\ref{cor:fo-dropped}, with a dropped set
  $\mathcal{D}$ and with the numbers formed at the gradient of the dropped
  objective and, in the second way, with part (d) of
  Lemma~\ref{lem:certified-bracket-enclosed} in place of its part (c):
  $\max\{0,\underline{h}^{\mathrm{cert}}_{\alpha}\}\le
  \underline{h}^{\mathcal{D}}_{\alpha}[\Decl^{\mathrm{ref}}]\le
  \hlow[\Decl^{\mathrm{ref}}]$, and where $\Decl^{\mathrm{ref}}$ satisfies
  Definition~\ref{def:admissible-declaration} and the hypotheses of
  Corollary~\ref{cor:keylength-bound} other than \eqref{eq:lower-bound-h}
  hold, the conclusions of
  Corollary~\ref{cor:keylength-bound} hold at
  $h=\underline{h}^{\mathrm{cert}}_{\alpha}$, for the protocol whose
  declaration is $\Decl^{\mathrm{ref}}$.
\end{enumerate}
\end{digestclaim}

\section*{22\quad The two minimisations and the arithmetic}\pdfbookmark[1]{22 The two minimisations and the arithmetic}{digest-group-meat-core/frank-wolfe/ass:lmo}

% claim meat-core/frank-wolfe/ass:lmo
\digestlabel{ass:lmo}{22.1}\pdfbookmark[2]{Assumption 22.1 Exact linear minimisation over the feasible set}{digest-meat-core/frank-wolfe/ass:lmo}
\SetShortId{A-LMO}
\renewcommand*{\RestatedName}{\hyperref{contract.pdf}{}{claim-meat-core/frank-wolfe/ass:lmo}{Assumption~22.1}}%
\begin{digestassumption}[Exact linear minimisation over the feasible set]
(i) At each iterate $x_t$ at which the engine accumulates a bracket of the
first-order bound, the block list $S_t$ it consumes, which is the backend's
return after any positivity repair the adapter applies to it, satisfies the
feasible-side inequality of \eqref{eq:lmo}
(Definition~\ref{def:first-order-setting}), its minimum taken over the
\emph{exact} $\Feas^{J}$ and not over the relaxed feasible set, at
$\lmoJ_t$, the block part of the gradient $(\lmoq_t,\lmoJ_t)$ the bracket
consumes at $x_t$, of the objective, or of that of a dropped program.
\end{digestassumption}

\section*{24\quad Requirements: reported numbers and withdrawal}\pdfbookmark[1]{24 Requirements: reported numbers and withdrawal}{digest-group-meat-core/frank-wolfe/req:first-order-withdrawal}

% claim meat-core/frank-wolfe/req:first-order-withdrawal
\digestlabel{req:first-order-withdrawal}{24.1}\pdfbookmark[2]{Requirement 24.1 Withdrawal of the first-order bound}{digest-meat-core/frank-wolfe/req:first-order-withdrawal}
\renewcommand*{\RestatedName}{\hyperref{contract.pdf}{}{claim-meat-core/frank-wolfe/req:first-order-withdrawal}{Requirement~24.1}}%
\begin{digestrequirement}[Withdrawal of the first-order bound]
An iterate at which any of the following holds is not a member of
$\Itset$: it contributes no bracket to the maximum in \eqref{eq:hlb}, and the
engine records the machine-readable reason.
\begin{enumerate}[label=(\roman*),leftmargin=2.4em]
\item (H-POS) fails (\iface{gradient_hypothesis_positivity});
\item (H-SUPP) fails (\iface{gradient_hypothesis_support});
\item a gradient perturbation was applied, that is, (H-NOPERT) fails
  (\iface{gradient_perturbation_applied});
\item the point consumed at that iterate fails one of the three checks of
  Requirement~\ref{req:iterate-checks} and is not repairable under
  Requirement~\ref{req:iterate-repair}, or its acceptance-side component
  $\mathbf q_t$ was not formed on the exact island under ($\alpha$)
  and ($\beta$) of Remark~\ref{C-rem:acc-side}, so that its membership in the
  product $\acSet\times\Feas^{J}(\varepsilon)$ is not established
  (\iface{iterate_not_feasible} in either case).
\end{enumerate}
Neither the point the oracle returned at that iterate nor its status
removes the iterate from $\Itset$.
If no iterate is left, $\Itset$ is empty: the engine withdraws $\hlb$ and
the tier falls to the estimate, with a reason from the same
machine-readable vocabulary.
\end{digestrequirement}

% claim meat-core/frank-wolfe/req:returned-point-measured
\digestlabel{req:returned-point-measured}{24.2}\pdfbookmark[2]{Requirement 24.2 Measurement of the returned point}{digest-meat-core/frank-wolfe/req:returned-point-measured}
\renewcommand*{\RestatedName}{\hyperref{contract.pdf}{}{claim-meat-core/frank-wolfe/req:returned-point-measured}{Requirement~24.2}}%
\begin{digestrequirement}[Measurement of the returned point]
Every run applies the three checks of Requirement~\ref{req:iterate-checks}
once more, at the end of the run, to the point it returned: per-block
positivity at exactly $0$, the coupling residual at \iface{iterate_coupling}
and each $\bowtie_k$ residual at \iface{constraint_rows}. It records whether
that point passes all three, together with the three quantities measured,
in \iface{optimizer_certificate}. The measurement enters neither the
first-order bound, nor the tier, nor the key length.
\end{digestrequirement}

% claim meat-core/frank-wolfe/req:first-order-clamp
\digestlabel{req:first-order-clamp}{24.3}\pdfbookmark[2]{Requirement 24.3 Clamping of the first-order bound at zero}{digest-meat-core/frank-wolfe/req:first-order-clamp}
\renewcommand*{\RestatedName}{\hyperref{contract.pdf}{}{claim-meat-core/frank-wolfe/req:first-order-clamp}{Requirement~24.3}}%
\begin{digestrequirement}[Clamping of the first-order bound at zero]
The engine passes $\max\{0,\hlb\}$ to the key length of
Corollary~\ref{cor:fo-keylength}, and records the clamp: it emits the
boolean \iface{first_order_clamped} beside \iface{h_lb}, and \iface{h_lb}
carries the unclamped value.
\end{digestrequirement}

% claim meat-core/frank-wolfe/req:negative-estimate-refused
\digestlabel{req:negative-estimate-refused}{24.4}\pdfbookmark[2]{Requirement 24.4 Refusal of a negative estimate}{digest-meat-core/frank-wolfe/req:negative-estimate-refused}
\renewcommand*{\RestatedName}{\hyperref{contract.pdf}{}{claim-meat-core/frank-wolfe/req:negative-estimate-refused}{Requirement~24.4}}%
\begin{digestrequirement}[Refusal of a negative estimate]
The engine refuses a run whose computed estimate, the objective at
the point the run's search on the declaration it searches returned at an
order the run evaluates, is negative, with the
machine-readable reason \iface{negative_objective} naming the value and the
point, and never clamps the estimate to zero. A bare exception is not a
conforming implementation of this rule; the reason must be readable by the
caller. The rule binds that search alone: the polish of
Requirement~\ref{req:bracket-on-reference}\,(i) and the solve of
Requirement~\ref{req:unreachable-acceptance-set}\,(i) form no estimate of
the run, and a negative objective value of theirs refuses nothing.
\end{digestrequirement}

% claim meat-core/frank-wolfe/req:bound-estimate-difference
\digestlabel{req:bound-estimate-difference}{24.5}\pdfbookmark[2]{Requirement 24.5 The recorded difference between the bound and the estimate}{digest-meat-core/frank-wolfe/req:bound-estimate-difference}
\renewcommand*{\RestatedName}{\hyperref{contract.pdf}{}{claim-meat-core/frank-wolfe/req:bound-estimate-difference}{Requirement~24.5}}%
\begin{digestrequirement}[The recorded difference between the bound and the estimate]
The engine records the difference $\hlb$ minus the reported
estimate, with its sign, as a diagnostic, and asserts nothing of it.
\end{digestrequirement}

% claim meat-core/frank-wolfe/req:no-imported-penalty
\digestlabel{req:no-imported-penalty}{24.6}\pdfbookmark[2]{Requirement 24.6 The program's own objective, with no imported penalty}{digest-meat-core/frank-wolfe/req:no-imported-penalty}
\renewcommand*{\RestatedName}{\hyperref{contract.pdf}{}{claim-meat-core/frank-wolfe/req:no-imported-penalty}{Requirement~24.6}}%
\begin{digestrequirement}[The program's own objective, with no imported penalty]
The engine imports no relative-entropy penalty onto $O^{J}_{\alpha}$ from any
source: it subtracts from it no penalty term taken from a
von-Neumann-entropy key-rate formulation.
\end{digestrequirement}

% claim meat-core/frank-wolfe/req:bracket-on-reference
\digestlabel{req:bracket-on-reference}{24.7}\pdfbookmark[2]{Requirement 24.7 The first-order bound evaluated on the reference companion}{digest-meat-core/frank-wolfe/req:bracket-on-reference}
\renewcommand*{\RestatedName}{\hyperref{contract.pdf}{}{claim-meat-core/frank-wolfe/req:bracket-on-reference}{Requirement~24.7}}%
\begin{digestrequirement}[The first-order bound evaluated on the reference companion]
Where the engine has admitted a compiled reference companion
$\Decl^{\mathrm{ref}}$ beside the declaration $\Decl$ it searched on
(Requirement~\ref{req:companion-admission}), it forms the run's first-order
number as follows, after the search on $\Decl$ has ended.
\begin{enumerate}[label=(\roman*),leftmargin=2.4em]
\item \emph{The points.} It evaluates the points of two sets. The first
  holds at most five points of the search on $\Decl$: the at most four
  members of the index set of $\Decl$'s bound
  (Requirement~\ref{req:first-order-withdrawal}) with the largest brackets
  on $\Decl$, ties broken by the order of the run, and the last point the
  search assessed; where $\Decl$'s index set is empty, the last point alone.
  The second holds the points kept by the \emph{polish}, a search on the
  Choi program of $\Decl^{\mathrm{ref}}$ that continues the search on
  $\Decl$ after it has ended. The polish reads the objective, the gradient,
  the maps and the entries of $\Decl^{\mathrm{ref}}$
  (Requirement~\ref{req:companion-admission}\,(i)) and no map generated from
  $\Decl$'s measured family; it starts at the point of the first set with
  the largest bracket on $\Decl$, at the last point where $\Decl$'s index
  set is empty, and forms its acceptance-side iterate as
  Requirement~\ref{req:acceptance-iterate-exact} requires, from that point's
  exact acceptance-side component. Where that component was not formed on
  the exact island (Requirement~\ref{req:acceptance-iterate-exact}),
  the polish does not run, the second set is empty and the record names
  that reason. It assesses each point it consumes by the
  gate of (ii) and keeps the at most four of them that contribute with the
  largest brackets on $\Decl^{\mathrm{ref}}$, ties broken by its own order,
  and the last point it assessed; a polish that assesses no point keeps
  none. It runs within the share \iface{reference_polish_share} of the
  wall budget of one evaluation of the run, where the run has one, set
  aside for it before the search on $\Decl$ begins, and ends where that
  share is spent, where an iteration cap it records is reached or where its
  own stopping rule holds, whichever comes first; a polish that fails or is
  cut short keeps, as above, the points it assessed until then. The share
  so set aside shortens the search on $\Decl$, which bears on the run's
  estimate, on $\hlb[\Decl]$, on a number certified on $\Decl$ in the
  fallback of Requirement~\ref{req:certified-evaluation}\,(i), and on the
  soundness of none of them. Its objective values are not the run's
  estimate: none is reported, and none refuses the run
  (Requirement~\ref{req:negative-estimate-refused} reads the estimate of the
  search on $\Decl$ alone).
  Each point $x_t=(\mathbf q_t,J_t)$ of either set is taken as the search
  or the polish consumed it, $\mathbf q_t$ included, and is a point of the
  Choi program of $\Decl^{\mathrm{ref}}$
  (Corollary~\ref{cor:fo-reference}\,(ii)).
\item \emph{The gate.} At each such point it forms the bracket of
  \eqref{eq:hlb} for $\Decl^{\mathrm{ref}}$: the objective of
  $\Decl^{\mathrm{ref}}$ at $x_t$, its gradient at $x_t$ with no
  perturbation, the acceptance-side minimisation of
  Requirements~\ref{req:acceptance-geometries}
  and~\ref{req:acceptance-iterate-exact}, and the feasible side of
  Requirement~\ref{req:verified-feasible-side}, from a fresh feasible-side
  oracle call on the linear form of that gradient over
  $\Decl^{\mathrm{ref}}$'s entries. It decides (H-POS) and (H-SUPP) on
  $\Decl^{\mathrm{ref}}$'s maps (Requirement~\ref{req:support-rank-check})
  and makes every check of Requirement~\ref{req:iterate-checks} against
  $\Decl^{\mathrm{ref}}$'s entries, on the point; a point at which
  Requirement~\ref{req:first-order-withdrawal} withdraws it on these terms
  contributes no bracket, with that requirement's reason.
\item \emph{The number.} Where some evaluated point contributes and the
  bracket is not withdrawn by (iv), the run's
  first-order number is $\hlb[\Decl^{\mathrm{ref}}]$, the largest bracket so
  formed, its clamp and key length are those of
  Requirement~\ref{req:first-order-clamp} and
  Corollary~\ref{cor:fo-keylength} at it, with the parameters of $\Decl$'s
  run, and its tier is \emph{first-order}. The bound on $\Decl$,
  $\hlb[\Decl]$, is recorded beside it with its brackets as a diagnostic,
  and no key length is formed from it. The difference
  $\hlb[\Decl^{\mathrm{ref}}]-\hlb[\Decl]$ is recorded with its sign, as a
  diagnostic, beside the difference of
  Requirement~\ref{req:bound-estimate-difference}, which is taken against
  the estimate of $\Decl$. A point of the second set that contributes
  counts as any evaluated point does: where no point of the first set
  contributes, it makes this item apply in place of the fallback of (iv),
  so that the run's first-order number moves from $\hlb[\Decl]$ to the
  companion.
\item \emph{The bracket withdrawn.} The bracket on the companion is
  withdrawn for one of two causes: no evaluated point contributes; or the
  certified evaluation finds the two enclosures of a block the companion
  carries inconsistent, the cause \iface{factor_enclosure_inconsistent}
  (Requirement~\ref{req:certified-evaluation}\,(iv)). An inconsistency shows
  that a module carrying the enclosures is defective, and the first-order
  number on the companion reads its entrywise enclosure, which may be the
  defective one; the second cause is decided after that number was formed,
  and the run then forms its first-order number again as this item says,
  before any key length is reported. Where the bracket is withdrawn, the
  run's first-order number is $\hlb[\Decl]$, formed as without a
  companion, with the reason \iface{reference_bracket_withdrawn}, its cause
  and the reasons of every evaluated point; no block is then evaluated off its snap.
  The record names in the group \emph{reference bracket withdrawn}, with the
  cause (\xref{foundations/req:snap-reference}\,(iii)(b)), every moved block
  on which a declaring instance whose link the companion replaces handed the
  reference and every moved class of each snapped emission of such a link
  whose normalizer the companion bypasses; it names as it does without a
  companion every other moved block or class that rests on
  \xref{foundations/ass:snapped-measurement}, so that a moved block on which a declaring
  instance whose link the companion does not replace handed a reference
  rests in the group \emph{reference bracket not evaluated}
  (\xref{foundations/req:snap-reference}\,(iii)(c)). A run
  that declares the strict option \iface{strict_snaps} refuses there instead
  and reports no key length (\xref{foundations/req:snap-reference}\,(iv)).
\item \emph{The blocks.} Where (iii) applies and the key length the run
  reports is not formed in the fallback of
  Requirement~\ref{req:certified-evaluation}\,(i), the record of the run names as
  evaluated off its snap (\xref{foundations/def:evaluated-off-snap})
  every moved block of the snap of each declaring instance whose
  link the companion replaces, on which the companion carries the handed
  reference, and every data sector of each snapped emission of such a link
  that holds a moved class and whose normalizer the companion bypasses; a kept block or sector is named as kept and is not
  counted among them. It names every moved block on which the companion
  carries the snapped block as not so evaluated, so that it rests on
  \xref{foundations/ass:snapped-measurement}
  (\xref{foundations/req:snap-reference}\,(iii)), and a run that declares
  \iface{strict_snaps} refuses there (\xref{foundations/req:snap-reference}\,(iv)).
  Where (iii) applies but the key length the run reports is formed from a
  number certified on $\Decl$ in the fallback of
  Requirement~\ref{req:certified-evaluation}\,(i), no number formed on the
  companion reaches that key length, and the number it is formed from reads
  the snapped blocks: the moved blocks on which the companion carries the
  handed reference and the data sectors whose normalizer it bypasses are
  then not evaluated off their snap, condition (ii) of
  \xref{foundations/def:evaluated-off-snap} failing for them. The record
  names them as resting on \xref{foundations/ass:snapped-measurement} in
  the group \emph{reference bracket withdrawn} with the cause
  \iface{companion_not_certified}
  (\xref{foundations/req:snap-reference}\,(iii)(b)), beside the moved blocks
  on which the companion carries the snapped block, which rest as above. The
  run decides this after the certified evaluation and before any key length
  is reported, and no block is then named as evaluated off its
  snap. A run that declares \iface{strict_snaps} forms no such
  number (Requirement~\ref{req:certified-evaluation}\,(i)).
\item \emph{The record.} The record carries, per evaluated point, the set
  of (i) it belongs to, its bracket on $\Decl^{\mathrm{ref}}$ and the
  withdrawal reasons there, and for a point of the first set its bracket on
  $\Decl$; for the polish, whether it ran and, where it did not, why, the
  point it started from, the share of the wall budget set aside for it and
  the wall time it took, its iteration cap, the number of points it
  assessed and the cause it stopped on; the measurements of
  Requirement~\ref{req:companion-admission}\,(iii), the largest radius of the
  companion's enclosure among them, which bear on the arithmetic of
  Remark~\ref{C-rem:fo-arithmetic} and on nothing else; and, per block,
  whether the companion
  carries the handed reference or the snapped block.
\end{enumerate}
A run with no admitted companion forms its first-order number on $\Decl$
alone, as the other requirements of this submodule state; a moved block on
which a reference was handed, and a moved class of an emission beside which a
companion that bypasses its normalizer was formed, then rest in the group
\emph{reference bracket not evaluated}
(\xref{foundations/req:snap-reference}\,(iii)).
\end{digestrequirement}

\section*{25\quad Requirements: checks and repairs of the iterates}\pdfbookmark[1]{25 Requirements: checks and repairs of the iterates}{digest-group-meat-core/frank-wolfe/req:iterate-checks}

% claim meat-core/frank-wolfe/req:iterate-checks
\digestlabel{req:iterate-checks}{25.1}\pdfbookmark[2]{Requirement 25.1 Feasibility checks at each iterate}{digest-meat-core/frank-wolfe/req:iterate-checks}
\renewcommand*{\RestatedName}{\hyperref{contract.pdf}{}{claim-meat-core/frank-wolfe/req:iterate-checks}{Requirement~25.1}}%
\begin{digestrequirement}[Feasibility checks at each iterate]
On the point it consumes at each iterate, after any repair, the engine
decides every condition defining $\Feas^{J}$ in \eqref{eq:choi-feas}, as the
relaxed feasible set relaxes it, all three and none of them optional:
\begin{enumerate}[label=(\roman*),leftmargin=2.4em]
\item per-block positive semidefiniteness, at exactly $0$ and at no window
  (\iface{iterate_positivity});
\item the coupling condition, its residual
  $\|\sum_s\tr_{\Qsys_s}J_s-\idon{A_0}\|_{\max}$, or, in
  untrusted-source mode, where the coupling condition is the
  normalisation $\sum_s\tr\omega_s=1$
  (Definition~\ref{def:untrusted-choi}\,(b)), its residual
  $|\sum_s\tr\omega_s-1|$, at the declared window \iface{iterate_coupling};
\item the declared constraint rows $\bowtie_k$, each residual at the declared
  window \iface{constraint_rows}.
\end{enumerate}
No fourth check is required, and none is required of the point the oracle
returned at that iterate. The
$\varepsilon$ at which a run reads the relaxed feasible set is the larger of
the two declared windows of (ii) and~(iii). No declared tolerance (\xref{foundations/def:tolerance}) is
that number: the windows govern conditions checked on an iterate and never on
a declaration.
\end{digestrequirement}

% claim meat-core/frank-wolfe/req:support-rank-check
\digestlabel{req:support-rank-check}{25.2}\pdfbookmark[2]{Requirement 25.2 Rank test for maximal support}{digest-meat-core/frank-wolfe/req:support-rank-check}
\renewcommand*{\RestatedName}{\hyperref{contract.pdf}{}{claim-meat-core/frank-wolfe/req:support-rank-check}{Requirement~25.2}}%
\begin{digestrequirement}[Rank test for maximal support]
The engine decides (H-SUPP) as the rank equality \eqref{C-eq:hsupp-check},
both sides computed and both read with the one declared threshold
\iface{support_cutoff}:
\begin{enumerate}[label=(\roman*),leftmargin=2.4em]
\item the left side is the count of eigenvalues of $X_{\hat c,s}$ above
  \iface{support_cutoff}, read relative to that block's own spectral radius;
\item the right side is $\operatorname{rank}\mathbf K_{\hat c,s}$ of
  \eqref{eq:Kstack}, read at the square root of \iface{support_cutoff}
  relative to that matrix's largest singular value, a window derived from
  the same threshold and not declared beside it; equivalently the count of
  eigenvalues of $\mathbf K_{\hat c,s}\mathbf K_{\hat c,s}^{\dagger}$ above
  \iface{support_cutoff} relative to its spectral radius. It is computed once
  per $(\hat c,s)$ from the declaration and $W_{\tau}$, and never from an
  iterate.
\end{enumerate}
\end{digestrequirement}

% claim meat-core/frank-wolfe/req:oracle-repair
\digestlabel{req:oracle-repair}{25.3}\pdfbookmark[2]{Requirement 25.3 Repair of an oracle return}{digest-meat-core/frank-wolfe/req:oracle-repair}
\renewcommand*{\RestatedName}{\hyperref{contract.pdf}{}{claim-meat-core/frank-wolfe/req:oracle-repair}{Requirement~25.3}}%
\begin{digestrequirement}[Repair of an oracle return]
The engine may reject outright an oracle return a block of which has a
negative spectrum, or project that block onto the positive semidefinite cone
within the declared window \iface{oracle_repair_window}, a number of its own,
declared beside $\tnum$ and never as an instance of it. It records every
such repair, with the block, the iterate and the measured depth, in
\iface{optimizer_certificate}, kept apart there from the iterate repairs of
Requirement~\ref{req:iterate-repair}, and uses the repaired return only as
a step of the search and in the diagnostic pairing of
Requirement~\ref{req:verified-feasible-side}\,(vi), never in a bracket and
never as a bound on anything.
\end{digestrequirement}

% claim meat-core/frank-wolfe/req:iterate-repair
\digestlabel{req:iterate-repair}{25.4}\pdfbookmark[2]{Requirement 25.4 Repair of a formed iterate}{digest-meat-core/frank-wolfe/req:iterate-repair}
\renewcommand*{\RestatedName}{\hyperref{contract.pdf}{}{claim-meat-core/frank-wolfe/req:iterate-repair}{Requirement~25.4}}%
\begin{digestrequirement}[Repair of a formed iterate]
When a block of the point consumed at iterate $t$ fails the positivity check
of Requirement~\ref{req:iterate-checks}, the engine either rejects it, in
which case $t$ contributes no bracket and is not a member of $\Itset$, or
projects it onto the positive semidefinite cone, provided its measured depth
is within the declared window \iface{iterate_repair_window}, a number of its
own, declared beside $\tnum$ and never as an instance of it. No other repair
is made on an iterate. A projection is governed by three rules:
\begin{enumerate}[label=(\roman*),leftmargin=2.4em]
\item the repaired point is $x_t$, the point at which the objective, the
  gradient and the pairings with $x_t$ of \eqref{eq:hlb} are evaluated;
\item the coupling and $\bowtie_k$ residuals of
  Requirement~\ref{req:iterate-checks} are measured again on the repaired
  point and never carried over from the point before it, and a residual that
  then fails its window refuses the iterate;
\item every repair is recorded, with the block, the iterate and the measured
  depth, in \iface{optimizer_certificate}, beside the oracle-return repairs
  of Requirement~\ref{req:oracle-repair} and distinguished from them.
\end{enumerate}
\end{digestrequirement}

% claim meat-core/frank-wolfe/req:per-sector-coupling
\digestlabel{req:per-sector-coupling}{25.5}\pdfbookmark[2]{Requirement 25.5 Coupling residual and record per Alice sector}{digest-meat-core/frank-wolfe/req:per-sector-coupling}
\renewcommand*{\RestatedName}{\hyperref{contract.pdf}{}{claim-meat-core/frank-wolfe/req:per-sector-coupling}{Requirement~25.5}}%
\begin{digestrequirement}[Coupling residual and record per Alice sector]
On a declaration in the general format the engine decides the coupling
condition of the relaxed feasible set of
Definition~\ref{def:relaxed-feasible-set} per Alice sector, one coupling
residual per Alice sector, each at the declared window
\iface{iterate_coupling}; and the record names the frame and the row form
per Alice sector (Requirements~\ref{req:frame-recorded}
and~\ref{req:row-form-record}).
\end{digestrequirement}

\section*{26\quad Requirements: the acceptance side}\pdfbookmark[1]{26 Requirements: the acceptance side}{digest-group-meat-core/frank-wolfe/req:acceptance-geometries}

% claim meat-core/frank-wolfe/req:acceptance-geometries
\digestlabel{req:acceptance-geometries}{26.1}\pdfbookmark[2]{Requirement 26.1 The admitted acceptance geometries}{digest-meat-core/frank-wolfe/req:acceptance-geometries}
\renewcommand*{\RestatedName}{\hyperref{contract.pdf}{}{claim-meat-core/frank-wolfe/req:acceptance-geometries}{Requirement~26.1}}%
\begin{digestrequirement}[The admitted acceptance geometries]
The engine solves the acceptance-side minimisation exactly for the two
geometries of Remark~\ref{C-rem:acc-side}, a singleton~(a) and a polytope~(b),
and admits an acceptance set in no other geometry: a declaration whose
$\acSet$ fits neither is refused at admission, before any solve, and no
first-order bound is formed for it.
\end{digestrequirement}

% claim meat-core/frank-wolfe/req:acceptance-iterate-exact
\digestlabel{req:acceptance-iterate-exact}{26.2}\pdfbookmark[2]{Requirement 26.2 Exact formation of the acceptance-side iterate}{digest-meat-core/frank-wolfe/req:acceptance-iterate-exact}
\renewcommand*{\RestatedName}{\hyperref{contract.pdf}{}{claim-meat-core/frank-wolfe/req:acceptance-iterate-exact}{Requirement~26.2}}%
\begin{digestrequirement}[Exact formation of the acceptance-side iterate]
The engine forms the acceptance-side iterate on the exact island,
under conditions ($\alpha$) and ($\beta$) of Remark~\ref{C-rem:acc-side}: it
initialises it at an exact rational point of $\acSet$, and updates it only by
convex combinations, with exact rational weights in $[0,1]$, of points
already established to lie in $\acSet$. An iterate whose acceptance-side
component was not so formed has not met the membership hypothesis and
contributes no bracket to $\hlb$, with the reason of
Requirement~\ref{req:first-order-withdrawal}\,(iv). The run's record states in
\iface{optimizer_certificate}, of the point consumed at each iterate, whether
its acceptance-side component was so formed.
\end{digestrequirement}

\section*{27\quad Requirements: the feasible side}\pdfbookmark[1]{27 Requirements: the feasible side}{digest-group-meat-core/frank-wolfe/req:verified-feasible-side}

% claim meat-core/frank-wolfe/req:verified-feasible-side
\digestlabel{req:verified-feasible-side}{27.1}\pdfbookmark[2]{Requirement 27.1 The verified feasible side of every bracket}{digest-meat-core/frank-wolfe/req:verified-feasible-side}
\renewcommand*{\RestatedName}{\hyperref{contract.pdf}{}{claim-meat-core/frank-wolfe/req:verified-feasible-side}{Requirement~27.1}}%
\begin{digestrequirement}[The verified feasible side of every bracket]
At every iterate $t$ at which the engine forms a bracket of \eqref{eq:hlb},
on the declaration it searched and on an admitted reference companion alike,
it uses in place of the pairing $\langle\lmoJ_t,S_t\rangle$ the number
$b^{\mathrm F}_t$ of Theorem~\ref{thm:firstorder}\,(iv), formed as follows,
and never the pairing with a point the backend returned, a dual objective
the backend reported, or the backend's status.
\begin{enumerate}[label=(\roman*),leftmargin=2.4em]
\item \emph{The groups and the data.} $b^{\mathrm F}_t$ is a sum over the
  groups of the program, one per Alice sector in trusted-source mode
  and one in untrusted-source mode, read as
  Lemma~\ref{lem:feasible-side-dual}\,(ii) reads them. Its $G$ is the exact
  Hermitian part of the gradient blocks as the engine computed them at
  $x_t$, the same gradient whose acceptance-side part the bracket uses; its
  coupling is the exact coupling condition of the group; its rows are the
  maps of Lemma~\ref{lem:feasible-side-dual}\,(ii), in the conjugated or
  frame-free form the engine imposes, built from the declared data taken
  exactly: $\Theta^{(\mathsf t)}_k$ as the declared block or the exact
  product $c_k\tau^{(\mathsf t)}_A$, the lifted operators, and
  $D^{(s_f)}_k=\hat O^{(s_f)}_k-c_k\idon{}$ formed exactly, all Hermitian by
  exact admissibility; and in trusted-source mode the frame in those maps is
  the binary64 $\widetilde W_{\mathsf t}$ the engine constructed for
  $W_{\mathsf t}$, read as the dyadic matrix it is. Every binary64 number is
  read as the dyadic rational it is.
\item \emph{The multipliers.} $Y$ and every $\Lambda_r$ are read off the
  dual the backend returned for that same solve, as the operators that the
  multipliers of the coupling condition and of row $r$ represent in the
  encoding the oracle used, each replaced by its exact Hermitian part. On a
  group with a one-dimensional coupling and no rows, and on any group whose
  solve returned no finite dual, whatever its status, $Y=0$ and every
  $\Lambda_r=0$; on an Alice sector for which (iv) forms no
  $\varepsilon_{\mathsf t}$, every $\Lambda_r$ of its rows is $0$ and $Y$ is
  kept. The multiplier of an inequality row is replaced by
  $\Lambda_r+\mu_r\id$, with $\mu_r$ a dyadic rational at least
  $\max\{0,-\ell_r\}$ and $\ell_r$ a lower bound on $\lmin(\Lambda_r)$ given
  by \xref{foundations/lem:least-eigenvalue-bound}, so that
  $\Lambda_r\succeq0$ holds exactly.
\item \emph{The eigenvalue.} The group's $\lambda$ is the least, over the
  blocks of the group, of the lower bounds on $\lmin(R_p)$ that
  \xref{foundations/lem:least-eigenvalue-bound} gives at a matrix $V$ of the
  engine's choosing, with $R_p$ formed from the rows of (i). The bound holds
  at every $V$ that lemma admits, the identity ($\theta=0$) always among
  them, and the choice bears only on how close $\lambda$ lies to
  $\min_p\lmin(R_p)$.
\item \emph{The frame.} For each Alice sector the engine forms a dyadic
  $\varepsilon_{\mathsf t}\ge\opn{\widetilde W_{\mathsf t}-W_{\mathsf t}}$,
  established in exact rational arithmetic, where it can: $0$ where
  $\widetilde W_{\mathsf t}=W_{\mathsf t}$ is decided exactly; and where
  $\tau^{(\mathsf t)}_A$ is diagonal, the frame is the full frame or
  the diagonal frame (Definition~\ref{def:support-frame}), and
  $\widetilde W_{\mathsf t}$ vanishes exactly wherever $W_{\mathsf t}$
  does, so that both have at most one non-zero entry in each row and each
  column and $\opn{\widetilde W_{\mathsf t}-W_{\mathsf t}}$ is the largest
  modulus of their entrywise difference: a dyadic at least
  $|\widetilde w-\sqrt{\tau_{ii}}|$ at every entry where $W_{\mathsf t}$
  carries $\sqrt{\tau_{ii}}$, $\widetilde w$ being the entry of
  $\widetilde W_{\mathsf t}$ there (for a real $\widetilde w>0$,
  $|\widetilde w^{2}-\tau_{ii}|/\widetilde w$ is one, since
  $|\widetilde w-\sqrt{\tau_{ii}}|=|\widetilde w^{2}-\tau_{ii}|/
  (\widetilde w+\sqrt{\tau_{ii}})$). Where every row of the sector is
  imposed in frame-free form, no deduction is formed: those rows read no
  frame, every $\kappa_r$ is $0$, and the sector's bound is that of
  Lemma~\ref{lem:feasible-side-dual}\,(i). In every other case no
  $\varepsilon_{\mathsf t}$ is formed. Where it is formed, the sector's bound
  is that of Lemma~\ref{lem:feasible-side-dual}\,(iv),
  $b(G_{\mathsf t};Y,\Lambda,\widetilde\lambda)-\sum_r\kappa_r$, with each
  $\kappa_r$ formed from dyadic upper bounds on
  $\opn{\widetilde W_{\mathsf t}}$, $\opn{\Lambda_r}$ and $\opn{O_{r,s}}$
  (for instance on their Frobenius norms, by an integer square root rounded
  up). Where it is not, the row multipliers are zero by (ii) and the
  sector's bound is that of Lemma~\ref{lem:feasible-side-dual}\,(i), which
  then reads no row. In untrusted-source mode there is no frame and no
  deduction.
\item \emph{The arithmetic.} The residuals, the shifted multipliers, every
  eigenvalue and norm bound, every deduction and $b^{\mathrm F}_t$ are
  computed in exact rational arithmetic on those dyadic data, as a
  certification of the exact island
  (\xref{foundations/def:exact-island}\,(iii)), the backend's dual entering
  it only as a candidate (\xref{foundations/req:exact-island}).
  $b^{\mathrm F}_t$ enters the binary64 bracket rounded toward $-\infty$, the
  direction that weakens a lower bound, as
  \xref{foundations/req:rounding-direction} rounds an emitted one.
\item \emph{The record.} Per bracket the engine records $b^{\mathrm F}_t$
  (\iface{feasible_side_bound}); the sum of the deductions $\kappa_r$
  (\iface{frame_deduction}) and the sum of the bounds on
  $\opn{\Lambda_r}$ it used (\iface{row_multiplier_norm}); where the backend
  returned a point, the pairing $\langle\lmoJ_t,S_t\rangle$ with it
  (\iface{oracle_pairing}) and the difference
  $\langle\lmoJ_t,S_t\rangle-b^{\mathrm F}_t$ with its sign
  (\iface{dual_gap}); and, per group, whether its multipliers were the
  backend's, zero, or the backend's with the row multipliers set to zero
  (\iface{feasible_side_multipliers}). None of these but
  $b^{\mathrm F}_t$ enters any number.
\end{enumerate}
\end{digestrequirement}

% claim meat-core/frank-wolfe/req:primal-pairing
\digestlabel{req:primal-pairing}{27.2}\pdfbookmark[2]{Requirement 27.2 The primal pairing in the first-order bracket}{digest-meat-core/frank-wolfe/req:primal-pairing}
\renewcommand*{\RestatedName}{\hyperref{contract.pdf}{}{claim-meat-core/frank-wolfe/req:primal-pairing}{Requirement~27.2}}%
\begin{digestrequirement}[The primal pairing in the first-order bracket]
In the bracket of \eqref{eq:hlb} the engine uses the primal pairing
$\langle\lmoJ_t,S_t\rangle$ at the returned point $S_t$ and nothing else; in
particular it never uses a dual objective reported by the backend in its
place.
\end{digestrequirement}

\section*{28\quad Requirements: the certified evaluation}\pdfbookmark[1]{28 Requirements: the certified evaluation}{digest-group-meat-core/frank-wolfe/req:certified-evaluation}

% claim meat-core/frank-wolfe/req:certified-evaluation
\digestlabel{req:certified-evaluation}{28.1}\pdfbookmark[2]{Requirement 28.1 The certified evaluation of the first-order bound}{digest-meat-core/frank-wolfe/req:certified-evaluation}
\renewcommand*{\RestatedName}{\hyperref{contract.pdf}{}{claim-meat-core/frank-wolfe/req:certified-evaluation}{Requirement~28.1}}%
\begin{digestrequirement}[The certified evaluation of the first-order bound]
On every run that forms a first-order number, unless it sets the option
\iface{certified_evaluation} to \iface{off}, its default being \iface{on},
the engine forms a certified number as follows, after the first-order number
is formed and before a key length is reported.
\begin{enumerate}[label=(\roman*),leftmargin=2.4em]
\item \emph{The declaration and the points.} It works on the declaration the
  run's first-order number is formed on, at the points that number is formed
  from: on an admitted reference companion where
  Requirement~\ref{req:bracket-on-reference}\,(iii) applies, at the at most
  five of the points part~(i) of that requirement evaluates, of the search
  and of the polish together, that contribute there with the largest
  brackets on the companion, a point that is in both sets counted once and
  ties broken by the sequence in which the search and then the polish
  consumed them, each with the solve part~(ii) of it
  makes there; otherwise on the declaration the run searched, at the at most
  four members of the index set of its bound
  (Requirement~\ref{req:first-order-withdrawal}) with the largest brackets,
  ties broken by the order of the run, and the last point the search
  assessed, each with the search's solve there. \emph{The fallback}: where it
  has worked on a companion by (ii)--(v) and certified no point there, it
  works once more, by (ii)--(v), on the declaration the run searched, at the
  points and with the solves of the second case, unless the run declares the
  strict option \iface{strict_snaps} of
  \xref{foundations/req:snap-reference}\,(iv); the companion's results are
  kept for the record (ix). The fallback runs nowhere else: not where a point
  of the companion is certified, not where no companion's bracket counts, not
  where the companion's bracket is withdrawn by (iv), and not where the
  arithmetic the evaluation needs is not available. Where the companion's
  bracket is withdrawn by (iv), the run's first-order number is formed on the
  declaration the run searched (Requirement~\ref{req:bracket-on-reference}\,(iv)),
  and the evaluation works once more there, by the second case above, as on
  any run whose first-order number is formed there. Each point
  is taken exactly: its acceptance-side component as the exact rational
  iterate formed on the exact island
  (Requirement~\ref{req:acceptance-iterate-exact}), a point whose component
  was not so formed being not certified
  (\iface{acceptance_iterate_not_exact}), and each Choi block as the exact
  Hermitian part of the dyadic matrix its binary64 entries are.
\item \emph{The data.} Every quantity is formed from the data of the
  declaration it works on, taken exactly, every binary64 number as the dyadic
  rational it is and the acceptance set as its exact rational data, and, for
  the blocks a reference companion replaces, from the exact centres and radii
  of the entrywise enclosures its compilation carries, recentred as
  \xref{foundations/lem:recentred-enclosure} recentres them, and from the
  exact centres and radii of the matrix box and the exact core of
  every factor enclosure its compilation carries
  (\xref{foundations/def:factor-enclosure}). No map, frame,
  gradient, support, rank, eigenvalue or check value the engine computed in
  binary64 enters it. In trusted-source mode the frame of each Alice
  sector is $W_{\mathsf t}$ of Theorem~\ref{thm:choi} for the
  support frame the engine constructed
  (Definition~\ref{def:support-frame}), held as an enclosure formed by a
  rigorously enclosed computation; where the engine forms none, every point
  is not certified (\iface{frame_not_enclosed}).
\item \emph{The maps in factor form.} At every joint pair that is not
  key-inert, every kept one under a dropped set, the branch
  operators and their pinched blocks are formed through factors
  $S'_{xv}$ with the Gram matrices of
  Lemma~\ref{lem:kraus-freedom}\,(v), built from factor forms
  $E^{\dagger}\widetilde XE$ of the elements of Alice's measurement or
  ensemble and of the measured family that enter them, $\widetilde X$
  positive definite and $E$ either an exact matrix of full row rank
  (\xref{foundations/cor:least-eigenvalue-over-enclosure}\,(iii)) or, on a
  block whose rank the factor rule of (iv) fixes, the adjoint $L^{\dagger}$
  of a factor $L$ of its factor enclosure, held through the box of the
  adjoints, of full row rank at every factor
  (\xref{foundations/cor:rank-over-factor-enclosure}\,(ii)), with
  $\widetilde X$ the exact core: on exact data from an exact factorisation,
  on a block a companion carries with a positive radius by the rule of (iv).
  In the general format an element of the measured family is the tensor
  product of its blocks, one per measured factor, and its factor form is
  the tensor product of theirs (Lemma~\ref{lem:certified-bracket-enclosed}\,(b)),
  $E$ and $\widetilde X$ the tensor products of the blocks' $E$ and
  $\widetilde X$, however many of its blocks a companion replaces and
  however many of them are held through boxes: $E$ is then held through an
  enclosure of the Kronecker product of their enclosures, which contains it
  at every member, and has full row rank at every member, $EE^{\dagger}$
  being the tensor product of the blocks' $EE^{\dagger}$, each positive
  definite. No rank is read off a computed spectrum, and a matrix function applied to
  an operator $A$, a branch operator $X$ or a pinched block
  $(Z_j\otimes\id)X^{p}(Z_j\otimes\id)$ of one, acts on its compression
  $Q^{\dagger}AQ$ by an isometry $Q$, held through an enclosure, as follows.
  \begin{enumerate}[label=(\alph*),leftmargin=2.2em]
  \item \emph{Containment.} The range of $Q$ is established to contain
    $\operatorname{supp}A$, exactly or by an argument on the factors and the
    enclosures that holds at every value they enclose. The range of the
    side-by-side listing of the pair's pass-class factors contains the support
    of every branch operator they form, and the pinching projector applied to
    a subspace containing $\operatorname{supp}X$ contains the support of the
    pinched block (Lemma~\ref{lem:kraus-freedom}\,(iv)(a)). Where the
    listing factors as $T_{\mathrm{out}}\widetilde B\,T_{\mathrm{in}}$, with
    $\widetilde B$ an exact matrix of full column rank, formed exactly,
    $T_{\mathrm{out}}$ the maps applied after it, certified injective by a
    Gram matrix certified positive definite, and $T_{\mathrm{in}}$ the maps
    applied before it, an orthonormal basis of
    $\operatorname{ran}(T_{\mathrm{out}}\widetilde B)$ is established from
    $\widetilde B$. On exact data $\widetilde B$ comes from an exact rank
    factorisation of the exact matrix of the factors of the
    pass classes, and $T_{\mathrm{in}}$ holds the rest of
    that factorisation and the frame $W_{\mathsf t}\otimes\id$ in the Choi
    form. Where an element is held through the box of the adjoints of the
    factors of one or more of its blocks, $\widetilde B$ is, for each pass class, the exact factor of
    Lemma~\ref{lem:certified-bracket-enclosed}\,(b), formed by exact rank
    factorisations of the rows whose factors are exact and, per tuple of the
    blocks held through boxes, of the stack of Alice's factors tensored with
    the exact factors of the elements whose blocks held through boxes are
    that tuple, the kind that lemma fixes for a block being the third for a
    block that the factor rule of (iv)(b) decides not zero, fixing its rank,
    the first for a block decided zero, and the second otherwise, an exact
    block included. The
    rows of such elements are grouped by their tuples, each block named by
    its factor, its sector and its letter, and never by the equality of
    their boxes: elements whose tuples name the same blocks share a group,
    whatever their boxes, and two distinct blocks with equal boxes are
    distinct entries of the tuples. Where the form the engine holds of such
    an element is not of this shape (its parts and the blocks naming them
    differing in number, a split into exact and boxed parts held where no
    block of the element is held through a box or missing where one is, or
    one tuple of blocks at two positions or of two dimensions), no claim
    admits it, and every point is not certified
    (\iface{boxed_form_inconsistent}), a defect of the engine and not of the
    declaration. $T_{\mathrm{in}}$ holds that lemma's
    matrix $C$, the stack of $C_0$ and of each tuple's
    $(C_{\beta}\otimes E^{\mathrm{box}}_{\beta})U_{\beta}$, the columns of
    each such block permuted by its own $U_{\beta}$ and by no other
    permutation, which lies in an enclosure containing it at every member,
    and the frame. That range contains the range of the listing at every member,
    and equals it at every member where the product of the maps in
    $T_{\mathrm{in}}$ is established to have full row rank at every member;
    otherwise the basis spans a subspace containing the range, which (a)
    admits and the floor of (b) decides. $T_{\mathrm{out}}$ is block
    diagonal over the outcome pairs, so the pinching projectors commute with
    it, and $(Z_j\otimes\id)\operatorname{ran}(T_{\mathrm{out}}\widetilde B)
    =\operatorname{ran}\big(T_{\mathrm{out}}(Z_j\otimes\id)\widetilde B\big)$,
    whose basis comes in the same way from an exact rank factorisation of
    $(Z_j\otimes\id)\widetilde B$. A point at which containment is not
    established is not certified (\iface{support_not_established}).
  \item \emph{The floor.} A positive lower bound on $\lmin(Q^{\dagger}AQ)$
    is then established, over the enclosures at every value they contain,
    which with (a) proves $\operatorname{ran}Q=\operatorname{supp}A$
    (Lemma~\ref{lem:kraus-freedom}\,(iv)(b)) at every member and makes the
    compression that of part (iv)(c) of that lemma. A point at which it is
    not is not certified (\iface{floor_not_certified}).
  \end{enumerate}
\item \emph{The rank rule over an enclosure.} A block a companion carries
  with a positive radius is decided, and its rank fixed, by the first of
  two rules that does so.
  \begin{enumerate}[label=(\alph*),leftmargin=2.2em]
  \item \emph{The entrywise rule.} The block is decided when every entry of
    its entrywise enclosure has centre and radius $0$, then zero, or some
    entry's centre exceeds its radius in modulus, then not zero. Its factor
    form is that of
    \xref{foundations/cor:least-eigenvalue-over-enclosure}\,(iv): the rows
    and columns whose every entry has centre and radius $0$ are dropped, and
    the rest must have a positive lower bound on its least eigenvalue over
    the enclosure.
  \item \emph{The factor rule}, where the compilation carries a factor
    enclosure of the block and the entrywise rule leaves the block
    undecided or fixes no rank for it. With no column the block is decided
    zero (\xref{foundations/cor:rank-over-factor-enclosure}\,(i)). With
    $k\ge1$ columns and a floor of part (ii) of that corollary established,
    $G$ and $\rho$ formed as its part (iii) forms them or by a rigorously
    enclosed computation of $L^{\dagger}L$ over the box, the block is
    decided not zero and its factor form is $L^{\dagger}$, held through the
    box of the adjoints, with the core as $\widetilde X$.
  \item \emph{The cross-check}, wherever the compilation carries a factor
    enclosure of the block, whichever rule fixes its rank. The engine
    encloses $LWL^{\dagger}$ entry by entry over the box of the factor
    enclosure, in exact rational arithmetic or by a rigorously enclosed
    computation, so that the enclosure contains $LWL^{\dagger}$ for every
    factor $L$, every member of the factor enclosure; and it decides, as (b)
    decides them, whether the factor enclosure has no column and, where it
    has columns, whether a floor of part (ii) of
    \xref{foundations/cor:rank-over-factor-enclosure} is established.
  \end{enumerate}
  A block that neither rule decides is undecided, and a block whose rank
  neither fixes has no certified rank. The two enclosures of a block are
  \emph{inconsistent} where the enclosure of $LWL^{\dagger}$ of (c) and the
  block's entrywise enclosure are disjoint at some entry, or where one of the
  two decides the block zero and the other decides it not zero, as (a) and
  (c) decide them. The companion's own block lies in its entrywise enclosure
  and, by (H-FAC) of Lemma~\ref{lem:certified-bracket-enclosed}, is a member
  of its factor enclosure, so it lies in both enclosures and both decisions
  hold of it: an inconsistency shows that one of the two was not formed or
  carried as the claims that establish it say. Where the two enclosures of a
  block are inconsistent, the factor rule is set aside for the block, every
  point of the companion is not certified
  (\iface{factor_enclosure_inconsistent}), whatever joint pairs the block is a
  factor at, key-inert or not, and whichever rule would fix its rank, and the
  companion's bracket is withdrawn with that cause
  (Requirement~\ref{req:bracket-on-reference}\,(iv)), its first-order number
  reading the entrywise enclosure, which may be the defective one. On a recentred
  enclosure, the entrywise rule leaves a block undecided exactly when its
  centre vanishes and some radius does not, its enclosure then deciding
  neither that it vanishes nor that it does not. Where an undecided block is
  a factor of an element of an outcome pair of a pass class of a
  key-live acceptance value whose element of Alice's is not zero, or
  a block with no certified rank is such a factor at a joint pair that is
  not key-inert, every point is not certified (\iface{rank_not_certified}),
  save where another block of that element is decided zero, an exact block
  that is zero or a block a rule decides zero: the element is then zero for
  every member, whatever its other blocks.
  What the rule establishes is the hypothesis of
  Lemma~\ref{lem:certified-bracket-enclosed} that such blocks are decided
  and of fixed rank, in the sense of that lemma; its hypothesis (H-FAC), that
  the companion's block is a member of the factor enclosure carried with it,
  is established by the claims of the modules that form and carry that
  enclosure, and decided by no rule here.
\item \emph{What is established at a point.} At each point the engine
  establishes, each step in exact rational arithmetic or in a rigorously
  enclosed computation at a working precision of at least $128$ bits:
  \begin{enumerate}[label=(\alph*),leftmargin=2.2em]
  \item every Choi block positive definite, by a positive lower bound of
    \xref{foundations/lem:least-eigenvalue-bound} on its exact Hermitian
    part (\iface{block_not_positive_definite}), which gives (H-SUPP) by
    Lemma~\ref{lem:branch-support}\,(i);
  \item (H-POS), with $q_t(\hat c)>0$ decided exactly and the lower ends of
    the enclosures of $P^{J}_{\hat c}(J_t)$ and $K^{J}_{\hat c}(J_t)$
    positive (\iface{positivity_not_certified}), on the exact
    acceptance support $\Chacc$ of \eqref{eq:chacc-def}, which is
    also where $\hat g$ of (c) lives: for a singleton the support of its
    one point, and for a polytope the union of the supports of its exact
    vertices (Remark~\ref{C-rem:acc-side});
  \item an enclosure $(\hat g,\delta)$ of the acceptance-side gradient and,
    for every block $p$ of every group, an exact Hermitian $M_p$ and a
    rational $\rho_p$ with $\|R_p-M_p\|_{\mathrm F}\le\rho_p$ for the
    residual of Lemma~\ref{lem:certified-bracket}\,(iii), both formed by a
    rigorously enclosed computation of the gradient of
    Corollary~\ref{cor:full-gradient} at the point (of the dropped objective
    under a dropped set) and of the frame: the enclosure of $R_p$ is made
    Hermitian-centred with symmetric radii by
    \xref{foundations/lem:recentred-enclosure}\,(iv), $M_p$ is its centre
    rounded entrywise to dyadic rationals and kept Hermitian, and $\rho_p$
    is a rational at least the radius that
    \xref{foundations/cor:least-eigenvalue-over-enclosure}\,(ii) gives for
    that centre plus a rational bound on the Frobenius norm of the
    rounding, so that the triangle inequality gives
    $\|R_p-M_p\|_{\mathrm F}\le\rho_p$. On a companion, the computation
    reads each block of the measured family through its entrywise enclosure
    where it forms $P^{J}_{\hat c}$, the term $P^{J}_{\hat c}$ of
    $K^{J}_{\hat c}$ and the $\Gamma^{(s')}_{\hat c}$, and, where it forms
    the deficits and the $\Delta^{(s')}_{\hat c}$, through the enclosure
    that holds the block's factor form of (iii), its entrywise enclosure
    where the entrywise rule of (iv) fixed the rank and its factor
    enclosure where the factor rule did, so that the enclosures of (b) and
    (c) hold at every pair of a set as in
    Lemma~\ref{lem:certified-bracket-enclosed}\,(c). The
    multipliers $Y$ and $\Lambda_r$ are read off the dual of that point's
    solve as Requirement~\ref{req:verified-feasible-side}\,(ii) reads them in
    these respects only: the operators the dual's multipliers represent,
    each replaced by its exact Hermitian part; zero where the solve returned
    no finite dual; and the exact shift that makes the multiplier of an
    inequality row positive semidefinite. The rule of that item that zeroes
    the row multipliers of an Alice sector without a bound on its frame's
    error is not read, the rows here being built from the enclosed exact
    frame;
  \item $b^{\mathrm A}_t$ as the exact minimum of
    $\langle\hat g,\cdot\rangle$ over $\acSet$ for a geometry of
    Requirement~\ref{req:acceptance-geometries} that it solves exactly, less
    $\delta$, a point of any other geometry being not certified
    (\iface{acceptance_geometry_not_certified}); and $b^{\mathrm F}_t$ by
    Lemma~\ref{lem:certified-bracket}\,(iii), each $\ell_p$ from
    \xref{foundations/lem:least-eigenvalue-bound}, in exact rational
    arithmetic;
  \end{enumerate}
  and the point's certified bracket $b^{\mathrm A}_t+b^{\mathrm F}_t$. A
  point at which any of these is not established is not certified, with the
  reason named, and contributes nothing.
\item \emph{The number, its key length and its tier.} Where some point is
  certified, the run's certified number
  $\underline{h}^{\mathrm{cert}}_{\alpha}$ (\iface{certified_bound}) is the
  largest certified bracket, \eqref{eq:hcert}, an exact rational, entering
  every binary64 use rounded toward $-\infty$ as
  \xref{foundations/req:rounding-direction} rounds an emitted bound; the run
  reports it at the certified tier (\xref{foundations/req:tiers}).
  The key length it reports is that of
  Corollary~\ref{cor:certified-keylength} (on a reference companion, for the
  protocol it declares, Corollary~\ref{cor:certified-keylength-reference}) at
  $h\coloneqq\max\{0,\underline{h}^{\mathrm{cert}}_{\alpha}\}$, with the
  parameters of the run, and it is the maximal key length of
  Definition~\ref{def:max-key-length} with $h$ in place of $\hlow$, computed
  exactly: $c\coloneqq n\,h-\lEC-\lceil\log_2(1/\eev)\rceil+1$ is formed
  exactly, a rational, $\lceil\log_2(1/\eev)\rceil$ being the least integer
  $k$ with $2^{k}\ge1/\eev$, decided by comparing the rational $1/\eev$ with
  powers of two; the order $\alpha$ is taken exactly, as the dyadic rational
  its binary64 value is, so that $\frac{\alpha}{\alpha-1}$ is a non-zero
  rational; the term
  $\tau\coloneqq\frac{\alpha}{\alpha-1}\log_2(1/\epa)$ exactly where
  $1/\epa$ is a power of two, $\tau$ being then rational, and otherwise
  through an enclosure $[L,U]\ni\tau$ formed in exact arithmetic or by a
  rigorously enclosed computation and refined until
  $\lfloor c-U\rfloor=\lfloor c-L\rfloor$, which is then
  $\lfloor c-\tau\rfloor$; the refinement ends, $\tau$ being irrational
  there, as a non-zero rational multiple of $\log_2(1/\epa)$, which is
  irrational for a dyadic rational $\epa$ whose reciprocal is not a power
  of two, so that $c-\tau$ is not an integer. The key length is
  $\max\{0,\lfloor c-\tau\rfloor\}$, and the rule of
  Requirement~\ref{req:reported-key-length} is met as written. The clamp is
  recorded as
  Requirement~\ref{req:first-order-clamp} records it. The
  first-order number and the key length formed from it, by
  Requirement~\ref{req:bracket-on-reference}\,(iii) or without a companion,
  are recorded beside it (\iface{first_order_key_length}). In the fallback
  of (i), the certified number is the largest certified bracket on the
  declaration the run searched, a lower bound on that declaration's value
  and not on the companion's, and its key length is that of
  Corollary~\ref{cor:certified-keylength} on that declaration, never that of
  Corollary~\ref{cor:certified-keylength-reference}; it is reported at the
  certified tier, and the blocks that
  Requirement~\ref{req:bracket-on-reference}\,(v) names for the fallback
  rest on \xref{foundations/ass:snapped-measurement}. The companion's
  first-order number stays the run's first-order number, recorded beside it
  with its key length, and no number is formed from the two together: the
  key length is that of the fallback's certified number alone, whether the
  companion's first-order number gives a larger key length or a smaller one.
  Where the companion's bracket is withdrawn by (iv), the run's first-order
  number is that of Requirement~\ref{req:bracket-on-reference}\,(iv) on the
  declaration the run searched, a number certified there is a lower bound on
  that declaration's value with the key length of
  Corollary~\ref{cor:certified-keylength}, and the blocks the companion
  carried rest on \xref{foundations/ass:snapped-measurement} with the cause
  \iface{factor_enclosure_inconsistent}; a run that declares
  \iface{strict_snaps} refuses there (\xref{foundations/req:snap-reference}\,(iv)).
  Where no point is certified, or the arithmetic the evaluation needs is not
  available to the engine (\iface{certifier_unavailable}), the run's number,
  key length and tier are those it has without this requirement, recorded
  with the reason \iface{not_certified} and the reasons of every point.
\item \emph{Premises and evaluations upstream of the declaration.} A run is
  not reported at the certified tier where (a) an assumption of the category
  numerical-trust is invoked on it, reason
  \iface{numerical_trust_premise}: an assumption of its assumption closure
  counts as invoked unless the module that owns it records that it was
  invoked nowhere on the run, as the receiver records its eigenvalue-slack
  assumption beside the settings at which it was invoked
  (\xref{receiver/req:eig-ledger-row}); or (b) the declaration it is formed
  on carries a weight routed from a receiver's high-photon weights
  (\xref{reductions/wpfss/def:weight-route}) into which a per-setting weight
  enters that the receiver's artifact records as evaluated in floating point
  (\xref{receiver/req:evaluation-label}), the uniform weight of a slot entry
  counting the weight of every setting as entering it, reason
  \iface{receiver_weight_floating_point}. Its certified number is then
  recorded with the reason, and the run's number, key length and tier are
  those it has without this requirement.
\item \emph{The strict option.} A run that sets
  \iface{certified_evaluation} to \iface{strict} and is not reported at the
  certified tier by (vi) and (vii) refuses
  (\iface{certification_refused}), naming the reasons, and reports no key
  length. A run reported at the certified tier through the fallback of (i)
  is so reported under this option too and does not refuse: the premise it
  rests on there, \xref{foundations/ass:snapped-measurement}, is named in its
  record, and the option that refuses a run resting on it is
  \iface{strict_snaps}, under which the fallback does not run.
\item \emph{The record.} The run records the option; per point, whether it
  was certified, $b^{\mathrm A}_t$, $b^{\mathrm F}_t$, its certified bracket,
  the largest radius of its enclosure of the gradient and its reasons; and
  per run $\underline{h}^{\mathrm{cert}}_{\alpha}$, the largest such radius
  over the certified points (\iface{max_gradient_radius}), the working
  precision (\iface{working_precision_bits}), the library of the
  rigorous arithmetic with its version (\iface{arithmetic_library}), and,
  as a diagnostic, $\underline{h}^{\mathrm{cert}}_{\alpha}$ less the
  first-order number on the same declaration, with its sign
  (\iface{certified_minus_first_order}), which in the fallback is the bound
  on the searched declaration that
  Requirement~\ref{req:bracket-on-reference}\,(iii) records and is absent
  where that declaration has none. It records further the declaration the
  per-point entries were evaluated on (\iface{declaration}:
  \iface{reference_companion} or \iface{searched}); the premise at the
  snap the certified number rests on (\iface{premises}): the short id of
  \xref{foundations/ass:snapped-measurement} wherever a moved block or class
  of the declaration the number is formed on rests on that assumption
  (\xref{foundations/req:snap-reference}\,(iii)), among them, on the
  declaration the run searched, every moved block or class it has, in the
  fallback, without a companion and where the companion's bracket is
  withdrawn alike, and, on a companion, every moved block on which it carries
  the snapped block and every moved block of another snap it shares with the
  declaration the run searched, such as that of another measured factor;
  and none otherwise; with it, the blocks and classes that
  rest on it (\iface{resting_blocks}), with their groups and causes as the
  record of the run names them where the key length is formed from a number
  on that declaration, those of
  Requirement~\ref{req:bracket-on-reference}\,(v) among them in the
  fallback, with the cause \iface{companion_not_certified}; and, wherever
  the fallback ran or the companion's bracket was withdrawn by (iv), the
  companion's points, each with whether it was certified and its reasons
  (\iface{companion_points}), which say why. The fields \iface{premises}
  and \iface{resting_blocks} describe the declaration that
  \iface{declaration} names, whatever the status: they name what a number
  certified there rests on, also where none is reported at the certified
  tier. In a fallback not reported at the certified tier, the status being
  \iface{withheld} or \iface{not_certified}, they therefore name the blocks
  of Requirement~\ref{req:bracket-on-reference}\,(v) with the cause
  \iface{companion_not_certified}, under which that item names them where
  the fallback's number forms the key length, while the run's key length is
  formed from the companion's first-order number and the record of the run
  names those blocks as evaluated off their snap. None of these but
  $\underline{h}^{\mathrm{cert}}_{\alpha}$ enters a number.
\end{enumerate}
\end{digestrequirement}

\bibliographystyle{plain}
\bibliography{../../../refs/references}

\end{document}
